\chapter{Relative Picard sheaf --- \texttt{Scheme.Modules.tensorObj} substrate (A.1.c.SubT)}
\label{chap:Picard_TensorObjSubstrate}
% archon:covers AlgebraicJacobian/Picard/TensorObjSubstrate.lean
% archon:covers AlgebraicJacobian/Picard/TensorObjSubstrate/PullbackTensorComp.lean
% archon:covers AlgebraicJacobian/Picard/TensorObjSubstrate/StalkTensor.lean
% archon:covers AlgebraicJacobian/Picard/TensorObjSubstrate/Vestigial.lean
% archon:covers AlgebraicJacobian/Picard/TensorObjSubstrate/DualInverse.lean
% archon:covers AlgebraicJacobian/Picard/TensorObjSubstrate/PresheafInternalHom.lean

\section{Motivation: the abelian-group-valued relative Picard sheaf}
\label{sec:tensorobj_motivation}

Fix a base field \(k\) and a smooth proper geometrically integral curve \(C\) over
\(k\). In Kleiman's framework (\cref{chap:Picard_RelPicFunctor},
\cref{def:rel_pic_sharp}), the relative Picard functor
\[
  \Pic_{C/k} \;\;\colon\;\; (\Sch/k)^{op} \;\longrightarrow\; \AddCommGrpCat,
  \qquad
  T \;\longmapsto\; \Pic(C \times_k T) \, / \, \pi_T^* \Pic(T),
\]
is, by definition, abelian-group-valued. The abelian-group structure on each
fiber \(\Pic(C \times_k T)\) comes from tensor product of line bundles
(Stacks tag 01CR / Hartshorne~II.\S 6), and the structure descends to the
quotient because \(\pi_T^*\) is a group homomorphism.

% SOURCE: [Kleiman], "The Picard scheme", \S 2 (FGA Explained Ch.~9 \S 9.2),
% Definitions df:aPf + df:Pfs (the absolute and relative Picard functors are
% abelian-group-valued; the relative one is defined as a quotient of abelian
% groups)
% (read from references/kleiman-picard-src/kleiman-picard.tex, L1274-L1318).
% SOURCE QUOTE:
% "\begin{dfn}\label{df:aPf}
%  The {\it absolute Picard functor} $\Pic_X$ is the functor from  the
%  category of (locally Noetherian) $S$-schemes $T$ to the category of
% abelian groups defined by the formula
% 	$$\Pic_X(T):=\Pic(X_T).$$
% \end{dfn}
%  [\dots]
% \begin{dfn}\label{df:Pfs}
%   The {\it relative Picard functor} $\Pic_{X/S}$ is defined by
% 	$$\Pic_{X/S}(T):=\Pic(X_T)\big/\Pic(T).$$
%  Denote its associated sheaves in the Zariski, \'etale, and fppf
% topologies  by
% 	$$\Pic_{(X/S)\zar},\ \Pic_{(X/S)\et},\
% 	\Pic_{(X/S)\fppf}.$$
%  \end{dfn}"
\textit{Source: [Kleiman], ``The Picard scheme'', \S 2, Defs.~df:aPf + df:Pfs.
The absolute and relative Picard functors take values in the category of
abelian groups; the relative functor is the quotient of one such group by
another, and the structure morphisms of the resulting presheaf are group
homomorphisms.}

Concretely, the abelian-group law on \(\Pic(C \times_k T)\) is the assignment
\(\bigl[\,\mathcal{L}\,\bigr] + \bigl[\,\mathcal{L}'\,\bigr] :=
\bigl[\,\mathcal{L} \otimes_{\mathcal{O}_{C \times_k T}} \mathcal{L}'\,\bigr]\),
with neutral element the class of the structure sheaf
\(\mathcal{O}_{C \times_k T}\) and inverse
\(- \bigl[\,\mathcal{L}\,\bigr] := \bigl[\,\mathcal{L}^{-1}\,\bigr]\), where
\(\mathcal{L}^{-1} = \mathcal{H}om_{\mathcal{O}_{C \times_k T}}(\mathcal{L},
\mathcal{O}_{C \times_k T})\) is the dual invertible sheaf. Pullback of
\(\mathcal{O}\)-modules along \(\pi_T : C \times_k T \to T\) preserves tensor
products and the structure sheaf, so the image
\(\pi_T^* \Pic(T) \subset \Pic(C \times_k T)\) is a subgroup, and the quotient
\(\Pic(C \times_k T) / \pi_T^* \Pic(T)\) inherits a canonical abelian-group
structure. This is the content of
\cref{lem:rel_pic_sharp_groupoid} in chapter \cref{chap:Picard_RelPicFunctor}.

The mathematics is straightforward; the obstacle to formalisation is purely
infrastructural, and --- crucially --- much smaller than a full
monoidal-category construction. The group law lives on \emph{isomorphism
classes} of line bundles: each group axiom is a \emph{proposition} asserting
the existence of an isomorphism (\(\mathrm{Nonempty}(\,\cdots \cong \cdots)\)),
and a proposition is insensitive to which isomorphism witnesses it. In
particular the coherence data of a monoidal category --- the pentagon,
triangle and hexagon identities, and a fortiori any closed structure --- is
\emph{never consumed} by a group law on iso-classes. What the law genuinely
requires on the carrier is therefore only:
\begin{enumerate}
  \item a binary tensor-product operation
    \(\otimes\, : \Scheme.\mathtt{Modules}\,X \times \Scheme.\mathtt{Modules}\,X
                  \to \Scheme.\mathtt{Modules}\,X\), well-defined on
    isomorphism classes;
  \item the structure sheaf \(\mathcal{O}_X\) as a designated unit object
    for \(\otimes\);
  \item the \(\otimes\)-invertibility predicate
    (\(\exists N,\ M \otimes_X N \cong \mathcal{O}_X\)) together with the three
    monoidal coherence isomorphisms of \(\otimes_X\): associativity
    \((M \otimes N) \otimes P \cong M \otimes (N \otimes P)\), unit
    \(\mathcal{O}_X \otimes M \cong M\) (and on the right), and commutativity
    \(M \otimes N \cong N \otimes M\). The inverse is carried by the
    predicate, not by a separate lemma.
\end{enumerate}
At Mathlib's pinned commit one might be tempted to obtain (3) by instantiating a
\emph{full} symmetric monoidal category on \(\Scheme.\mathtt{Modules}\,X\) and
reading the three isomorphisms off it. Two Mathlib facts make such a route
conceivable. First, the category \(\mathtt{PresheafOfModules}\,R\) of presheaves
of modules over a presheaf of commutative rings \(R\) is already a (symmetric)
monoidal category in Mathlib (\(\mathtt{PresheafOfModules.monoidalCategory}\)),
with the sectionwise relative tensor \((M \otimes_R N)(U) = M(U) \otimes_{R(U)}
N(U)\); this applies verbatim to the \emph{varying} structure sheaf
\(R = \mathcal{O}_X\). Second, Mathlib's localization-of-monoidal-categories API
(\(\mathtt{CategoryTheory.Localization.Monoidal.LocalizedMonoidal}\)) takes a
monoidal category \(C\), a localization functor \(L : C \to D\) at a morphism
property \(W\), and an instance \(W.\mathtt{IsMonoidal}\), and produces a
\emph{full} monoidal-category structure on the localization \(D\); sheafification
of presheaves of modules is the localization \(L = \mathtt{sheafification}\) at the
locally-bijective class \(W = J.W\), so a full monoidal structure on
\(\Scheme.\mathtt{Modules}\,X\) would follow once \(J.W.\mathtt{IsMonoidal}\) were
supplied.

This full-monoidal-category route is, however, \emph{not} the goal, and is not
pursued. The goal is the commutative \emph{group} on locally-trivial iso-classes
(à la \(\mathtt{CommRing.Pic}\)), and that group law consumes only propositions, so
it never needs a coherent monoidal category --- neither the
pentagon/triangle/hexagon identities nor the \(J.W.\mathtt{IsMonoidal}\) instance
that the \(\mathtt{LocalizedMonoidal}\) API requires. What it needs is exactly the
\emph{existence} of the three coherence isomorphisms of (3), and these are produced
directly: the associator by gluing canonical local isomorphisms
(\cref{lem:tensorobj_assoc_iso}, via the single linchpin
\cref{lem:tensorobj_restrict_iso}), the unitors by the cheap \(\mathtt{mapIso}\)
pattern (\cref{lem:tensorobj_unit_iso}), and the braiding likewise
(\cref{lem:tensorobj_comm_iso}). With those three existence facts the group law
assembles \emph{by hand} the commutative group of \(\otimes\)-invertible
\(\otimes\)-iso-classes --- the Picard group
(\cref{lem:tensorobj_isoclass_commgroup}) --- each group-axiom field fed a single
existence-of-isomorphism. This is the scheme-theoretic analogue of Mathlib's
fixed-ring \(\mathtt{CommRing.Pic}\,R := \mathrm{Skeleton}(\mathtt{Module.Invertible}\,R)\),
whose \(\mathtt{CommGroup}\) is read off the \(\mathrm{Skeleton}\) of the coherent
monoidal category \(\mathtt{MonoidalCategory}(\mathtt{Module.Invertible}\,R)\) that
Mathlib supplies over a fixed ring; the present construction cannot take that
\(\mathrm{Skeleton}\) route, since the coherent monoidal category on
\(\Scheme.\mathtt{Modules}\,X\) for the varying structure sheaf is exactly what it
avoids. The
\(\otimes\)-invertibility predicate (\cref{def:scheme_modules_isinvertible},
Mathlib's \(\mathtt{Module.Invertible}\) idiom) carries each inverse existentially.
The \emph{full} \(\mathtt{LocalizedMonoidal}\) / \(J.W.\mathtt{IsMonoidal}\) coherent
monoidal category is therefore not built for the group law and is off the critical
path; what the group law does still need from this circle of ideas is the
\emph{existence} of the associator (\cref{lem:tensorobj_assoc_iso}), realized as a
sheafification transport whose single open left-whisker obligation is closed by the
stalk--tensor commutation (d.2). The stalk apparatus (d.1/d.2) is accordingly
\emph{on} the critical path, while the coherent monoidal category itself is not.

\section{Mathlib API survey}
\label{sec:tensorobj_api_survey}

The relevant pieces of the Mathlib API at the pinned commit
(\texttt{b80f227}) were originally assembled into a single \emph{realization route}
(route~(e)): obtaining a full monoidal structure on the substrate by instantiating
Mathlib's abstract localization-of-monoidal-categories API rather than
hand-assembling the associator. That route is \emph{superseded}. The group law on
iso-classes never consumes a coherent monoidal category, and the associator is
obtained by direct gluing through \cref{lem:tensorobj_restrict_iso}, so the
\(J.W.\mathtt{IsMonoidal}\) instance on which route~(e) hinges is never required.
The survey below is retained as a Mathlib-API reference and to record why the route
is vestigial; the four Mathlib facts are all verified against the on-disk Mathlib
source.

\paragraph{(i) The presheaf-of-modules monoidal category.} Mathlib's
\texttt{Mathlib.Algebra.Category.ModuleCat.Presheaf.Monoidal} constructs a
monoidal-category instance
\(\mathtt{PresheafOfModules.monoidalCategory}\) on
\(\mathtt{PresheafOfModules}\,(R \circ \mathtt{forget}_2\,\mathtt{CommRingCat}\,
\mathtt{RingCat})\), with the sectionwise relative tensor
\[
  (M \otimes_R N)(U) \;\;=\;\; M(U) \otimes_{R(U)} N(U)
\]
% SOURCE: [Mathlib], `PresheafOfModules.monoidalCategoryStruct` /
% `PresheafOfModules.monoidalCategory`, in
% `Mathlib/Algebra/Category/ModuleCat/Presheaf/Monoidal.lean` (L104, L125)
% (read from .lake/packages/mathlib/Mathlib/Algebra/Category/ModuleCat/Presheaf/Monoidal.lean).
% SOURCE QUOTE:
% "noncomputable instance monoidalCategoryStruct :
%     MonoidalCategoryStruct (PresheafOfModules.{u} (R ⋙ forget₂ _ _)) where
%   tensorObj := tensorObj
%   whiskerLeft _ _ _ g := tensorHom (𝟙 _) g
%   whiskerRight f _ := tensorHom f (𝟙 _)
%   ...
%   tensorUnit := unit _
%   associator M₁ M₂ M₃ := isoMk (fun _ ↦ α_ _ _ _) ...
% noncomputable instance monoidalCategory :
%     MonoidalCategory (PresheafOfModules.{u} (R ⋙ forget₂ _ _)) where ..."
\textit{Source: [Mathlib], \texttt{PresheafOfModules.monoidalCategoryStruct} /
\texttt{PresheafOfModules.monoidalCategory}, in
\texttt{Mathlib/Algebra/Category/ModuleCat/Presheaf/Monoidal.lean}, L104--L125.}
together with associator, left/right unitors and braiding making it symmetric
monoidal. The base presheaf of rings \(R \circ \mathtt{forget}_2\,
\mathtt{CommRingCat}\,\mathtt{RingCat}\) is exactly the project's
\(\mathtt{Sheaf.val}\,X.\mathtt{ringCatSheaf}\) carrier (equal by \(\mathtt{rfl}\)),
so this monoidal structure applies verbatim to the \emph{varying} structure sheaf
\(\mathcal{O}_X\) --- no project work. This is the presheaf-level analogue of
Stacks tag 03DM (relative tensor product of \(\mathcal{O}_X\)-modules).

\paragraph{(ii) The localization-of-monoidal-categories API.} Mathlib's
\texttt{Mathlib.CategoryTheory.Localization.Monoidal.Basic} provides the class
\(\mathtt{MorphismProperty.IsMonoidal}\) and the construction
\(\mathtt{Localization.Monoidal.LocalizedMonoidal}\). The class records exactly
the whiskering-stability data:
% SOURCE: [Mathlib], `CategoryTheory.MorphismProperty.IsMonoidal` and
% `CategoryTheory.Localization.Monoidal.LocalizedMonoidal`, in
% `Mathlib/CategoryTheory/Localization/Monoidal/Basic.lean` (L44, L86, L440)
% (read from .lake/packages/mathlib/Mathlib/CategoryTheory/Localization/Monoidal/Basic.lean).
% SOURCE QUOTE:
% "class IsMonoidal : Prop extends W.IsMultiplicative where
%   whiskerLeft (X : C) {Y₁ Y₂ : C} (g : Y₁ ⟶ Y₂) (hg : W g) : W (X ◁ g)
%   whiskerRight {X₁ X₂ : C} (f : X₁ ⟶ X₂) (hf : W f) (Y : C) : W (f ▷ Y)
%  ...
% def LocalizedMonoidal (L : C ⥤ D) (W : MorphismProperty C)
%     [W.IsMonoidal] [L.IsLocalization W] {unit : D} (_ : L.obj (𝟙_ C) ≅ unit) := D
%  ...
% noncomputable instance :
%     MonoidalCategory (LocalizedMonoidal L W ε) where ..."
\textit{Source: [Mathlib], \texttt{MorphismProperty.IsMonoidal} and
\texttt{Localization.Monoidal.LocalizedMonoidal}, in
\texttt{Mathlib/CategoryTheory/Localization/Monoidal/Basic.lean}, L44--L440.}
Given \([\MonoidalCategory\,C]\), \([W.\mathtt{IsMonoidal}]\), a localization
functor \(L : C \to D\) with \([L.\mathtt{IsLocalization}\,W]\) and an isomorphism
\(\varepsilon : L(\mathbf{1}_C) \cong \mathtt{unit}\), the localized category
\(\mathtt{LocalizedMonoidal}\,L\,W\,\varepsilon\) is a \(\MonoidalCategory\) with
\emph{all} coherence (associator, unitors, pentagon, triangle, naturalities)
derived. There is \emph{no} \(\mathtt{MonoidalClosed}\) anywhere in the
hypothesis chain.

\paragraph{(iii) The proof template for $W$-monoidality.} Mathlib's
\texttt{Mathlib.CategoryTheory.Sites.Point.IsMonoidalW} proves
\((J.W\,(A := A)).\mathtt{IsMonoidal}\) \emph{stalkwise} for presheaves valued in a
\emph{fixed} monoidal concrete category \(A\) (i.e.\ \(C^{op} \to A\)), via the
conservative-family characterisation \(\mathtt{hP.W\_iff}\),
\(\mathtt{Functor.Monoidal.map\_tensor}\) and typeclass inference:
% SOURCE: [Mathlib], `CategoryTheory.ObjectProperty.IsConservativeFamilyOfPoints.isMonoidal_W`
% and the `(J.W).IsMonoidal` instance, in
% `Mathlib/CategoryTheory/Sites/Point/IsMonoidalW.lean` (L48, L57)
% (read from .lake/packages/mathlib/Mathlib/CategoryTheory/Sites/Point/IsMonoidalW.lean).
% SOURCE QUOTE:
% "lemma ObjectProperty.IsConservativeFamilyOfPoints.isMonoidal_W
%     [J.HasSheafCompose (forget A)] :
%     (J.W (A := A)).IsMonoidal :=
%   .mk' _ (fun f g hf hg ↦ by
%     simp only [hP.W_iff (A := A)] at hf hg ⊢
%     intro Φ
%     rw [Functor.Monoidal.map_tensor]
%     infer_instance)"
\textit{Source: [Mathlib],
\texttt{ObjectProperty.IsConservativeFamilyOfPoints.isMonoidal\_W}, in
\texttt{Mathlib/CategoryTheory/Sites/Point/IsMonoidalW.lean}, L48--L57.}
This is the \emph{template} for the proof technique used in route~(e), but it does
\emph{not} apply directly to \(\mathtt{PresheafOfModules}\,R\): modules over a
\emph{varying} ring are not of the form \(C^{op} \to A\) for a fixed monoidal
\(A\).

\paragraph{The gap and route~(e), superseded.} The substrate
\(\Scheme.\mathtt{Modules}.\mathtt{tensorObj}\) of this chapter is the
sheafification of \(\mathtt{PresheafOfModules.Monoidal.tensorObj}\), and
sheafification of presheaves of modules is exactly the localization
\(L = \mathtt{sheafification}\) at the locally-bijective class \(W = J.W\). Route~(e)
proposed to obtain a \emph{full} monoidal structure on
\(\Scheme.\mathtt{Modules}\,X = \mathtt{SheafOfModules}\,X.\mathtt{ringCatSheaf}\)
by applying \(\mathtt{LocalizedMonoidal}\) (ii) to the already-monoidal category (i),
which would require the instance \(J.W.\mathtt{IsMonoidal}\) (its two
whiskering-stability fields \cref{lem:whisker_of_W}, with load-bearing
local-injectivity half \cref{lem:islocallyinjective_whisker_of_W}). This is
\emph{not} how the group law is realized. The group law consumes only the
\emph{existence} of the three coherence isomorphisms, and the associator is built
directly by gluing canonical local isomorphisms through
\cref{lem:tensorobj_restrict_iso} (\cref{lem:tensorobj_assoc_iso}); it never invokes
the \emph{full} \(J.W.\mathtt{IsMonoidal}\) instance or a coherent monoidal category.
The \emph{full monoidal-category} route~(e) is therefore off the critical path. The
associator does, however, consume a \emph{single} left-whisker existence lemma
\cref{lem:islocallyinjective_whiskerleft_via_stalk}, closed by the d.2 stalk--tensor
commutation; so the stalk apparatus (d.1/d.2) itself is on the critical path, even
though the full coherent monoidal category built from it is not.

The stalk-infrastructure analysis that route~(e) turned on is recorded here for
reference only, since that route is superseded. The \emph{stalk module structure}
is already present in Mathlib: the file
\texttt{Mathlib.Algebra.Category.ModuleCat.Stalk} supplies, for a topological
space \(X\), a presheaf of commutative rings \(R\) on \(X\) and a point
\(x \in X\), the \(R.\mathtt{stalk}\,x\)-module structure on the stalk
\(\mathtt{stalk}\,M.\mathtt{presheaf}\,x\) of a presheaf of modules \(M\),
together with the germ--scalar compatibility \(\mathtt{germ\_smul}\); a linearity
packaging of the induced stalk map of a morphism was built project-side
(\cref{lem:stalk_linear_map}, the \emph{linearity half} of ingredient (d.1)).
Route~(e) additionally required a (d.1)-\emph{bridge} between the site-level
\(J.W\) and the topological stalkwise-isomorphism criterion on
\(\mathtt{Opens}\,X\), and (d.2) the commutation \((A \otimes^{p} B)_x \cong A_x
\otimes_{R_x} B_x\) of the stalk with the relative module-presheaf tensor; there is
moreover no monoidal \(\mathtt{SheafOfModules}\) in Mathlib. The full \(J.W.\mathtt{IsMonoidal}\) coherent monoidal category that route~(e) would
build is not pursued; but the (d.1)-linearity packaging and the (d.2) stalk--tensor
commutation are \emph{on} the critical path, since the associator's single open
left-whisker obligation \cref{lem:islocallyinjective_whiskerleft_via_stalk} is closed
precisely by (d.2) (\cref{sec:tensorobj_stalk_tensor}). What is superseded is only the
\emph{pair}-of-whiskering-fields \(J.W.\mathtt{IsMonoidal}\) packaging and the
\(\mathtt{LocalizedMonoidal}\) monoidal category, not the stalk infrastructure.

The project's policy on this gap is \emph{project-side}: the substrate is
built inside \texttt{AlgebraicJacobian/Picard/} (the file
\(\texttt{AlgebraicJacobian/Picard/TensorObjSubstrate.lean}\) is its home)
and the consumer instance is closed against it. A parallel Mathlib upstream
PR with the same content is welcome but not on the critical path for the
project; the consumer is closed against the project-side substrate
independently of any upstream timing.

\section{The substrate \(\Scheme.\mathtt{Modules}.\mathtt{tensorObj}\)}
\label{sec:tensorobj_substrate}

\begin{definition}


\leanok
[The substrate operation $\otimes$ on $\Scheme.\mathtt{Modules}\,X$]
  \label{def:scheme_modules_tensorobj}
  \lean{AlgebraicGeometry.Scheme.Modules.tensorObj}
  Let \(X\) be a scheme. For
  \(M, N \in \Scheme.\mathtt{Modules}\,X\), define
  \[
    M \otimes_X N \;\;\in\;\; \Scheme.\mathtt{Modules}\,X
  \]
  as follows. Choose an affine open cover \(\{U_i = \Spec A_i\}_{i \in I}\)
  of \(X\). On each \(U_i\), the restriction
  \(M|_{U_i}\) corresponds (via the equivalence between
  \(\mathcal{O}_{\Spec A_i}\)-modules and \(A_i\)-modules) to an
  \(A_i\)-module \(M_i\), and similarly
  \(N|_{U_i}\) corresponds to an \(A_i\)-module \(N_i\). Define
  \((M \otimes_X N)|_{U_i}\) as the sheafification of the presheaf
  \(U_i \mapsto M_i \otimes_{A_i} N_i\). On overlaps
  \(U_i \cap U_j = \Spec A_{ij}\) the resulting modules agree via the
  canonical compatibility
  \((M_i \otimes_{A_i} N_i) \otimes_{A_i} A_{ij}
     \;\cong\;
     (M_j \otimes_{A_j} N_j) \otimes_{A_j} A_{ij}\),
  so the data \(\{(M \otimes_X N)|_{U_i}\}\) glue to a single object
  \(M \otimes_X N\) in \(\Scheme.\mathtt{Modules}\,X\). The result is
  canonically independent of the affine cover chosen.

  Equivalently: the underlying presheaf of \(M \otimes_X N\) is the
  presheaf-of-modules tensor product
  \(\mathtt{PresheafOfModules.Monoidal.tensorObj}\) of the underlying
  presheaves of \(M\) and \(N\), composed with the sheafification
  functor on the small Zariski site of \(X\). The substrate operation
  \(\otimes_X\) is, by construction, a \(\Scheme.\mathtt{Modules}\,X\)
  object reflecting the relative tensor product of
  \(\mathcal{O}_X\)-modules described in Stacks tag 03DM.
\end{definition}

\begin{lemma}


\leanok
[Functoriality of $\otimes_X$]
  \label{lem:scheme_modules_tensorobj_functoriality}
  \lean{AlgebraicGeometry.Scheme.Modules.tensorObj_functoriality}
  \uses{def:scheme_modules_tensorobj}
  Let \(X\) be a scheme. The assignment
  \((M, N) \mapsto M \otimes_X N\) of
  \cref{def:scheme_modules_tensorobj} extends to a bifunctor
  \[
    \otimes_X \;\;\colon\;\; \Scheme.\mathtt{Modules}\,X \;\times\;
                          \Scheme.\mathtt{Modules}\,X \;\;\longrightarrow\;\;
                          \Scheme.\mathtt{Modules}\,X
  \]
  natural in both arguments: a pair of morphisms
  \(f : M \to M'\) and \(g : N \to N'\) determines a morphism
  \(f \otimes g : M \otimes_X N \to M' \otimes_X N'\) compatible with
  identities and composition. This bifunctorial action on morphisms is
  the only datum the present lemma provides. The further coherence data
  of a monoidal structure --- the unitors
  \[
    \lambda_M \;\colon\; \mathcal{O}_X \otimes_X M \;\xrightarrow{\sim}\; M,
    \qquad
    \rho_M \;\colon\; M \otimes_X \mathcal{O}_X \;\xrightarrow{\sim}\; M,
  \]
  the associator
  \(\alpha_{M, N, P} : (M \otimes_X N) \otimes_X P \xrightarrow{\sim}
   M \otimes_X (N \otimes_X P)\), and the braiding
  \(\beta_{M, N} : M \otimes_X N \xrightarrow{\sim} N \otimes_X M\) ---
  are \emph{not} outputs of this lemma; they are off the critical path,
  recorded in \cref{rem:scheme_modules_monoidal_off_path}.
\end{lemma}

\begin{proof}
  \uses{def:scheme_modules_tensorobj}
  Pick an affine open cover \(\{U_i = \Spec A_i\}\) of \(X\). On each
  affine, the assignment is the standard tensor product of modules over a
  commutative ring \(A_i\), which is bifunctorial. The bifunctorial action
  on morphisms glues across the cover because the underlying gluing data of
  \cref{def:scheme_modules_tensorobj} respect the canonical morphisms on
  overlaps; equivalently, the morphism action is inherited from
  \(\mathtt{PresheafOfModules.Monoidal.tensorObj}\) (Mathlib module
  \texttt{Mathlib.CategoryTheory.Monoidal.PresheafOfModules}) under
  sheafification. The unitor/associator/braiding coherence data are
  \emph{not} constructed here; they are off the critical path
  (\cref{rem:scheme_modules_monoidal_off_path}).
\end{proof}

\begin{remark}
[The full monoidal category on all of $\Scheme.\mathtt{Modules}\,X$ is off-path; the group law consumes only propositions]
  \label{rem:scheme_modules_monoidal_off_path}
  \uses{def:scheme_modules_tensorobj, lem:scheme_modules_tensorobj_functoriality,
        lem:whisker_of_W}
  One could ask for the bifunctor \(\otimes_X\) of
  \cref{lem:scheme_modules_tensorobj_functoriality}, together with
  \(\mathcal{O}_X\), \(\alpha\), \(\lambda\), \(\rho\) and \(\beta\), to
  assemble into a full symmetric \(\MonoidalCategory\) structure on
  \emph{all} of \(\Scheme.\mathtt{Modules}\,X\), discharging the pentagon,
  triangle and hexagon coherence identities. This is explicitly \emph{off-path}
  and \emph{not pursued}. One could in principle obtain it as the output of
  \(\mathtt{Localization.Monoidal.LocalizedMonoidal}\) applied to the
  already-monoidal presheaf category \(\mathtt{PresheafOfModules}\,\mathcal{O}_X\)
  and the sheafification localizer \(J.W\), supplying the instance
  \(J.W.\mathtt{IsMonoidal}\) (\cref{lem:whisker_of_W}); but the group law never
  consumes a coherent monoidal category, so this construction is superseded and not
  carried out. The substrate provides only the \emph{existence} of the three
  coherence isomorphisms that the group law actually consumes.

  Were that superseded route pursued, the instance \(J.W.\mathtt{IsMonoidal}\)
  would need \emph{no} closed structure: its
  \(\mathtt{whiskerLeft}\) field \(W\,g \Rightarrow W(F \triangleleft g)\) for an
  \emph{arbitrary} module \(F\) is the flatness-free stalkwise statement
  \cref{lem:islocallyinjective_whisker_of_W} (a \(J.W\)-morphism \(g\) is a
  stalkwise isomorphism, so \((F \triangleleft g)_x = \mathrm{id}_{F_x} \otimes
  g_x\) is again an isomorphism for any \(F\)), the same technique Mathlib uses for
  fixed-base presheaves in
  \(\mathtt{CategoryTheory.Sites.Point.IsMonoidalW}\). It would use neither
  \(\mathtt{MonoidalClosed}(\mathtt{PresheafOfModules}\,R)\) nor the
  sectionwise-flatness hypothesis of \cref{lem:flat_whisker_localizer} (which is
  false over a non-affine open and is likewise off the critical path). The \emph{full}
  \(J.W.\mathtt{IsMonoidal}\) instance and the \(\mathtt{LocalizedMonoidal}\) monoidal
  category are not on the critical path: the group law is realized without them. The
  single left-whisker existence lemma and its underlying d.1/d.2 stalk apparatus, by
  contrast, \emph{are} on the critical path, since the associator's sheafification
  transport consumes exactly that one whisker obligation
  (\cref{lem:islocallyinjective_whiskerleft_via_stalk}, closed via d.2).

  Independently of how the monoidal structure is built, the relative Picard
  group law of \cref{lem:rel_pic_sharp_groupoid} is a structure on
  \emph{isomorphism classes} of line bundles, and every group axiom there is a
  \emph{proposition} of the form \(\mathrm{Nonempty}(\cdots \cong \cdots)\).
  Such a proposition is insensitive to which isomorphism witnesses it, so the
  coherence data of the monoidal category --- pentagon, triangle, hexagon, and
  a fortiori any monoidal-closed structure --- is never \emph{consumed} by the
  group law. The group law needs only the existence of the three coherence
  isomorphisms \cref{lem:tensorobj_assoc_iso,lem:tensorobj_unit_iso,%
  lem:tensorobj_comm_iso}, supplied directly: the associator by gluing canonical
  local isomorphisms through the restriction-compatibility isomorphism
  \cref{lem:tensorobj_restrict_iso} (\cref{lem:tensorobj_assoc_iso}), and the
  unitors and braiding by the cheap \(\mathtt{mapIso}\) pattern
  (\cref{lem:tensorobj_unit_iso,lem:tensorobj_comm_iso}). The carrier --- iso-classes
  of the substrate tensor \(\otimes_X\) --- is the scheme-theoretic analogue of
  Mathlib's Picard group \(\mathtt{CommRing.Pic}\,R := \mathrm{Skeleton}(
  \mathtt{Module.Invertible}\,R)\) (\texttt{Mathlib.RingTheory.PicardGroup}), but it
  is built \emph{differently}: Mathlib reads its \(\mathtt{CommGroup}\) off the
  \(\mathrm{Skeleton}\) of the coherent monoidal category
  \(\mathtt{MonoidalCategory}(\mathtt{Module.Invertible}\,R)\) over a fixed ring,
  whereas here --- with no coherent monoidal category on
  \(\Scheme.\mathtt{Modules}\,X\) available for the varying structure sheaf --- the
  group is assembled by hand from the existence-of-isomorphism facts.
\end{remark}

\section{The lift through \(\texttt{LineBundle.OnProduct}\)}
\label{sec:tensorobj_onproduct_lift}

The consumer in chapter \cref{chap:Picard_RelPicFunctor} works not with
\(\Scheme.\mathtt{Modules}(C \times_k T)\) directly but with the subtype
\(\texttt{LineBundle.OnProduct}\,\pi_C\,\pi_T\) of locally-trivial
modules on the fiber product \(C \times_k T\), introduced in
\cref{chap:Picard_LineBundlePullback}. This section records the
\emph{existence-of-isomorphism facts on line bundles} (associativity, unit,
commutativity and inverse) that the relative Picard group law consumes. These facts
are supplied \emph{directly}: associativity by the direct-gluing associator
\cref{lem:tensorobj_assoc_iso} (through the linchpin
\cref{lem:tensorobj_restrict_iso}), unit and commutativity by the cheap
\(\mathtt{mapIso}\) pattern (\cref{lem:tensorobj_unit_iso,lem:tensorobj_comm_iso}),
and the inverse existentially from \(\otimes\)-invertibility. They are \emph{not}
read off any \(\MonoidalCategory\) structure: the route that would derive them from
\cref{lem:jw_ismonoidal} (route~(e)) is superseded and off-path. The section also
records the closure of \(\otimes\) on the locally-trivial
\(\texttt{LineBundle.OnProduct}\) subtype (\cref{lem:tensorobj_lift_onproduct}) and
the interaction with the pullback functor \(\pi_T^*\). The open-immersion
restriction-compatibility isomorphism \cref{lem:tensorobj_restrict_iso} is on the
critical path precisely because it is consumed by the direct-gluing associator
\cref{lem:tensorobj_assoc_iso}; it is \emph{not} tied to the superseded
whiskering-stability lemma \cref{lem:islocallyinjective_whisker_of_W}, which belongs
to the off-path route~(e).

\begin{lemma}

\leanok
[Sectionwise lax-monoidal structure on \(\mathtt{restrictScalars}\,\varphi\)]
  \label{lem:restrictscalars_laxmonoidal}
  \lean{PresheafOfModules.restrictScalarsLaxMonoidal}
  \uses{def:scheme_modules_tensorobj}
  % NOTE: this CommRingCat-level lax-monoidal instance
  % (`PresheafOfModules.restrictScalarsLaxMonoidal`, with helpers
  % `restrictScalarsLaxε`, `restrictScalarsLaxμ`) is a self-contained
  % supplement. It is NOT used by `lem:tensorobj_restrict_iso`, whose proof
  % (open-immersion sectionwise base change) identifies restriction along an
  % open immersion with base change along a structure-sheaf isomorphism and
  % needs no lax-monoidal pushforward. Retained for potential reuse; off the
  % critical path for the tensor group law.
  Let \(\varphi : R \to S\) be a morphism of presheaves of \emph{commutative}
  rings on a site (so that both \(R\) and \(S\) factor through
  \(\mathtt{CommRingCat}\) and the category of presheaves of modules over each
  is monoidal). Then the presheaf-of-modules restriction-of-scalars functor
  \(\mathtt{PresheafOfModules.restrictScalars}\,\varphi\) is lax monoidal.
  Consequently, composing with the already-monoidal Mathlib functor
  \(\mathtt{pushforward}_0\mathtt{OfCommRingCat}\) via
  \(\mathtt{Functor.LaxMonoidal.comp}\), the full pushforward
  \[
    \mathtt{pushforward}\,\varphi
      \;\;:=\;\;
    \mathtt{pushforward}_0\,F\,R \;\circ\; \mathtt{restrictScalars}\,\varphi
  \]
  is lax monoidal. This lemma is a self-contained
  \(\mathtt{CommRingCat}\)-level lax-monoidal supplement: it stands on its own
  and is retained for potential reuse, but it is \emph{off the critical path}
  for the tensor group law. In particular it is \emph{not} an ingredient of
  \cref{lem:tensorobj_restrict_iso}, whose proof identifies restriction along an
  open immersion with base change along a structure-sheaf isomorphism and
  consequently needs no lax-monoidal structure on the pushforward.
\end{lemma}

\begin{proof}


\leanok
  This is a \emph{sectionwise} lift of the existing fact that, for a ring map
  \(f : R \to S\), the module restriction-of-scalars functor
  \(\mathtt{ModuleCat.restrictScalars}\,f\) is lax monoidal
  (\texttt{Mathlib.Algebra.Category.ModuleCat.Monoidal.Adjunction}, where it is
  obtained as the \(\mathtt{rightAdjointLaxMonoidal}\) of the
  extension/restriction adjunction \(\mathtt{extendRestrictScalarsAdj}\) ---
  i.e.\ Mathlib already uses the same mate idiom one level down). The presheaf
  monoidal structure is computed sectionwise,
  \((M_1 \otimes M_2).\mathrm{obj}\,X = M_1.\mathrm{obj}\,X \otimes M_2.\mathrm{obj}\,X\),
  and the presheaf functor \(\mathtt{restrictScalars}\,\varphi\) acts
  sectionwise as \(\mathtt{ModuleCat.restrictScalars}\,(\varphi.\mathrm{app}\,X)\).
  Hence the lax structure maps \(\mu\) (laxity) and \(\varepsilon\) (unit)
  assemble from the per-section \(\mathtt{ModuleCat}\) ones, the required
  naturality squares being the naturality of the sectionwise maps against the
  presheaf restriction morphisms. Packaging these as a
  \(\mathtt{Functor.CoreLaxMonoidal}\) yields the lax-monoidal structure; no
  new diagram chase is needed beyond transporting the sectionwise coherence.
  The composite \(\mathtt{pushforward}\,\varphi\) is then lax monoidal by
  \(\mathtt{Functor.LaxMonoidal.comp}\) together with the monoidality of
  \(\mathtt{pushforward}_0\mathtt{OfCommRingCat}\) (Mathlib), modulo a definitional
  transport across the unfolding of \(\mathtt{pushforward}\,\varphi\).

  The commutativity hypothesis is essential: the presheaf-of-modules monoidal
  category requires \(\mathtt{CommRingCat}\)-valued presheaves of rings, so the
  instance must be stated over \(\mathtt{CommRingCat}\)-factored \(R, S\). In
  the project this holds: the relevant \(\varphi\) is
  \((\mathtt{Scheme.Hom.toRingCatSheafHom}\,f).\mathrm{hom}\), which is
  structure-sheaf-valued and hence commutative-ring-valued.
\end{proof}

\begin{lemma}


\leanok
[Tensor product commutes with restriction along an open immersion]
  \label{lem:tensorobj_restrict_iso}
  \lean{AlgebraicGeometry.Scheme.Modules.tensorObj_restrict_iso}
  \uses{def:scheme_modules_tensorobj, lem:restrictscalars_ringiso_tensorequiv}
  Let \(X\) be a scheme, let \(M, N \in \Scheme.\mathtt{Modules}\,X\) be
  \emph{arbitrary} \(\mathcal{O}_X\)-modules, and let
  \(f : U \hookrightarrow X\) be an open immersion, with \((-)|_f\) the
  restriction of \(\mathcal{O}_X\)-modules along \(f\). The canonical comparison
  morphism
  \[
    (M \otimes_X N)\big|_f
      \;\xrightarrow{\;\;}\;
    M\big|_f \;\otimes_U\; N\big|_f
  \]
  is an isomorphism. No local-freeness, flatness, or line-bundle hypothesis on
  \(M, N\) is required: the statement holds for all \(\mathcal{O}_X\)-modules.
  The reason is structural to open immersions --- restriction along \(f\) is
  base change along the structure-sheaf \emph{isomorphism} induced by \(f\),
  not along a general ring map --- and base change along a ring isomorphism is
  trivially invertible and commutes with tensor products.
\end{lemma}

\begin{proof}
  \uses{def:scheme_modules_tensorobj}
  Write \(\varphi\) for the morphism of structure-sheaf-valued presheaves of
  rings induced by the open immersion \(f\) (concretely
  \(\varphi = (\mathtt{Scheme.Hom.toRingCatSheafHom}\,f).\mathrm{hom}\)), so that
  restriction along \(f\) is computed by the inverse-image (pullback) functor on
  presheaves of modules. The proof exploits the defining feature of an open
  immersion: the comparison of structure-sheaf rings over the image is an
  \emph{isomorphism}, not merely a ring map. For an open \(V \subseteq U\) the
  section ring \(\Gamma\bigl(X, f(V)\bigr)\) of \(\mathcal{O}_X\) over the image
  of \(V\) is identified with \(\Gamma(U, V)\) by the canonical \emph{ring
  isomorphism}
  \(f.\mathtt{appIso}\,V : \Gamma\bigl(X, f(V)\bigr) \xrightarrow{\sim}
  \Gamma(U, V)\)
  (the \(\mathtt{CommRingCat}\) isomorphism
  \(\mathtt{AlgebraicGeometry.Scheme.Hom.appIso}\)). Restriction along \(f\) is
  therefore extension of scalars along
  this ring isomorphism, and the proof needs no flatness anywhere.

  \emph{Step 1 (reduce restriction to the abstract pullback).} The restriction
  \((-)|_f\) of \(\mathcal{O}_X\)-modules along the open immersion \(f\) is
  naturally isomorphic to the abstract presheaf-of-modules pullback
  \((\mathtt{PresheafOfModules.pullback}\,\varphi).\mathrm{obj}(-)\), via the
  present comparison \(\mathtt{Scheme.Modules.restrictFunctorIsoPullback}\). It
  therefore suffices to produce the comparison and prove it an isomorphism for
  the abstract pullback functor.

  \emph{Step 2 (move the pullback inside the sheafification).} The substrate
  \(\otimes_X\) is sheafification of the presheaf-of-modules tensor product
  (\cref{def:scheme_modules_tensorobj}). The pullback commutes with
  sheafification: there is a natural isomorphism
  \(\mathtt{sheafification} \circ \mathtt{pullback}\,\varphi
   \;\cong\;
   \mathtt{PresheafOfModules.pullback}\,\varphi \circ \mathtt{sheafification}\)
  (the present \(\mathtt{SheafOfModules.sheafificationCompPullback}\)). Hence the
  whole question descends to the presheaf level, where it is the base-change
  comparison
  \[
    (\mathtt{pullback}\,\varphi).\mathrm{obj}(A \otimes B)
      \;\longrightarrow\;
    (\mathtt{pullback}\,\varphi).\mathrm{obj}\,A
      \;\otimes\;
    (\mathtt{pullback}\,\varphi).\mathrm{obj}\,B
  \]
  for the underlying presheaves of modules \(A, B\) of \(M\) and \(N\).

  % SOURCE: [Stacks Project], Modules on Ringed Spaces, Lemma
  % \ref{lemma-exactness-pushforward-pullback} (the pair $(f^*, f_*)$ of
  % pullback/pushforward of $\mathcal{O}$-modules along a morphism of ringed
  % spaces is an adjoint pair)
  % (read from references/stacks-modules.tex, L252-273).
  % SOURCE QUOTE:
  % "\begin{lemma}
  % \label{lemma-exactness-pushforward-pullback}
  % Let $f : (X, \mathcal{O}_X) \to (Y, \mathcal{O}_Y)$
  % be a morphism of ringed spaces.
  % \begin{enumerate}
  % \item The functor
  % $f_* : \textit{Mod}(\mathcal{O}_X) \to \textit{Mod}(\mathcal{O}_Y)$
  % is left exact. In fact it commutes with all limits.
  % \item The functor
  % $f^* : \textit{Mod}(\mathcal{O}_Y) \to \textit{Mod}(\mathcal{O}_X)$
  % is right exact. In fact it commutes with all colimits.
  %  [\dots]
  % \end{enumerate}
  % \end{lemma}
  %  [\dots]
  % Parts (1) and (2) hold because $(f^*, f_*)$ is an adjoint pair
  % of functors, see
  % Sheaves, Lemma \ref{sheaves-lemma-adjoint-pullback-pushforward-modules}"
  \textit{Source: [Stacks Project], Modules on Ringed Spaces, Lemma
  \texttt{lemma-exactness-pushforward-pullback}: pullback and pushforward of
  $\mathcal{O}$-modules along a morphism form an adjoint pair $(f^*, f_*)$, the
  abstract universal property on which the comparison below rests.}

  \emph{Step 3 (the genuine residual: H1 alone).} Completing this comparison at the
  presheaf level requires a sectionwise handle on the abstract functor
  \(\mathtt{PresheafOfModules.pullback}\,\varphi\), which is defined as the left
  adjoint of pushforward (the presheaf-of-modules avatar of the
  \(f^* \dashv f_*\) adjunction of the cited Stacks lemma) and admits no
  sectionwise formula in Mathlib at the pinned commit. The closure passes through
  two presheaf-level ingredients, of which exactly one now remains open.

  The first ingredient --- a strong-monoidal upgrade of
  \(\mathtt{restrictScalars}\) along the structure-sheaf ring \emph{isomorphism}
  \(f.\mathtt{appIso}\), so that base change along the isomorphism commutes with
  \(\otimes\) --- is \emph{closed in this chapter} as
  \cref{lem:restrictscalars_ringiso_tensorequiv}
  (\(\mathtt{restrictScalarsRingIsoTensorEquiv}\)): for a ring isomorphism
  \(e : R \simeq S\) the canonical map \(a \otimes_R b \mapsto a \otimes_S b\) is an
  \(R\)-linear equivalence \(\mathtt{restrictScalars}_e(M) \otimes_R
  \mathtt{restrictScalars}_e(N) \xrightarrow{\sim} \mathtt{restrictScalars}_e(M
  \otimes_S N)\).

  The \emph{sole remaining residual} is therefore the presheaf-level identification
  \[
    \mathtt{pushforward}\,\beta \;\cong\; \mathtt{pullback}\,\varphi
  \]
  via \(\mathtt{Adjunction.leftAdjointUniq}\). One exhibits a
  presheaf-level adjunction
  \(\mathtt{pushforward}\,\beta \dashv \mathtt{pushforward}\,\varphi\) and concludes
  that \(\mathtt{pushforward}\,\beta\), being a left adjoint of
  \(\mathtt{pushforward}\,\varphi\), is canonically isomorphic to the abstract left
  adjoint \(\mathtt{pullback}\,\varphi\). Here \(\beta\) is the structure map of the
  restriction functor, so that \((M|_f).\mathrm{val} =
  (\mathtt{pushforward}\,\beta).\mathrm{obj}\,M.\mathrm{val}\) definitionally;
  composing the resulting comparison with
  \cref{lem:restrictscalars_ringiso_tensorequiv} yields the asserted isomorphism.

  This presheaf-level adjunction is the presheaf analogue of the sheaf-level
  \(\mathtt{SheafOfModules.pushforwardPushforwardAdj}\), and must be built from
  presheaf-level \(\mathtt{pushforwardNatTrans}\) and
  \(\mathtt{pushforwardCongr}\). These two building blocks are \emph{Mathlib-absent}
  at the pinned commit: only the \emph{sheaf}-level versions exist
  (\texttt{Mathlib/.../SheafOfModules/Sheaf/PushforwardContinuous.lean}), whereas
  \(\mathtt{pushforwardId}\) and \(\mathtt{pushforwardComp}\) do already exist at the
  presheaf level. This residual is being constructed \emph{directly}, mirroring the
  existing sheaf-level template app-for-app at the presheaf level (a Mathlib-gradient
  build, not a deferral to an upstream PR).

  This lemma is the single substrate linchpin of the relative Picard group law for
  \emph{arbitrary} \(\mathcal{O}_X\)-modules \(M, N\): the associator
  \cref{lem:tensorobj_assoc_iso} consumes the comparison
  \((M \otimes_X N)|_f \cong M|_f \otimes_U N|_f\) on its sole consumer
  \cref{lem:tensorobj_preserves_locally_trivial}, where the restriction is commuted
  past the \emph{sheafified} tensor \(M \otimes_X N\) \emph{before} any local
  trivialisation of \(M, N\) is introduced --- so the comparison is required for the
  global, untrivialised modules, and no free-cover specialisation can stand in for it.
  Consequently the H1 residual \(\mathtt{pushforward}\,\beta \cong
  \mathtt{pullback}\,\varphi\) is genuinely on the critical path, and closing it is
  the \emph{single open substrate obligation} of this chapter. The superseded
  route~(e) \((J.W).\mathtt{IsMonoidal} \to \mathtt{Sheaf.monoidalCategory}\)
  (\cref{lem:jw_ismonoidal}) remains named only as a last-resort fallback should the
  presheaf-level adjunction itself bottom out; it is not the plan.
\end{proof}

\begin{lemma}
[Presheaf-level pushforward adjunction and strong-monoidal restriction (H1/H2 substrate)]
  \label{lem:presheaf_pushforward_adj_substrate}
  \lean{PresheafOfModules.pushforwardNatTrans,
        PresheafOfModules.pushforwardCongr,
        PresheafOfModules.pushforwardPushforwardAdj,
        PresheafOfModules.isIso_of_isIso_app,
        PresheafOfModules.restrictScalarsMonoidalOfBijective}
  \uses{def:scheme_modules_tensorobj}
  These five presheaf-level declarations are the two ingredients on which the
  linchpin \cref{lem:tensorobj_restrict_iso} rests; both are absent from Mathlib at
  the pinned commit and are supplied project-side.

  \emph{H1 (the pushforward adjunction).} For a morphism \(\beta\) of presheaves of
  rings, \(\mathtt{pushforwardNatTrans}\) and \(\mathtt{pushforwardCongr}\) are the
  presheaf-level de-sheafifications of Mathlib's sheaf-level
  \(\mathtt{SheafOfModules.pushforwardNatTrans}\) and
  \(\mathtt{SheafOfModules.pushforwardCongr}\)
  (\texttt{Mathlib/Algebra/Category/ModuleCat/Sheaf/PushforwardContinuous.lean}),
  reproduced app-for-app one level down at the presheaf of modules. From them
  \(\mathtt{pushforwardPushforwardAdj}\) assembles the presheaf-level adjunction
  \(\mathtt{pushforward}\,\beta \dashv \mathtt{pushforward}\,\varphi\), the
  presheaf analogue of the sheaf-level
  \(\mathtt{SheafOfModules.pushforwardPushforwardAdj}\). Pairing this against the
  existing \(\mathtt{PresheafOfModules.pullbackPushforwardAdjunction}\) through
  \(\mathtt{Adjunction.leftAdjointUniq}\) yields the canonical isomorphism
  \(\mathtt{pushforward}\,\beta \cong \mathtt{pullback}\,\varphi\) of two left
  adjoints of the same pushforward --- exactly the residual identified in Step~3 of
  the proof of \cref{lem:tensorobj_restrict_iso}.

  \emph{H2 (strong-monoidal restriction).} \(\mathtt{isIso\_of\_isIso\_app}\) is the
  sectionwise isomorphism criterion: a morphism of presheaves of modules is an
  isomorphism once each of its section maps is. Built on it,
  \(\mathtt{restrictScalarsMonoidalOfBijective}\) equips the presheaf-level
  restriction-of-scalars functor
  \(\mathtt{PresheafOfModules.restrictScalars}\,\alpha\) with a \emph{strong}-monoidal
  structure whenever the base-change comparison \(\alpha\) is \emph{sectionwise
  bijective} --- the presheaf-level lift of the ModuleCat-level H2 core
  \cref{lem:restrictscalars_ringiso_strongmonoidal}, assembled section-by-section via
  \(\mathtt{Functor.Monoidal.ofLaxMonoidal}\). It is to be distinguished from the
  module-level, ring-equivalence form
  \(\mathtt{restrictScalarsMonoidalOfRingEquiv}\) already pinned in
  \cref{lem:restrictscalars_ringiso_strongmonoidal}: the present declaration is the
  \emph{presheaf}-level, \emph{bijective}-comparison form, the one the substrate lift
  actually consumes (where the comparison map
  \((\alpha.\mathtt{app}\,X).\mathtt{hom}\) is a bijective \(\mathtt{CommRingCat}\)
  morphism but is not presented as a \((\,\cdot\,)\!.\mathtt{toRingHom}\)).

  Together H1 and H2 are the two Mathlib-absent ingredients the linchpin needed; they
  are self-contained presheaf-level results and are natural upstream-PR candidates.
\end{lemma}

\begin{lemma}

\leanok
[Base change along a ring isomorphism commutes with $\otimes$]
  \label{lem:restrictscalars_ringiso_tensorequiv}
  \lean{restrictScalarsRingIsoTensorEquiv}
  Let \(e : R \simeq S\) be an isomorphism of commutative rings, and let \(M, N\) be
  \(S\)-modules. Restriction of scalars along \(e\) commutes with the tensor
  product: the canonical map
  \[
    \mathtt{restrictScalars}_e(M) \otimes_R \mathtt{restrictScalars}_e(N)
      \;\xrightarrow{\ \sim\ }\;
    \mathtt{restrictScalars}_e\bigl(M \otimes_S N\bigr),
    \qquad a \otimes_R b \;\longmapsto\; a \otimes_S b,
  \]
  is an \(R\)-linear equivalence. Equivalently, \(\mathtt{restrictScalars}_e\) is
  \emph{strong} monoidal, not merely lax, when \(e\) is a ring \emph{isomorphism}:
  base change along an isomorphism \(R \simeq S\) identifies the two ground rings, so
  the lax structure map of \(\mathtt{ModuleCat.restrictScalars}\) is invertible. The
  forward map is \(R\)-linear, but its inverse \(a \otimes_S b \mapsto a \otimes_R b\)
  is only \emph{additive} (it is not \(S\)-linear), so the inverse is constructed as
  a homomorphism of additive groups and then verified to be the two-sided inverse of
  the forward \(R\)-linear map. This closes the strong-monoidal ingredient of
  \cref{lem:tensorobj_restrict_iso}, leaving the presheaf-level pushforward
  adjunction as the sole open residual there.
\end{lemma}

\begin{lemma}

[Strong-monoidal upgrade of \(\mathtt{restrictScalars}\) along a ring isomorphism (ModuleCat-level H2 core)]
  \label{lem:restrictscalars_ringiso_strongmonoidal}
  \lean{restrictScalars_isIso_μ, restrictScalars_isIso_ε,
        restrictScalarsMonoidalOfRingEquiv,
        restrictScalars_isIso_μ_of_bijective,
        restrictScalars_isIso_ε_of_bijective}
  \uses{lem:restrictscalars_ringiso_tensorequiv}
  Let \(e : R \simeq S\) be an isomorphism of commutative rings and let
  \(\mathtt{ModuleCat.restrictScalars}\,e.\mathtt{toRingHom} :
  \mathtt{ModuleCat}\,S \to \mathtt{ModuleCat}\,R\) be the associated
  restriction-of-scalars functor, equipped with its canonical \emph{lax}-monoidal
  structure (the lax tensorator \(\mu\) and lax unit \(\varepsilon\) supplied by
  Mathlib). Then:
  \begin{enumerate}
    \item the lax tensorator
      \(\mu_{M, N} : \mathtt{restrictScalars}_e M \otimes_R \mathtt{restrictScalars}_e N
      \to \mathtt{restrictScalars}_e(M \otimes_S N)\) is an isomorphism
      (\(\mathtt{restrictScalars\_isIso\_\mu}\));
    \item the lax unit
      \(\varepsilon : \mathbf{1}_R \to \mathtt{restrictScalars}_e(\mathbf{1}_S)\)
      is an isomorphism (\(\mathtt{restrictScalars\_isIso\_\varepsilon}\));
    \item consequently \(\mathtt{restrictScalars}\,e.\mathtt{toRingHom}\) is
      \emph{strong} monoidal: it carries the packaged structure
      \((\mathtt{ModuleCat.restrictScalars}\,e.\mathtt{toRingHom}).\mathtt{Monoidal}\)
      (\(\mathtt{restrictScalarsMonoidalOfRingEquiv}\)).
  \end{enumerate}
  The same two invertibility statements hold, more flexibly, under the hypothesis that
  the underlying ring homomorphism is merely \emph{bijective} rather than literally of
  the form \(e.\mathtt{toRingHom}\) for a packaged ring equivalence \(e\)
  (\(\mathtt{restrictScalars\_isIso\_\mu\_of\_bijective}\),
  \(\mathtt{restrictScalars\_isIso\_\varepsilon\_of\_bijective}\)); these bijective
  forms are the ones consumed by the presheaf-level lift, where the relevant
  comparison map \((\alpha.\mathtt{app}\,X).\mathtt{hom}\) is a bijective
  \(\mathtt{CommRingCat}\) morphism but is not presented as a
  \((\,\cdot\,)\!.\mathtt{toRingHom}\).
\end{lemma}

\begin{proof}
  \uses{lem:restrictscalars_ringiso_tensorequiv}
  Statements (1) and (2) are the categorical repackaging of
  \cref{lem:restrictscalars_ringiso_tensorequiv}: when the ground-ring map is a ring
  \emph{isomorphism} the two ground rings are identified, so the lax structure maps
  \(\mu\) and \(\varepsilon\) of \(\mathtt{ModuleCat.restrictScalars}\) --- which are
  in general only one-directional --- become invertible. For (1) the inverse of
  \(\mu_{M, N}\) is the additive map \(a \otimes_S b \mapsto a \otimes_R b\) of
  \cref{lem:restrictscalars_ringiso_tensorequiv}, verified to be two-sided inverse to
  the canonical \(R\)-linear \(\mu_{M, N}\); for (2) the unit comparison
  \(\mathbf{1}_R = R \to \mathtt{restrictScalars}_e S\) is the additive identification
  induced by \(e\). Strong monoidality (3) is then the Mathlib upgrade of a
  lax-monoidal functor all of whose structure maps are isomorphisms
  (\(\mathtt{Functor.Monoidal.ofLaxMonoidal}\)), packaging \(\mu^{-1}, \varepsilon^{-1}\)
  into \((\mathtt{ModuleCat.restrictScalars}\,e.\mathtt{toRingHom}).\mathtt{Monoidal}\).
  The bijective-hypothesis forms are proved identically, replacing ``\(e\) is a ring
  equivalence'' by ``the underlying ring homomorphism is bijective'' and reconstructing
  the inverse ground-ring map from bijectivity; this is what lets the presheaf-level
  lift feed in a section-wise-bijective comparison map directly.

  These declarations constitute the \emph{ModuleCat-level H2}: strong
  \(\mathtt{restrictScalars}\) along a ring isomorphism commutes with \(\otimes\). The
  one \emph{remaining} H2 step is the bounded, no-Mathlib-gap \emph{presheaf} lift:
  assembling \((\mathtt{PresheafOfModules.restrictScalars}\,\alpha).\mathtt{Monoidal}\)
  section-by-section via \(\mathtt{Functor.Monoidal.ofLaxMonoidal}\), using a
  \(\mathtt{CommRingCat}\)-iso\,\(\Rightarrow\)\,bijective bridge to invoke the
  bijective forms above at each section, together with the
  \(\mathtt{PresheafOfModules}\) iso-iff-app-iso reflection to lift section-wise
  invertibility to the presheaf level; once that lift is in place,
  \(\mathtt{pushforward}\,\beta\) is monoidal. This presheaf H2 lift is bounded and
  has no Mathlib-absent ingredient; the sole genuinely Mathlib-absent piece of
  \cref{lem:tensorobj_restrict_iso} is the presheaf-level pushforward adjunction
  (``H1'').
\end{proof}

\begin{lemma}


\leanok
[Tensor product of locally-trivial modules is locally trivial]
  \label{lem:tensorobj_preserves_locally_trivial}
  \lean{AlgebraicGeometry.Scheme.Modules.tensorObj_isLocallyTrivial}
  \uses{def:scheme_modules_tensorobj}
  Let \(X\) be a scheme, and let
  \(L, L' \in \Scheme.\mathtt{Modules}\,X\) be locally-trivial modules of
  rank one (i.e.\ line bundles). The tensor product \(L \otimes_X L'\) of
  \cref{def:scheme_modules_tensorobj} is again locally trivial of rank
  one. Consequently the bifunctor \(\otimes_X\) restricts to a bifunctor
  on the full subcategory \(\texttt{LineBundle}\,X \subset \Scheme.\mathtt{Modules}\,X\)
  of locally-trivial rank-one modules.
\end{lemma}

\begin{proof}




  Pick a common affine open cover \(\{U_i\}\) of \(X\) on which both
  \(L|_{U_i}\) and \(L'|_{U_i}\) are free of rank one. On each \(U_i\)
  there are trivialisations \(L|_{U_i} \cong \mathcal{O}_{U_i}\) and
  \(L'|_{U_i} \cong \mathcal{O}_{U_i}\). The tensor product of two free
  rank-one modules over a commutative ring is again free of rank one;
  hence \((L \otimes_X L')|_{U_i} \cong \mathcal{O}_{U_i}\). So
  \(L \otimes_X L'\) is locally trivial of rank one, i.e.\ an object of
  \(\texttt{LineBundle}\,X\).
\end{proof}

\emph{Closed standalone result --- off the critical path.} The following lemma
\emph{is} formalized and sorry-free; it is not, however, consumed by the relative
Picard group law. The associator \cref{lem:tensorobj_assoc_iso} is realized
\emph{unconditionally} through the varying-ring stalk--tensor commutation
\cref{lem:islocallyinjective_whiskerleft_via_stalk}
(\cref{sec:tensorobj_stalk_tensor}), which closes the single left-whiskering
obligation for \emph{all} modules with no flatness and no local-triviality
hypothesis. The sectionwise-flatness whiskering statement below is therefore
neither the live whiskering-stability route nor on the critical path: it is a valid,
self-contained result about flat whiskering of sheafification-localizer morphisms,
retained for potential reuse. (Unlike the genuinely abandoned
\cref{lem:whisker_of_W}, it is a closed lemma, not a stub to be skipped.)

\begin{lemma}

\leanok
% Closed, sorry-free standalone result (flat whiskering). Off the critical path:
% the associator is realized via the d.2 stalk-tensor route
% (lem:islocallyinjective_whiskerleft_via_stalk), not via this flat lemma.

[Flat whiskering preserves the sheafification localizer]
  \label{lem:flat_whisker_localizer}
  \lean{PresheafOfModules.W_whiskerLeft_of_flat}
  \uses{def:scheme_modules_tensorobj}
  Let \(R\) be a presheaf of commutative rings on a site carrying a Grothendieck
  topology \(J\), and let \(J.W\) be the class of morphisms of presheaves of
  \(R\)-modules that are inverted by sheafification --- equivalently, the
  \emph{locally bijective} morphisms, i.e.\ those that are both locally injective
  and locally surjective for \(J\) (the characterisation
  \(\mathtt{WEqualsLocallyBijective}\)). Let \(F\) be a presheaf of \(R\)-modules
  that is \emph{flat} (sectionwise: \(F(U)\) is a flat \(R(U)\)-module for every
  object \(U\) of the site). Then for every morphism \(g\) in \(J.W\) the
  left-whiskered morphism
  \[
    F \triangleleft g \;\;=\;\; \mathrm{id}_F \otimes g
  \]
  (the action of \(F \otimes^{p} -\) on \(g\), with \(\otimes^{p}\) the
  presheaf-of-modules tensor product) again lies in \(J.W\). By the symmetry of
  the presheaf-of-modules monoidal structure, the right-whiskered morphism
  \(g \triangleright F\) likewise lies in \(J.W\) when \(F\) is flat.

  Both halves --- the left-whiskering statement and its braiding-conjugate
  right-whiskering statement --- are established. They are valid standalone
  results about flat whiskering of sheafification-localizer morphisms. They are,
  however, \emph{off the critical path}: the live whiskering-stability obligation
  for the unconditional associator is closed for an \emph{arbitrary} whiskering
  object \(F\) by the stalk--tensor commutation
  \cref{lem:islocallyinjective_whiskerleft_via_stalk}
  (\cref{sec:tensorobj_stalk_tensor}), which never invokes the
  sectionwise-flatness hypothesis of the present lemma. The sectionwise flatness
  condition --- \(F(U)\) flat over \(R(U)\) for \emph{every} object \(U\) of the
  site, including non-affine opens --- is not supplied by \(\otimes\)-invertibility
  or local triviality, which are affine-local properties; the present (flat) lemma
  is therefore retained as an independent, closed result, superseded on the critical
  path by the d.2 route of \cref{sec:tensorobj_stalk_tensor}.
\end{lemma}

\begin{proof}
  \uses{def:scheme_modules_tensorobj}
  \leanok
  By the locally-bijective characterisation of the sheafification localizer
  (\(J.W\) is exactly the class of morphisms that are \(J\)-locally injective and
  \(J\)-locally surjective), it suffices to show that \(F \triangleleft g\) is
  both locally injective and locally surjective whenever \(g\) is.

  The presheaf-of-modules tensor product is computed sectionwise,
  \((F \otimes^{p} G)(U) = F(U) \otimes_{R(U)} G(U)\), so for
  \(g : A \to B\) the whiskered morphism acts on sections over \(U\) as the
  \(R(U)\)-linear map
  \[
    \mathrm{id}_{F(U)} \otimes g(U) \;\colon\;
      F(U) \otimes_{R(U)} A(U) \;\longrightarrow\; F(U) \otimes_{R(U)} B(U).
  \]

  \emph{Local surjectivity} is automatic and needs no hypothesis on \(F\):
  tensoring is right exact, so applying \(F \otimes^{p} -\) carries the cokernel
  of \(g\) to the cokernel of \(F \triangleleft g\); since \(g\) is locally
  surjective its cokernel is locally zero, hence so is the cokernel of
  \(F \triangleleft g\), and \(F \triangleleft g\) is locally surjective.

  \emph{Local injectivity} is where flatness enters. Because each \(F(U)\) is a
  flat \(R(U)\)-module, the functor \(F(U) \otimes_{R(U)} -\) preserves injective
  linear maps (Mathlib
  \(\mathtt{Module.Flat.lTensor\_preserves\_injective\_linearMap}\)). Thus
  whenever \(g(U)\) is injective on a covering sieve, so is
  \((\mathrm{id}_{F(U)} \otimes g(U))\); the local injectivity of \(g\) is
  therefore transported sectionwise to \(F \triangleleft g\).

  Locally injective and locally surjective together give
  \(F \triangleleft g \in J.W\). The right-whiskered statement
  \(g \triangleright F\) follows by conjugating \(F \triangleleft g\) with the
  braiding of the symmetric presheaf-of-modules monoidal structure, which is an
  isomorphism and hence preserves membership in \(J.W\); the flatness of \(F\) is
  the only input, and it is symmetric in the two tensor factors.
\end{proof}

\emph{Superseded route --- off path, not to be formalized.} The following lemma
belongs to the abandoned coherent-monoidal-category / sheafification-whiskering
approach to the associator. The group law on iso-classes consumes only the
\emph{existence} of the coherence isomorphisms (\cref{lem:tensorobj_isoclass_commgroup}),
and the associator is realized as the sheafification transport of the presheaf
associator (\cref{lem:tensorobj_assoc_iso}), whose single open left-whisker obligation
is closed by the d.2 route of \cref{lem:islocallyinjective_whiskerleft_via_stalk}. The
\emph{full} \((J.W).\mathtt{IsMonoidal}\) / \(\mathtt{LocalizedMonoidal}\)
monoidal-category packaging is not required. This block is retained for the historical
record only and must \emph{not} be formalized.

\begin{lemma}
% LOAD-BEARING (NOT pending deletion): the Lean lemma is invoked by the associator `tensorObj_assoc_iso`, by `pullbackUnitIso`, and by `isIso_sheafifyDelta_unitPair_of_isIso_sheafifyEta`. It is wired into the graph as a dependency of \cref{lem:tensorobj_assoc_iso}.

\leanok
[Sheafification inverts a localizer morphism of presheaves of modules]
  \label{lem:isiso_sheafification_map_of_W}
  \lean{PresheafOfModules.isIso_sheafification_map_of_W}
  Let \(R_0\) be a presheaf of commutative rings on a site carrying a
  Grothendieck topology \(J\), let \(R_{sh}\) be its sheafification, and let
  \(\alpha : R_0 \to R_{sh}.\mathrm{obj}\) be the (locally bijective) comparison
  map of presheaves of rings. Let \(f : A \to B\) be a morphism of presheaves of
  \(R_0\)-modules whose underlying morphism of presheaves of abelian groups
  \((\mathtt{toPresheaf}\,R_0).\mathrm{map}\,f\) lies in the sheafification
  localizer class \(J.W\) --- equivalently, is locally bijective for \(J\). Then
  the presheaf-of-modules sheafification functor
  \(\mathtt{PresheafOfModules.sheafification}\,\alpha\) sends \(f\) to an
  isomorphism:
  \[
    (\mathtt{sheafification}\,\alpha).\mathrm{map}\,f
    \quad\text{is an isomorphism.}
  \]
  This is a closed go/no-go bridge --- it converts ``\,a morphism lies in
  \(J.W\)\,'' into ``\,its sheafification is invertible\,''. Under route~(e) the
  associator is the API-derived one (\cref{lem:jw_ismonoidal}), so this bridge is
  retained as a supplement rather than as the associator's load-bearing step.
\end{lemma}

\begin{proof}

\leanok
  This is a one-morphism reading of the structural identity that the
  sheafification of presheaves of \(R_0\)-modules \emph{is} the localization of
  \(\mathtt{PresheafOfModules}\,R_0\) at the class
  \(J.W.\mathtt{inverseImage}\,(\mathtt{toPresheaf}\,R_0)\) of morphisms whose
  underlying map of presheaves of abelian groups is locally bijective. Concretely,
  the underlying \(\mathtt{AddCommGrp}\)-valued sheafification inverts exactly the
  locally bijective morphisms, and Mathlib's
  \(\mathtt{PresheafOfModules.inverseImage\_W\_toPresheaf\_eq\_inverseImage\_%
  isomorphisms}\) identifies the preimage of \(J.W\) under
  \(\mathtt{toPresheaf}\,R_0\) with the preimage of the class of isomorphisms under
  the module-level sheafification functor. Reading that equality of morphism
  classes at the single morphism \(f\): the hypothesis
  \(J.W\,\bigl((\mathtt{toPresheaf}\,R_0).\mathrm{map}\,f\bigr)\) places \(f\) in
  that preimage, hence \((\mathtt{sheafification}\,\alpha).\mathrm{map}\,f\) is an
  isomorphism. No new construction is required beyond instantiating the existing
  Mathlib identity at one morphism; this lemma is a closed thin wrapper, retained
  as a supplement under route~(e).
\end{proof}

\section{Route~(e): the \(J.W\)-monoidality instance and the localized monoidal structure}
\label{sec:tensorobj_route_e}

\emph{Full-monoidal-category construction --- off path, not to be formalized.} The
construction below realizes a \emph{full} symmetric monoidal category on all of
\(\Scheme.\mathtt{Modules}\,X\) by instantiating Mathlib's
localization-of-monoidal-categories API. It was the original plan for supplying the
associator. The relative Picard group law, however, is a structure on
\emph{isomorphism classes} of invertible sheaves and consumes only the
\emph{existence} of the associator, unitor and braiding isomorphisms
(\cref{lem:tensorobj_isoclass_commgroup}, built by hand as the scheme-theoretic
analogue of Mathlib's \(\mathtt{CommRing.Pic}\), not off a coherent monoidal
category's \(\mathtt{Skeleton}\)). The associator is obtained as the sheafification
transport of the presheaf associator (\cref{lem:tensorobj_assoc_iso}), whose
\emph{single} open left-whisker obligation is closed by the d.2 stalk--tensor
commutation (\cref{lem:islocallyinjective_whiskerleft_via_stalk,%
lem:stalk_tensor_commutation}); it uses \emph{no} \emph{full}
\((J.W).\mathtt{IsMonoidal}\) instance and \emph{no} \(\mathtt{LocalizedMonoidal}\)
monoidal category. Accordingly the \emph{full monoidal-category} target of this
section --- the \emph{pair} of whiskering-stability fields (\cref{lem:whisker_of_W})
and the \(\mathtt{LocalizedMonoidal}\) instance (\cref{lem:jw_ismonoidal}) --- is off
path and must not be re-attempted; it is kept for the historical record. The d.1/d.2
stalk apparatus it shares with the live single whisker obligation is \emph{not}
vestigial, however: it is on the critical path through
\cref{lem:islocallyinjective_whiskerleft_via_stalk}.

The monoidal structure on the substrate is realized by \emph{instantiating}
Mathlib's localization-of-monoidal-categories API
(\cref{sec:tensorobj_api_survey}): apply
\(\mathtt{Localization.Monoidal.LocalizedMonoidal}\) to the already-monoidal
presheaf category \(\mathtt{PresheafOfModules}\,\mathcal{O}_X\) and the
sheafification localizer \(J.W\). The single new input the API requires is the
instance \(J.W.\mathtt{IsMonoidal}\), whose two fields are the
whiskering-stability lemmas below.

\paragraph{What the group law consumes.} The group law is assembled directly on
isomorphism classes of \emph{locally trivial} (invertible) sheaves
(\cref{lem:tensorobj_isoclass_commgroup}), mirroring Mathlib's ring-level Picard
group \(\mathtt{Module.Invertible}\) / \(\mathtt{CommRing.Pic}\)
(\texttt{Mathlib/RingTheory/PicardGroup.lean}). It needs only the \emph{existence} of
the associator \cref{lem:tensorobj_assoc_iso}, the unitors \cref{lem:tensorobj_unit_iso}
and the braiding \cref{lem:tensorobj_comm_iso}, together with \(\otimes\)-invertibility
supplying inverses; it uses \emph{no} \(J.W.\mathtt{IsMonoidal}\) instance, \emph{no}
\(\mathtt{LocalizedMonoidal}\) monoidal category, and \emph{no} \(\mathtt{Skeleton}\).
It does, however, consume the associator \cref{lem:tensorobj_assoc_iso}, which is
realized as the sheafification transport of the presheaf associator and whose
\emph{single} open left-whisker obligation \cref{lem:islocallyinjective_whiskerleft_via_stalk}
is closed by the stalk--tensor commutation (d.2, \cref{lem:stalk_tensor_commutation}).
Hence the stalk ingredients (d.1)/(d.2) are \emph{on} the critical path of the group
law --- through the associator's single whisker obligation --- even though the full
\(J.W.\mathtt{IsMonoidal} \to \mathtt{LocalizedMonoidal}\) monoidal category on all of
\(\mathtt{SheafOfModules}\) (the subject of the present section) is \emph{not} built:
that full monoidal generality, with its \emph{pair} of whiskering-stability fields
(\cref{lem:whisker_of_W}) and all derived coherence, is what the group law does not
consume. The single open obligation of the present section's full-monoidal target is
thus the local-injectivity half packaged for an arbitrary whisker object
(\cref{lem:islocallyinjective_whisker_of_W}), itself a consequence of the same d.2
commutation.

The following lemma is the \(R_x\)-linear stalk-map packaging (the ``d.1''
ingredient). It is consumed by the stalk--tensor commutation of
\cref{sec:tensorobj_stalk_tensor} (\cref{lem:stalk_tensor_commutation,
lem:islocallyinjective_whiskerleft_via_stalk}) and is therefore a live ingredient
of the unconditional associator route.

\begin{lemma}

\leanok
% Live ingredient: consumed by lem:stalk_tensor_commutation and
% lem:islocallyinjective_whiskerleft_via_stalk (sec:tensorobj_stalk_tensor).

[The $R.\mathtt{stalk}\,x$-linear stalk map of a morphism of presheaves of modules]
  \label{lem:stalk_linear_map}
  \lean{PresheafOfModules.stalkLinearMap}
  \uses{def:scheme_modules_tensorobj}
  Let \(X\) be a topological space, \(R\) a presheaf of commutative rings on
  \(X\), and \(x \in X\) a point. For a morphism \(g : M \to N\) of presheaves of
  \(R\)-modules, the induced map on stalks
  \[
    g_x \;\colon\; M_x \;\longrightarrow\; N_x
  \]
  is \(R.\mathtt{stalk}\,x\)-linear (\(\mathtt{stalkLinearMap}\)), and is
  characterised on germs by
  \(g_x(\mathrm{germ}_x\,s) = \mathrm{germ}_x(g_U\,s)\)
  (\(\mathtt{stalkLinearMap\_germ}\)). When the underlying map of
  \(\mathtt{Ab}\)-valued stalks is an isomorphism, \(g_x\) is bijective
  (\(\mathtt{stalkLinearMap\_bijective\_of\_isIso}\)); in that case it bundles to
  an \(R.\mathtt{stalk}\,x\)-linear equivalence
  \(M_x \xrightarrow{\sim} N_x\) (\(\mathtt{stalkLinearEquivOfIsIso}\)).

  These four declarations are the \emph{(d.1)-partial implementation} --- the
  \(R_x\)-linearity packaging of stalk maps --- consumed by the
  \(\mathrm{id} \otimes g_x\) step in the proofs of
  \cref{lem:stalk_tensor_commutation} and
  \cref{lem:islocallyinjective_whiskerleft_via_stalk}. They build on the
  stalk module structure of \texttt{Mathlib.Algebra.Category.ModuleCat.Stalk};
  the linearity packaging itself is project-original.
\end{lemma}

\begin{proof}
  \uses{def:scheme_modules_tensorobj}
  \leanok
  The stalk \(M_x = \mathtt{stalk}\,M.\mathtt{presheaf}\,x\) carries its
  \(R.\mathtt{stalk}\,x\)-module structure from
  \texttt{Mathlib.Algebra.Category.ModuleCat.Stalk}. The underlying additive map
  \(g_x\) on stalks is the filtered colimit of the section maps \(g_U\); since
  each section map is \(R(U)\)-linear and the scalar action on the stalk is the
  colimit of the section actions, \(g_x\) respects the
  \(R.\mathtt{stalk}\,x\)-action, which gives the linear map together with its
  germ characterisation. Bijectivity is read off the underlying
  \(\mathtt{Ab}\)-stalk isomorphism (a bijective linear map is a linear
  equivalence), and bundling the inverse yields the linear equivalence.
\end{proof}

\emph{Off-path duplicate --- the full-generality packaging, not to be formalized.}
The following lemma is the load-bearing whiskering-stability field of the \emph{full}
\((J.W).\mathtt{IsMonoidal}\) / \(\mathtt{LocalizedMonoidal}\) monoidal-category route
to the associator. That \emph{full monoidal category} is not built: the group law on
iso-classes consumes only the \emph{existence} of the associator
(\cref{lem:tensorobj_isoclass_commgroup}). The associator instead arises as the
sheafification transport whose single open left-whisker obligation is closed by the
live lemma \cref{lem:islocallyinjective_whiskerleft_via_stalk}. The stalk ingredients
(d.1)/(d.2) themselves --- the \(R_x\)-linear stalk map \cref{lem:stalk_linear_map}
and the stalk--tensor commutation \cref{lem:stalk_tensor_commutation} --- are
\emph{on} the critical path through that live lemma; what is off path is only the
\emph{full-generality} \(J.W.\mathtt{IsMonoidal}\) packaging of the following
duplicate, retained for the historical record.

\begin{lemma}
% SUPERSEDED route (route-(d)/(e) whiskering/stalk machinery); pending deletion once the assoc re-route (morphism-level descent) lands. Declaration still present and still backing the current `tensorObj_assoc_iso` proof.

\leanok
[Local injectivity of a whiskered localizer morphism (stalkwise, flatness-free)]
  \label{lem:islocallyinjective_whisker_of_W}
  % \lean: this superseded, off-critical-path obligation is the SAME statement
  % (left-whiskering of a J.W-morphism by an arbitrary F is locally injective) as
  % the LIVE lemma lem:islocallyinjective_whiskerleft_via_stalk, and is realised by
  % the same sorry-free decl. Both blueprint nodes legitimately share the pin
  % (leandag flags only Lean-side name collisions, not shared blueprint pins).
  \lean{PresheafOfModules.isLocallyInjective_whiskerLeft_of_W}
  \uses{def:scheme_modules_tensorobj, lem:stalk_linear_map,
        lem:tensorobj_restrict_iso, lem:tensorobj_unit_iso}
  Let \(R\) be a presheaf of commutative rings on a site carrying a Grothendieck
  topology \(J\), and let \(J.W\) be the sheafification localizer (the locally
  bijective morphisms). Concretely the site carries the
  \(\mathtt{WEqualsLocallyBijective}\) property for \(\mathtt{Ab}\), so \(J.W\) is
  exactly the class of \(J\)-locally bijective morphisms; this is recorded as a
  typeclass hypothesis \([J.\mathtt{WEqualsLocallyBijective}\,\mathtt{Ab}]\) on
  the statement. Let \(F\) be an \emph{arbitrary} presheaf of
  \(R\)-modules, and let \(g : M \to N\) be a morphism of presheaves of
  \(R\)-modules whose underlying morphism of presheaves of abelian groups
  \((\mathtt{toPresheaf}\,R).\mathrm{map}\,g\) lies in \(J.W\). Then the
  left-whiskered morphism
  \[
    F \triangleleft g \;\;=\;\; \mathrm{id}_F \otimes g
  \]
  is \emph{\(J\)-locally injective}. No flatness and no local-triviality
  hypothesis on \(F\) is required.

  This is the \(\mathtt{whiskerLeft}\)-field content (the injectivity half) of the
  \emph{full} instance \(J.W.\mathtt{IsMonoidal}\) (\cref{lem:jw_ismonoidal}), the
  load-bearing field of the off-path \(\mathtt{LocalizedMonoidal}\) monoidal-category
  path (arbitrary \(F\)). Its content for an arbitrary \(F\) is exactly the live lemma
  \cref{lem:islocallyinjective_whiskerleft_via_stalk}, which the associator's
  sheafification transport consumes and which is closed \emph{unconditionally} by the
  d.2 stalk--tensor commutation (\cref{lem:stalk_tensor_commutation}); the d.1/d.2
  ingredients are therefore on the critical path. The two proofs recorded below --- a
  sheaf-level argument specialised to locally trivial \(F\), and the general stalkwise
  argument via (d.1)/(d.2) --- are retained for the historical record; the live group
  law closes this obligation in its general (d.2) form through
  \cref{lem:islocallyinjective_whiskerleft_via_stalk}.
\end{lemma}

\begin{proof}
  \uses{def:scheme_modules_tensorobj, lem:stalk_linear_map,
        lem:tensorobj_restrict_iso, lem:tensorobj_unit_iso}


\leanok
  \emph{PRIMARY route (\(F\) locally trivial --- the consumer's case; sheaf-level,
  no stalks, no (d.2)).} Suppose \(F\) is locally trivial
  (\(\mathtt{LineBundle.IsLocallyTrivial}\,F\)), which is exactly the hypothesis its
  sole consumer --- the associator \cref{lem:tensorobj_assoc_iso} --- supplies.
  \(J\)-local injectivity of a morphism is \emph{local on a covering family}: it
  suffices to produce a cover of \(X\) on each member of which \(F \triangleleft g\)
  restricts to a locally injective morphism. Take a trivialising open cover
  \(\{V\}\) of \(X\) for \(F\), so that \(F|_V \cong \mathcal{O}_V\) on each \(V\).
  On such a \(V\):
  \begin{itemize}
    \item the substrate--restriction compatibility \cref{lem:tensorobj_restrict_iso}
      gives a canonical isomorphism
      \[
        (F \triangleleft g)\big|_V \;\cong\; (F|_V) \triangleleft (g|_V)
      \]
      --- restriction along the open immersion \(V \hookrightarrow X\) commutes with
      \(\otimes\), and the whiskering \(F \triangleleft g = \mathrm{id}_F \otimes g\)
      is \(\otimes\) with an identity, so it restricts to \(\mathrm{id}_{F|_V}
      \otimes (g|_V) = (F|_V) \triangleleft (g|_V)\);
    \item the trivialisation \(F|_V \cong \mathcal{O}_V\) together with the left
      unitor \cref{lem:tensorobj_unit_iso} (\(\mathcal{O}_V \otimes_V (-) \cong
      (-)\)) identifies \((F|_V) \triangleleft (g|_V) \cong g|_V\).
  \end{itemize}
  Composing, \((F \triangleleft g)|_V\) is isomorphic to \(g|_V\). Now \(g \in J.W\),
  and membership in \(J.W\) (local bijectivity) is stable under restriction along an
  open immersion, so \(g|_V \in J.W\); in particular \(g|_V\) is \(J\)-locally
  injective. An isomorphism preserves local injectivity, hence \((F \triangleleft
  g)|_V\) is locally injective on \(V\). As the \(\{V\}\) cover \(X\) and local
  injectivity is local on the cover, \(F \triangleleft g\) is \(J\)-locally
  injective. \emph{No stalks and no (d.2) stalk--tensor commutation enter}; the
  single comparison lemma \cref{lem:tensorobj_restrict_iso} (a presheaf-level
  comparison, decomposed into its two open-immersion base-change halves) carries the
  argument, and thereby moves \emph{onto} the critical path.

  This is \emph{not} the section-level route that tensors a bare section map \(g(V)\)
  --- only injective --- and stalls on a \(\mathrm{Tor}_1\)/flatness obligation. Here
  the argument is sheaf-level and uses that \(g\) is locally \emph{bijective} (a
  \(J.W\) morphism), reduced through the unitor to \(g|_V\) (again a \(J.W\)
  morphism); a mere injection is never tensored, so no flatness is needed.

  \emph{FALLBACK route (arbitrary \(F\); stalkwise, required only for the
  full-generality route~(e) \(\mathtt{LocalizedMonoidal}\) path).} For an
  \emph{arbitrary} (not necessarily locally trivial) \(F\), local injectivity is
  proved stalkwise, by the argument Mathlib uses for fixed-base presheaves in
  \(\mathtt{CategoryTheory.Sites.Point.IsMonoidalW}\)
  (\cref{sec:tensorobj_api_survey}), ported to \(\mathtt{PresheafOfModules}\). A
  morphism in \(J.W\) is, on the topological site of \(X\), a \emph{stalkwise
  isomorphism}: at each point \(x\) the induced map on stalks \(g_x\) is an
  isomorphism. The presheaf-of-modules tensor is sectionwise, so it commutes with
  the stalk (a filtered colimit), giving
  \[
    (F \triangleleft g)_x
      \;=\; \mathrm{id}_{F_x} \otimes_{R_x} g_x
      \;\colon\; F_x \otimes_{R_x} M_x \;\longrightarrow\; F_x \otimes_{R_x} N_x .
  \]
  Tensoring an isomorphism \(g_x\) by the identity \(\mathrm{id}_{F_x}\) yields an
  isomorphism --- no flatness is needed, since \(g_x\) is invertible, not merely
  injective. Hence \(F \triangleleft g\) is a stalkwise isomorphism, in particular
  locally bijective, in particular locally injective, for \emph{any} \(F\).

  The stalkwise computation rests on three ingredients with distinct status at the
  pinned commit:
  \begin{enumerate}
    \item[(d.1-done)] the \(R.\mathtt{stalk}\,x\)-linear stalk-map packaging
      (\(\mathtt{stalkLinearMap}\) and its companions, \cref{lem:stalk_linear_map}):
      a morphism of presheaves of modules induces an \(R.\mathtt{stalk}\,x\)-linear
      map on stalks, characterised on germs, that becomes a linear equivalence once
      the underlying \(\mathtt{Ab}\)-stalk map is an isomorphism. This is built
      project-side, axiom-clean, and furnishes the \(\mathrm{id} \otimes g_x\) step
      above.
    \item[(d.1-bridge)] (remaining) the site-\(J.W\) \(\Longleftrightarrow\)
      stalkwise-iso characterisation on \(\mathtt{Opens}\,X\): that a morphism lies
      in \(J.W\) iff its stalk maps are isomorphisms at every point. The concrete
      Mathlib assembly route uses the \(\mathtt{WEqualsLocallyBijective}\) structure
      together with the topological stalk criteria --- local injectivity
      \(\Longleftrightarrow\) injectivity on stalks and local surjectivity
      \(\Longleftrightarrow\) surjectivity on stalks --- bridging the site-level
      \(\mathtt{Presheaf.IsLocallyInjective}\) / \(\mathtt{IsLocallySurjective}\)
      to the \(\mathtt{TopCat.Presheaf}\) stalk versions. This half is still
      Mathlib-absent at the \(\mathtt{PresheafOfModules}\) level.
    \item[(d.2)] (remaining) the commutation \((A \otimes^{p} B)_x \cong A_x
      \otimes_{R_x} B_x\) of the stalk (a filtered colimit) with the relative
      module-presheaf tensor, under which \((F \triangleleft g)_x\) is identified
      with \(\mathrm{lTensor}\,F_x\,(g_x) = \mathrm{id}_{F_x} \otimes g_x\) (using
      that \(\mathtt{tensorLeft}\) / \(\mathtt{tensorRight}\) preserve filtered
      colimits over a module category). This is genuinely Mathlib-absent over a
      \emph{varying} ring and is the largest remaining piece.
  \end{enumerate}
  Mathlib supplies the analogues for presheaves valued in a fixed monoidal concrete
  category, but not for modules over a varying ring; porting (d.1-bridge) and (d.2)
  to \(\mathtt{PresheafOfModules}\) is the remaining work. Once (d.2) lands, the
  conclusion is immediate and flatness-free: by
  \(\mathtt{stalkLinearMap\_bijective\_of\_isIso}\) the stalk map \(g_x\) is a
  linear equivalence, and the linear-equivalence \(\mathrm{lTensor}\) by \(F_x\)
  carries it to an isomorphism, finishing the proof.

  This stalkwise fallback is required \emph{only} if the full-generality
  \(J.W.\mathtt{IsMonoidal}\) / \(\mathtt{LocalizedMonoidal}\) route~(e) --- the
  monoidal category on \emph{all} of \(\Scheme.\mathtt{Modules}\,X\), for an
  arbitrary whisker object \(F\) --- is later wanted. For the relative Picard group
  law the whisker objects are locally trivial, and the PRIMARY route above suffices,
  trading the Mathlib-absent (d.2) for the single comparison
  \cref{lem:tensorobj_restrict_iso}.
\end{proof}

\emph{Superseded route --- off path, not to be formalized.} The following lemma
supplies the two whiskering-stability fields of the abandoned
\((J.W).\mathtt{IsMonoidal}\) instance (\cref{lem:jw_ismonoidal}). The group law
needs no such instance: the associator is built directly by gluing canonical local
isomorphisms (\cref{lem:tensorobj_assoc_iso}, via \cref{lem:tensorobj_restrict_iso}).
This block is retained for the historical record only and must \emph{not} be
formalized.

\begin{lemma}

\leanok
% SUPERSEDED route (route-(d)/(e) whiskering/stalk machinery); pending deletion once the assoc re-route (morphism-level descent) lands. Declaration still present and still backing the current `tensorObj_assoc_iso` proof.
[Flatness-free whiskering preserves the sheafification localizer]
  \label{lem:whisker_of_W}
  \lean{PresheafOfModules.W_whiskerLeft_of_W}
  \uses{def:scheme_modules_tensorobj, lem:islocallyinjective_whisker_of_W}
  With notation as in \cref{lem:islocallyinjective_whisker_of_W}: for an
  \emph{arbitrary} presheaf of \(R\)-modules \(F\) and a morphism \(g\) whose
  underlying additive-presheaf map lies in \(J.W\), both whiskerings again lie in
  \(J.W\):
  \[
    g \in J.W \;\Longrightarrow\; F \triangleleft g \in J.W,
    \qquad
    g \in J.W \;\Longrightarrow\; g \triangleright F \in J.W .
  \]
  These two statements are precisely the \(\mathtt{whiskerLeft}\) and
  \(\mathtt{whiskerRight}\) data required by the class
  \(\mathtt{MorphismProperty.IsMonoidal}\) (\cref{sec:tensorobj_api_survey}) for
  the morphism property \(J.W\); together they furnish the instance
  \(J.W.\mathtt{IsMonoidal}\) of \cref{lem:jw_ismonoidal}. No flatness and no
  local triviality is used --- they are the flatness-free replacements for the
  flat variants of \cref{lem:flat_whisker_localizer}.
\end{lemma}

\begin{proof}
  \uses{lem:islocallyinjective_whisker_of_W}
  \leanok
  By the locally-bijective characterisation of \(J.W\), it suffices to show
  \(F \triangleleft g\) is both locally injective and locally surjective.
  \emph{Local surjectivity} is free and needs no hypothesis on \(F\): the presheaf
  tensor is right exact, so \(F \otimes^{p} -\) carries the locally-zero cokernel
  of \(g\) to the cokernel of \(F \triangleleft g\), which is therefore locally
  zero (Mathlib \(\mathtt{isLocallySurjective\_whiskerLeft}\)). \emph{Local
  injectivity} is \cref{lem:islocallyinjective_whisker_of_W}. The two together
  give \(F \triangleleft g \in J.W\). The right-whiskered statement
  \(g \triangleright F\) follows by conjugating \(F \triangleleft g\) with the
  braiding of the symmetric presheaf monoidal structure, which is an isomorphism
  and hence preserves membership in \(J.W\).
\end{proof}

\emph{Superseded route --- off path, not to be formalized.} The following lemma is
the route~(e) target: a full \(\MonoidalCategory\) structure on all of
\(\Scheme.\mathtt{Modules}\,X\) from \(\mathtt{LocalizedMonoidal}\). It is not
needed. The group law consumes only the existence of the associator, unitor and
braiding isomorphisms (\cref{lem:tensorobj_isoclass_commgroup}, assembled by hand
rather than off a coherent monoidal category's \(\mathtt{Skeleton}\)), and the
associator is now built directly by
gluing (\cref{lem:tensorobj_assoc_iso}, via \cref{lem:tensorobj_restrict_iso}). This
block is retained for the historical record only and must \emph{not} be formalized.

\begin{lemma}
% SUPERSEDED route (route-(d)/(e) whiskering/stalk machinery); pending deletion once the assoc re-route (morphism-level descent) lands. Declaration still present and still backing the current `tensorObj_assoc_iso` proof.
[The \(J.W\)-monoidality instance and the localized monoidal substrate (route~(e) target)]
  \label{lem:jw_ismonoidal}
  % \lean: TODO placeholder per DAG integrity rule 1 — the full
  % (J.W).IsMonoidal / LocalizedMonoidal route-(e) packaging is stated but
  % deliberately NOT formalized (the group law needs only existence of the
  % coherence isos), so no real Lean declaration backs this node.
  \lean{AlgebraicGeometry.TODO.jw_ismonoidal}
  \uses{def:scheme_modules_tensorobj, lem:whisker_of_W}
  Let \(X\) be a scheme. The sheafification localizer \(J.W\) on
  \(\mathtt{PresheafOfModules}\,\mathcal{O}_X\) is \emph{monoidal}: it carries an
  instance
  \[
    J.W.\mathtt{IsMonoidal},
  \]
  whose two fields are the whiskering-stability lemmas \cref{lem:whisker_of_W}
  (and which is multiplicative because sheafification inverts identities and
  composites). Consequently, by Mathlib's
  \(\mathtt{Localization.Monoidal.LocalizedMonoidal}\)
  (\cref{sec:tensorobj_api_survey}) applied to the already-monoidal category
  \(\mathtt{PresheafOfModules}\,\mathcal{O}_X\), the localization functor
  \(L = \mathtt{PresheafOfModules.sheafification}\) (which is the localization at
  \(J.W\)), and the unit isomorphism \(L(\mathbf{1}) \cong \mathcal{O}_X\), the
  category
  \[
    \Scheme.\mathtt{Modules}\,X
      \;=\; \mathtt{SheafOfModules}\,X.\mathtt{ringCatSheaf}
  \]
  acquires a full \(\MonoidalCategory\) structure whose tensor is the substrate
  \(\otimes_X\) (\cref{def:scheme_modules_tensorobj}) and whose associator,
  left/right unitors, braiding and \emph{all} coherence (pentagon, triangle,
  naturalities) are derived by the API. This is the route~(e) target: the
  associator \cref{lem:tensorobj_assoc_iso}, the unitors
  \cref{lem:tensorobj_unit_iso} and the braiding \cref{lem:tensorobj_comm_iso} are
  then components of this derived monoidal structure rather than hand-assembled
  composites.

  This block states the route~(e) construction; it depends on the sole open obligation
  \cref{lem:islocallyinjective_whisker_of_W}. It mirrors Mathlib's
  \(\mathtt{CommRing.Pic}\) construction one categorical level up.
\end{lemma}

\begin{proof}
  \uses{lem:whisker_of_W}
  Multiplicativity of \(J.W\) (closure under identities and composition) is part
  of its being a localizer (\(\mathtt{WEqualsLocallyBijective}\)). The two
  whiskering-stability fields of \(\mathtt{MorphismProperty.IsMonoidal}\) are
  exactly the left- and right-whiskering statements of \cref{lem:whisker_of_W}.
  Hence \(J.W.\mathtt{IsMonoidal}\) holds. Sheafification of presheaves of
  modules is a localization at \(J.W\) (the underlying identification of
  sheafification with localization at the locally-bijective class), and the unit
  presheaf sheafifies to \(\mathcal{O}_X\); feeding these to
  \(\mathtt{LocalizedMonoidal}\) produces the \(\MonoidalCategory\) structure on
  \(\Scheme.\mathtt{Modules}\,X\) with all coherence derived, as in
  \cref{sec:tensorobj_api_survey}. No \(\mathtt{MonoidalClosed}\) structure enters
  at any point.
\end{proof}

\begin{lemma}


\leanok
[Associator for $\otimes_X$ (unconditional)]
  \label{lem:tensorobj_assoc_iso}
  \lean{AlgebraicGeometry.Scheme.Modules.tensorObj_assoc_iso}
  \uses{def:scheme_modules_tensorobj,
        lem:islocallyinjective_whiskerleft_via_stalk,
        lem:isiso_sheafification_map_of_W}
  Let \(X\) be a scheme and let
  \(M, N, P \in \Scheme.\mathtt{Modules}\,X\). There is an isomorphism
  \[
    (M \otimes_X N) \otimes_X P
      \;\xrightarrow{\ \sim\ }\;
    M \otimes_X (N \otimes_X P)
  \]
  in \(\Scheme.\mathtt{Modules}\,X\), natural in \(M, N, P\). The group law on
  \(\otimes\)-iso-classes (\cref{lem:tensorobj_isoclass_commgroup}) consumes only
  the \emph{existence} of this isomorphism, the proposition
  \(\mathrm{Nonempty}\bigl((M \otimes_X N) \otimes_X P \cong M \otimes_X (N
  \otimes_X P)\bigr)\); no pentagon, triangle, or naturality coherence of it is
  ever required.

  The Lean pin is unconditional: it holds for arbitrary
  \(\mathcal{O}_X\)-modules \(M, N, P\), with no local-triviality or
  invertibility hypothesis. (Earlier drafts carried vestigial
  \(\mathtt{IsLocallyTrivial}\) hypotheses to match the carrier; these were
  dropped once the construction below was seen to use none of them.)
\end{lemma}

\begin{proof}
  \uses{def:scheme_modules_tensorobj,
        lem:islocallyinjective_whiskerleft_via_stalk}
  \textbf{Realization (sheafification transport, unconditional via d.2).} The
  associator is the sheafification transport of the presheaf-of-modules associator
  \(\mathtt{PresheafOfModules.monoidalCategoryStruct.associator}\), which exists for
  \emph{all} presheaves of modules over the varying ring. The transport requires the
  single left-whisker \(F \triangleleft \mathtt{toSheafify}\) to lie in the
  sheafification localizer \(J.W\) for the relevant \(F\); that single open
  left-whiskering obligation \(\mathtt{isLocallyInjective\_whiskerLeft\_of\_W}\) is
  closed \emph{unconditionally} --- for an arbitrary whiskering object, with no
  flatness and no local triviality --- by the d.2 stalk--tensor commutation
  \cref{lem:islocallyinjective_whiskerleft_via_stalk} (\cref{sec:tensorobj_stalk_tensor}):
  a \(J.W\)-morphism \(g\) is a stalkwise isomorphism, and by the d.2 commutation
  \(\bigl(F \triangleleft g\bigr)_x \cong \mathrm{id}_{F_x} \otimes g_x\) is again an
  isomorphism, so \(F \triangleleft g \in J.W\). Hence the associator is itself
  unconditional and axiom-clean once d.2 (\cref{lem:stalk_tensor_commutation}) is
  formalized; it does \emph{not} require the full \(J.W.\mathtt{IsMonoidal}\) instance
  or the \(\mathtt{LocalizedMonoidal}\) monoidal category --- only the one left-whisker
  lemma --- and no sectionwise flatness enters. The local-gluing description below
  records the equivalent concrete local realization through the
  restriction-compatibility comparison \cref{lem:tensorobj_restrict_iso}.

  Concretely, the same isomorphism is described by gluing canonical local isomorphisms
  through the restriction-compatibility comparison \cref{lem:tensorobj_restrict_iso}.
  Fix a
  common open cover \(\{U\}\) of \(X\). On each member \(U\) form the canonical
  composite
  \[
    \bigl((M \otimes_X N) \otimes_X P\bigr)\big|_U
      \;\xrightarrow{\ \sim\ }\;
    \bigl(M|_U \otimes_U N|_U\bigr) \otimes_U P|_U
      \;\xrightarrow{\ \sim\ }\;
    M|_U \otimes_U \bigl(N|_U \otimes_U P|_U\bigr)
      \;\xrightarrow{\ \sim\ }\;
    \bigl(M \otimes_X (N \otimes_X P)\bigr)\big|_U ,
  \]
  in which:
  \begin{itemize}
    \item the first arrow is two applications of the restriction-compatibility
      isomorphism \cref{lem:tensorobj_restrict_iso} along the open immersion
      \(U \hookrightarrow X\): first \((M \otimes_X N)|_U \cong M|_U \otimes_U
      N|_U\) tensored with \(P|_U\), then \(\bigl((M \otimes_X N) \otimes_X
      P\bigr)|_U \cong (M \otimes_X N)|_U \otimes_U P|_U\);
    \item the middle arrow is the \emph{canonical presheaf-level associator}
      \(\alpha\) of the (symmetric) monoidal category
      \(\mathtt{PresheafOfModules}\,\mathcal{O}_U\)
      (\(\mathtt{PresheafOfModules.monoidalCategoryStruct.associator}\),
      \cref{sec:tensorobj_api_survey}(i)), read on the restricted modules;
    \item the third arrow is again two applications of
      \cref{lem:tensorobj_restrict_iso}, in the opposite direction.
  \end{itemize}
  Each of the three arrows is a \emph{canonical} isomorphism: the
  restriction-compatibility comparison \cref{lem:tensorobj_restrict_iso} is induced
  by the open immersion with no choices, hence natural in its module arguments, and
  the presheaf associator \(\alpha\) is a fixed natural transformation. Therefore,
  on a double overlap \(U \cap U'\), restricting the \(U\)-composite and the
  \(U'\)-composite to \(U \cap U'\) yields the same morphism: each of the three
  arrows commutes with the further restriction to \(U \cap U'\) by the naturality of
  \cref{lem:tensorobj_restrict_iso} and of the presheaf associator, so the two local
  composites agree on the overlap.

  Because the internal hom \(\mathcal{H}om(-,-)\) of sheaves of
  \(\mathcal{O}_X\)-modules is itself a sheaf on \(X\), a family of morphisms defined
  on the members of an open cover and agreeing on the overlaps glues to a single
  global morphism. Applying this both to the local composites above and to their
  local inverses produces a global morphism
  \((M \otimes_X N) \otimes_X P \to M \otimes_X (N \otimes_X P)\) whose restriction
  to each \(U\) is the local composite, together with a two-sided inverse glued from
  the local inverses; hence the glued morphism is a global isomorphism --- the
  associator. The single non-formal input of this re-route is the
  restriction-compatibility comparison \cref{lem:tensorobj_restrict_iso}; the presheaf
  associator (\cref{sec:tensorobj_api_survey}(i)) and the sheaf-of-modules
  internal-hom gluing are already available.

  Note that, unlike the pointwise-existence statement
  \cref{lem:tensorobj_preserves_locally_trivial} --- a \emph{proposition} proved
  cover-pointwise with \emph{no} overlap obligation --- the associator is
  \emph{global data}, a single isomorphism of sheaves on all of \(X\). Assembling it
  from local pieces is therefore genuinely stronger than that pattern: the
  overlap-naturality compatibility of the local composites must be discharged before
  gluing, and is discharged above. The locally-trivial hypotheses of the statement
  are not consumed; the comparison \cref{lem:tensorobj_restrict_iso} is applied to the
  global, untrivialised restrictions \((M \otimes_X N)|_U\), so no free-cover
  specialisation can replace it. What this proof does \emph{not} require is the
  \emph{full} \((J.W).\mathtt{IsMonoidal}\) instance or the
  \(\mathtt{LocalizedMonoidal}\) monoidal category
  (\cref{lem:whisker_of_W,lem:jw_ismonoidal}): only the single left-whisker obligation
  \cref{lem:islocallyinjective_whiskerleft_via_stalk}, closed via the d.2 stalk--tensor
  commutation \cref{lem:stalk_tensor_commutation}, is consumed by the sheafification
  transport. The group law on \(\otimes\)-iso-classes
  consumes only the \emph{existence} of the resulting isomorphism, never its pentagon,
  triangle, or naturality coherence.
\end{proof}

\begin{lemma}
[Left and right unitors for $\otimes_X$]
  \label{lem:tensorobj_unit_iso}
  \lean{AlgebraicGeometry.Scheme.Modules.tensorObj_left_unitor, AlgebraicGeometry.Scheme.Modules.tensorObj_right_unitor}
  \uses{def:scheme_modules_tensorobj}
  Let \(X\) be a scheme and \(M \in \Scheme.\mathtt{Modules}\,X\) an
  \emph{arbitrary} \(\mathcal{O}_X\)-module. With
  \(\mathcal{O}_X = \mathtt{SheafOfModules.unit}\,X.\mathtt{ringCatSheaf}\) the
  structure-sheaf unit object, there are natural isomorphisms
  \[
    \mathcal{O}_X \otimes_X M \;\xrightarrow{\ \sim\ }\; M,
    \qquad
    M \otimes_X \mathcal{O}_X \;\xrightarrow{\ \sim\ }\; M .
  \]
\end{lemma}

\begin{proof}
  \uses{def:scheme_modules_tensorobj}
  The presheaf monoidal structure on \(\mathtt{PresheafOfModules}\,R\) supplies
  the left and right unitors \(\lambda, \rho\) at the presheaf level. Their
  sheafifications \(\mathtt{sheafification.mapIso}(\lambda)\),
  \(\mathtt{sheafification.mapIso}(\rho)\), composed with the sheafification
  counit identifying the sheafified structure-sheaf presheaf with
  \(\mathcal{O}_X\) (the step already used in the structure-sheaf-unit
  comparison \(\mathtt{tensorObj\_unit\_iso}\) in the substrate file), give the
  two isomorphisms. This route is the cheap \(\mathtt{mapIso}\) pattern and uses no
  abstract pullback.
\end{proof}

\begin{lemma}


\leanok
[Braiding for $\otimes_X$]
  \label{lem:tensorobj_comm_iso}
  \lean{AlgebraicGeometry.Scheme.Modules.tensorObj_braiding}
  \uses{def:scheme_modules_tensorobj}
  Let \(X\) be a scheme and \(M, N \in \Scheme.\mathtt{Modules}\,X\)
  \emph{arbitrary} \(\mathcal{O}_X\)-modules. There is a natural isomorphism
  \[
    M \otimes_X N \;\xrightarrow{\ \sim\ }\; N \otimes_X M .
  \]
\end{lemma}

\begin{proof}
  \uses{def:scheme_modules_tensorobj}
  The presheaf-of-modules monoidal category is symmetric, so it carries a
  braiding \(\beta\) at the presheaf level
  (\texttt{Mathlib.CategoryTheory.Monoidal.PresheafOfModules}). Its
  sheafification \(\mathtt{sheafification.mapIso}(\beta)\) is the asserted
  isomorphism, again by the cheap \(\mathtt{mapIso}\) pattern.
\end{proof}

\begin{lemma}


\leanok
[Inverse of an invertible module]
  \label{lem:tensorobj_inverse_invertible}
  \lean{AlgebraicGeometry.Scheme.Modules.exists_tensorObj_inverse}
  \uses{def:scheme_modules_tensorobj, lem:tensorobj_preserves_locally_trivial,
        lem:tensorobj_restrict_iso}
  % NOTE: ON the critical path under the PRIMARY locally-trivial group law (it
  % supplies inverses to lem:tensorobj_isoclass_commgroup, in whose uses it sits).
  % Only under the route-(e) fallback is it dispensable: there the group law uses
  % the tensor-invertibility predicate IsInvertible (def:scheme_modules_isinvertible),
  % which carries the inverse existentially, so no separate ``a tensor inverse
  % exists'' lemma is required.
  % SOURCE: [Stacks Project], Tag 01CR (Modules on Ringed Spaces, \S Invertible
  % modules), Lemma~lemma-constructions-invertible (item~2: the dual of an
  % invertible is invertible and the contraction map is an isomorphism) +
  % Definition~definition-pic (the Picard group as the abelian group of
  % iso-classes of invertibles)
  % (read from references/stacks-modules.tex, L4200--L4213 and L4350--L4357).
  % SOURCE QUOTE:
  % "\begin{lemma}
  % \label{lemma-constructions-invertible}
  % Let $(X, \mathcal{O}_X)$ be a ringed space.
  % ...
  % \item If $\mathcal{L}$ is an invertible $\mathcal{O}_X$-module, then so is
  % $\SheafHom_{\mathcal{O}_X}(\mathcal{L}, \mathcal{O}_X)$ and the evaluation map
  % $\mathcal{L} \otimes_{\mathcal{O}_X}
  % \SheafHom_{\mathcal{O}_X}(\mathcal{L}, \mathcal{O}_X) \to \mathcal{O}_X$
  % is an isomorphism.
  % ...
  % \begin{definition}
  % \label{definition-pic}
  % Let $(X, \mathcal{O}_X)$ be a ringed space.
  % The {\it Picard group} $\Pic(X)$ of $X$ is the
  % abelian group whose elements are isomorphism classes of
  % invertible $\mathcal{O}_X$-modules, with addition
  % corresponding to tensor product.
  % \end{definition}"
  \textit{Source: [Stacks Project], Tag~01CR, Lemma~\texttt{lemma-constructions-%
  invertible} (item~2, the dual of an invertible is invertible and the contraction
  is an isomorphism) and Definition~\texttt{definition-pic} (the Picard group is
  the abelian group of iso-classes of invertible \(\mathcal{O}_X\)-modules under
  tensor product).}
  Let \(X\) be a scheme and let \(L \in \texttt{LineBundle}\,X\) be a line
  bundle. The dual
  \[
    L^{-1} \;\;:=\;\; \mathcal{H}om_{\mathcal{O}_X}(L, \mathcal{O}_X)
  \]
  is again a line bundle on \(X\), and there is a canonical natural
  isomorphism
  \(\varepsilon_L : L \otimes_X L^{-1} \xrightarrow{\sim} \mathcal{O}_X\).
  Consequently every line bundle has a two-sided tensor inverse under
  \(\otimes_X\), and the monoid \(\bigl(\texttt{LineBundle}\,X / \sim, \otimes_X,
  \mathcal{O}_X\bigr)\) of isomorphism classes is an abelian group ---
  the Picard group \(\Pic(X)\) (Stacks tag 01CR).
\end{lemma}

\begin{proof}
  \uses{lem:tensorobj_preserves_locally_trivial, lem:tensorobj_restrict_iso,
        lem:tensorobj_unit_iso, lem:internal_hom_eval, lem:dual_isLocallyTrivial,
        lem:topresheaf_map_hommk}
  \textbf{Infrastructure-blocked.} This statement is \emph{not} realizable with the
  \(\mathcal{O}_X\)-module infrastructure presently available: its construction
  depends on a \emph{sheaf internal hom of \(\mathcal{O}_X\)-modules}
  \(\mathcal{H}om_{\mathcal{O}_X}(L, \mathcal{O}_X)\), an object that does not yet
  exist at the presheaf, sheaf, or general-categorical level (there is no
  monoidal-closed / internal-hom / dual structure on
  \(\mathtt{PresheafOfModules}\) or \(\mathtt{SheafOfModules}\), and no object-level
  descent assembling a global \(\mathcal{O}_X\)-module from local pieces). The dual
  object \(L^{-1}\) therefore cannot even be \emph{named} until that infrastructure
  is built. The construction of the required internal hom, its dual, and its
  evaluation is decomposed into individually-formalizable sub-steps in
  \cref{sec:tensorobj_dual_infra} (\cref{def:presheaf_internal_hom} through
  \cref{lem:dual_isLocallyTrivial}); the present lemma is realized once those land,
  and its Lean body is a placeholder until then.

  We record the intended mathematical route, all three steps of which are correct
  modulo the missing object. \emph{The route is as follows.} Take
  \(L^{-1} := \mathcal{H}om_{\mathcal{O}_X}(L, \mathcal{O}_X)\) to be the
  sheaf-level dual of \cref{lem:internal_hom_isSheaf}, and let
  \(\varepsilon_L : L \otimes_X L^{-1} \to \mathcal{O}_X\) be the evaluation of
  \cref{lem:internal_hom_eval}.

  \emph{Step 1 (\(L^{-1}\) is a line bundle).} Local triviality of \(L\)
  (\(\mathtt{LineBundle.IsLocallyTrivial}\,L\)) furnishes an open cover \(\{U\}\) of
  \(X\) with trivialisations \(L|_U \cong \mathcal{O}_U\). The internal hom commutes
  with restriction along the open immersion \(U \hookrightarrow X\), so on each \(U\)
  \[
    L^{-1}\big|_U \;\cong\; \mathcal{H}om_{\mathcal{O}_U}(L|_U, \mathcal{O}_U)
                 \;\cong\; \mathcal{H}om_{\mathcal{O}_U}(\mathcal{O}_U, \mathcal{O}_U)
                 \;\cong\; \mathcal{O}_U ,
  \]
  the last identification being the evaluation-at-\(1\) isomorphism for the dual of a
  free rank-one module. Hence \(L^{-1}\) is locally trivial of rank one, i.e.\ an
  object of \(\texttt{LineBundle}\,X\).

  \emph{Step 2 (the contraction is a local isomorphism).} The evaluation/contraction
  morphism
  \[
    \varepsilon_L \;\colon\; L \otimes_X L^{-1}
      \;=\; L \otimes_X \mathcal{H}om_{\mathcal{O}_X}(L, \mathcal{O}_X)
      \;\longrightarrow\; \mathcal{O}_X,
    \qquad s \otimes \phi \longmapsto \phi(s),
  \]
  is a morphism in \(\Scheme.\mathtt{Modules}\,X\). Restricting to a trivialising
  open \(U\), the restriction-compatibility isomorphism \cref{lem:tensorobj_restrict_iso}
  commutes \(\otimes_X\) past \((-)|_U\):
  \[
    (L \otimes_X L^{-1})\big|_U \;\cong\; L|_U \otimes_U L^{-1}|_U
      \;\cong\; \mathcal{O}_U \otimes_U \mathcal{O}_U ,
  \]
  using the trivialisations of Step~1. Under these identifications \(\varepsilon_L|_U\)
  becomes the contraction \(\mathcal{O}_U \otimes_U \mathcal{O}_U \to \mathcal{O}_U\) of
  the free rank-one module with its dual, which is precisely the left unitor
  \cref{lem:tensorobj_unit_iso} (\(\mathcal{O}_U \otimes_U (-) \cong (-)\)) and hence an
  isomorphism. Thus \(\varepsilon_L|_U\) is an isomorphism on each member of the cover.

  \emph{Step 3 (glue to a global isomorphism).} The local trivialisations are induced
  by the open immersions with no choices, so the local contraction isomorphisms
  \(\varepsilon_L|_U\) agree on the overlaps \(U \cap U'\) by the naturality of
  \cref{lem:tensorobj_restrict_iso} and of the unitor under further restriction. As
  ``being an isomorphism'' is local on \(X\) --- a morphism of sheaves of
  \(\mathcal{O}_X\)-modules whose restriction to every member of an open cover is an
  isomorphism is itself an isomorphism --- the globally-defined \(\varepsilon_L\) is an
  isomorphism \(L \otimes_X L^{-1} \xrightarrow{\sim} \mathcal{O}_X\). Equivalently, the
  local inverses glue (they agree on overlaps for the same reason) to a two-sided
  inverse of \(\varepsilon_L\).

  Consequently, \emph{once \(L^{-1}\) and \(\varepsilon_L\) are available from
  \cref{sec:tensorobj_dual_infra}}, every line bundle \(L\) has the two-sided tensor
  inverse \(L^{-1}\) under \(\otimes_X\), and the monoid of \(\otimes\)-iso-classes of
  line bundles is an abelian group --- the Picard group \(\Pic(X)\) (Stacks tag 01CR).
  The local-triviality of \(L^{-1}\) (Step~1) is \cref{lem:dual_isLocallyTrivial} and
  the local-iso\(\Rightarrow\)global-iso passage (Step~3) is the already-closed
  \cref{lem:tensorobj_restrict_iso}; the single non-available input is the dual object
  and its evaluation of \cref{sec:tensorobj_dual_infra}.
\end{proof}

\begin{lemma}




\leanok
[Closure of \(\otimes\) on the locally-trivial carrier
  \(\texttt{LineBundle.OnProduct}\)]
  \label{lem:tensorobj_lift_onproduct}
  \lean{AlgebraicGeometry.Scheme.Modules.tensorObjOnProduct}
  \uses{def:scheme_modules_tensorobj,
        lem:tensorobj_preserves_locally_trivial}
  Let \(C \to S\) and \(T \to S\) be \(S\)-schemes. The carrier is the
  \emph{locally-trivial} subtype
  \[
    \texttt{LineBundle.OnProduct}\,\pi_C\,\pi_T
    \;\;=\;\;
    \bigl\{\,M \in \Scheme.\mathtt{Modules}(C \times_S T)
            \,\bigm|\, \mathtt{IsLocallyTrivial}\,M \,\bigr\}
  \]
  introduced in \cref{chap:Picard_LineBundlePullback}. The bifunctor
  \(\otimes_{C \times_S T}\) of
  \cref{lem:scheme_modules_tensorobj_functoriality} restricts to this subtype:
  if \(\mathtt{IsLocallyTrivial}\,M\) and \(\mathtt{IsLocallyTrivial}\,M'\) then
  \(\mathtt{IsLocallyTrivial}(M \otimes_{C \times_S T} M')\). Hence \(\otimes\)
  closes on \(\texttt{LineBundle.OnProduct}\,\pi_C\,\pi_T\), supplying the binary
  operation the relative group law multiplies with.
\end{lemma}

\begin{proof}
  \uses{lem:tensorobj_preserves_locally_trivial}
  \leanok
  Closure is exactly
  \cref{lem:tensorobj_preserves_locally_trivial}: the tensor product of two
  locally-trivial rank-one modules is again locally trivial of rank one
  (\(\mathtt{tensorObj\_isLocallyTrivial}\)). The Lean construction
  \(\mathtt{tensorObjOnProduct}\) applies this directly to a pair of objects of
  the \(\mathtt{IsLocallyTrivial}\) subtype \(\texttt{LineBundle.OnProduct}\),
  returning the tensor with its locally-trivial witness; no
  \(\otimes\)-invertibility predicate and no coherence isomorphism enters this
  closure statement. (The \(\otimes\)-invertibility carrier
  \(\mathtt{IsInvertible}\) of \cref{def:scheme_modules_isinvertible} is a
  separate predicate, connected to \(\mathtt{IsLocallyTrivial}\) only by an
  off-critical-path equivalence; the units-group packaging is
  \cref{lem:tensorobj_isoclass_commgroup}, not this lemma.)
\end{proof}

\begin{lemma}
[\(\pi_T^*\) is a tensor-functor]
  \label{lem:pullback_compatible_with_tensorobj}
  \lean{AlgebraicGeometry.LineBundle.OnProduct.pullback_tensorObj_iso}
  % NOTE: \lean{} pin target absent from Lean source as of iter-271 (verified by grep).
  % Forward reference to a not-yet-formalized declaration — NOT a stale rename; there is
  % no existing Lean decl under any name to repin to. The `unmatched_lean` status is the
  % normal "unformalized target" state, not a pin defect.
  \uses{lem:tensorobj_lift_onproduct}
  % NOTE: pullback preserves the structure sheaf and tensor products (and hence
  % tensor-invertibility) by the standard pullback-tensor compatibility of
  % \(\mathcal{O}\)-modules (Stacks \texttt{lemma-tensor-product-pullback} /
  % \texttt{lemma-pullback-invertible}), independently of the open-immersion
  % restriction-compatibility iso.
  With notation as in \cref{lem:tensorobj_lift_onproduct}, the pullback
  functor
  \[
    \pi_T^{*} \;\;\colon\;\; \texttt{LineBundle}\,T \;\longrightarrow\;
                          \texttt{LineBundle.OnProduct}\,\pi_C\,\pi_T
  \]
  of \cref{def:pullback_along_projection} is a (strong) monoidal
  functor: there is a canonical natural isomorphism
  \(\pi_T^*(\mathcal{N} \otimes_T \mathcal{N'})
     \xrightarrow{\sim}
   \pi_T^* \mathcal{N} \otimes_{C \times_S T} \pi_T^* \mathcal{N'}\)
  for \(\mathcal{N}, \mathcal{N'} \in \texttt{LineBundle}\,T\), and a
  canonical isomorphism
  \(\pi_T^* \mathcal{O}_T \xrightarrow{\sim} \mathcal{O}_{C \times_S T}\),
  both inherited from the analogous data on the underlying
  \(\Scheme.\mathtt{Modules}\) functors via the
  pullback-tensor-compatibility of \(\mathcal{O}\)-modules.
\end{lemma}

\begin{proof}
  \uses{lem:tensorobj_lift_onproduct}
  On any affine open of \(T\) over which \(\mathcal{N}\) and
  \(\mathcal{N'}\) trivialise simultaneously, the assertion reduces to
  the standard fact that the extension of scalars
  \(B \otimes_A (M \otimes_A M')
   \;\cong\; (B \otimes_A M) \otimes_B (B \otimes_A M')\)
  for a ring map \(A \to B\) and \(A\)-modules \(M, M'\). This affine
  identity glues across an affine trivialising cover of \(T\): on each
  affine open where \(\mathcal{N}\) and \(\mathcal{N}'\) are free of rank one
  over \(A\), both sides evaluate to the free rank-one \(B\)-module, so
  the comparison is an isomorphism. No flatness of \(A \to B\) is required
  at this level, as the modules involved are locally free. The compatibility
  with the unit, the associator, and the unitors is verified pointwise on the
  same cover.
\end{proof}

\section{The varying-ring stalk--tensor commutation (d.2) and the unconditional
left-whiskering}
\label{sec:tensorobj_stalk_tensor}

This section supplies the one genuinely Mathlib-absent ingredient on which the
\emph{unconditional} associator \cref{lem:tensorobj_assoc_iso} rests: the commutation
of the stalk functor with the relative-ring tensor product of presheaves of modules
(``d.2''). Together with the already-built \(R_x\)-linear stalk comparison
(\cref{lem:stalk_linear_map}, ``d.1'') it closes the single open left-whiskering
obligation for an \emph{arbitrary} whiskering object --- with no flatness and no
local-triviality hypothesis --- which is exactly what makes the sheafification
transport of the presheaf associator work for all modules. The construction is to be
formalized in the dedicated Lean file
\(\mathtt{AlgebraicJacobian/Picard/TensorObjSubstrate/StalkTensor.lean}\).

\begin{lemma}

\leanok
[Stalk of a tensor product (d.2): varying-ring stalk--tensor commutation]
  \label{lem:stalk_tensor_commutation}
  \lean{PresheafOfModules.stalkTensorIso}
  \uses{def:scheme_modules_tensorobj, lem:stalk_linear_map}
  % SOURCE: [Stacks Project], Tag (Modules on Ringed Spaces, \S Tensor product),
  % Lemma lemma-stalk-tensor-product (stalk of a tensor product of modules is the
  % tensor product of the stalks over the stalk ring)
  % (read from references/stacks-modules.tex, L2332-L2344).
  % SOURCE QUOTE:
  % "\begin{lemma}
  % \label{lemma-stalk-tensor-product}
  % Let $(X, \mathcal{O}_X)$ be a ringed space.
  % Let $\mathcal{F}$, $\mathcal{G}$ be $\mathcal{O}_X$-modules.
  % Let $x \in X$. There is a canonical isomorphism
  % of $\mathcal{O}_{X, x}$-modules
  % $$
  % (\mathcal{F} \otimes_{\mathcal{O}_X} \mathcal{G})_x
  % =
  % \mathcal{F}_x \otimes_{\mathcal{O}_{X, x}} \mathcal{G}_x
  % $$
  % functorial in $\mathcal{F}$ and $\mathcal{G}$.
  % \end{lemma}"
  \textit{Source: [Stacks Project], Modules on Ringed Spaces,
  Lemma~\texttt{lemma-stalk-tensor-product}.}
  Let \(R : (\mathtt{Opens}\,X)^{op} \to \mathtt{CommRingCat}\) be the structure
  presheaf of a scheme \(X\), and let \(A, B\) be presheaves of
  \(R\)-modules (objects of \(\mathtt{PresheafOfModules}\,(R \circ
  \mathtt{forget_2}\,\mathtt{CommRingCat}\,\mathtt{RingCat})\)). For each point
  \(x \in X\) there is a canonical isomorphism of \((R \circ
  \mathtt{forget_2}).\mathtt{stalk}\,x\)-modules
  \[
    (A \otimes^{p} B).\mathtt{stalk}\,x
      \;\xrightarrow{\ \sim\ }\;
    A.\mathtt{stalk}\,x \;\otimes_{R.\mathtt{stalk}\,x}\; B.\mathtt{stalk}\,x,
  \]
  natural in \(A\) and \(B\), where \(\otimes^{p}\) is the presheaf-of-modules tensor
  product \(\mathtt{PresheafOfModules.Monoidal.tensorObj}\).
\end{lemma}

\begin{proof}
  \uses{def:scheme_modules_tensorobj, lem:stalk_linear_map}
  \leanok
  The stalk at \(x\) is the filtered colimit over the neighbourhoods \(U \ni x\) of
  the sections functor, and sectionwise the tensor product is computed pointwise:
  \((A \otimes^{p} B)(U) = A(U) \otimes_{R(U)} B(U)\)
  (\(\mathtt{PresheafOfModules.Monoidal.tensorObj\_obj}\)). The neighbourhood system
  \(\{U \ni x\}\) is cofiltered under inclusion, so the colimit is filtered. The
  isomorphism is assembled in five named stages. Stages (i)--(ii) construct the
  forward additive comparison map and its germ characterisation; stages (iii)--(v)
  upgrade it to an \(R_x\)-linear map, construct the reverse map, and check mutual
  inversion.

  \emph{(i) Per-neighbourhood balanced bilinear map and its section-level descent}
  (\(\mathtt{stalkTensorBilin}\), \(\mathtt{stalkTensorDescU}\)). For a fixed
  neighbourhood \(U \ni x\) the assignment
  \[
    (a, b) \;\longmapsto\; \mathrm{germ}_x\,a \,\otimes\, \mathrm{germ}_x\,b,
    \qquad a \in A(U),\ b \in B(U),
  \]
  is \(R(U)\)-balanced bilinear into \(A_x \otimes_{R_x} B_x\): it is additive in each
  variable, and for \(r \in R(U)\) the two ways of moving \(r\) across the tensor agree
  after passing to germs, since \(\mathrm{germ}_x(r \cdot a) = (\mathrm{germ}_x r)
  \cdot \mathrm{germ}_x a\) and the stalk tensor is balanced over
  \(R_x = R.\mathtt{stalk}\,x\) (the germ--scalar compatibility of
  \cref{lem:stalk_linear_map}). By the universal property of the relative tensor
  product it descends to an additive map on the section module
  \[
    \mathtt{stalkTensorDescU}_U \;\colon\;
      (A \otimes^{p} B)(U) = A(U) \otimes_{R(U)} B(U)
        \;\longrightarrow\; A_x \otimes_{R_x} B_x,
    \qquad a \otimes b \longmapsto \mathrm{germ}_x a \otimes \mathrm{germ}_x b.
  \]

  \emph{(ii) Colimit descent to the forward comparison map}
  (\(\mathtt{stalkTensorDesc}\), \(\mathtt{stalkTensorDesc\_germ\_tmul}\)). The
  maps \(\mathtt{stalkTensorDescU}_U\) are compatible with the restriction maps of the
  presheaf \(A \otimes^{p} B\) as \(U\) shrinks toward \(x\) (a smaller neighbourhood
  computes the same germs), so they form a cocone over the filtered diagram
  \(\{(A \otimes^{p} B)(U)\}_{U \ni x}\). The induced map out of the colimit is the
  forward additive comparison map
  \[
    \mathtt{stalkTensorDesc} \;\colon\;
      (A \otimes^{p} B).\mathtt{stalk}\,x \;\longrightarrow\; A_x \otimes_{R_x} B_x,
  \]
  characterised on germs of simple tensors by
  \(\mathtt{stalkTensorDesc}\bigl(\mathrm{germ}_x(a \otimes b)\bigr) =
  \mathrm{germ}_x a \otimes \mathrm{germ}_x b\)
  (\(\mathtt{stalkTensorDesc\_germ\_tmul}\)). This is the forward half of the asserted
  isomorphism (\cref{lem:stalk_tensor_desc_forward}).

  \emph{(iii) \(R_x\)-linearity packaging} (\(\mathtt{stalkTensorLinearMap}\)). The
  additive map of (ii) is upgraded to an \(R.\mathtt{stalk}\,x\)-linear map: on a germ
  \(\mathrm{germ}_x r\) of a section \(r \in R(U)\) one checks
  \(\mathtt{stalkTensorDesc}\bigl((\mathrm{germ}_x r)\cdot \xi\bigr) =
  (\mathrm{germ}_x r)\cdot \mathtt{stalkTensorDesc}(\xi)\) on germs of simple tensors
  and extends by additivity. The one subtlety is a \emph{carrier-duality obstacle}:
  the section-level tensor \(A(U) \otimes_{R(U)} B(U)\) is a module over the
  \(\mathtt{RingCat}\) carrier \((R \circ \mathtt{forget_2})(U)\), whereas the natural
  scalar \(r : R(U)\) lives in the \(\mathtt{CommRingCat}\) carrier \(R(U)\); the two
  carriers are canonically identified but not definitionally equal, so the
  scalar-pulling lemma \(\mathtt{TensorProduct.smul\_tmul'}\) applies only after
  inserting a canonical \(\mathtt{RingEquiv}\) bridge between the
  \(\mathtt{CommRingCat}\) and \(\mathtt{RingCat}\) carriers. Once that bridge is
  in place the linearity computation mirrors the germ-representative pattern used
  for the d.1 stalk map \(\mathtt{stalkLinearMap}\) (\cref{lem:stalk_linear_map}).

  \emph{(iv) The reverse map.} The map
  \[
    A_x \otimes_{R_x} B_x \;\longrightarrow\; (A \otimes^{p} B).\mathtt{stalk}\,x
  \]
  is built by the universal property of the tensor product over \(R_x\) from the
  \(R_x\)-bilinear map of the two stalks
  \[
    (\mathrm{germ}_x a,\ \mathrm{germ}_x b) \;\longmapsto\;
      \mathrm{germ}_x\bigl(a|_{U \cap V} \otimes b|_{U \cap V}\bigr),
  \]
  where \(a \in A(U)\) and \(b \in B(V)\) are chosen representatives; this is a
  \emph{nested} colimit descent, descending first the \(A\)-variable out of the
  filtered colimit \(\{A(U)\}_{U \ni x}\) and then the \(B\)-variable out of
  \(\{B(V)\}_{V \ni x}\), passing to the common refinement \(U \cap V\) to form the
  simple tensor and its germ. Well-definedness on germs (independence of the chosen
  representatives and of the neighbourhoods \(U, V\)) is exactly filteredness of the
  neighbourhood system.

  It remains to check that the reverse map is \(R_x\)-bilinear, and in particular that
  its underlying bilinear map is \emph{balanced} over \(R_x\), i.e. that pulling a scalar
  out of the left variable agrees with pulling it out of the right. The key point is to
  establish this balancing \emph{at the stalk level}, with every scalar living in the
  stalk ring \(R_x = R.\mathtt{stalk}\,x\) throughout. Concretely, for a germ
  \(\mathrm{germ}_x r\) of a section \(r \in R(W)\) one verifies the identity
  \[
    \mathtt{revBihom}\bigl((\mathrm{germ}_x r)\cdot \mathrm{germ}_x a,\ \mathrm{germ}_x b\bigr)
      \;=\; (\mathrm{germ}_x r)\cdot \mathtt{revBihom}\bigl(\mathrm{germ}_x a,\ \mathrm{germ}_x b\bigr),
  \]
  and symmetrically on the right variable, where the scalar \(\mathrm{germ}_x r \in R_x\)
  never leaves the stalk ring. This holds because of the germ--scalar compatibility
  \(\mathrm{germ}_x(r \cdot s) = (\mathrm{germ}_x r)\cdot(\mathrm{germ}_x s)\) of the germ
  map \(R(W) \to R_x\) (the Lean lemma \(\mathtt{germ\_smul}\)) — the same germ--scalar
  compatibility already used in stages (i) and (iii). Multiplying a representative
  \(r \cdot a\) and then taking germs produces \((\mathrm{germ}_x r)\cdot \mathrm{germ}_x a\),
  so the action of \(r\) is transparently the \(R_x\)-action on the stalk and may be
  moved freely across the tensor of stalks, where balancedness over \(R_x\) is built into
  \(A_x \otimes_{R_x} B_x\). Because the whole computation takes place over the single
  ring \(R_x\), no carrier identification intervenes and the balancing follows directly.

  One must \emph{not} instead reduce this balancing to a section-level identity. If the
  scalar were pushed through the presheaf restriction to a common neighbourhood \(W\) and
  the balancing read off inside the section module \(A(W) \otimes_{R(W)} B(W)\), then the
  scalar action appearing there is the one produced by the presheaf-of-modules structure
  of \(A\) via restriction of scalars: it lives over the \(\mathtt{RingCat}\) carrier
  \((R \circ \mathtt{forget_2})(W)\), \emph{not} over the \(\mathtt{CommRingCat}\) carrier
  \(R(W)\) that annotates the section tensor. This is precisely the
  \(\mathtt{CommRingCat}\)-vs-\(\mathtt{RingCat}\) carrier-duality obstacle already met in
  stage (iii), here compounded by the extra restriction-of-scalars wrapper, so the naive
  ``move the scalar across the tensor'' step (\(\mathtt{TensorProduct.smul\_tmul'}\)) does
  not apply at the section level without first inserting the carrier bridge. Working at the
  stalk level as above sidesteps this wrapper entirely.

  Where a \(\mathtt{CommRingCat}\)-vs-\(\mathtt{RingCat}\) carrier identification is
  genuinely unavoidable, it is the \emph{same} canonical bridge already invoked for the
  stage-(iii) \(R_x\)-linearity packaging (the \(\mathtt{stalkTensorDescU\_smul}\)/%
  \(\mathtt{stalkTensorLinearMap}\) step, \cref{lem:stalk_tensor_linear_map}): the two
  carriers are canonically identified but not definitionally equal, and the same
  \(\mathtt{RingEquiv}\) bridge resolves the obstruction here. Thus the stage-(iv)
  balancing reuses the technique already resolved in stage (iii) rather than re-deriving
  it.

  \emph{(v) Mutual inversion on germ generators} (\(\mathtt{stalkTensorIso}\)). The two
  composites are checked to be identities on generators. The composite
  \(A_x \otimes_{R_x} B_x \to (A \otimes^{p} B).\mathtt{stalk}\,x \to A_x
  \otimes_{R_x} B_x\) is the identity on a simple tensor \(\mathrm{germ}_x a \otimes
  \mathrm{germ}_x b\) by the germ characterisation \(\mathtt{stalkTensorDesc\_germ\_tmul}\)
  of (ii), and extends to all of \(A_x \otimes_{R_x} B_x\) by
  \(\mathtt{TensorProduct.induction\_on}\). The opposite composite
  \((A \otimes^{p} B).\mathtt{stalk}\,x \to A_x \otimes_{R_x} B_x \to (A \otimes^{p}
  B).\mathtt{stalk}\,x\) is the identity on a germ \(\mathrm{germ}_x(a \otimes b)\) of
  a simple tensor by the same characterisation; since germs generate the stalk and the
  germ maps are jointly epimorphic (\(\mathtt{stalk\_hom\_ext}\)), together with
  \(\mathtt{TensorProduct.induction\_on}\) applied to each section, the composite is the
  identity. Bundling the \(R_x\)-linear forward map of (iii) with the reverse map of
  (iv) and these two inversion identities yields the asserted isomorphism
  \(\mathtt{stalkTensorIso}\).

  Naturality in \(A, B\) is inherited from naturality of the sectionwise tensor and of
  the germ maps. This is the presheaf-of-modules, varying-ring analogue of the Stacks
  lemma above; the only subtlety beyond the ringed-space statement is that the ground
  ring varies with \(U\), handled by the ``tensor of filtered colimits over the colimit
  ring'' identification together with the carrier-duality bridge of (iii). Conceptually
  the tensor product of modules commutes with filtered colimits in each variable and
  the ground ring \(R(U)\) itself colimits to the stalk ring \(R.\mathtt{stalk}\,x\),
  which is the reason the forward map (ii) and the reverse map (iv) are mutually
  inverse.
\end{proof}

\begin{lemma}

\leanok
[Forward stalk--tensor comparison map (d.2, partial)]
  \label{lem:stalk_tensor_desc_forward}
  \lean{PresheafOfModules.stalkTensorDesc}
  \uses{def:scheme_modules_tensorobj}
  With \(R\), \(A\), \(B\), \(x\) as in \cref{lem:stalk_tensor_commutation}, there is
  a canonical natural \emph{additive} comparison map
  \[
    \mathtt{stalkTensorDesc} \;\colon\;
      (A \otimes^{p} B).\mathtt{stalk}\,x \;\longrightarrow\;
      A.\mathtt{stalk}\,x \;\otimes_{R.\mathtt{stalk}\,x}\; B.\mathtt{stalk}\,x,
  \]
  characterised on germs of simple tensors by
  \[
    \mathtt{stalkTensorDesc}\bigl(\mathrm{germ}_x(a \otimes b)\bigr)
      \;=\; \mathrm{germ}_x a \,\otimes\, \mathrm{germ}_x b .
  \]
  This is the \emph{forward half} of the isomorphism
  \cref{lem:stalk_tensor_commutation}: the additive comparison map of
  stages (i)--(ii) of that lemma's proof, prior to the \(R_x\)-linearity upgrade
  (iii) and the construction of the reverse map (iv).
\end{lemma}

\begin{lemma}

\leanok
[\(R_x\)-linear stalk--tensor comparison map (d.2, stage (iii))]
  \label{lem:stalk_tensor_linear_map}
  \lean{PresheafOfModules.stalkTensorLinearMap}
  \uses{def:scheme_modules_tensorobj, lem:stalk_tensor_desc_forward, lem:stalk_linear_map}
  With \(R\), \(A\), \(B\), \(x\) as in \cref{lem:stalk_tensor_commutation}, the
  additive forward map \(\mathtt{stalkTensorDesc}\) of
  \cref{lem:stalk_tensor_desc_forward} upgrades to an \(R.\mathtt{stalk}\,x\)-\emph{linear}
  map
  \[
    \mathtt{stalkTensorLinearMap} \;\colon\;
      (A \otimes^{p} B).\mathtt{stalk}\,x \;\longrightarrow_{R.\mathtt{stalk}\,x}\;
      A.\mathtt{stalk}\,x \;\otimes_{R.\mathtt{stalk}\,x}\; B.\mathtt{stalk}\,x,
  \]
  agreeing with \(\mathtt{stalkTensorDesc}\) on underlying additive maps, and hence
  still characterised on germs of simple tensors by
  \[
    \mathtt{stalkTensorLinearMap}\bigl(\mathrm{germ}_x(a \otimes b)\bigr)
      \;=\; \mathrm{germ}_x a \,\otimes\, \mathrm{germ}_x b .
  \]
  This is the stage-(iii) landing of the construction narrated in
  \cref{lem:stalk_tensor_commutation}, and the direct input to stage (iv) (the reverse
  map).
\end{lemma}

The proof of the live whiskering lemma below factors through two named
sub-lemmas: the \emph{d.1-bridge} relating membership in the sheafification
localizer \(J.W\) to stalkwise isomorphism, and the \emph{second-factor
naturality} of the d.2 commutation isomorphism \cref{lem:stalk_tensor_commutation}.
We state both before the lemma that consumes them.

\begin{lemma}

\leanok
[d.1-bridge: locally bijective implies stalkwise iso]
  \label{lem:W_implies_stalkwise_iso}
  \lean{PresheafOfModules.isIso_stalkFunctor_map_of_W}
  \uses{lem:stalk_linear_map}
  Let \(X\) be a topological space and \(R\) a presheaf of commutative rings on the
  small site \(\mathtt{Opens}\,X\), equipped with its canonical Grothendieck topology
  \(J\) satisfying \(J.\mathtt{WEqualsLocallyBijective}\,\mathtt{Ab}\). Let
  \(g : M \to N\) be a morphism of presheaves of \(R\)-modules whose underlying
  morphism of \(\mathtt{Ab}\)-valued presheaves lies in \(J.W\). Then for every point
  \(x \in X\) the induced \(\mathtt{Ab}\)-stalk map
  \[
    (\mathtt{stalkFunctor}\,\mathtt{Ab}\,x).\mathtt{map}\,
      \bigl((\mathtt{toPresheaf}).\mathtt{map}\,g\bigr)
    \;\colon\; M_x \;\longrightarrow\; N_x
  \]
  is an isomorphism. Conversely, if the \(\mathtt{Ab}\)-stalk map is an isomorphism at
  every point, then \((\mathtt{toPresheaf}).\mathtt{map}\,g \in J.W\).
\end{lemma}

\begin{proof}
  \uses{lem:stalk_linear_map}
  \leanok
  Membership in \(J.W = J.\mathtt{WEqualsLocallyBijective}\,\mathtt{Ab}\) unfolds to
  local bijectivity: \(g\) is both \(J\)-locally injective and \(J\)-locally
  surjective. On the topological site \(\mathtt{Opens}\,X\) the Grothendieck-topology
  notions of local injectivity and local surjectivity coincide with the classical
  ``locally on a cover'' conditions, and these are detected stalkwise: a morphism of
  \(\mathtt{Ab}\)-valued presheaves is locally injective iff every stalk map is
  injective and locally surjective iff every stalk map is surjective. Concretely, for
  the topological site one identifies local bijectivity of
  \((\mathtt{toPresheaf}).\mathtt{map}\,g\) with the associated sheaf map being an
  isomorphism and applies the stalkwise criterion
  \(\mathtt{TopCat.Presheaf.isIso\_iff\_stalkFunctor\_map\_iso}\); equivalently one
  combines \(\mathtt{GrothendieckTopology.W.isLocallyInjective}\) with the injective
  stalk criterion \(\mathtt{TopCat.Presheaf.app\_injective\_iff\_stalkFunctor\_map\_injective}\)
  and the matching locally-surjective-on-stalks characterisation. A stalk map that is
  both injective and surjective is an isomorphism of \(\mathtt{Ab}\)-stalks, which is
  the forward implication; the converse runs the same equivalences in reverse, since a
  family of stalk isomorphisms forces local bijectivity. (The \(R_x\)-linear refinement
  of these stalk maps, used downstream, is \cref{lem:stalk_linear_map}.)
\end{proof}

\begin{lemma}
% NOTE: iter-238 (review agent) — dangling `\lean{}` pin REMOVED. This B-naturality argument
% was never formalized as a standalone declaration: the prover (iter-237) inlined it as the
% `have key` step inside `isLocallyInjective_whiskerLeft_of_W` (Vestigial.lean), to avoid
% re-plumbing the +/0 carrier-duality at the module-morphism boundary twice. No decl named
% `PresheafOfModules.stalkTensorIso_naturality_right` exists, so the pin was deleted. The block
% is retained (unformalized) as an auxiliary statement because its `\label` is still consumed by
% `\uses`/`\cref` downstream; its mathematical content is discharged inside the whiskering lemma.
[B-naturality of the stalk--tensor comparison]
  \label{lem:stalk_tensor_commutation_naturality_right}
  % \lean: TODO placeholder per DAG integrity rule 1. As the NOTE above records,
  % no standalone decl exists (the B-naturality content was inlined into
  % PresheafOfModules.isLocallyInjective_whiskerLeft_of_W); the placeholder keeps
  % the node rule-1 compliant without falsely claiming a real declaration.
  \lean{AlgebraicGeometry.TODO.stalk_tensor_commutation_naturality_right}
  \uses{lem:stalk_tensor_commutation, lem:stalk_linear_map}
  Let \(R\), \(A\) and \(x \in X\) be as in \cref{lem:stalk_tensor_commutation}, and let
  \(g : M \to N\) be a morphism of presheaves of \(R\)-modules. The isomorphism
  \(\mathtt{stalkTensorIso}\) of \cref{lem:stalk_tensor_commutation} is natural in the
  \emph{second} tensor factor: the square
  \[
    \begin{array}{ccc}
      (A \otimes^{p} M).\mathtt{stalk}\,x
        & \xrightarrow{\ (A \triangleleft g)_x\ }
        & (A \otimes^{p} N).\mathtt{stalk}\,x \\[4pt]
      \big\downarrow{\scriptstyle \mathtt{stalkTensorIso}\,A\,M}
        & & \big\downarrow{\scriptstyle \mathtt{stalkTensorIso}\,A\,N} \\[4pt]
      A_x \otimes_{R_x} M_x
        & \xrightarrow{\ \mathrm{id}_{A_x} \otimes\, g_x\ }
        & A_x \otimes_{R_x} N_x
    \end{array}
  \]
  commutes, where \(g_x\) is the \(R_x\)-linear stalk map of
  \cref{lem:stalk_linear_map} and \(\mathrm{id}_{A_x} \otimes\, g_x =
  \mathtt{LinearMap.lTensor}\,A_x\,g_x\). Equivalently, under \(\mathtt{stalkTensorIso}\)
  the stalk of the left-whiskered morphism \(A \triangleleft g = \mathrm{id}_A \otimes g\)
  is identified with \(\mathtt{LinearMap.lTensor}\,A_x\,(\mathtt{stalkLinearMap}\,g\,x)\).
\end{lemma}

\begin{proof}
  \uses{lem:stalk_tensor_commutation, lem:stalk_linear_map}
  Both composites in the square are \(R_x\)-linear maps out of
  \((A \otimes^{p} M).\mathtt{stalk}\,x\), so by \(\mathtt{TensorProduct.induction\_on}\)
  together with the joint epimorphy of the germ maps \(\mathtt{stalk\_hom\_ext}\) it
  suffices to check that they agree on germs of simple tensors
  \(\mathrm{germ}_x(a \otimes m)\), with \(a \in A(U)\) and \(m \in M(U)\). On such a
  generator the left-then-bottom composite sends \(\mathrm{germ}_x(a \otimes m)\) first
  to \(\mathrm{germ}_x a \otimes \mathrm{germ}_x m\) by the germ characterisation
  \(\mathtt{stalkTensorLinearMap\_germ\_tmul}\) of \(\mathtt{stalkTensorIso}\,A\,M\)
  (stage (iii) of \cref{lem:stalk_tensor_commutation}), and then to
  \(\mathrm{germ}_x a \otimes g_x(\mathrm{germ}_x m)\) by definition of
  \(\mathtt{LinearMap.lTensor}\). The germ--functoriality characterisation
  \(\mathtt{stalkLinearMap\_germ}\) of \cref{lem:stalk_linear_map} gives
  \(g_x(\mathrm{germ}_x m) = \mathrm{germ}_x(g_U\,m)\), so this equals
  \(\mathrm{germ}_x a \otimes \mathrm{germ}_x(g_U\,m)\). The top-then-right composite
  sends \(\mathrm{germ}_x(a \otimes m)\) first to
  \(\mathrm{germ}_x\bigl((A \triangleleft g)_U(a \otimes m)\bigr) =
  \mathrm{germ}_x(a \otimes g_U\,m)\) --- the whiskered map acts as
  \(\mathrm{id}_A \otimes g\) sectionwise --- and then, again by
  \(\mathtt{stalkTensorLinearMap\_germ\_tmul}\) for \(\mathtt{stalkTensorIso}\,A\,N\),
  to \(\mathrm{germ}_x a \otimes \mathrm{germ}_x(g_U\,m)\). The two results agree, and
  the inverse direction (the analogous identity for the reverse map
  \(\mathtt{stalkTensorRev}\) via \(\mathtt{stalkTensorRev\_germ\_tmul}\)) is checked the
  same way. Hence the square commutes.
\end{proof}

\begin{lemma}

\leanok
[Unconditional left-whiskering of the localizer, via d.2]
  \label{lem:islocallyinjective_whiskerleft_via_stalk}
  \lean{PresheafOfModules.isLocallyInjective_whiskerLeft_of_W}
  \uses{def:scheme_modules_tensorobj, lem:stalk_tensor_commutation, lem:stalk_linear_map}
  Let \(R\), \(J\) be as above with \(J.\mathtt{WEqualsLocallyBijective}\,\mathtt{Ab}\),
  let \(F\) be an \emph{arbitrary} presheaf of \(R\)-modules, and let \(g : M \to N\)
  be a morphism whose underlying morphism lies in \(J.W\). Then the left-whiskered
  morphism \(F \triangleleft g = \mathrm{id}_F \otimes g\) is \(J\)-locally injective
  (indeed lies in \(J.W\)). No flatness and no local-triviality hypothesis on \(F\) is
  required. This discharges the open left-whiskering obligation
  \cref{lem:islocallyinjective_whisker_of_W} that the unconditional associator
  \cref{lem:tensorobj_assoc_iso} consumes.
\end{lemma}

% SOURCE QUOTE PROOF: (Archon-original assembly from the Stacks stalk-tensor lemma
% above and the stalkwise criterion for membership in the sheafification localizer;
% no single external source proof.)
\begin{proof}
  \leanok
  \uses{lem:stalk_tensor_commutation, lem:stalk_linear_map,
    lem:W_implies_stalkwise_iso, lem:stalk_tensor_commutation_naturality_right}
  The proof runs in three movements: the d.1-bridge turns membership in \(J.W\) into a
  stalkwise statement, the second-factor naturality of the d.2 isomorphism transports
  the whiskered stalk into a tensor of the stalk maps, and a flatness-free
  ``tensoring by an equivalence'' argument concludes.

  \emph{(a) The d.1-bridge.} Specialise to the topological site \(\mathtt{Opens}\,X\),
  where stalks exist. By the hypothesis \(g \in J.W\) and the d.1-bridge
  \cref{lem:W_implies_stalkwise_iso}, the \(\mathtt{Ab}\)-stalk map \(g_x : M_x \to N_x\)
  is an isomorphism at every point \(x \in X\). By \cref{lem:stalk_linear_map} this
  stalk map carries an \(R_x\)-linear structure, and the bijectivity just obtained
  upgrades it to an \(R_x\)-linear equivalence
  \(\mathtt{stalkLinearEquivOfIsIso} : M_x \xrightarrow{\sim} N_x\).

  \emph{(b) Identifying the whiskered stalk.} By the second-factor naturality of the
  stalk--tensor comparison \cref{lem:stalk_tensor_commutation_naturality_right}, the
  isomorphism \(\mathtt{stalkTensorIso}\) of \cref{lem:stalk_tensor_commutation}
  identifies the stalk \((F \triangleleft g)_x\) of the whiskered morphism with the
  map
  \[
    \mathrm{id}_{F_x} \otimes g_x \;=\; \mathtt{LinearMap.lTensor}\,F_x\,g_x
      \;\colon\; F_x \otimes_{R_x} M_x \longrightarrow F_x \otimes_{R_x} N_x .
  \]

  \emph{(c) Conclusion.} Tensoring the \(R_x\)-linear equivalence
  \(g_x : M_x \xrightarrow{\sim} N_x\) of (a) on the left by the identity of \(F_x\)
  yields an \(R_x\)-linear equivalence
  \(\mathtt{LinearMap.lTensor}\,F_x\,g_x = \mathrm{id}_{F_x} \otimes g_x\)
  (this is \(\mathtt{LinearEquiv.lTensor}\); it needs \emph{no} flatness of \(F_x\),
  because lifting an isomorphism through \(F_x \otimes_{R_x} (-)\) only uses the
  functoriality of the tensor product, and the inverse is supplied by the inverse
  equivalence). By (b), \((F \triangleleft g)_x\) is therefore an isomorphism of
  \(\mathtt{Ab}\)-stalks for every \(x\). Applying the converse half of the d.1-bridge
  \cref{lem:W_implies_stalkwise_iso} to the morphism \(F \triangleleft g\) gives
  \(F \triangleleft g \in J.W\), and in particular it is \(J\)-locally injective
  (\(\mathtt{GrothendieckTopology.W.isLocallyInjective}\)). No flatness, no local
  triviality, and the whiskering object \(F\) is arbitrary.

  The statement holds over any Grothendieck site \(J\) satisfying
  \(J.\mathtt{WEqualsLocallyBijective}\,\mathtt{Ab}\), and is applied in particular
  at the topological site \(\mathtt{Opens}\,X\) of a scheme, where the stalk
  arguments of (a)--(c) are available and through which the unconditional associator
  \cref{lem:tensorobj_assoc_iso} is realized.
\end{proof}

\subsection{Stalk--tensor comparison helper declarations}
\label{subsec:tensorobj_stalktensor_helpers}

The following are the internal Lean helper declarations assembling the forward
comparison map \cref{lem:stalk_tensor_desc_forward}, its \(R_x\)-linear packaging
\cref{lem:stalk_tensor_linear_map}, the reverse map, and the mutual-inversion germ
identities feeding the stalk--tensor isomorphism \cref{lem:stalk_tensor_commutation}.
Each is proved directly in Lean and carries the dependency edges of the construction.

\begin{definition}\leanok
[Balanced bilinear comparison map]
  \label{def:stalk_tensor_bilin}
  \lean{PresheafOfModules.stalkTensorBilin}
  \uses{def:scheme_modules_tensorobj}
  For a neighbourhood \(U \ni x\), the \(R(U)\)-bilinear additive map
  \(A(U) \times B(U) \to A_x \otimes_{R_x} B_x\) sending \((a,b)\) to
  \(\mathrm{germ}_x a \otimes \mathrm{germ}_x b\).
\end{definition}
\begin{proof}
  \uses{def:scheme_modules_tensorobj}
  \leanok
  Proved directly in Lean, as the germ-of-tensor pairing built from the
  \(\mathtt{TensorProduct.mk}\) bilinear map post-composed with the germ maps.
\end{proof}

\begin{lemma}\leanok
[Balancing identity of the comparison map]
  \label{lem:stalk_tensor_bilin_balanced}
  \lean{PresheafOfModules.stalkTensorBilin_balanced}
  \uses{def:stalk_tensor_bilin}
  The bilinear map \(\mathtt{stalkTensorBilin}\) is \(R(U)\)-balanced:
  \((r \cdot a, b)\) and \((a, r \cdot b)\) have the same image.
\end{lemma}
\begin{proof}
  \uses{def:stalk_tensor_bilin}
  \leanok
  Proved directly in Lean, from germ--scalar compatibility and the balancedness of
  the stalk tensor over \(R_x\).
\end{proof}

\begin{definition}\leanok
[Per-neighbourhood descended comparison map]
  \label{def:stalk_tensor_desc_u}
  \lean{PresheafOfModules.stalkTensorDescU}
  \uses{def:stalk_tensor_bilin, lem:stalk_tensor_bilin_balanced}
  The descent of \(\mathtt{stalkTensorBilin}\) through the relative tensor product to
  an additive map \((A \otimes^{p} B)(U) = A(U) \otimes_{R(U)} B(U) \to A_x
  \otimes_{R_x} B_x\).
\end{definition}
\begin{proof}
  \uses{def:stalk_tensor_bilin, lem:stalk_tensor_bilin_balanced}
  \leanok
  Proved directly in Lean, via \(\mathtt{TensorProduct.liftAddHom}\) applied to the
  balanced bilinear map.
\end{proof}

\begin{lemma}\leanok
[Per-neighbourhood descent on a simple tensor]
  \label{lem:stalk_tensor_desc_u_tmul}
  \lean{PresheafOfModules.stalkTensorDescU_tmul}
  \uses{def:stalk_tensor_desc_u}
  The descent \(\mathtt{stalkTensorDescU}\) sends \(a \otimes b\) to
  \(\mathrm{germ}_x a \otimes \mathrm{germ}_x b\).
\end{lemma}
\begin{proof}
  \uses{def:stalk_tensor_desc_u}
  \leanok
  Proved directly in Lean, unfolding the lift on a simple tensor.
\end{proof}

\begin{lemma}\leanok
[Scalar compatibility of the per-neighbourhood descent]
  \label{lem:stalk_tensor_desc_u_smul}
  \lean{PresheafOfModules.stalkTensorDescU_smul}
  \uses{def:stalk_tensor_desc_u}
  For \(r \in R(U)\) and a section \(z\), \(\mathtt{stalkTensorDescU}\,(r \cdot z) =
  (\mathrm{germ}_x r) \cdot \mathtt{stalkTensorDescU}\,z\).
\end{lemma}
\begin{proof}
  \uses{def:stalk_tensor_desc_u}
  \leanok
  Proved directly in Lean by induction on \(z\), moving the scalar across the tensor
  via germ--scalar compatibility.
\end{proof}

\begin{lemma}\leanok
[Factoring through the germ]
  \label{lem:germ_stalk_tensor_desc}
  \lean{PresheafOfModules.germ_stalkTensorDesc}
  \uses{lem:stalk_tensor_desc_forward, def:stalk_tensor_desc_u}
  The comparison map \(\mathtt{stalkTensorDesc}\) factors the per-neighbourhood map
  \(\mathtt{stalkTensorDescU}\) through the germ, by the colimit universal property.
\end{lemma}
\begin{proof}
  \uses{lem:stalk_tensor_desc_forward, def:stalk_tensor_desc_u}
  \leanok
  Proved directly in Lean, as the colimit \(\iota\)-descent identity.
\end{proof}

\begin{lemma}\leanok
[Element form of the germ factorisation]
  \label{lem:stalk_tensor_desc_germ}
  \lean{PresheafOfModules.stalkTensorDesc_germ}
  \uses{lem:stalk_tensor_desc_forward, def:stalk_tensor_desc_u}
  The comparison map sends the germ of a section \(w\) over \(U\) to its
  per-neighbourhood descent \(\mathtt{stalkTensorDescU}\,w\).
\end{lemma}
\begin{proof}
  \uses{lem:stalk_tensor_desc_forward, def:stalk_tensor_desc_u}
  \leanok
  Proved directly in Lean, the elementwise reading of
  \cref{lem:germ_stalk_tensor_desc}.
\end{proof}

\begin{lemma}\leanok
[Germ characterisation on a simple tensor]
  \label{lem:stalk_tensor_desc_germ_tmul}
  \lean{PresheafOfModules.stalkTensorDesc_germ_tmul}
  \uses{lem:stalk_tensor_desc_forward, lem:stalk_tensor_desc_u_tmul}
  The comparison map sends the germ of \(a \otimes b\) to
  \(\mathrm{germ}_x a \otimes \mathrm{germ}_x b\).
\end{lemma}
\begin{proof}
  \uses{lem:stalk_tensor_desc_forward, lem:stalk_tensor_desc_u_tmul}
  \leanok
  Proved directly in Lean, combining the germ factorisation with the simple-tensor
  value of \(\mathtt{stalkTensorDescU}\).
\end{proof}

\begin{lemma}\leanok
[Germ characterisation of the linear comparison map]
  \label{lem:stalk_tensor_linear_map_germ_tmul}
  \lean{PresheafOfModules.stalkTensorLinearMap_germ_tmul}
  \uses{lem:stalk_tensor_linear_map, lem:stalk_tensor_desc_germ_tmul}
  The \(R_x\)-linear comparison map \(\mathtt{stalkTensorLinearMap}\) sends the germ of
  \(a \otimes b\) to \(\mathrm{germ}_x a \otimes \mathrm{germ}_x b\).
\end{lemma}
\begin{proof}
  \uses{lem:stalk_tensor_linear_map, lem:stalk_tensor_desc_germ_tmul}
  \leanok
  Proved directly in Lean: the linear repackaging agrees with the additive map on
  germs.
\end{proof}

\begin{lemma}\leanok
[Germ of a restricted simple tensor]
  \label{lem:germ_tensor_obj_map_tmul}
  \lean{PresheafOfModules.germ_tensorObj_map_tmul}
  \uses{def:scheme_modules_tensorobj}
  For \(j : Z \to W\) and sections \(a, b\) over \(W\), the germ over \(Z\) of
  \((a|_Z) \otimes (b|_Z)\) equals the germ over \(W\) of \(a \otimes b\).
\end{lemma}
\begin{proof}
  \uses{def:scheme_modules_tensorobj}
  \leanok
  Proved directly in Lean, folding the restriction through
  \(\mathtt{tensorObj\_map\_tmul}\) and using germ--restriction compatibility.
\end{proof}

\begin{definition}\leanok
[Inner cocone leg of the reverse map]
  \label{def:rev_inner_leg}
  \lean{PresheafOfModules.revInnerLeg}
  \uses{def:scheme_modules_tensorobj}
  For a fixed \(A\)-section \(a\) over \(U\) and a \(B\)-neighbourhood \(V \ni x\), the
  additive map \(B(V) \to (A \otimes^{p} B)_x\) sending \(b\) to the germ of
  \((a|_{U \cap V}) \otimes (b|_{U \cap V})\).
\end{definition}
\begin{proof}
  \uses{def:scheme_modules_tensorobj}
  \leanok
  Proved directly in Lean, the germ-of-tensor pairing over the refinement
  \(U \cap V\).
\end{proof}

\begin{lemma}\leanok
[Evaluation of the inner leg]
  \label{lem:rev_inner_leg_apply}
  \lean{PresheafOfModules.revInnerLeg_apply}
  \uses{def:rev_inner_leg}
  \(\mathtt{revInnerLeg}\) sends \(b\) to the germ of
  \((a|_{U \cap V}) \otimes (b|_{U \cap V})\).
\end{lemma}
\begin{proof}
  \uses{def:rev_inner_leg}
  \leanok
  Proved directly in Lean, by definitional unfolding.
\end{proof}

\begin{definition}\leanok
[Inner descent of the reverse map]
  \label{def:rev_inner}
  \lean{PresheafOfModules.revInner}
  \uses{def:rev_inner_leg}
  The additive map \(B_x \to (A \otimes^{p} B)_x\) descending \(\mathtt{revInnerLeg}\)
  through the \(B\)-stalk colimit.
\end{definition}
\begin{proof}
  \uses{def:rev_inner_leg}
  \leanok
  Proved directly in Lean, via the colimit universal property over the
  \(B\)-neighbourhood diagram.
\end{proof}

\begin{lemma}\leanok
[Factoring the inner descent through the germ]
  \label{lem:germ_rev_inner}
  \lean{PresheafOfModules.germ_revInner}
  \uses{def:rev_inner, def:rev_inner_leg}
  The inner descent \(\mathtt{revInner}\) factors \(\mathtt{revInnerLeg}\) through the
  \(B\)-germ, by the colimit universal property.
\end{lemma}
\begin{proof}
  \uses{def:rev_inner, def:rev_inner_leg}
  \leanok
  Proved directly in Lean, the colimit \(\iota\)-descent identity.
\end{proof}

\begin{lemma}\leanok
[Germ characterisation of the inner descent]
  \label{lem:rev_inner_germ}
  \lean{PresheafOfModules.revInner_germ}
  \uses{def:rev_inner, lem:rev_inner_leg_apply}
  \(\mathtt{revInner}\) sends \(\mathrm{germ}_x b\) to the germ of
  \((a|_{U \cap V}) \otimes (b|_{U \cap V})\).
\end{lemma}
\begin{proof}
  \uses{def:rev_inner, lem:rev_inner_leg_apply}
  \leanok
  Proved directly in Lean, from the germ factorisation and the inner-leg value.
\end{proof}

\begin{definition}\leanok
[Outer cocone leg of the reverse map]
  \label{def:rev_outer_leg}
  \lean{PresheafOfModules.revOuterLeg}
  \uses{def:rev_inner, lem:rev_inner_germ}
  The additive map \(A(U) \to (B_x \to (A \otimes^{p} B)_x)\) sending an \(A\)-section
  \(a\) to the inner descent \(\mathtt{revInner}\) it determines.
\end{definition}
\begin{proof}
  \uses{def:rev_inner, lem:rev_inner_germ}
  \leanok
  Proved directly in Lean; additivity in \(a\) is checked on germ generators of
  \(B_x\) by bilinearity of the section tensor.
\end{proof}

\begin{lemma}\leanok
[Evaluation of the outer leg]
  \label{lem:rev_outer_leg_apply}
  \lean{PresheafOfModules.revOuterLeg_apply}
  \uses{def:rev_outer_leg}
  The outer leg evaluated at \(a\) is the inner descent \(\mathtt{revInner}\).
\end{lemma}
\begin{proof}
  \uses{def:rev_outer_leg}
  \leanok
  Proved directly in Lean, by definitional unfolding.
\end{proof}

\begin{definition}\leanok
[Outer descent: the additive bilinear comparison]
  \label{def:rev_bihom}
  \lean{PresheafOfModules.revBihom}
  \uses{def:rev_outer_leg}
  The additive map \(A_x \to (B_x \to (A \otimes^{p} B)_x)\) descending
  \(\mathtt{revOuterLeg}\) through the \(A\)-stalk colimit.
\end{definition}
\begin{proof}
  \uses{def:rev_outer_leg}
  \leanok
  Proved directly in Lean, via the colimit universal property over the
  \(A\)-neighbourhood diagram.
\end{proof}

\begin{lemma}\leanok
[Factoring the outer descent through the germ]
  \label{lem:germ_rev_bihom}
  \lean{PresheafOfModules.germ_revBihom}
  \uses{def:rev_bihom, def:rev_outer_leg}
  The outer descent \(\mathtt{revBihom}\) factors \(\mathtt{revOuterLeg}\) through the
  \(A\)-germ.
\end{lemma}
\begin{proof}
  \uses{def:rev_bihom, def:rev_outer_leg}
  \leanok
  Proved directly in Lean, the colimit \(\iota\)-descent identity.
\end{proof}

\begin{lemma}\leanok
[Germ characterisation of the additive bilinear comparison]
  \label{lem:rev_bihom_germ_tmul}
  \lean{PresheafOfModules.revBihom_germ_tmul}
  \uses{def:rev_bihom, lem:rev_outer_leg_apply, lem:rev_inner_germ}
  \(\mathtt{revBihom}(\mathrm{germ}_x a)(\mathrm{germ}_x b)\) is the germ of
  \((a|_{U \cap V}) \otimes (b|_{U \cap V})\).
\end{lemma}
\begin{proof}
  \uses{def:rev_bihom, lem:rev_outer_leg_apply, lem:rev_inner_germ}
  \leanok
  Proved directly in Lean, from the outer-germ factorisation and the inner-descent
  germ value.
\end{proof}

\begin{lemma}\leanok
[Section-level balancing identity]
  \label{lem:rev_bihom_balanced_germ}
  \lean{PresheafOfModules.revBihom_balanced_germ}
  \uses{def:scheme_modules_tensorobj}
  Over a self-intersection, restricting a scalar-twisted simple tensor commutes
  between the two factors: \(\mathrm{germ}\bigl((r_0 \cdot a)|_Z \otimes b|_Z\bigr) =
  \mathrm{germ}\bigl(a|_Z \otimes (r_0 \cdot b)|_Z\bigr)\).
\end{lemma}
\begin{proof}
  \uses{def:scheme_modules_tensorobj}
  \leanok
  Proved directly in Lean, transporting the clean \(W\)-level identity
  \((r_0 \cdot a) \otimes b = a \otimes (r_0 \cdot b)\) down through the restriction.
\end{proof}

\begin{lemma}\leanok
[Stalk-level balancing of the bilinear comparison]
  \label{lem:rev_bihom_balanced}
  \lean{PresheafOfModules.revBihom_balanced}
  \uses{def:rev_bihom, lem:rev_bihom_germ_tmul, lem:rev_bihom_balanced_germ}
  The bilinear comparison \(\mathtt{revBihom}\) is \(R_x\)-balanced, the condition
  feeding \(\mathtt{TensorProduct.liftAddHom}\) to produce the reverse map.
\end{lemma}
\begin{proof}
  \uses{def:rev_bihom, lem:rev_bihom_germ_tmul, lem:rev_bihom_balanced_germ}
  \leanok
  Proved directly in Lean: on germ generators over a common neighbourhood the scalar
  stays in \(R_x\) and the two germs reduce to the section-level identity.
\end{proof}

\begin{definition}\leanok
[The reverse map of the stalk--tensor comparison]
  \label{def:stalk_tensor_rev}
  \lean{PresheafOfModules.stalkTensorRev}
  \uses{def:rev_bihom, lem:rev_bihom_balanced}
  The descent of the \(R_x\)-balanced bilinear map \(\mathtt{revBihom}\) through the
  relative tensor product over \(R_x\), giving \(A_x \otimes_{R_x} B_x \to
  (A \otimes^{p} B)_x\) --- the candidate inverse to \(\mathtt{stalkTensorLinearMap}\).
\end{definition}
\begin{proof}
  \uses{def:rev_bihom, lem:rev_bihom_balanced}
  \leanok
  Proved directly in Lean, via \(\mathtt{TensorProduct.liftAddHom}\) applied to the
  balanced comparison.
\end{proof}

\begin{lemma}\leanok
[Germ characterisation of the reverse map]
  \label{lem:stalk_tensor_rev_germ_tmul}
  \lean{PresheafOfModules.stalkTensorRev_germ_tmul}
  \uses{def:stalk_tensor_rev, lem:rev_bihom_germ_tmul}
  The reverse map sends \(\mathrm{germ}_x a \otimes \mathrm{germ}_x b\) to the germ of
  \((a|_{U \cap V}) \otimes (b|_{U \cap V})\).
\end{lemma}
\begin{proof}
  \uses{def:stalk_tensor_rev, lem:rev_bihom_germ_tmul}
  \leanok
  Proved directly in Lean, unfolding the lift on a simple tensor.
\end{proof}

\section{\(\otimes\)-invertibility and the commutative-monoid group-law engine}
\label{sec:tensorobj_invertibility}

The relative Picard group law is assembled from \emph{\(\otimes\)-invertibility}
and the monoidal coherence isomorphisms of
\cref{lem:tensorobj_assoc_iso,lem:tensorobj_unit_iso,lem:tensorobj_comm_iso},
following Mathlib's Picard-group idiom (\(\mathtt{Module.Invertible}\),
\(\mathtt{CommRing.Pic} = \mathrm{Units}(\mathrm{Skeleton}(\mathtt{ModuleCat}\,R))\),
\(\mathtt{instCommGroupPic}\); \texttt{Mathlib.RingTheory.PicardGroup}) as the
design template. Under the route~(e) fallback the inverse of an invertible object
is carried by the predicate itself, so neither the open-immersion
restriction-compatibility isomorphism \cref{lem:tensorobj_restrict_iso} nor a
separate tensor-inverse lemma is needed there. Under the PRIMARY locally-trivial
route, by contrast, both \cref{lem:tensorobj_restrict_iso} (via the PRIMARY proof
of \cref{lem:islocallyinjective_whisker_of_W}) and the tensor-inverse lemma
\cref{lem:tensorobj_inverse_invertible} (supplying inverses to
\cref{lem:tensorobj_isoclass_commgroup}) are on the critical path.

\begin{definition}\leanok
[\(\otimes\)-invertibility of an \(\mathcal{O}_X\)-module]
  \label{def:scheme_modules_isinvertible}
  \lean{AlgebraicGeometry.Scheme.Modules.IsInvertible}
  \uses{def:scheme_modules_tensorobj}
  % SOURCE: [Stacks Project], Tag 01CR (Modules on Ringed Spaces, \S Invertible
  % modules), Definition 01CS (definition-invertible) + Lemma 0B8K
  % (lemma-invertible, the "$\exists\,\mathcal{N}$" characterisation)
  % (read from references/stacks-modules.tex, L4046-L4079).
  % SOURCE QUOTE:
  % "\begin{definition}
  % \label{definition-invertible}
  % Let $(X, \mathcal{O}_X)$ be a ringed space. An
  % {\it invertible $\mathcal{O}_X$-module} is a sheaf
  % of $\mathcal{O}_X$-modules $\mathcal{L}$ such that
  % the functor
  % $$
  % \textit{Mod}(\mathcal{O}_X) \longrightarrow \textit{Mod}(\mathcal{O}_X),\quad
  % \mathcal{F} \longmapsto \mathcal{L} \otimes_{\mathcal{O}_X} \mathcal{F}
  % $$
  % is an equivalence of categories. We say that $\mathcal{L}$ is
  % {\it trivial} if it is isomorphic as an $\mathcal{O}_X$-module
  % to $\mathcal{O}_X$.
  % \end{definition}
  % [\dots]
  % \begin{lemma}
  % \label{lemma-invertible}
  % Let $(X, \mathcal{O}_X)$ be a ringed space. Let $\mathcal{L}$
  % be an $\mathcal{O}_X$-module. Equivalent are
  % \begin{enumerate}
  % \item $\mathcal{L}$ is invertible, and
  % \item there exists an $\mathcal{O}_X$-module $\mathcal{N}$
  % such that
  % $\mathcal{L} \otimes_{\mathcal{O}_X} \mathcal{N} \cong \mathcal{O}_X$.
  % \end{enumerate}
  % In this case $\mathcal{L}$ is locally a direct summand of a finite free
  % $\mathcal{O}_X$-module and the module $\mathcal{N}$ in (2) is isomorphic to
  % $\SheafHom_{\mathcal{O}_X}(\mathcal{L}, \mathcal{O}_X)$.
  % \end{lemma}"
  \textit{Source: [Stacks Project], Tag 01CR (Modules on Ringed Spaces,
  ``Invertible modules''), Definition~01CS and Lemma~0B8K. An invertible
  \(\mathcal{O}_X\)-module is one for which \(\mathcal{L} \otimes_{\mathcal{O}_X}
  (-)\) is an equivalence; equivalently there exists \(\mathcal{N}\) with
  \(\mathcal{L} \otimes_{\mathcal{O}_X} \mathcal{N} \cong \mathcal{O}_X\), and
  then \(\mathcal{N} \cong \mathcal{H}om_{\mathcal{O}_X}(\mathcal{L},
  \mathcal{O}_X)\).}
  Let \(X\) be a scheme and \(M \in \Scheme.\mathtt{Modules}\,X\). We say \(M\)
  is \emph{\(\otimes\)-invertible} when it admits a tensor inverse: there is an
  object \(N \in \Scheme.\mathtt{Modules}\,X\) together with an isomorphism
  \[
    M \otimes_X N \;\cong\;
      \mathtt{SheafOfModules.unit}\,X.\mathtt{ringCatSheaf}
      \;=\; \mathcal{O}_X ,
  \]
  i.e.\ \(\mathtt{IsInvertible}(M) :\equiv \exists\, N,\
  \mathrm{Nonempty}\bigl(M \otimes_X N \cong
  \mathtt{SheafOfModules.unit}\,X.\mathtt{ringCatSheaf}\bigr)\), with
  \(\otimes_X\) the substrate of \cref{def:scheme_modules_tensorobj}. The
  predicate is a \emph{proposition} carrying its inverse witness existentially;
  consequently the dual/inverse needed by the group law is definitional and no
  separate ``a tensor inverse exists'' lemma is consumed downstream. This is the
  scheme-level analogue of Mathlib's \(\mathtt{Module.Invertible}\), transported
  one categorical level up to \(\Scheme.\mathtt{Modules}\,X\) with the substrate
  \(\otimes_X\).
\end{definition}

\begin{lemma}[The commutative monoid of \(\otimes\)-iso-classes and its units]
  \label{lem:tensorobj_isoclass_commgroup}
  % NOTE: the monolithic `tensorObjIsoclassCommMonoid` of the original design was
  % refactored into the carrier `PicGroup` (def:pic_carrier) and its group law
  % `picCommGroup` (thm:pic_commgroup); pin corrected to those two decls (review iter-244).
  \lean{AlgebraicGeometry.Scheme.Modules.PicGroup, AlgebraicGeometry.Scheme.Modules.picCommGroup}
  \uses{def:scheme_modules_tensorobj, def:scheme_modules_isinvertible,
        lem:scheme_modules_tensorobj_functoriality, lem:tensorobj_lift_onproduct,
        lem:tensorobj_inverse_invertible, lem:tensorobj_assoc_iso,
        lem:tensorobj_unit_iso, lem:tensorobj_comm_iso,
        def:tensorobjisoofiso, def:tensorobj_unit_self_iso, def:tensorobj_middlefour,
        def:pic_setoid, def:pic_mul, def:pic_inv}
  % SOURCE: [Stacks Project], Tag 01CR (Modules on Ringed Spaces, \S Invertible
  % modules), Definition 01CX (definition-pic) + Lemma lemma-constructions-
  % invertible (tensor of invertibles is invertible; dual is the inverse)
  % (read from references/stacks-modules.tex, L4200-L4213 and L4350-L4357).
  % SOURCE QUOTE:
  % "\begin{lemma}
  % \label{lemma-constructions-invertible}
  % Let $(X, \mathcal{O}_X)$ be a ringed space.
  % \begin{enumerate}
  % \item If $\mathcal{L}$, $\mathcal{N}$ are invertible
  % $\mathcal{O}_X$-modules, then so is
  % $\mathcal{L} \otimes_{\mathcal{O}_X} \mathcal{N}$.
  % \item If $\mathcal{L}$ is an invertible $\mathcal{O}_X$-module, then so is
  % $\SheafHom_{\mathcal{O}_X}(\mathcal{L}, \mathcal{O}_X)$ and the evaluation map
  % $\mathcal{L} \otimes_{\mathcal{O}_X}
  % \SheafHom_{\mathcal{O}_X}(\mathcal{L}, \mathcal{O}_X) \to \mathcal{O}_X$
  % is an isomorphism.
  % \end{enumerate}
  % \end{lemma}
  % [\dots]
  % \begin{definition}
  % \label{definition-pic}
  % Let $(X, \mathcal{O}_X)$ be a ringed space.
  % The {\it Picard group} $\Pic(X)$ of $X$ is the
  % abelian group whose elements are isomorphism classes of
  % invertible $\mathcal{O}_X$-modules, with addition
  % corresponding to tensor product.
  % \end{definition}"
  % SOURCE: [Mathlib], `CommRing.Pic` definition, `CommGroup` instance,
  % `mul_eq_tensor` and `inv_eq_dual`, in
  % `Mathlib/RingTheory/PicardGroup.lean` (L405--L412, L492, L495)
  % (read from .lake/packages/mathlib/Mathlib/RingTheory/PicardGroup.lean).
  % SOURCE QUOTE:
  % "/-- The Picard group of a commutative semiring R consists of the invertible R-modules,
  % up to isomorphism. -/
  % def CommRing.Pic (R : Type u) [CommSemiring R] : Type u :=
  %   Shrink (Skeleton <| SemimoduleCat.{u} R)ˣ
  % ...
  % noncomputable instance : CommGroup (Pic R) := fast_instance% (equivShrink _).symm.commGroup
  % ...
  % theorem inv_eq_dual (M : Pic R) : M⁻¹ = Pic.mk R (Dual R M) := ...
  % theorem mul_eq_tensor (M N : Pic R) : M * N = Pic.mk R (M ⊗[R] N) := ..."
  \textit{Source: [Stacks Project], Tag 01CR (Modules on Ringed Spaces,
  ``Invertible modules''), Definition~01CX and Lemma~\texttt{lemma-constructions-%
  invertible}. The isomorphism classes of invertible \(\mathcal{O}_X\)-modules
  form an abelian group \(\Pic(X)\) under \(\otimes\), with unit
  \(\mathcal{O}_X\) and inverse the dual; the tensor product of invertibles is
  invertible and the dual is a tensor inverse.}
  Let \(X\) be a scheme. The isomorphism classes of \(\otimes\)-invertible
  objects of \(\Scheme.\mathtt{Modules}\,X\)
  (\cref{def:scheme_modules_isinvertible}) form a commutative monoid under the
  substrate tensor product \(\otimes_X\) in which every element is a unit ---
  equivalently, an abelian group --- the Picard group \(\Pic(X)\) (Stacks tag
  01CR).

  \textbf{Primary construction (directly on locally-trivial iso-classes, à la
  \(\mathtt{Module.Invertible}\)).} The commutative group is built directly on
  isomorphism classes of \emph{locally-trivial} (equivalently \(\otimes\)-invertible)
  sheaves, mirroring Mathlib's ring-level Picard group
  \(\mathtt{Module.Invertible}\) / \(\mathtt{CommRing.Pic}\)
  (\texttt{Mathlib/RingTheory/PicardGroup.lean}). The group operation is
  \([L]\cdot[M] := [L \otimes_X M]\), well-defined on iso-classes by the
  bifunctoriality of \(\otimes_X\)
  (\cref{lem:scheme_modules_tensorobj_functoriality}) and again locally trivial by
  the closure lemma \cref{lem:tensorobj_lift_onproduct}; the identity is
  \([\mathcal{O}_X]\); and the inverse of \([L]\) is the class of the dual
  \([L^{-1}]\) supplied by \cref{lem:tensorobj_inverse_invertible} (with
  \(L \otimes_X L^{-1} \cong \mathcal{O}_X\)).

  The four group-axiom fields are defined \emph{by hand}, each discharged by a
  \emph{single} existence-of-isomorphism read as an equality of iso-classes --- passage
  to iso-classes turns an isomorphism of sheaves into an equality of classes:
  \[
    \mathtt{one\_mul} \;\leftarrow\; \langle \lambda_X \rangle,
    \quad
    \mathtt{mul\_one} \;\leftarrow\; \langle \rho_X \rangle,
    \quad
    \mathtt{mul\_assoc} \;\leftarrow\; \langle \alpha_{X,Y,Z} \rangle,
    \quad
    \mathtt{mul\_comm} \;\leftarrow\; \langle \beta_{X,Y} \rangle,
  \]
  where the unit laws are the unitors \cref{lem:tensorobj_unit_iso}, associativity is
  the associator \cref{lem:tensorobj_assoc_iso}, and commutativity is the braiding
  \cref{lem:tensorobj_comm_iso}, each consumed only as the proposition
  \(\mathrm{Nonempty}(\cdots \cong \cdots)\); the pentagon, triangle and hexagon
  coherence carried by a \(\MonoidalCategory\) typeclass is \emph{never} invoked in
  any of the four axiom proofs.

  This by-hand assembly does \emph{not} go through
  \(\mathtt{CategoryTheory.monoidOfSkeletalMonoidal}\) or \(\mathtt{Skeleton}\)
  (\texttt{Mathlib/CategoryTheory/Monoidal/Skeleton.lean}): those constructions
  \emph{require} a coherent \([\mathtt{MonoidalCategory}]\) (and, for the commutative
  law, \([\mathtt{BraidedCategory}]\)) on the carrier category --- precisely the
  coherent monoidal category on \(\Scheme.\mathtt{Modules}\,X\) over the \emph{varying}
  structure sheaf that this construction explicitly avoids. Consequently the group law
  requires \emph{no} \(J.W.\mathtt{IsMonoidal}\), \emph{no}
  \(\mathtt{LocalizedMonoidal}\) and \emph{no} \(\mathtt{Skeleton}\) instance on
  \(\Scheme.\mathtt{Modules}\,X\). The genuine precedent is Mathlib's fixed-ring
  Picard group \(\mathtt{CommRing.Pic}\,R := \mathrm{Skeleton}(
  \mathtt{Module.Invertible}\,R)\) (\texttt{Mathlib/RingTheory/PicardGroup.lean}):
  there the \(\mathtt{CommGroup}\) \emph{is} read off the \(\mathrm{Skeleton}\) of the
  coherent monoidal category \(\mathtt{MonoidalCategory}(\mathtt{Module.Invertible}\,R)\)
  that Mathlib supplies over a fixed ring. The present construction \emph{cannot} take
  that \(\mathrm{Skeleton}\) route --- no coherent monoidal category on
  \(\Scheme.\mathtt{Modules}\,X\) for the varying structure sheaf is available --- so
  its \(\mathtt{CommGroup}\) is built directly from the existence-of-isomorphism
  propositions above, the construction ultimately resting, as
  \(\mathtt{CommRing.Pic}\) does, on iso-classes of invertible modules. The pinned
  carrier \(\mathtt{tensorObjIsoclassCommMonoid}\) realizes this commutative-group
  structure; the scheme-level construction is genuinely absent from Mathlib at the
  pinned commit.

  This is the \emph{lower-risk} path: every ingredient is either already proved
  (the unitors \cref{lem:tensorobj_unit_iso}, the braiding
  \cref{lem:tensorobj_comm_iso}) or locally-trivial-scoped and assembled (the
  associator \cref{lem:tensorobj_assoc_iso}). The only remaining obligations on this
  critical path are the restriction-compatibility comparison
  \cref{lem:tensorobj_restrict_iso} and the
  dual/inverse \cref{lem:tensorobj_inverse_invertible}. The associator
  \cref{lem:tensorobj_assoc_iso} is supplied by the unconditional route of
  \cref{sec:tensorobj_stalk_tensor} and transitively invokes the d.2 stalk--tensor
  commutation \cref{lem:stalk_tensor_commutation}.

  \textbf{Fallback (full route~(e)).} Alternatively, once
  \cref{lem:jw_ismonoidal} lands, \(\Scheme.\mathtt{Modules}\,X =
  \mathtt{SheafOfModules}\,X.\mathtt{ringCatSheaf}\) is a full (symmetric)
  \(\MonoidalCategory\) via \(\mathtt{Localization.Monoidal.LocalizedMonoidal}\),
  and the Picard group is the units of its skeleton,
  \(\Pic(X) = \mathrm{Units}(\mathrm{Skeleton}(\Scheme.\mathtt{Modules}\,X))\),
  exactly as Mathlib builds \(\mathtt{CommRing.Pic}\,R =
  \mathrm{Units}(\mathrm{Skeleton}(\mathtt{ModuleCat}\,R))\) and transports
  \(\mathtt{instCommGroupPic}\). This fallback yields the group on \emph{all}
  invertible modules and is retained for the full monoidal generality, but it is
  not required for the relative Picard group law, which the primary construction
  already supplies.
\end{lemma}

\begin{proof}
  \uses{def:scheme_modules_tensorobj, def:scheme_modules_isinvertible,
        lem:scheme_modules_tensorobj_functoriality, lem:tensorobj_lift_onproduct,
        lem:tensorobj_inverse_invertible, lem:tensorobj_assoc_iso,
        lem:tensorobj_unit_iso, lem:tensorobj_comm_iso,
        def:tensorobjisoofiso, def:tensorobj_unit_self_iso, def:tensorobj_middlefour,
        def:pic_setoid, def:pic_mul, def:pic_inv}
  Work directly on isomorphism classes of locally-trivial sheaves, as in Mathlib's
  \(\mathtt{Module.Invertible}\) / \(\mathtt{CommRing.Pic}\)
  (\texttt{Mathlib/RingTheory/PicardGroup.lean}). The binary operation
  \([L]\cdot[M] := [L \otimes_X M]\) is well-defined on iso-classes by the
  bifunctoriality of \(\otimes_X\)
  (\cref{lem:scheme_modules_tensorobj_functoriality}) --- an isomorphism in either
  argument tensors to an isomorphism --- and its value is again locally trivial by
  the closure lemma \cref{lem:tensorobj_lift_onproduct}, so the operation closes on
  the carrier. The identity element is \([\mathcal{O}_X]\). The four group-axiom
  fields are each defined \emph{by hand} from a \emph{single} existence-of-isomorphism,
  read as an equality of iso-classes,
  \[
    \mathtt{mul\_assoc} \leftarrow \langle\alpha\rangle,
    \quad
    \mathtt{one\_mul} \leftarrow \langle\lambda\rangle,
    \quad
    \mathtt{mul\_one} \leftarrow \langle\rho\rangle,
    \quad
    \mathtt{mul\_comm} \leftarrow \langle\beta\rangle:
  \]
  associativity is \cref{lem:tensorobj_assoc_iso}, the unit laws are the unitors
  \cref{lem:tensorobj_unit_iso}, and commutativity is the braiding
  \cref{lem:tensorobj_comm_iso}. No pentagon, triangle or hexagon identity is
  invoked, since each law asserts only \(\mathrm{Nonempty}(\cdots \cong \cdots)\).
  This is the same one-isomorphism-per-law shape that Mathlib's
  \(\mathtt{monoidOfSkeletalMonoidal}\)
  (\texttt{Mathlib/CategoryTheory/Monoidal/Skeleton.lean}) uses on a skeletal
  monoidal category via its skeletality map \(\mathtt{hC}\), but we do \emph{not}
  instantiate that construction: it requires a coherent \([\mathtt{MonoidalCategory}]\)
  (and \([\mathtt{BraidedCategory}]\)) on \(\Scheme.\mathtt{Modules}\,X\), which is
  exactly the structure for the varying structure sheaf that is unavailable.
  Finally every element is invertible: the dual \(L^{-1}\) of
  \cref{lem:tensorobj_inverse_invertible} satisfies \(L \otimes_X L^{-1} \cong
  \mathcal{O}_X\), so \([L^{-1}]\) is a two-sided inverse of \([L]\); equivalently,
  \([L]\) is a unit precisely when \(L\) is \(\otimes\)-invertible
  (\cref{def:scheme_modules_isinvertible}), the inverse witness carried
  existentially by \(\mathtt{IsInvertible}\). This assembles the abelian group ---
  the Picard group of \(X\) (Stacks tag 01CR) --- with \emph{no}
  \(J.W.\mathtt{IsMonoidal}\), \(\mathtt{LocalizedMonoidal}\) or \(\mathtt{Skeleton}\),
  exactly as \(\mathtt{CommRing.Pic}\) carries its \(\mathtt{CommGroup}\) on
  iso-classes of invertible modules (there off a fixed-ring \(\mathtt{Skeleton}\),
  here by hand).

  Alternatively (fallback), once \cref{lem:jw_ismonoidal} provides the full
  \(\MonoidalCategory\) structure on \(\Scheme.\mathtt{Modules}\,X\) via
  \(\mathtt{LocalizedMonoidal}\), the group is the units of the skeleton,
  \(\mathrm{Units}(\mathrm{Skeleton}(\Scheme.\mathtt{Modules}\,X))\), exactly the
  Mathlib \(\mathtt{Units}(\mathrm{Skeleton}(\cdots))\) template used for
  \(\mathtt{CommRing.Pic}\); this route is not needed for the group law.
\end{proof}

\section{Pullback-monoidality: invertibility under a general scheme morphism}
\label{sec:tensorobj_pullback_monoidality}

The relative Picard consumer (\cref{chap:Picard_RelPicFunctor}) carries its line
bundles on the \emph{local-triviality} predicate \(\mathtt{IsLocallyTrivial}\)
(\cref{def:IsLocallyTrivial}): its carrier \(\mathtt{OnProduct}\) is the
subobject \(\{\,L \in \Scheme.\mathtt{Modules}\,(C \times_S T) \mid
\mathtt{IsLocallyTrivial}\,L\,\}\) of locally-trivial \(\mathcal{O}\)-modules, and
its functorial structure maps are \emph{pullback} maps of \(\mathcal{O}\)-modules:
the projection pullback \(\pi_T^*\) along \(C \times_S T \to T\) and the
base-change pullback along \(C \times_S T' \to C \times_S T\). For these maps to
descend to the carrier and to be \emph{additive} (an \(\mathtt{AddMonoidHom}\) of
isomorphism classes) one needs the comparison isomorphism
\[
  f^*(L \otimes_X L')
    \;\cong\;
  f^*L \otimes_Y f^*L'
\]
for \emph{locally-trivial} \(L, L'\) and a \emph{general} morphism of schemes
\(f : Y \to X\) --- the projection and the base-change morphism are neither open
immersions nor flat. This restricted comparison isomorphism, on locally-trivial
pairs, is the load-bearing result of the section
(\cref{lem:pullback_tensor_iso_loctriv}); the invertibility corollary
\cref{lem:isinvertible_pullback} is then a three-line consequence. The comparison
is needed \emph{only} on locally-trivial pairs --- the carrier never presents a
general \(\mathcal{O}\)-module --- so no comparison for arbitrary modules is
required.

This is a genuinely different fact from the open-immersion compatibility
\cref{lem:tensorobj_restrict_iso}, and not a corollary of it. That lemma handles
restriction along an open immersion, and it does so by routing through the
\emph{lax right} adjoint \(\mathtt{pushforward}/\mathtt{restrictScalars}\), which is
strong monoidal there \emph{only} because an open immersion's structure-sheaf ring
map is a local \emph{isomorphism} (\(\mathtt{appIso}\,f\)). A general scheme
morphism has no such isomorphism, so the right-adjoint route is unavailable.

\paragraph{The reduction to two facts about pullback.} The corollary
\cref{lem:isinvertible_pullback} is a three-line consequence of two facts about
pullback, exactly as in the Stacks proof of \texttt{lemma-pullback-invertible}. The
first is that pullback commutes with the tensor product on the \emph{locally-trivial}
pairs the consumer presents: the canonical comparison
\(f^*(L \otimes_X L') \cong f^*L \otimes_Y f^*L'\) is an isomorphism whenever
\(L, L'\) are locally trivial, for every \(f\)
(\cref{lem:pullback_tensor_iso_loctriv}). The second is that pullback preserves the
structure sheaf, \(f^*\mathcal{O}_X \cong \mathcal{O}_Y\)
(\cref{lem:pullback_unit_iso}, already in hand). Given a locally-trivial inverse
\(N\) with \(L \otimes_X N \cong \mathcal{O}_X\), pulling back yields
\[
  f^*L \otimes_Y f^*N
    \;\cong\; f^*(L \otimes_X N)
    \;\cong\; f^*\mathcal{O}_X
    \;\cong\; \mathcal{O}_Y,
\]
so \(f^*N\) is a tensor inverse of \(f^*L\). The locally-trivial comparison
\cref{lem:pullback_tensor_iso_loctriv} is the load-bearing input.

\paragraph{Why the apparent obstruction does not apply.} It is tempting to argue that
the comparison can never be an isomorphism because an (op)lax monoidal functor need
not preserve invertible objects --- global sections \(\Gamma\) is lax monoidal, yet
\(\Gamma(\mathbb{P}^1, \mathcal{O}(1)) = 0\) is not invertible. But that
counterexample concerns the \emph{lax} tensorator of \emph{pushforward}: \(\Gamma\)
is a right adjoint (global sections), and its comparison
\(\Gamma(\mathcal{F}) \otimes \Gamma(\mathcal{G}) \to \Gamma(\mathcal{F} \otimes
\mathcal{G})\) genuinely fails to be an isomorphism. The comparison we need is the
\emph{oplax} \(\delta\) of \emph{pullback}, a \emph{left} adjoint, running the other
way, \(f^*(M \otimes N) \to f^*M \otimes f^*N\). Pullback \emph{provably preserves
local triviality} (\cref{lem:IsLocallyTrivial_pullback}): if \(M\) is trivialised
by \(\mathcal{O}_U\) on a cover \(\{U_i\}\), then \(f^*M\) is trivialised by
\(\mathcal{O}_{f^{-1}U_i}\) on \(\{f^{-1}U_i\}\) --- using only
\(f^*\mathcal{O} \cong \mathcal{O}\) and the compatibility of pullback with
restriction along open immersions. Consequently, on locally-trivial objects, both
sides of \(\delta\) are locally trivial, and \(\delta\) restricts to a comparison
between the structure sheaf and itself on each chart, where it is forced to be an
isomorphism. The bare oplax structure supplies \(\delta\) as a \emph{morphism} for
free (\cref{lem:pullback_tensor_map}); the entire remaining content is the
chart-local check that, on a trivialising cover, \(\delta\) is an isomorphism.

\paragraph{The construction route.} The isomorphism is obtained by \emph{upgrading
the already-built oplax comparison map} \(\delta_{\mathrm{sheaf}} =
\mathtt{pullbackTensorMap}\) (\cref{lem:pullback_tensor_map}) to an isomorphism
through a chart-chase --- the same shape as the proven local-triviality results
\cref{lem:IsLocallyTrivial_pullback} and
\cref{lem:tensorobj_preserves_locally_trivial}, and \emph{not} via any concrete
inverse-image model or filtered-colimit construction. The argument decomposes into
four steps, carried out in the sub-lemmas cited by
\cref{lem:pullback_tensor_iso_loctriv}:
\begin{itemize}
  \item[(D1$'$)] \emph{Naturality of \(\delta\) in \((M,N)\).} The sheaf-level
    comparison \(\delta_{\mathrm{sheaf}}\) is natural in both arguments, assembled from
    the naturality of the abstract oplax \(\delta\) (Mathlib's
    \(\mathtt{Functor.OplaxMonoidal.}\delta\mathtt{\_natural}\)) and of the
    sheafification reconciliations used to build it
    (\cref{lem:pullback_tensor_map_natural}). This naturality is what lets a chart
    replace \(M|_{U_i}, N|_{U_i}\) by the trivial bundle \(\mathcal{O}\) without
    changing whether \(\delta\) is an isomorphism there.
  \item[(D2$'$)] \emph{\(\delta\) on the unit pair \((\mathcal{O},\mathcal{O})\) is an
    isomorphism.} On the trivial bundle the comparison
    \(\delta_{(\mathcal{O},\mathcal{O})} : f^*(\mathcal{O} \otimes \mathcal{O}) \to
    f^*\mathcal{O} \otimes f^*\mathcal{O}\) is reduced, by the \(\delta\)-wrapping
    \(\mathtt{isIso\_sheafifyDelta\_unitPair\_of\_isIso\_sheafifyEta}\) (the presheaf
    oplax left-unitality factored through sheafification, using
    \(\mathtt{W\_of\_isIso\_sheafification}\) and the flatness-free right-whiskering
    \(\mathtt{W\_whiskerRight\_of\_W}\)), to the single \emph{unit square}
    \(a_Y.\mathrm{map}(\eta\,F) \circ \mathtt{sheafifyUnitIso} = \mathtt{pullbackValIso}
    \circ \mathtt{pullbackObjUnitToUnit}\) (\cref{lem:eta_bridge_unit_square}). That
    square, transposed across the sheaf pullback--pushforward adjunction, bottoms out in
    the structure-sheaf comparison \(\mathtt{pullbackUnitIso}\)
    (\cref{lem:pullback_unit_iso}, an isomorphism for \emph{all} \(f\)); hence
    \(\delta_{(\mathcal{O},\mathcal{O})}\) is an isomorphism
    (\cref{lem:pullback_tensor_iso_unit}).
  \item[(D3$'$)] \emph{Composition coherence of \(\delta\)} (the genuinely new
    ingredient). For composable scheme morphisms \(h, f\) the comparison \(\delta\) for
    the composite \(h \circ f\) factors through the comparisons for \(f\) and for \(h\),
    conjugated by the pullback pseudofunctoriality isomorphism
    \(\mathtt{pullbackComp}\,h\,f\) (\cref{lem:pullback_tensor_map_basechange}).
    Specialising to \(h := j_i'\) at an open \(U_i\) of \(X\) with preimage
    \(f^{-1}U_i\) --- with \(g_i : f^{-1}U_i \to U_i\) the restriction of \(f\) and
    \(j_i : U_i \hookrightarrow X\), \(j_i' : f^{-1}U_i \hookrightarrow Y\) the open
    immersions, so \(f \circ j_i' = j_i \circ g_i\) --- this intertwines the
    restriction-to-\(f^{-1}U_i\) of \(\delta\) for \(f\) with the \(\delta\) for
    \(g_i\). Because \(\mathtt{pullbackTensorMap}\) is a hand-built four-fold composite
    rather than an adjunction transpose, the unit-coherence mate calculus of
    \cref{lem:pullbackObjUnitToUnit_comp} does \emph{not} transfer; the coherence is
    instead established as a paste of four composition-coherence squares, decomposing
    \(\delta\) of the composite pullback via Mathlib's
    \(\mathtt{Functor.OplaxMonoidal.comp\_}\delta\) and the conjugation
    \(\mathtt{conjugateEquiv\_pullbackComp\_inv}\) of the pseudofunctor isomorphisms.
  \item[(D4$'$)] \emph{Chart-chase assembly.} Choose a common affine cover \(\{U_i\}\)
    of \(X\) trivialising both \(M\) and \(N\) (the common refinement, exactly as in
    \cref{lem:tensorobj_preserves_locally_trivial}). On each \(f^{-1}U_i\), the
    base-change coherence (D3$'$) aligns the restriction of \(\delta\) with the
    \(\delta\) for \(g_i\); naturality (D1$'$) replaces the trivialised
    \(M|_{U_i}, N|_{U_i}\) by \(\mathcal{O}\); and (D2$'$) makes the resulting
    comparison on the unit pair an isomorphism. Thus \(\delta\) restricts to an
    isomorphism on every member of the cover \(\{f^{-1}U_i\}\) of \(Y\), and the
    restriction-detects-isomorphism criterion
    \cref{lem:isiso_of_isiso_restrict} promotes it to a global isomorphism. The
    structure mirrors the chase of \cref{lem:IsLocallyTrivial_pullback}.
\end{itemize}
The only genuinely new sub-step is (D3$'$); (D1$'$) is Mathlib naturality, (D2$'$)
is the unit coherence backed by \(\mathtt{pullbackUnitIso}\), and (D4$'$) reuses the
already-proven cover machinery
(\cref{lem:tensorobj_preserves_locally_trivial},
\cref{lem:isiso_of_isiso_restrict}). No concrete inverse-image model, no
filtered-colimit/\(\otimes\) interchange, and no general comparison for arbitrary
modules is constructed.

\begin{lemma}

\leanok
[The presheaf pushforward is lax monoidal]
  \label{lem:presheaf_pushforward_laxmonoidal}
  \lean{AlgebraicGeometry.Scheme.Modules.presheafPushforwardLaxMonoidal}
  Let \(\varphi\) be a morphism of \(\mathtt{CommRingCat}\)-valued ring presheaves,
  presenting a pushforward of presheaves of modules. Then the presheaf pushforward
  \(\mathtt{PresheafOfModules.pushforward}\,\varphi\) is \emph{lax monoidal}: it
  carries the comparison
  \[
    \mu :
      (\mathtt{pushforward}\,\varphi).\mathrm{obj}\,A
        \otimes (\mathtt{pushforward}\,\varphi).\mathrm{obj}\,B
      \;\longrightarrow\;
      (\mathtt{pushforward}\,\varphi).\mathrm{obj}\,(A \otimes B),
  \]
  obtained as the composite of the strong-monoidal pushforward of presheaves over a
  fixed ring presheaf (\(\mathtt{pushforward}_0\mathtt{OfCommRingCat}\)) with the lax
  monoidal restriction of scalars \(\mathtt{restrictScalars}\,\varphi\) along the
  ring-presheaf map. This is a statement at the \emph{presheaf} level of
  \(\mathcal{O}\)-modules; it does not require a monoidal structure on
  \(\mathtt{SheafOfModules}\).
\end{lemma}

\begin{proof}



  Restriction of scalars along a homomorphism of ring presheaves is lax monoidal --- the
  comparison sends \((a \otimes b)\) to itself under the identification of the
  restricted tensor product with the source tensor product --- and the pushforward of
  presheaves of modules over a \emph{fixed} ring presheaf is strong monoidal. The
  presheaf pushforward along \(\varphi\) factors as the composite of these two functors,
  so its lax monoidal structure is the composite of the strong-monoidal structure on the
  first factor with the lax structure on the second; the comparison \(\mu\) is the
  whiskered composite of the two comparisons.
\end{proof}

\begin{lemma}

\leanok
[The presheaf pullback is oplax monoidal; the comparison map $\delta$]
  \label{lem:presheaf_pullback_oplaxmonoidal}
  \lean{AlgebraicGeometry.Scheme.Modules.presheafPullbackOplaxMonoidal}
  \uses{lem:presheaf_pushforward_laxmonoidal}
  In the same setting, the presheaf pullback \(\mathtt{PresheafOfModules.pullback}\,\varphi\)
  --- the left adjoint of the lax-monoidal presheaf pushforward of
  \cref{lem:presheaf_pushforward_laxmonoidal} --- is \emph{oplax monoidal}. In
  particular it carries the canonical comparison \emph{morphism}
  \[
    \delta :
      (\mathtt{pullback}\,\varphi).\mathrm{obj}\,(A \otimes B)
      \;\longrightarrow\;
      (\mathtt{pullback}\,\varphi).\mathrm{obj}\,A
        \otimes (\mathtt{pullback}\,\varphi).\mathrm{obj}\,B,
  \]
  produced by the doctrinal adjunction transfer
  \(\mathtt{Adjunction.leftAdjointOplaxMonoidal}\): the oplax structure on a left
  adjoint is the adjunction mate of the lax structure on its right adjoint. The
  comparison \(\delta\) is at the \emph{presheaf} level, and it is only a
  \emph{morphism}: upgrading \(\delta\) to an isomorphism is the remaining content, and
  is exactly what the chart-chase of \cref{lem:pullback_tensor_iso_loctriv}
  supplies for a locally-trivial pair.
\end{lemma}

\begin{proof}
  \uses{lem:presheaf_pushforward_laxmonoidal}
  The presheaf pullback is by construction the left adjoint of the presheaf pushforward,
  which is lax monoidal by \cref{lem:presheaf_pushforward_laxmonoidal}. Mathlib's
  \(\mathtt{Adjunction.leftAdjointOplaxMonoidal}\) turns the lax-monoidal structure on a
  right adjoint into an oplax-monoidal structure on the left adjoint, with the oplax
  comparison \(\delta\) defined as the mate of the lax comparison \(\mu\) under the
  adjunction. Applying this to the pullback--pushforward adjunction yields the asserted
  oplax structure and the comparison morphism \(\delta\).
\end{proof}

% NOTE: ABANDONED general route, OFF the critical path. This lemma asserts the
% comparison iso for ARBITRARY modules M,N. The consumer (the relative Picard carrier
% OnProduct = { L // IsLocallyTrivial L }) never presents a general module: it needs
% the comparison only on LOCALLY-TRIVIAL pairs, which is supplied by the chart-chase
% lem:pullback_tensor_iso_loctriv WITHOUT any concrete inverse-image / filtered-colimit
% build. The general statement here is retained for the record but is no longer
% load-bearing; its proof (Steps D1-D4 below, via lem:pullback_lan_decomposition and
% lem:pullback0_tensor_iso) is the Mathlib-scale build the project has pivoted off.
% lem:isinvertible_pullback no longer \uses this lemma.
\begin{lemma}
[Pullback commutes with $\otimes$ for a general scheme morphism]
  \label{lem:pullback_tensor_iso}
  \lean{AlgebraicGeometry.Scheme.Modules.pullbackTensorIso}
  % NOTE: \lean{} pin target absent from Lean source as of iter-271 (verified by grep).
  % Forward reference to a not-yet-formalized (and, per the ABANDONED-route NOTE above,
  % off-critical-path) declaration — NOT a stale rename; no existing Lean decl to repin to.
  \uses{def:scheme_modules_tensorobj, lem:tensorobj_restrict_iso,
        lem:pullback_lan_decomposition, lem:pullback0_tensor_iso,
        lem:pullback_tensor_map}
  Let \(f : Y \to X\) be a morphism of schemes and let
  \(M, N \in \Scheme.\mathtt{Modules}\,X\) be \emph{arbitrary} \(\mathcal{O}_X\)-modules.
  Writing \(f^* = \mathtt{Scheme.Modules.pullback}\,f\) for the inverse-image functor
  of \(\mathcal{O}\)-modules, the canonical comparison
  \[
    f^*(M \otimes_X N)
      \;\xrightarrow{\;\sim\;}\;
    (f^*M) \otimes_Y (f^*N)
  \]
  is an isomorphism, natural in \(M\) and \(N\). No local-freeness, flatness, or
  open-immersion hypothesis is required: it holds for all \(\mathcal{O}_X\)-modules
  and all morphisms \(f\). This is the general strong-monoidality of pullback; its
  underlying morphism is the comparison map \cref{lem:pullback_tensor_map}, here
  upgraded to an isomorphism.
  % SOURCE: [Stacks Project], Modules on Ringed Spaces, Lemma
  % \ref{lemma-tensor-product-pullback} (pullback of $\mathcal{O}$-modules along a
  % morphism of ringed spaces commutes with tensor product, functorially)
  % (read from references/stacks-modules.tex, L2392-2400).
  % SOURCE QUOTE:
  % "\begin{lemma}
  % \label{lemma-tensor-product-pullback}
  % Let $f : (X, \mathcal{O}_X) \to (Y, \mathcal{O}_Y)$ be
  % a morphism of ringed spaces. Let $\mathcal{F}$, $\mathcal{G}$
  % be $\mathcal{O}_Y$-modules. Then
  % $f^*(\mathcal{F} \otimes_{\mathcal{O}_Y} \mathcal{G})
  % = f^*\mathcal{F} \otimes_{\mathcal{O}_X} f^*\mathcal{G}$
  % functorially in $\mathcal{F}$, $\mathcal{G}$.
  % \end{lemma}"
  \textit{Source: [Stacks Project], Modules on Ringed Spaces, Lemma
  \texttt{lemma-tensor-product-pullback}: for a morphism of ringed spaces
  \(f : (X, \mathcal{O}_X) \to (Y, \mathcal{O}_Y)\) and \(\mathcal{O}_Y\)-modules
  \(\mathcal{F}, \mathcal{G}\), one has \(f^*(\mathcal{F} \otimes \mathcal{G}) =
  f^*\mathcal{F} \otimes f^*\mathcal{G}\) functorially.}
\end{lemma}

\begin{proof}
  \uses{def:scheme_modules_tensorobj, lem:tensorobj_restrict_iso,
        lem:pullback_lan_decomposition, lem:pullback0_tensor_iso,
        lem:pullback_tensor_map}
  Mathematically, pullback is extension of scalars and extension of scalars is strong
  monoidal, so the comparison is an isomorphism for every \(f\). The only work is to
  carry this through the formalisation's \emph{abstract} pullback
  \(f^* = (\mathtt{pushforward}\,f).\mathtt{leftAdjoint}\), which exposes no sectionwise
  value. We do not evaluate \(f^*\) sectionwise; we equip a \emph{concrete} model of the
  pullback with a strong-monoidal structure and transport it to \(f^*\) along the
  uniqueness of left adjoints. This mirrors Mathlib's construction of the
  extension-of-scalars monoidal structure, where the
  \emph{left} adjoint's strong-monoidal structure is built explicitly (via the
  base-change distributivity isomorphism) and the right adjoint's lax structure is
  derived from the adjunction. The argument has four steps.

  \emph{Step 1 (D1: decomposition).} The presheaf pushforward factors as a fixed-ring
  pushforward followed by restriction of scalars,
  \(\mathtt{pushforward}\,\varphi = \mathtt{pushforward}_0\,F\,R \circ
  \mathtt{restrictScalars}\,\varphi\). Taking left adjoints reverses the composite, so the
  presheaf pullback factors as
  \[
    \mathtt{pullback}\,\varphi \;\cong\; \mathtt{extendScalars}\,\varphi \circ
    \mathtt{pullback}_0, \qquad
    \mathtt{pullback}_0 = (\mathtt{pushforward}_0\,F\,R).\mathtt{leftAdjoint},
  \]
  with \(\mathtt{pullback}_0\) the topological inverse image
  (\cref{lem:pullback_lan_decomposition}).

  \emph{Step 2 (D2: the scalar half is strong, for free).} Extension of scalars
  \(\mathtt{extendScalars}\,\varphi\) is strong monoidal: its tensorator is the
  base-change distributivity isomorphism
  \(\mathtt{TensorProduct.AlgebraTensorModule.distribBaseChange}\),
  \[
    (M \otimes_R N) \otimes_R S
      \;\cong\;
    (M \otimes_R S) \otimes_S (N \otimes_R S),
  \]
  an isomorphism for \emph{every} ring map \(R \to S\), with no flatness. Composing a
  strong-monoidal functor with \(\mathtt{pullback}_0\), the comparison \(\delta\) for
  \(\mathtt{pullback}\,\varphi\) is therefore an isomorphism \emph{if and only if} the
  comparison \(\delta_0\) for \(\mathtt{pullback}_0\) is. This isolates the topological
  half as the sole genuine obligation.

  \emph{Step 3 (D3: the topological half --- the genuine build).} The functor
  \(\mathtt{pullback}_0\) is the left Kan extension
  \(\mathrm{Lan}\,(\mathtt{Opens.map}\,f.\mathtt{base})^{\mathrm{op}}\) on underlying
  presheaves (its right adjoint \(\mathtt{pushforward}_0\) is restriction along
  \(F^{\mathrm{op}}\)). Pointwise, with \(F = (\mathtt{Opens.map}\,f.\mathtt{base})^{\mathrm{op}}\),
  \[
    \mathtt{pullback}_0(M \otimes N)(V)
      = \operatorname*{colim}_{(F \downarrow V)} \bigl(M(U) \otimes N(U)\bigr),
    \quad
    (\mathtt{pullback}_0 M \otimes \mathtt{pullback}_0 N)(V)
      = \Bigl(\operatorname*{colim} M(U)\Bigr) \otimes \Bigl(\operatorname*{colim} N(U)\Bigr).
  \]
  The comma category \((F \downarrow V) = \{U : f^{-1}V \subseteq U\}\) is up-directed
  --- \(U_1 \cup U_2\) is an upper bound of any \(U_1, U_2\) --- and the tensor product
  commutes with filtered colimits (the diagonal \((F\downarrow V) \to (F\downarrow V)
  \times (F \downarrow V)\) is final). Hence the comparison \(\delta_0\) is an
  isomorphism (\cref{lem:pullback0_tensor_iso}). This is the genuine content of the lemma.

  \emph{Step 4 (D4: sheafify and transport).} Steps 1--3 furnish a strong-monoidal
  structure on the concrete presheaf-level inverse image. Sheafify the resulting presheaf
  isomorphism to the sheaf level --- the unit and tensor factors are reconciled with
  their sheafified forms by the sheafification-monoidality brick
  \(\mathtt{sheafifyTensorUnitIso}\), the same \(\mathtt{sheafificationCompPullback}\)
  shape used in \cref{lem:tensorobj_restrict_iso}. By definition
  \(\mathtt{Scheme.Modules.pullback}\,f = (\mathtt{pushforward}\,f).\mathtt{leftAdjoint}\)
  is \emph{a} left adjoint of the pushforward; the uniqueness of left adjoints
  \(\mathtt{Adjunction.leftAdjointUniq}\) supplies a canonical natural isomorphism from it
  to the concrete model, along which the strong-monoidal structure travels. Its
  tensorator becomes the asserted isomorphism
  \[
    f^*(M \otimes_X N)
      \;\xrightarrow{\;\sim\;}\;
    (f^*M) \otimes_Y (f^*N),
  \]
  natural in \(M\) and \(N\), refining the comparison map
  \(\delta_{\mathrm{sheaf}} = \mathtt{pullbackTensorMap}\) of
  \cref{lem:pullback_tensor_map}; the abstract pullback is never evaluated sectionwise.

  No shortcut from the bare oplax map is available: the presheaf pullback is oplax
  monoidal (\cref{lem:presheaf_pullback_oplaxmonoidal}) and so supplies \(\delta\) as a
  morphism for free, but an (op)lax monoidal functor does not preserve invertible objects
  --- global sections \(\Gamma\) is lax monoidal yet \(\Gamma(\mathbb{P}^1,
  \mathcal{O}(1)) = 0\) is not invertible --- so the isomorphism must be proved, which is
  exactly the work of Steps~1--4.
\end{proof}

% NOTE: OFF the critical path. General presheaf-pullback infrastructure, retained
% (and proven), but no longer load-bearing: it served only the abandoned general
% comparison lem:pullback_tensor_iso. The committed loc-triv route
% (lem:pullback_tensor_iso_loctriv) does not use it.
\begin{lemma}

\leanok
[Presheaf pullback factors as topological inverse image then extension of scalars]
  \label{lem:pullback_lan_decomposition}
  \lean{AlgebraicGeometry.Scheme.Modules.pullbackLanDecomposition}
  \uses{lem:presheaf_pushforward_laxmonoidal, def:pullback0, def:pullback0_adjunction,
        def:extend_scalars, def:extend_scalars_adjunction}
  Let \(\varphi\) present a pushforward of presheaves of \(\mathcal{O}\)-modules along a
  functor \(F\) of ring presheaves. The presheaf pushforward factors as the fixed-ring
  pushforward followed by restriction of scalars,
  \[
    \mathtt{PresheafOfModules.pushforward}\,\varphi
      \;=\;
    \mathtt{pushforward}_0\,F\,R \circ \mathtt{restrictScalars}\,\varphi .
  \]
  Consequently its left adjoint, the presheaf pullback, factors as
  \[
    \mathtt{PresheafOfModules.pullback}\,\varphi
      \;\cong\;
    \mathtt{extendScalars}\,\varphi \circ \mathtt{pullback}_0,
    \qquad
    \mathtt{pullback}_0 := (\mathtt{pushforward}_0\,F\,R).\mathtt{leftAdjoint},
  \]
  where \(\mathtt{pullback}_0\) is the topological inverse image --- the left Kan
  extension along \((\mathtt{Opens.map}\,f.\mathtt{base})^{\mathrm{op}}\) --- and
  \(\mathtt{extendScalars}\,\varphi\) is sectionwise extension of scalars along the
  structure-sheaf ring maps.
\end{lemma}

\begin{proof}
  \uses{lem:presheaf_pushforward_laxmonoidal}
  The pushforward factorisation is the definitional decomposition of the presheaf
  pushforward into a fixed-ring pushforward and a restriction of scalars, the same
  decomposition that supplies the lax structure of
  \cref{lem:presheaf_pushforward_laxmonoidal}. The left adjoint of a composite is the
  composite of the left adjoints in the reverse order; the left adjoint of
  \(\mathtt{restrictScalars}\,\varphi\) is \(\mathtt{extendScalars}\,\varphi\), and the
  left adjoint of \(\mathtt{pushforward}_0\,F\,R\) is by definition
  \(\mathtt{pullback}_0\). The displayed isomorphism is the resulting decomposition of the
  pullback, canonical because each factor's adjoint is unique up to canonical
  isomorphism. That \(\mathtt{pullback}_0\) is the left Kan extension along
  \((\mathtt{Opens.map}\,f.\mathtt{base})^{\mathrm{op}}\) follows because its right adjoint
  \(\mathtt{pushforward}_0\) is, on underlying presheaves, restriction along that functor.
\end{proof}

% NOTE: ABANDONED general route, OFF the critical path. This is the Mathlib-scale
% filtered-colimit/tensor interchange (the irreducible piece D3 of the abandoned general
% comparison lem:pullback_tensor_iso). The consumer never needs it: the loc-triv
% comparison lem:pullback_tensor_iso_loctriv is obtained by a chart-chase upgrading the
% already-built oplax map, with no concrete inverse-image colimit model. Retained for the
% record only.
\begin{lemma}
[The topological inverse image commutes with $\otimes$]
  \label{lem:pullback0_tensor_iso}
  \lean{AlgebraicGeometry.Scheme.Modules.pullback0TensorIso}
  % NOTE: \lean{} pin target absent from Lean source as of iter-271 (verified by grep).
  % Forward reference to a not-yet-formalized (and, per the ABANDONED-route NOTE above,
  % off-critical-path) declaration — NOT a stale rename; no existing Lean decl to repin to.
  \uses{lem:pullback_lan_decomposition}
  Let \(f : Y \to X\) be a morphism of schemes and let \(\mathtt{pullback}_0\) be the
  topological inverse image of presheaves of \(\mathcal{O}\)-modules of
  \cref{lem:pullback_lan_decomposition} --- the left Kan extension along
  \((\mathtt{Opens.map}\,f.\mathtt{base})^{\mathrm{op}}\). For arbitrary presheaves of
  modules \(M, N\), the canonical comparison
  \[
    \delta_0 :
    \mathtt{pullback}_0(M \otimes N)
      \;\xrightarrow{\;\sim\;}\;
    \mathtt{pullback}_0 M \otimes \mathtt{pullback}_0 N
  \]
  is an isomorphism, natural in \(M, N\).
\end{lemma}

\begin{proof}
  \uses{lem:pullback_lan_decomposition}
  Evaluate both sides pointwise. As a left Kan extension, \(\mathtt{pullback}_0\) has the
  pointwise colimit formula
  \[
    \mathtt{pullback}_0(P)(V)
      \;=\;
    \operatorname*{colim}_{(F \downarrow V)} P(U),
    \qquad F = (\mathtt{Opens.map}\,f.\mathtt{base})^{\mathrm{op}},
  \]
  over the comma category \((F \downarrow V)\), whose objects are the opens \(U\) of
  \(X\) with \(f^{-1}V \subseteq U\). Since the presheaf tensor product is sectionwise,
  \((M \otimes N)(U) = M(U) \otimes N(U)\), one has
  \[
    \mathtt{pullback}_0(M \otimes N)(V)
      = \operatorname*{colim}_{(F \downarrow V)} \bigl(M(U) \otimes N(U)\bigr),
    \quad
    \bigl(\mathtt{pullback}_0 M \otimes \mathtt{pullback}_0 N\bigr)(V)
      = \Bigl(\operatorname*{colim} M(U)\Bigr) \otimes \Bigl(\operatorname*{colim} N(U)\Bigr),
  \]
  and \(\delta_0(V)\) is the canonical comparison between them. The comma category
  \((F \downarrow V)\) is up-directed: it is nonempty (the whole space \(X\) contains
  \(f^{-1}V\)) and any two members \(U_1, U_2\) have the upper bound \(U_1 \cup U_2\),
  still containing \(f^{-1}V\). Over a filtered (here up-directed) index the tensor
  product commutes with the colimit --- the diagonal \((F\downarrow V) \to (F\downarrow V)
  \times (F\downarrow V)\) is final, so the colimit of the pointwise tensors agrees with
  the tensor of the colimits --- whence \(\delta_0(V)\) is an isomorphism for every \(V\),
  and \(\delta_0\) is an isomorphism of presheaves, natural in \(M, N\).

  The pointwise left-Kan-extension model of \(\mathtt{pullback}_0\) and the
  filtered-colimit/tensor interchange for \(\mathcal{O}\)-module-valued presheaves are not
  in Mathlib; this lemma is therefore established bottom-up from the
  left-Kan-extension colimit formula and the filtered-colimit preservation of the
  tensor product.
\end{proof}

\begin{lemma}

\leanok
[The sheaf-level pullback--tensor comparison map $\delta_{\mathrm{sheaf}}$]
  \label{lem:pullback_tensor_map}
  \lean{AlgebraicGeometry.Scheme.Modules.pullbackTensorMap}
  \uses{lem:presheaf_pullback_oplaxmonoidal, def:scheme_modules_tensorobj,
        def:sheafify_tensor_unit_iso, def:pullback_val_iso}
  Let \(f : Y \to X\) be a morphism of schemes and let
  \(M, N \in \Scheme.\mathtt{Modules}\,X\) be \emph{arbitrary} \(\mathcal{O}_X\)-modules.
  There is a canonical comparison \emph{morphism}
  \[
    \delta_{\mathrm{sheaf}} :
      (\mathtt{Scheme.Modules.pullback}\,f).\mathrm{obj}\,(M \otimes_X N)
      \;\longrightarrow\;
      (\mathtt{Scheme.Modules.pullback}\,f).\mathrm{obj}\,M
        \otimes_Y (\mathtt{Scheme.Modules.pullback}\,f).\mathrm{obj}\,N
  \]
  of \(\mathcal{O}_Y\)-modules, natural in \(M\) and \(N\). It is a \emph{map only}: in
  general it is not asserted to be an isomorphism.
\end{lemma}

\begin{proof}
  \uses{lem:presheaf_pullback_oplaxmonoidal, def:scheme_modules_tensorobj,
        def:sheafify_tensor_unit_iso, def:pullback_val_iso}
  The presheaf pullback is oplax monoidal (\cref{lem:presheaf_pullback_oplaxmonoidal}),
  so at the presheaf level it carries the comparison morphism
  \[
    \delta : (\mathtt{pullback}\,\varphi).\mathrm{obj}\,(A \otimes B)
      \;\longrightarrow\;
      (\mathtt{pullback}\,\varphi).\mathrm{obj}\,A
        \otimes (\mathtt{pullback}\,\varphi).\mathrm{obj}\,B,
  \]
  where \(\varphi\) is the ring-presheaf map presenting \(f\). One transports \(\delta\)
  through sheafification to the sheaf level, using the same
  \(\mathtt{sheafificationCompPullback}\) device that supplies Steps 1--2 of
  \cref{lem:tensorobj_restrict_iso}: the abstract pullback of \(\mathcal{O}\)-modules
  agrees, up to canonical natural isomorphism, with the sheafification of the presheaf
  pullback, and likewise the sheafified tensor product \(\otimes_Y\) is the
  sheafification of the presheaf tensor. The right-hand side reconciliation of the unit
  and tensor factors with their sheafified forms is supplied by the
  sheafification-monoidality brick \(\mathtt{sheafifyTensorUnitIso}\). Whiskering and
  composing these canonical isomorphisms with the sheafification of \(\delta\) produces
  the asserted sheaf-level morphism \(\delta_{\mathrm{sheaf}}\), natural in \(M, N\)
  because each ingredient is. No claim of invertibility is made: \(\delta_{\mathrm{sheaf}}\)
  is the carrier on which the invertible-case isomorphism argument of
  \cref{lem:isinvertible_pullback} acts.
\end{proof}

\begin{lemma}

\leanok
[Reduction of $\mathtt{pullbackTensorMap}$ iso-ness to the sheafified presheaf $\delta$]
  \label{lem:isiso_pullbacktensormap_of_sheafifydelta}
  \lean{AlgebraicGeometry.Scheme.Modules.isIso_pullbackTensorMap_of_isIso_sheafifyDelta}
  \uses{lem:pullback_tensor_map, lem:isiso_sheafify_tensorhom_pullbackvaliso}
  Let \(f : Y \to X\) be a morphism of schemes and let
  \(M, N \in \Scheme.\mathtt{Modules}\,X\). Write
  \[
    \varphi' :
      (X.\mathtt{presheaf} \circ \mathtt{forget}_2)
      \;\longrightarrow\;
      (\mathtt{Opens.map}\,f.\mathtt{base})^{\mathrm{op}}
        \circ (Y.\mathtt{presheaf} \circ \mathtt{forget}_2)
  \]
  for the induced presheaf-of-modules map presenting \(f\), and let
  \(a_Y = \mathtt{PresheafOfModules.sheafification}\,(\mathbf{1}_{Y.\mathtt{ringCatSheaf}})\)
  denote sheafification on \(Y\). If the sheafified presheaf-level oplax comparison
  \[
    a_Y.\mathrm{map}\bigl(
      \delta\,(\mathtt{PresheafOfModules.pullback}\,\varphi')\,M.\mathtt{val}\,N.\mathtt{val}
    \bigr)
  \]
  is an isomorphism, then the sheaf-level comparison
  \(\mathtt{pullbackTensorMap}\,f\,M\,N : f^*(M \otimes_X N) \to f^*M \otimes_Y f^*N\)
  of \cref{lem:pullback_tensor_map} is an isomorphism.
\end{lemma}

\begin{proof}
  \uses{lem:pullback_tensor_map, lem:isiso_sheafify_tensorhom_pullbackvaliso}
  Unfolding \(\mathtt{pullbackTensorMap}\) (\cref{lem:pullback_tensor_map}) exhibits it as a
  four-fold composite
  \[
    \mathtt{pullbackTensorMap}\,f\,M\,N
      \;=\;
    p_1 \;\mathbin{;}\; a_Y.\mathrm{map}\,\delta \;\mathbin{;}\; p_3 \;\mathbin{;}\; p_4,
  \]
  where:
  \(p_1 = (\mathtt{sheafificationCompPullback}).\mathrm{hom}\), reconciling the abstract
  \(\mathcal{O}\)-module pullback with the sheafification of the presheaf pullback;
  \(a_Y.\mathrm{map}\,\delta\) is the sheafification of the presheaf-level oplax comparison
  \(\delta\); \(p_3 = (\mathtt{sheafifyTensorUnitIso}).\mathrm{hom}\), the
  sheafification-monoidality brick reconciling the sheafified tensor with the sheaf tensor;
  and \(p_4 = a_Y.\mathrm{map}\) of the tensor \(\otimes_{\mathrm{m}}\) of the two
  \(\mathtt{pullbackValIso}\) comparisons (one for \(M\), one for \(N\)) identifying the
  sheafified pullback value with the pullback.

  The outer three factors are isomorphisms unconditionally: \(p_1\) and \(p_3\) are the
  \(.\mathrm{hom}\) components of isomorphisms, and \(p_4\) is the image under the additive
  functor \(a_Y.\mathrm{map}\) of a tensor product of the two \(\mathtt{pullbackValIso}\)
  \emph{isomorphisms}, hence itself an isomorphism. The hypothesis supplies the one
  conditional middle factor \(a_Y.\mathrm{map}\,\delta\). A composite of isomorphisms is an
  isomorphism, so \(\mathtt{pullbackTensorMap}\,f\,M\,N\) is an isomorphism.
\end{proof}

\begin{lemma}

\leanok
[Composition coherence of the pushforward unit comparison]
  \label{lem:unitToPushforwardObjUnit_comp}
  \lean{AlgebraicGeometry.Scheme.Modules.unitToPushforwardObjUnit_comp}
  Let \(h : Z \to Y\) and \(f : Y \to X\) be composable morphisms of schemes. Write
  \(\mathrm{upu}\) for the pushforward (right-adjoint) unit comparison
  \(\mathtt{SheafOfModules.unitToPushforwardObjUnit}\), the canonical map
  \(\mathcal{O} \to f_*(\mathcal{O})\) comparing the structure sheaf of the target with
  the pushforward of the structure sheaf of the source. Then \(\mathrm{upu}\) satisfies
  the composition coherence
  \[
    \mathrm{upu}(h \mathbin{;} f)
      \;=\;
    \mathrm{upu}\,f \;\mathbin{;}\;
      (\mathtt{pushforward}\,f).\mathrm{map}\bigl(\mathrm{upu}\,h\bigr) \;\mathbin{;}\;
      (\mathtt{pushforwardComp}\,h\,f).\mathrm{hom},
  \]
  where \(\mathtt{pushforwardComp}\,h\,f : (h \mathbin{;} f)_* \cong f_* \circ h_*\) is the
  pseudofunctoriality isomorphism of the pushforward. Sectionwise both sides are nothing
  but the functoriality of the structure-sheaf ring maps, so the identity holds on the
  nose after the comparison of sections. This is the right-adjoint half of the unit
  composition coherence, and it is the input consumed by the mate transport that produces
  the pullback-side coherence \cref{lem:pullbackObjUnitToUnit_comp}.
\end{lemma}

\begin{proof}



  Both sides are morphisms of \(\mathcal{O}\)-modules; it suffices to compare them on
  sections over each open. Each comparison \(\mathrm{upu}\) acts on sections as the
  structure-sheaf ring homomorphism induced by the relevant morphism of schemes, and the
  pseudofunctor coherence \(\mathtt{pushforwardComp}\) acts as the identification of
  \((h \mathbin{;} f)\)-pushforward sections with iterated \(f_*\,h_*\)-pushforward
  sections. The composition coherence is therefore the functoriality identity
  \((h \mathbin{;} f)^\sharp = h^\sharp \circ f^\sharp\) of the structure-sheaf ring maps,
  which holds definitionally; the two sides agree section by section.
\end{proof}

\begin{lemma}

\leanok
[Composition coherence of the pullback unit comparison]
  \label{lem:pullbackObjUnitToUnit_comp}
  \lean{AlgebraicGeometry.Scheme.Modules.pullbackObjUnitToUnit_comp}
  \uses{lem:unitToPushforwardObjUnit_comp}
  Let \(h : Z \to Y\) and \(f : Y \to X\) be composable morphisms of schemes. Write
  \(\mathrm{pbu}\) for the pullback (left-adjoint) unit comparison
  \(\mathtt{SheafOfModules.pullbackObjUnitToUnit}\), the canonical map
  \(f^*\mathcal{O} \to \mathcal{O}\). Then \(\mathrm{pbu}\) satisfies the composition
  coherence
  \[
    \mathrm{pbu}(h \mathbin{;} f)
      \;=\;
    (\mathtt{pullbackComp}\,h\,f).\mathrm{inv} \;\mathbin{;}\;
      (\mathtt{pullback}\,h).\mathrm{map}\bigl(\mathrm{pbu}\,f\bigr) \;\mathbin{;}\;
      \mathrm{pbu}\,h,
  \]
  where \(\mathtt{pullbackComp}\,h\,f : (h \mathbin{;} f)^* \cong h^* \circ f^*\) is the
  pseudofunctoriality isomorphism of the pullback. This is the genuinely-new ingredient
  for \cref{lem:pullback_unit_iso}: the left-adjoint pullback is defined abstractly, with
  no sectionwise value, so this identity is \emph{not} a sectionwise statement and cannot
  be checked on sections.
\end{lemma}

\begin{proof}
  \uses{lem:unitToPushforwardObjUnit_comp}
  The pullback \(f^*\) is the left adjoint of the pushforward \(f_*\), and likewise for
  \(h\) and \(h \mathbin{;} f\); the comparison \(\mathrm{pbu}\,f\) is the adjunction
  \emph{mate} of the pushforward comparison \(\mathrm{upu}\,f\)
  (\cref{lem:unitToPushforwardObjUnit_comp}) across the pullback--pushforward adjunction,
  and similarly for \(h\) and \(h \mathbin{;} f\). Since the mate correspondence is a
  bijection, the asserted identity follows once both sides are transposed under the
  adjunction \((h \mathbin{;} f)^* \dashv (h \mathbin{;} f)_*\).

  Under this transpose the left-hand side \(\mathrm{pbu}(h \mathbin{;} f)\) becomes the
  pushforward comparison \(\mathrm{upu}(h \mathbin{;} f)\). The right-hand side, a
  composite involving \((\mathtt{pullbackComp}\,h\,f).\mathrm{inv}\) and the mates of
  \(\mathrm{pbu}\,f\) and \(\mathrm{pbu}\,h\), transposes --- using the conjugation of the
  pseudofunctoriality isomorphisms (\(\mathtt{pullbackComp}\,h\,f\) on the left adjoints
  corresponds to \(\mathtt{pushforwardComp}\,h\,f\) on the right adjoints) and the unit of
  the composite adjunction \((f^* \dashv f_*) \circ (h^* \dashv h_*)\) --- exactly into the
  right-hand side of the pushforward coherence
  \cref{lem:unitToPushforwardObjUnit_comp}, namely
  \(\mathrm{upu}\,f \mathbin{;} f_*(\mathrm{upu}\,h) \mathbin{;}
  (\mathtt{pushforwardComp}\,h\,f).\mathrm{hom}\). Thus the two transposed sides coincide
  by \cref{lem:unitToPushforwardObjUnit_comp}, and the original identity follows by the
  injectivity of the transpose.
\end{proof}

\begin{lemma}

\leanok
[Pullback preserves the structure sheaf]
  \label{lem:pullback_unit_iso}
  \lean{AlgebraicGeometry.Scheme.Modules.pullbackUnitIso}
  \uses{def:scheme_modules_tensorobj, def:pullback_obj_unit_to_unit_iso}
  Let \(f : Y \to X\) be a morphism of schemes. The pullback of the structure sheaf
  is the structure sheaf: there is a canonical isomorphism
  \[
    f^*\bigl(\mathtt{SheafOfModules.unit}\,X.\mathtt{ringCatSheaf}\bigr)
      \;\xrightarrow{\;\sim\;}\;
    \mathtt{SheafOfModules.unit}\,Y.\mathtt{ringCatSheaf},
    \qquad\text{i.e.}\qquad
    f^* \mathcal{O}_X \;\cong\; \mathcal{O}_Y,
  \]
  where \(\mathtt{SheafOfModules.unit}\,X.\mathtt{ringCatSheaf} = \mathcal{O}_X\) is
  the monoidal unit of \(\Scheme.\mathtt{Modules}\,X\).
\end{lemma}

\begin{proof}
  \uses{def:scheme_modules_tensorobj, def:pullback_obj_unit_to_unit_iso}
  This is a one-step consequence of a Mathlib instance, with no chart-chase. The
  canonical comparison map is the genuine Mathlib morphism
  \[
    \mathtt{SheafOfModules.pullbackObjUnitToUnit}\,f :
      f^*\mathcal{O}_X \;\longrightarrow\; \mathcal{O}_Y,
  \]
  the unit comparison \(\eta\) of \(f^*\); abbreviate it \(\mathrm{pbu}\,f\). Mathlib
  proves this map is an isomorphism whenever the comparison functor
  \(\mathtt{Opens.map}\,f.\mathtt{base}\) on opens is \(\mathtt{Final}\): the instance
  \(\mathtt{SheafOfModules.instIsIsoPullbackObjUnitToUnitOfFinal}\).

  The point is that the finality hypothesis holds \emph{unconditionally}, for every
  morphism of schemes \(f\). The preimage functor \(\mathtt{Opens.map}\,f.\mathtt{base}
  : \mathtt{Opens}\,X \to \mathtt{Opens}\,Y\) is a frame homomorphism: it preserves
  arbitrary unions and finite intersections, hence preserves all finite limits. A
  functor between small categories that preserves finite limits is representably flat,
  so \(\mathtt{final\_of\_representablyFlat}\) gives \((\mathtt{Opens.map}\,f.\mathtt{base}).\mathtt{Final}\)
  with no hypothesis on \(f\) whatsoever. (There is no need to pass to affine charts:
  finality is not a special feature of open immersions but a general property of
  preimage-on-opens.)

  Consequently the instance \(\mathtt{instIsIsoPullbackObjUnitToUnitOfFinal}\) fires
  directly, and \(\mathrm{pbu}\,f\) is an isomorphism for every \(f\). The asserted
  isomorphism \(f^*\mathcal{O}_X \cong \mathcal{O}_Y\) is the bundled iso
  \(\mathtt{pullbackUnitIso}\,f\) of this comparison.
\end{proof}

\begin{lemma}

\leanok
[Naturality of the sheaf-level pullback--tensor comparison $\delta_{\mathrm{sheaf}}$ (D1$'$)]
  \label{lem:pullback_tensor_map_natural}
  \lean{AlgebraicGeometry.Scheme.Modules.pullbackTensorMap_natural}
  \uses{lem:pullback_tensor_map, lem:presheaf_pullback_oplaxmonoidal,
        lem:sheafify_tensor_unit_iso_natural, lem:pullback_val_iso_natural}
  Let \(f : Y \to X\) be a morphism of schemes. The sheaf-level comparison map
  \(\delta_{\mathrm{sheaf}} = \mathtt{pullbackTensorMap}\) of
  \cref{lem:pullback_tensor_map} is \emph{natural} in both module arguments: for
  morphisms \(a : M \to M'\) and \(b : N \to N'\) in \(\Scheme.\mathtt{Modules}\,X\)
  the square
  \[
    \begin{array}{ccc}
      f^*(M \otimes_X N)
        & \xrightarrow{\;\delta_{\mathrm{sheaf}}\;}
        & f^*M \otimes_Y f^*N \\[2pt]
      \big\downarrow{\scriptstyle f^*(a \otimes b)}
        & & \big\downarrow{\scriptstyle f^*a \,\otimes\, f^*b} \\[4pt]
      f^*(M' \otimes_X N')
        & \xrightarrow{\;\delta_{\mathrm{sheaf}}\;}
        & f^*M' \otimes_Y f^*N'
    \end{array}
  \]
  commutes.
\end{lemma}

\begin{proof}
  \uses{lem:pullback_tensor_map, lem:presheaf_pullback_oplaxmonoidal,
        lem:sheafify_tensor_unit_iso_natural, lem:pullback_val_iso_natural}
  The map \(\delta_{\mathrm{sheaf}}\) is built (\cref{lem:pullback_tensor_map}) by
  transporting the abstract oplax comparison \(\delta\) of the presheaf pullback
  (\cref{lem:presheaf_pullback_oplaxmonoidal}) through the sheafification
  reconciliations. The asserted square is the pasting of four naturality squares, one for
  each ingredient of this composite; we record them in order, the only delicate one being the
  naturality of the oplax comparison itself.

  \emph{Square 2: naturality of the oplax comparison \(\delta\).} The comparison \(\delta\) is
  the oplax-monoidal comparison of the presheaf-level pullback functor
  \(F = \mathtt{PresheafOfModules.pullback}\,\varphi'\), whose defining ring map
  \(\varphi' = (\mathtt{toRingCatSheafHom}\,f).\mathrm{hom}\) has, as the source of its domain,
  the ring presheaf of \(X\). For the monoidal structure on \(\mathtt{PresheafOfModules}\) to
  be available on this domain the ring presheaf must be presented in its \emph{canonical}
  form as the composite
  \(\mathcal{O}_X \circ \mathtt{forget}_2\,\mathtt{CommRingCat}\,\mathtt{RingCat}\) --- the
  presentation on which that monoidal structure is defined. When the comparison map's
  construction is unfolded, however, the domain ring re-presents in a form definitionally
  equal, but not explicitly identical, to the canonical one, and on that form the
  monoidal structure is not directly available. The naturality of \(\delta\),
  \[
    \delta \circ F(a \otimes b) \;=\; (Fa \otimes Fb) \circ \delta,
  \]
  which holds for any oplax monoidal functor, is therefore applied by re-establishing the
  canonical domain-ring spelling at the functor argument itself: one ascribes the defining
  ring map \(\varphi'\) explicitly to the canonical composite-functor source, so that the
  domain monoidal structure resolves, and the naturality identity then matches the goal
  definitionally. This is a structural device --- it changes only the presentation of
  \(F\)'s defining ring map. No restatement of \(\mathtt{pullbackTensorMap}\) or of its helper
  isomorphisms is needed, and the unit-case lemma \(\mathtt{pullbackTensorMap\_unit\_isIso}\)
  (D2\(^\prime\)) is unaffected. After this step the comparison \(\delta\) commutes past
  \(F(a \otimes b)\), producing the factor \(Fa \otimes Fb\).

  \emph{Squares 1, 3, 4 and bifunctoriality.} The remaining ingredients of
  \(\delta_{\mathrm{sheaf}}\) are honest natural transformations and contribute already-closed
  naturality squares: the sheafification, the tensor/unit reconciliation
  \(\mathtt{sheafifyTensorUnitIso}\) (Square 3, \cref{lem:sheafify_tensor_unit_iso_natural}),
  and the \(\mathtt{pullbackValIso}\) reconciliation (Square 4,
  \cref{lem:pullback_val_iso_natural}). Finally, the tensor of the two component morphisms
  factors as \(f^*a \otimes f^*b\) by the bifunctoriality of \(\otimes_Y\) --- the interchange
  law \((f_1 \otimes f_2) \circ (g_1 \otimes g_2) = (f_1 \circ g_1) \otimes (f_2 \circ g_2)\),
  here \(\mathtt{tensorHom\_comp\_tensorHom}\). Pasting the four squares together with this
  bifunctoriality identity yields the asserted commuting square.
\end{proof}

\begin{lemma}

\leanok
  [Section-level naturality of the tensor/unit reconciliation
   \(\mathtt{sheafifyTensorUnitIso}\)]
  \label{lem:sheafify_tensor_unit_iso_natural}
  \lean{AlgebraicGeometry.Scheme.Modules.sheafifyTensorUnitIso_hom_natural}
  \uses{lem:pullback_tensor_map, def:sheafify_tensor_unit_iso,
        lem:sheafify_tensor_unit_iso_hom_eq_prime, lem:restrictscalarsid_map}
  The tensor/unit reconciliation isomorphism \(\mathtt{sheafifyTensorUnitIso}\) entering the
  fourth naturality square of \cref{lem:pullback_tensor_map_natural} is natural in both
  module arguments. Equivalently: after descending to sections over each open \(U\) and
  passing to the underlying module components, the naturality square is the
  bilinear, sectionwise naturality of the sheafification unit \(\eta\), verified one pure
  tensor \(p \otimes q\) at a time, using bilinearity.
\end{lemma}

\begin{proof}
  \uses{lem:pullback_tensor_map}
  Write \(\Phi\) for the natural transformation whose naturality square (in the two module
  arguments \(M, N\), along morphisms \(a : M \to M'\) and \(b : N \to N'\)) is to be shown
  to commute; \(\Phi\) is assembled from the reconciliation isomorphism
  \(\mathtt{sheafifyTensorUnitIso}\) (the intermediate construction inside
  \cref{lem:pullback_tensor_map}), so its components are morphisms of presheaves of modules.
  The square is closed at the categorical (tensor-of-morphisms) level, without any descent to
  sections or to elements.

  \emph{Step 1: a single tensor-of-morphisms presentation.} The \(\mathrm{hom}\) leg of
  \(\mathtt{sheafifyTensorUnitIso}\) is rewritten as a \emph{single} tensor of morphisms of the
  form \(a.\mathrm{map}(\eta_P \otimes \eta_Q)\), where
  \(\eta = (\mathtt{sheafificationAdjunction}).\mathrm{unit}\) is the unit
  \(\mathbf{1} \Rightarrow (\text{sheafify}) \circ (\text{inclusion})\) of the sheafification
  adjunction and \(a\) is the sheafification functor. The essential point is that all factors
  are kept inside \emph{one} monoidal structure on \(\mathtt{PresheafOfModules}\) --- the
  canonical presentation as a composite with the forgetful functor \(\mathtt{forget}_2\) ---
  so that the bifunctoriality law of \(\otimes\) can be applied consistently within one
  monoidal structure.

  \emph{Step 2: merge the two functor-image factors.} Both composites around the square are
  images, under the sheafification functor, of a tensor of morphisms. The two functor-image
  factors are merged into a single image of a composite by functoriality,
  \(F(g) \circ F(h) = F(g \circ h)\).

  \emph{Step 3: bifunctoriality and unit naturality.} The resulting equality of two tensors of
  morphisms is closed by the bifunctoriality (interchange) law
  \[
    (f_1 \otimes f_2) \circ (g_1 \otimes g_2) \;=\; (f_1 \circ g_1) \otimes (f_2 \circ g_2),
  \]
  applied within the single chosen monoidal structure, together with the two
  single-component naturality squares of the sheafification unit \(\eta\), one in each module
  argument:
  \[
    p \circ \eta_{M'} = \eta_{M} \circ (f^*a),
    \qquad
    q \circ \eta_{N'} = \eta_{N} \circ (f^*b).
  \]
  Each of these holds because \(\eta\) is a natural transformation. Combining the two
  single-component squares through the interchange law equates the two tensors of morphisms,
  which is exactly the naturality square of \(\mathtt{sheafifyTensorUnitIso}\).
\end{proof}

\begin{lemma}

\leanok
  [Naturality of the \(\mathtt{pullbackValIso}\) comparison]
  \label{lem:pullback_val_iso_natural}
  \lean{AlgebraicGeometry.Scheme.Modules.pullbackValIso_hom_natural}
  \uses{def:pullback_val_iso}
  The \(\mathtt{pullbackValIso}\) comparison --- the isomorphism reconciling the underlying
  presheaf of a pullback sheaf of modules with the pullback of its underlying presheaf ---
  is natural in its module argument. This is one of the (closed) input squares of the
  pasting that proves \cref{lem:pullback_tensor_map_natural}.
\end{lemma}

\begin{lemma}

\leanok
[The comparison on the unit pair $(\mathcal{O},\mathcal{O})$ is an isomorphism (D2$'$)]
  \label{lem:pullback_tensor_iso_unit}
  \lean{AlgebraicGeometry.Scheme.Modules.pullbackTensorMap_unit_isIso}
  \uses{lem:pullback_tensor_map, lem:pullback_unit_iso, lem:isiso_pullbacktensormap_of_sheafifydelta}
  Let \(f : Y \to X\) be a morphism of schemes. On the monoidal unit pair
  \((\mathcal{O}_X, \mathcal{O}_X)\) the sheaf-level comparison
  \[
    \delta_{(\mathcal{O},\mathcal{O})} :
      f^*(\mathcal{O}_X \otimes_X \mathcal{O}_X)
      \;\longrightarrow\;
      f^*\mathcal{O}_X \otimes_Y f^*\mathcal{O}_X
  \]
  of \cref{lem:pullback_tensor_map} is an isomorphism.
\end{lemma}

\begin{proof}
  \uses{lem:pullback_tensor_map, lem:isiso_pullbacktensormap_of_sheafifydelta, lem:isiso_sheafifyeta_of_unitsquare, lem:eta_bridge_unit_square}
  By \cref{lem:isiso_pullbacktensormap_of_sheafifydelta} applied on the unit pair,
  it suffices to verify that the sheafified presheaf comparison
  \(a_Y.\mathrm{map}\,(\delta\,(\mathtt{pullback}\,\varphi')\,\mathbf{1}\,\mathbf{1})\)
  is an isomorphism; the underlying presheaf of the structure-sheaf unit coincides with
  the presheaf monoidal unit \(\mathbf{1}\).

  On the unit pair the presheaf-level oplax comparison \(\delta\) is governed by Mathlib's
  oplax-monoidal \emph{unit coherence}: \(\mathtt{Functor.OplaxMonoidal.left\_unitality\_hom}\)
  expresses \(\delta\,(\mathtt{pullback}\,\varphi')\,\mathbf{1}\,\mathbf{1}\) as a composite of
  the functorial unitor \(F.\mathrm{map}\,(\lambda_{\mathbf{1}}).\mathrm{hom}\), the target
  unitor \(\lambda_{F.\mathrm{obj}\,\mathbf{1}}\), and the whiskered presheaf unit comparison
  \(\eta\,(\mathtt{pullback}\,\varphi') \rhd F.\mathrm{obj}\,\mathbf{1}\). The unitors are
  isomorphisms, so after sheafification by \(a_Y\) the genuine remaining content is
  \[
    \mathtt{IsIso}\bigl(a_Y.\mathrm{map}\,(\eta\,(\mathtt{pullback}\,\varphi'))\bigr),
  \]
  the sheafified presheaf unit comparison (the \(\delta\)-wrapping reduction). By \cref{lem:isiso_sheafifyeta_of_unitsquare} it suffices, in turn, to prove the
  \emph{unit square}
  \[
    (\mathtt{pullbackValIso}\,f\,\mathcal{O}_X).\mathrm{inv} \mathbin{;}
      a_Y.\mathrm{map}\,(\eta\,(\mathtt{pullback}\,\varphi')) \mathbin{;}
      \mathtt{sheafifyUnitIso}.\mathrm{hom}
      \;=\; \mathtt{pullbackObjUnitToUnit}\,\varphi,
  \]
  which is \cref{lem:eta_bridge_unit_square}. Granting that square,
  \cref{lem:isiso_sheafifyeta_of_unitsquare} yields
  \(\mathtt{IsIso}\bigl(a_Y.\mathrm{map}\,(\eta\,(\mathtt{pullback}\,\varphi'))\bigr)\) --- the
  right-hand side \(\mathtt{pullbackObjUnitToUnit}\,\varphi = \mathtt{pullbackUnitIso}\,f :
  f^*\mathcal{O}_X \xrightarrow{\sim} \mathcal{O}_Y\) being an isomorphism for \emph{every}
  \(f\) (\cref{lem:pullback_unit_iso}) --- and the \(\delta\)-wrapping reduction above then
  delivers that \(\delta_{(\mathcal{O},\mathcal{O})}\) is an isomorphism.
\end{proof}

\subsection{The unit square (D2$'$): a mate-calculus telescope}

The proof of \cref{lem:pullback_tensor_iso_unit} bottoms out in the single morphism equation
labelled the \emph{unit square} below. Its proof is a telescope across three nested adjunction
layers (sheafification, presheaf pullback--pushforward, sheaf pullback--pushforward). This
subsection isolates the load-bearing intermediate identities as named lemmas so each can be
discharged independently, then assembles them.

We fix throughout a morphism of schemes \(f : Y \to X\) and write
\(\varphi = f.\mathtt{toRingCatSheafHom}\), \(\varphi' = \varphi.\mathrm{hom}\). Let \(a_X, a_Y\)
denote sheafification of presheaves of modules on \(X\), \(Y\); \(\mathbf{1}^{\mathrm{p}}\) the
presheaf monoidal unit; \(F = \mathtt{pullback}\,\varphi'\) the presheaf pullback;
\(\eta\,F : F.\mathrm{obj}\,\mathbf{1}^{\mathrm{p}} \to \mathbf{1}^{\mathrm{p}}\) its oplax unit
comparison; and \(\varepsilon(G)\) the lax unit comparison of a lax monoidal functor \(G\). The
isomorphisms \(\mathtt{pullbackValIso}\,f\,\mathcal{O}_X\) and \(\mathtt{sheafifyUnitIso}\) are
the file's comparison morphisms relating the sheaf-level pullback value to the sheafification of
the presheaf-level pullback. Write
\(\mathtt{sheafCounit}_\bullet = (\mathtt{sheafificationAdjunction}\,(\mathbf{1}_{\bullet.\mathtt{ringCatSheaf}.\mathrm{val}})).\mathrm{counit}\).

\begin{lemma}\leanok
[Presheaf-side unit mate identity]
  \label{lem:presheaf_unit_comp_map_eta}
  \lean{AlgebraicGeometry.Scheme.Modules.presheafUnit_comp_map_eta}
  \uses{lem:presheaf_pushforward_laxmonoidal, lem:presheaf_pullback_oplaxmonoidal}
  At the presheaf level the pullback--pushforward adjunction unit at the monoidal unit,
  post-composed with the pushforward of the oplax unit comparison, recovers the lax unit:
  \[
    \mathtt{presheafAdj}.\mathrm{unit}.\mathrm{app}\,\mathbf{1}^{\mathrm{p}}
      \;\mathbin{;}\;
      (\mathtt{pushforward}\,\varphi').\mathrm{map}\,(\eta\,F)
      \;=\; \varepsilon(\mathtt{pushforward}\,\varphi'),
  \]
  where \(\mathtt{presheafAdj} = \mathtt{pullbackPushforwardAdjunction}\,\varphi'\) is the
  presheaf pullback--pushforward adjunction and \(F = \mathtt{pullback}\,\varphi'\).
\end{lemma}

\begin{proof}
  \uses{lem:presheaf_pushforward_laxmonoidal, lem:presheaf_pullback_oplaxmonoidal}
  This is Mathlib's general adjunction-mate identity
  \(\mathtt{Adjunction.unit\_app\_unit\_comp\_map\_}\eta\) instantiated at the presheaf
  pullback--pushforward adjunction \(\mathtt{pullbackPushforwardAdjunction}\,\varphi'\). That
  identity holds for any adjunction whose left adjoint is oplax monoidal and whose right adjoint
  is lax monoidal in a mate-compatible way. The compatible pair is
  \(\mathtt{presheafPushforwardLaxMonoidal}\)
  (\cref{lem:presheaf_pushforward_laxmonoidal}) and \(\mathtt{presheafPullbackOplaxMonoidal}\)
  (\cref{lem:presheaf_pullback_oplaxmonoidal}), so the
  mate identity applies to the concrete adjunction and yields the stated equation.
\end{proof}

\begin{lemma}\leanok
[IsIso plumbing from the unit square]
  \label{lem:isiso_sheafifyeta_of_unitsquare}
  \lean{AlgebraicGeometry.Scheme.Modules.isIso_sheafifyEta_of_unitSquare}
  \uses{lem:pullback_unit_iso}
  Suppose the \emph{unit square}
  \[
    (\mathtt{pullbackValIso}\,f\,\mathcal{O}_X).\mathrm{inv} \;\mathbin{;}\;
      a_Y.\mathrm{map}\,(\eta\,F) \;\mathbin{;}\;
      \mathtt{sheafifyUnitIso}.\mathrm{hom}
      \;=\; \mathtt{pullbackObjUnitToUnit}\,\varphi
  \]
  commutes. Then \(a_Y.\mathrm{map}\,(\eta\,F)\) is an isomorphism.
\end{lemma}

\begin{proof}
  \uses{lem:pullback_unit_iso}
  The right-hand side \(\mathtt{pullbackObjUnitToUnit}\,\varphi : f^*\mathcal{O}_X \to
  \mathcal{O}_Y\) is an isomorphism for \emph{every} \(f\): it is \(\mathtt{pullbackUnitIso}\,f\)
  (\cref{lem:pullback_unit_iso}), the iso-ness coming from the finality of
  \(\mathtt{Opens.map}\,f.\mathrm{base}\). The morphisms \(\mathtt{pullbackValIso}\,f\,\mathcal{O}_X\)
  and \(\mathtt{sheafifyUnitIso}\) are isomorphisms by construction. Solving the square for the
  middle factor exhibits
  \[
    a_Y.\mathrm{map}\,(\eta\,F)
      \;=\;
      (\mathtt{pullbackValIso}\,f\,\mathcal{O}_X).\mathrm{hom} \;\mathbin{;}\;
      \mathtt{pullbackObjUnitToUnit}\,\varphi \;\mathbin{;}\;
      \mathtt{sheafifyUnitIso}.\mathrm{hom}^{-1},
  \]
  a composite of three isomorphisms, hence an isomorphism.
\end{proof}

\begin{lemma}\leanok
[The unit square]
  \label{lem:eta_bridge_unit_square}
  \lean{AlgebraicGeometry.Scheme.Modules.pullbackEtaUnitSquare}
  \uses{lem:presheaf_unit_comp_map_eta, lem:comp_homequiv_factor_sheafify_pullback,
        lem:leftadjointuniq_app_unit_eta, lem:epsilon_presheaf_to_sheaf_unit,
        lem:pullbackObjUnitToUnit_comp, def:sheafify_unit_iso, def:pullback_val_iso,
        lem:pullback_sheafify_unit_eta_triangle, lem:restrictscalarsid_map}
  The unit square commutes:
  \[
    (\mathtt{pullbackValIso}\,f\,\mathcal{O}_X).\mathrm{inv} \;\mathbin{;}\;
      a_Y.\mathrm{map}\,(\eta\,F) \;\mathbin{;}\;
      \mathtt{sheafifyUnitIso}.\mathrm{hom}
      \;=\; \mathtt{pullbackObjUnitToUnit}\,\varphi.
  \]
\end{lemma}

\begin{proof}
  \uses{lem:presheaf_unit_comp_map_eta, lem:comp_homequiv_factor_sheafify_pullback,
        lem:leftadjointuniq_app_unit_eta, lem:epsilon_presheaf_to_sheaf_unit}
  Transpose the square across the \emph{sheaf} pullback--pushforward adjunction
  \(\mathtt{pullbackPushforwardAdjunction}\,\varphi\) by applying
  \(\mathtt{homEquiv}.\mathrm{injective}\) and then \(\mathtt{homEquiv\_unit}\),
  \(\mathtt{homEquiv\_counit}\). Since
  \(\mathtt{homEquiv}(\mathtt{pullbackObjUnitToUnit}\,\varphi) =
  \mathtt{unitToPushforwardObjUnit}\,\varphi\) (the mate identity for the structure-sheaf unit
  comparison), the square reduces to the concrete pushforward-side identity
  \begin{equation}
    \mathtt{sheafAdj}.\mathrm{unit}.\mathrm{app}\,\mathcal{O}_X \;\mathbin{;}\;
      (\mathtt{pushforward}\,\varphi).\mathrm{map}\,(m)
      \;=\; \mathtt{unitToPushforwardObjUnit}\,\varphi,
    \tag{$\ast\ast$}
  \end{equation}
  where \(m\) is the four-fold composite
  \[
    m \;=\;
      (\mathtt{pullback}\,\varphi).\mathrm{map}\,(\mathtt{sheafCounit}_X.\mathrm{inv})
      \;\mathbin{;}\;
      (\mathtt{sheafificationCompPullback}\,\varphi).\mathrm{hom}.\mathrm{app}\,\mathbf{1}^{\mathrm{p}}
      \;\mathbin{;}\;
      a_Y.\mathrm{map}\,(\eta\,F)
      \;\mathbin{;}\;
      \mathtt{sheafCounit}_Y.\mathrm{hom},
  \]
  \(\mathtt{sheafAdj} = \mathtt{pullbackPushforwardAdjunction}\,\varphi\), and
  \(\mathtt{sheafificationCompPullback}\,\varphi = \mathtt{Adjunction.leftAdjointUniq}\,A\,B\) is
  the canonical comparison of the two composite adjunctions
  \begin{itemize}
    \item \(A\) with left adjoint \(a_X \circ \mathtt{pullback}_{\mathrm{sheaf}}\,\varphi\),
      namely \(A = (\mathtt{sheafificationAdjunction}\,(\mathbf{1}_X)).\mathrm{comp}\,
      (\mathtt{pullbackPushforwardAdjunction}_{\mathrm{sheaf}}\,\varphi)\);
    \item \(B\) with left adjoint \(F \circ a_Y\), namely
      \(B = (\mathtt{pullbackPushforwardAdjunction}\,\varphi').\mathrm{comp}\,
      (\mathtt{sheafificationAdjunction}\,(\mathbf{1}_Y))\).
  \end{itemize}
  We prove $(\ast\ast)$ in seven steps.

  \begin{enumerate}
    \item \textbf{Distribute.} By functoriality of \(\mathtt{pushforward}\,\varphi\)
      (\(\mathtt{Functor.map\_comp}\)), expand
      \((\mathtt{pushforward}\,\varphi).\mathrm{map}\,(m)\) over the four factors of \(m\), so
      that the left side of $(\ast\ast)$ becomes
      \(\mathtt{sheafAdj}.\mathrm{unit}.\mathrm{app}\,\mathcal{O}_X\) followed by the four
      pushed-forward factors.
    \item \textbf{Unit naturality.} Pull the leading factor
      \((\mathtt{pushforward}\,\varphi).\mathrm{map}\bigl((\mathtt{pullback}\,\varphi).\mathrm{map}\,(\mathtt{sheafCounit}_X.\mathrm{inv})\bigr)\)
      past the unit using the naturality square \(\mathtt{Adjunction.unit\_naturality}\) of
      \(\mathtt{sheafAdj}\). This rewrites the head as
      \[
        \mathtt{sheafCounit}_X.\mathrm{inv} \;\mathbin{;}\;
          \mathtt{sheafAdj}.\mathrm{unit}.\mathrm{app}\,(a_X.\mathrm{obj}\,\mathbf{1}^{\mathrm{p}})
          \;\mathbin{;}\;
          (\mathtt{pushforward}\,\varphi).\mathrm{map}\bigl(
            (\mathtt{sheafificationCompPullback}\,\varphi).\mathrm{hom}.\mathrm{app}\,\mathbf{1}^{\mathrm{p}}\bigr)
          \;\mathbin{;}\;\cdots,
      \]
      leaving the genuinely adjunction-theoretic head
      \(\mathtt{sheafAdj}.\mathrm{unit}.\mathrm{app}\,(a_X.\mathrm{obj}\,\mathbf{1}^{\mathrm{p}})
      \mathbin{;} (\mathtt{pushforward}\,\varphi).\mathrm{map}\,(g)\) with
      \(g = (\mathtt{sheafificationCompPullback}\,\varphi).\mathrm{hom}.\mathrm{app}\,\mathbf{1}^{\mathrm{p}}\).
    \item \textbf{Composite-adjunction homEquiv factorisation}
      (\cref{lem:comp_homequiv_factor_sheafify_pullback}). The head of step~2 is
      \((\mathtt{sheafAdj}_X.\mathrm{homEquiv})^{-1}\bigl(A.\mathrm{homEquiv}\,g\bigr)\),
      where \(A.\mathrm{homEquiv}\) is the hom-set bijection of the composite adjunction \(A\).
    \item \textbf{leftAdjointUniq transport}
      (\cref{lem:leftadjointuniq_app_unit_eta}). Applying
      \(\mathtt{Adjunction.homEquiv\_leftAdjointUniq\_hom\_app}\,A\,B\,\mathbf{1}^{\mathrm{p}}\)
      rewrites \(A.\mathrm{homEquiv}\,g = B.\mathrm{unit}.\mathrm{app}\,\mathbf{1}^{\mathrm{p}}\),
      and expanding \(B.\mathrm{unit}\) by \(\mathtt{Adjunction.comp\_unit\_app}\) on \(B\) gives
      \[
        B.\mathrm{unit}.\mathrm{app}\,\mathbf{1}^{\mathrm{p}}
          \;=\;
          \mathtt{presheafAdj}.\mathrm{unit}.\mathrm{app}\,\mathbf{1}^{\mathrm{p}}
          \;\mathbin{;}\;
          (\mathtt{pushforward}\,\varphi').\mathrm{map}\bigl(\mathtt{toSheafify}_Y(F.\mathrm{obj}\,\mathbf{1}^{\mathrm{p}})\bigr).
      \]
    \item \textbf{Composite-adjunction unit expansion.} The expansion of \(B.\mathrm{unit}\) used in step~4
      follows from \(\mathtt{Adjunction.comp\_unit\_app}\) applied to the composite
      \(B = (\mathtt{pullbackPushforwardAdjunction}\,\varphi').\mathrm{comp}\,
      (\mathtt{sheafificationAdjunction}\,(\mathbf{1}_Y))\): the unit of a composite adjunction
      is the pasting of the two factor units.
    \item \textbf{Presheaf-side mate identity}
      (\cref{lem:presheaf_unit_comp_map_eta}). The oplax unit comparison \(\eta\,F\) enters
      through the factor \(a_Y.\mathrm{map}\,(\eta\,F)\) of \(m\); the surrounding
      sheafification counits and the \(\mathtt{toSheafify}_Y\) of step~4 cancel through
      \(\mathtt{sheafCounit}_Y\), exposing the presheaf composite
      \(\mathtt{presheafAdj}.\mathrm{unit}.\mathrm{app}\,\mathbf{1}^{\mathrm{p}} \mathbin{;}
      (\mathtt{pushforward}\,\varphi').\mathrm{map}\,(\eta\,F)\). By
      \cref{lem:presheaf_unit_comp_map_eta} this equals the presheaf-level lax unit
      \(\varepsilon_{\mathrm{pre}} = \varepsilon(\mathtt{pushforward}\,\varphi')\).
    \item \textbf{$\varepsilon$ reconciliation}
      (\cref{lem:epsilon_presheaf_to_sheaf_unit}). Finally identify the presheaf-level
      \(\varepsilon_{\mathrm{pre}}\) with the sheaf-level structure-sheaf comparison
      \(\mathtt{unitToPushforwardObjUnit}\,\varphi\) at the \((-).\mathrm{val}\)
      (underlying-presheaf) level --- both act on sections as \(\varphi.\mathrm{hom}.\mathrm{app}\,X\)
      --- using the pushforward \(\leftrightarrow\) sheafification compatibility that defines the two
      adjunctions. This completes $(\ast\ast)$, hence the unit square.
  \end{enumerate}
\end{proof}

\begin{lemma}\leanok
[Composite-adjunction homEquiv factorisation]
  \label{lem:comp_homequiv_factor_sheafify_pullback}
  \lean{AlgebraicGeometry.Scheme.Modules.compHomEquivFactor}
  \uses{lem:presheaf_pushforward_laxmonoidal}
  Let \(\mathtt{adj}_1 : L_1 \dashv R_1\) and \(\mathtt{adj}_2 : L_2 \dashv R_2\) be composable
  adjunctions and \(A = \mathtt{adj}_1.\mathrm{comp}\,\mathtt{adj}_2 : L_1 \circ L_2 \dashv
  R_2 \circ R_1\) their composite. Then the hom-set bijection of the composite factors as the
  composite of the two factor bijections: for every \(g\),
  \[
    A.\mathrm{homEquiv}\,g
      \;=\;
    \mathtt{adj}_1.\mathrm{homEquiv}\bigl(\mathtt{adj}_2.\mathrm{homEquiv}\,g\bigr).
  \]
  Consequently, for the composite adjunction \(A\) of the unit square (left adjoint
  \(a_X \circ \mathtt{pullback}_{\mathrm{sheaf}}\,\varphi\)), the head of step~2 of
  \cref{lem:eta_bridge_unit_square} equals
  \((\mathtt{sheafAdj}_X.\mathrm{homEquiv})^{-1}\bigl(A.\mathrm{homEquiv}\,g\bigr)\) with
  \(g = (\mathtt{sheafificationCompPullback}\,\varphi).\mathrm{hom}.\mathrm{app}\,\mathbf{1}^{\mathrm{p}}\).
\end{lemma}

\begin{proof}
  \uses{lem:presheaf_pushforward_laxmonoidal}
  The composite adjunction \(\mathtt{Adjunction.comp}\) is defined so that its unit is the
  pasting \(\mathtt{adj}_1.\mathrm{unit} \mathbin{;} R_1(\mathtt{adj}_2.\mathrm{unit})_{L_1(-)}\)
  and its counit dually; equivalently its hom-set bijection \(A.\mathrm{homEquiv}\) is the
  composite of the two factor bijections in the order \(\mathtt{adj}_2.\mathrm{homEquiv}\) then
  \(\mathtt{adj}_1.\mathrm{homEquiv}\). This is the standard factorisation
  \(\mathtt{Adjunction.comp\_homEquiv}\) (the naturality of \(\mathtt{homEquiv}\) under
  composition of adjunctions): a morphism \(L_1 L_2\,c \to d\) is transposed first across
  \(\mathtt{adj}_1\) to \(L_2\,c \to R_1\,d\) and then across \(\mathtt{adj}_2\) to
  \(c \to R_2 R_1\,d\). Applying this with \(\mathtt{adj}_1 =
  \mathtt{sheafificationAdjunction}\,(\mathbf{1}_X)\), \(\mathtt{adj}_2 =
  \mathtt{pullbackPushforwardAdjunction}_{\mathrm{sheaf}}\,\varphi\) and transposing back along
  \(\mathtt{adj}_1\) (i.e.\ applying \((\mathtt{sheafAdj}_X.\mathrm{homEquiv})^{-1}\)) gives the
  stated rewrite of the head of step~2.
\end{proof}

\begin{lemma}

\leanok
[The sheafify-then-pullback comparison is the left-adjoint-uniqueness isomorphism]
  \label{lem:sheafification_comp_pullback_eq_leftadjointuniq}
  \lean{AlgebraicGeometry.Scheme.Modules.sheafificationCompPullback_eq_leftAdjointUniq}
  Write \(A\) for the composite adjunction obtained by following the sheafification adjunction
  on \(X\) with the sheaf-level pullback--pushforward adjunction of \(\varphi\), namely
  \(A = (\mathtt{sheafificationAdjunction}\,(\mathbf{1}_X)).\mathrm{comp}\,
  (\mathtt{pullbackPushforwardAdjunction}_{\mathrm{sheaf}}\,\varphi)\), whose left adjoint is
  \(a_X \circ \mathtt{pullback}_{\mathrm{sheaf}}\,\varphi\); and write \(B\) for the composite
  adjunction obtained by following the presheaf-level pullback--pushforward adjunction of
  \(\varphi'\) with the sheafification adjunction on \(Y\), namely
  \(B = (\mathtt{pullbackPushforwardAdjunction}\,\varphi').\mathrm{comp}\,
  (\mathtt{sheafificationAdjunction}\,(\mathbf{1}_Y))\), whose left adjoint is \(F \circ a_Y\).
  These two left adjoints are canonically isomorphic, and the canonical ``sheafify-then-pullback''
  comparison agrees on the nose with the left-adjoint-uniqueness isomorphism:
  \[
    \mathtt{sheafificationCompPullback}\,\varphi
      \;=\;
    \mathtt{Adjunction.leftAdjointUniq}\,A\,B.
  \]
\end{lemma}

\begin{proof}


  The two composite adjunctions \(A\) and \(B\) share the same right adjoint:
  pushing forward along \(\varphi\) and then forgetting to presheaves of modules on \(X\) is, on
  the nose, the same as forgetting to presheaves on \(Y\) and then pushing forward along
  \(\varphi'\) --- the pushforward commutes with the forgetful functor to presheaves by
  definition. Hence \(A\) and \(B\) are two adjunctions sharing the same right adjoint, so the
  left-adjoint-uniqueness isomorphism \(\mathtt{Adjunction.leftAdjointUniq}\,A\,B\) between their
  left adjoints is defined. Unfolding the canonical comparison
  \(\mathtt{sheafificationCompPullback}\,\varphi\) and unfolding
  \(\mathtt{Adjunction.leftAdjointUniq}\,A\,B\) produces the same morphism --- both sides are the
  composite isomorphism assembled from the two factor adjunctions in the same order. The equality
  therefore holds by reflexivity.
\end{proof}

\begin{lemma}\leanok
[leftAdjointUniq transport of the composite unit]
  \label{lem:leftadjointuniq_app_unit_eta}
  \lean{AlgebraicGeometry.Scheme.Modules.leftAdjointUniqUnitEta}
  \uses{lem:comp_homequiv_factor_sheafify_pullback, lem:sheafification_comp_pullback_eq_leftadjointuniq}
  Let \(A\), \(B\) be the two composite adjunctions of \cref{lem:eta_bridge_unit_square} (with
  left adjoints \(a_X \circ \mathtt{pullback}_{\mathrm{sheaf}}\,\varphi\) and \(F \circ a_Y\)
  respectively, which agree up to the canonical iso
  \(\mathtt{sheafificationCompPullback}\,\varphi = \mathtt{Adjunction.leftAdjointUniq}\,A\,B\)),
  and let \(g = (\mathtt{sheafificationCompPullback}\,\varphi).\mathrm{hom}.\mathrm{app}\,\mathbf{1}^{\mathrm{p}}\).
  Then
  \[
    A.\mathrm{homEquiv}\,g
      \;=\;
    B.\mathrm{unit}.\mathrm{app}\,\mathbf{1}^{\mathrm{p}}
      \;=\;
    \mathtt{presheafAdj}.\mathrm{unit}.\mathrm{app}\,\mathbf{1}^{\mathrm{p}}
      \;\mathbin{;}\;
    (\mathtt{pushforward}\,\varphi').\mathrm{map}\bigl(\mathtt{toSheafify}_Y(F.\mathrm{obj}\,\mathbf{1}^{\mathrm{p}})\bigr).
  \]
\end{lemma}

\begin{proof}
  \uses{lem:comp_homequiv_factor_sheafify_pullback, lem:sheafification_comp_pullback_eq_leftadjointuniq}
  The comparison \(\mathtt{sheafificationCompPullback}\,\varphi\) is by definition
  \(\mathtt{Adjunction.leftAdjointUniq}\,A\,B\)
  (\cref{lem:sheafification_comp_pullback_eq_leftadjointuniq}), the canonical isomorphism between two left
  adjoints of the same functor. Mathlib's
  \(\mathtt{Adjunction.homEquiv\_leftAdjointUniq\_hom\_app}\,A\,B\) states that applying
  \(A.\mathrm{homEquiv}\) to the component of this comparison at an object \(c\) returns
  \(B.\mathrm{unit}.\mathrm{app}\,c\); taking \(c = \mathbf{1}^{\mathrm{p}}\) gives the first
  equality. For the second, expand \(B.\mathrm{unit}\) by
  \(\mathtt{Adjunction.comp\_unit\_app}\) on \(B = (\mathtt{pullbackPushforwardAdjunction}\,
  \varphi').\mathrm{comp}\,(\mathtt{sheafificationAdjunction}\,(\mathbf{1}_Y))\): the unit of a
  composite adjunction at \(\mathbf{1}^{\mathrm{p}}\) is the presheaf unit
  \(\mathtt{presheafAdj}.\mathrm{unit}.\mathrm{app}\,\mathbf{1}^{\mathrm{p}}\) followed by the
  pushforward of the sheafification unit \(\mathtt{toSheafify}_Y\) at
  \(F.\mathrm{obj}\,\mathbf{1}^{\mathrm{p}}\), which is the displayed right-hand side.
\end{proof}

\begin{lemma}\leanok
[leftAdjointUniq transport of the composite unit ($P$-general form)]
  \label{lem:leftadjointuniq_app_unit_eta_general}
  \lean{AlgebraicGeometry.Scheme.Modules.leftAdjointUniqUnitEta_app}
  \uses{lem:leftadjointuniq_app_unit_eta, lem:sheafification_comp_pullback_eq_leftadjointuniq}
  This is the \(P\)-general twin of \cref{lem:leftadjointuniq_app_unit_eta}: the latter is the
  specialisation to the monoidal unit \(P = \mathbf{1}^{\mathrm{p}}\), whereas the present form
  holds for an \emph{arbitrary} presheaf \(P\) over \(X\). Fix a morphism of schemes
  \(\varphi\) (with structure-ring map \(\varphi'\)), and let
  \[
    A_{\varphi}
      \;=\;
    (\mathtt{PresheafPullbackPushforwardAdj}\,\varphi').\mathrm{comp}\,\mathtt{sheafAdj}
  \]
  be the two-layer composite adjunction (presheaf pullback--pushforward composed with the
  sheafification adjunction), and let \(B_{\varphi}\) denote its sheaf-level counterpart (the
  composite adjunction whose left adjoint is the sheaf-level pullback). Writing
  \(\mathtt{sheafCompPb} = \mathtt{SheafOfModules.sheafificationCompPullback}\) for the
  canonical comparison, the lemma asserts
  \[
    A_{\varphi}.\mathrm{homEquiv}\,P\,\_\,\bigl((\mathtt{sheafCompPb}\,\varphi).\mathrm{hom}.\mathrm{app}\,P\bigr)
      \;=\;
    B_{\varphi}.\mathrm{unit}.\mathrm{app}\,P.
  \]
  It is the key \textbf{step-1} brick of the tail assembly of
  \cref{lem:pullback_tensor_map_basechange}: applying \(A_{\varphi}.\mathrm{homEquiv}\) to the
  comparison component recovers the composite-adjunction unit at \(P\). Its proof is identical
  to that of \cref{lem:leftadjointuniq_app_unit_eta} with \(\mathbf{1}^{\mathrm{p}}\) replaced
  by \(P\): \(\mathtt{sheafCompPb}\,\varphi\) is by definition
  \(\mathtt{Adjunction.leftAdjointUniq}\,A_{\varphi}\,B_{\varphi}\)
  (\cref{lem:sheafification_comp_pullback_eq_leftadjointuniq}), so Mathlib's
  \(\mathtt{Adjunction.homEquiv\_leftAdjointUniq\_hom\_app}\) at \(P\) returns
  \(B_{\varphi}.\mathrm{unit}.\mathrm{app}\,P\).
\end{lemma}

\begin{lemma}\leanok
[$\varepsilon$ reconciliation: presheaf unit to sheaf unit on underlying presheaves]
  \label{lem:epsilon_presheaf_to_sheaf_unit}
  \lean{AlgebraicGeometry.Scheme.Modules.epsilonPresheafToSheafUnit}
  \uses{lem:presheaf_pushforward_laxmonoidal, lem:unitToPushforwardObjUnit_comp}
  The presheaf-level lax-monoidal unit comparison
  \(\varepsilon(\mathtt{pushforward}\,\varphi')\) and the sheaf-level structure-sheaf unit
  comparison \(\mathtt{unitToPushforwardObjUnit}\,\varphi\) agree as morphisms of
  \emph{underlying presheaves} of modules: after applying the forgetful functor
  \((-).\mathrm{val}\) that sends a sheaf of modules to its underlying presheaf, their
  components coincide. Concretely, on every section both maps act as the ring map
  \(\varphi.\mathrm{hom}.\mathrm{app}\) underlying \(\varphi\):
  \[
    (\mathtt{unitToPushforwardObjUnit}\,\varphi).\mathrm{val}.\mathrm{app}\,X\,a
      \;=\; \varphi.\mathrm{hom}.\mathrm{app}\,X\,a,
  \]
  and the presheaf unit comparison \(\varepsilon(\mathtt{pushforward}\,\varphi')\) takes the
  same value on sections. The asserted identity is the equality of these two presheaf-level
  morphisms (the \((-).\mathrm{val}\) images), \emph{not} a sheaf-level lax-monoidal-unit
  equation.

  The exact form of the identity is the \((-).\mathrm{val}\)-level (underlying-presheaf)
  reconciliation, because the sheaf-level pushforward carries no lax monoidal structure in this
  setup: the presheaf lax unit
  \(\varepsilon(\mathtt{pushforward}\,\varphi')\) is the one that exists, and
  it is reconciled with the sheaf-level \(\mathtt{unitToPushforwardObjUnit}\,\varphi\) through
  the sheafification \(\leftrightarrow\) pushforward compatibility.
\end{lemma}

\begin{proof}
  \uses{lem:presheaf_pushforward_laxmonoidal, lem:unitToPushforwardObjUnit_comp}
  Equality of presheaf morphisms is checked sectionwise. Over each open both presheaf-level
  maps are the single ring map \(\varphi.\mathrm{hom}.\mathrm{app}\,X\): the presheaf lax unit
  of a pushforward (\cref{lem:presheaf_pushforward_laxmonoidal}) is the structure-map
  comparison \(\mathcal{O} \to \varphi_*\mathcal{O}\), and the underlying presheaf component of
  \(\mathtt{unitToPushforwardObjUnit}\,\varphi\) is the same comparison by construction --- its
  value on a section is \(\varphi.\mathrm{hom}.\mathrm{app}\,X\) (the same comparison whose
  composition coherence is \cref{lem:unitToPushforwardObjUnit_comp}). Since the two presheaf
  morphisms agree on every section, they are equal. Tracing the sheafification reconciliation
  backwards then identifies this presheaf-level identity with the step~6 presheaf head
  \(\varepsilon(\mathtt{pushforward}\,\varphi')\) of \cref{lem:eta_bridge_unit_square}, closing
  $(\ast\ast)$.
\end{proof}

% NOTE: Composition coherence of the pullback tensorator \(\delta\) for composable
% morphisms \(h : Z \to Y\), \(f : Y \to X\). The open-immersion base-change square
% of \cref{lem:pullback_tensor_iso_loctriv} is the specialisation \(h := j'\).
% \(\mathtt{pullbackTensorMap}\) is not an adjunction transpose (it is a four-fold
% composite), so the mate-calculus argument for the unit coherence does not apply
% directly; the genuine proof pastes four squares directly.
\begin{lemma}

\leanok
[Composition coherence of the pullback tensorator $\delta$ (D3$'$)]
  \label{lem:pullback_tensor_map_basechange}
  \lean{AlgebraicGeometry.Scheme.Modules.pullbackTensorMap_restrict}
  \uses{lem:pullback_tensor_map, lem:pullback_tensor_map_natural,
        lem:presheaf_pullback_oplaxmonoidal, def:tensorobjisoofiso,
        lem:toringcatsheafhom_comp_hom_reconcile, lem:pullbackcomp_delta,
        lem:sheafificationcomppullback_comp, def:pullback_val_iso,
        lem:pullbackValIso_comp_coherence,
        def:sheafify_tensor_unit_iso}
  Let \(h : Z \to Y\) and \(f : Y \to X\) be composable morphisms of schemes and let
  \(M, N \in \Scheme.\mathtt{Modules}\,X\). Write \(\delta_{\mathrm{sheaf}} =
  \mathtt{pullbackTensorMap}\) for the sheaf-level pullback--tensor comparison of
  \cref{lem:pullback_tensor_map}, \(h^{*}(-) = (\mathtt{pullback}\,h).\mathrm{map}\) for
  the pullback functor, and \(\mathtt{pullbackComp}\,h\,f : h^{*}\circ f^{*} \xrightarrow{\sim} (h \circ f)^{*}\)
  for the pullback pseudofunctoriality isomorphism. Then the
  comparison for the composite \(h \circ f\) factors through the comparisons for \(f\)
  and for \(h\), conjugated by \(\mathtt{pullbackComp}\,h\,f\):
  \[
    \delta_{\mathrm{sheaf}}^{\,h \circ f}(M, N)
      \;=\;
      (\mathtt{pullbackComp}\,h\,f).\mathrm{inv}_{M \otimes_X N} \;\mathbin{;}\;
      h^{*}\!\bigl(\delta_{\mathrm{sheaf}}^{\,f}(M, N)\bigr) \;\mathbin{;}\;
      \delta_{\mathrm{sheaf}}^{\,h}(f^{*}M,\, f^{*}N) \;\mathbin{;}\;
      \mathtt{tensorObjIsoOfIso}\bigl((\mathtt{pullbackComp}\,h\,f)_M,\,
        (\mathtt{pullbackComp}\,h\,f)_N\bigr).\mathrm{hom},
  \]
  an equality of morphisms \((h \circ f)^{*}(M \otimes_X N) \to (h^{*}f^{*}M)
  \otimes_Z (h^{*}f^{*}N)\) in \(\Scheme.\mathtt{Modules}\,Z\), where
  \(\mathtt{tensorObjIsoOfIso}\) is the functoriality of \(\otimes\) in a pair of
  isomorphisms.

  \emph{Remark (the open-immersion base-change square).} The base-change form consumed
  downstream by \cref{lem:pullback_tensor_iso_loctriv} is the \emph{specialisation} of
  this coherence. Given \(f : Y \to X\), an open \(U \subseteq X\) with preimage
  \(f^{-1}U\), the restriction \(g : f^{-1}U \to U\) and the open immersions
  \(j : U \hookrightarrow X\), \(j' : f^{-1}U \hookrightarrow Y\), the morphism
  \(f^{-1}U \to X\) admits the two equal factorisations \(j' \circ f = g \circ j\)
  (i.e. \(j' \mathbin{;} f = g \mathbin{;} j\)). Applying the displayed coherence to
  each factorisation and equating the two results identifies the restriction
  \((j')^{*}\bigl(\delta^{f}_{\mathrm{sheaf}}(M,N)\bigr)\) with
  \(\delta^{g}_{\mathrm{sheaf}}(j^{*}M, j^{*}N)\), through the pseudofunctoriality
  isomorphisms; in particular each is an isomorphism if and only if the other is. This
  is the form used in the chart-chase, where \(M, N\) are trivialised on \(U\).
\end{lemma}

\begin{proof}
  \uses{lem:pullback_tensor_map, lem:pullback_tensor_map_natural,
        lem:presheaf_pullback_oplaxmonoidal,
        lem:pullbackObjUnitToUnit_comp, def:tensorobjisoofiso,
        lem:toringcatsheafhom_comp_hom_reconcile, lem:pullbackcomp_delta,
        lem:sheafificationcomppullback_comp, def:pullback_val_iso,
        lem:pullbackValIso_comp_coherence, lem:cmp_leg,
        def:sheafify_tensor_unit_iso}
  Unfolding \(\delta_{\mathrm{sheaf}} = \mathtt{pullbackTensorMap}\)
  (\cref{lem:pullback_tensor_map}) on both sides exposes each as the same four-fold
  composite
  \[
    S_1 \;\mathbin{;}\; a.\mathrm{map}\,\delta \;\mathbin{;}\; S_3 \;\mathbin{;}\; S_4,
  \]
  where \(S_1\) is the sheafification--pullback comparison
  \(\mathtt{sheafificationCompPullback}\), \(a.\mathrm{map}\,\delta\) is the
  sheafification of the presheaf-level oplax comparison \(\delta\) of
  \cref{lem:presheaf_pullback_oplaxmonoidal}, \(S_3\) is the monoidal-unit
  sheafification isomorphism \(\mathtt{sheafifyTensorUnitIso}\), and \(S_4\) is the
  tensor of the connecting isomorphisms \(\mathtt{pullbackValIso}\). The identity to
  prove is therefore a paste of \emph{four commuting squares}, conjugated throughout by
  the pullback pseudofunctoriality isomorphism \(\mathtt{pullbackComp}\,h\,f\). Unlike
  the \emph{naturality} paste of \cref{lem:pullback_tensor_map_natural} (D1$'$), this is
  a \emph{composition}-coherence paste: the left-hand side is the composite-morphism
  (\(h \circ f\)) instance, while the right-hand side interleaves the \(f\)-instance
  (transported by \(h^{*}\)) with the \(h\)-instance (on the pulled-back modules
  \(f^{*}M, f^{*}N\)).

  We record \emph{why} the unit-comparison mirror does not apply. The unit coherence
  \cref{lem:pullbackObjUnitToUnit_comp} succeeds because
  \(\mathtt{pullbackObjUnitToUnit}\) is, by definition, an adjunction transpose, so its
  composition coherence is obtained by transposing across the pullback--pushforward
  adjunction. The tensorator \(\mathtt{pullbackTensorMap}\) is \emph{not} a transpose ---
  it is the four-fold composite above --- and there is no analogous
  \(\mathtt{homEquiv}\)-bridge, so the transpose-and-inject opening of the unit analog
  leaves an un-evaluable transpose of a concrete composite. The proof instead
  pastes the four squares directly.

  \emph{Sq2 (the \(\delta\) core).} The presheaf-level oplax comparison \(\delta\) of the
  \emph{composite} pullback decomposes by the Mathlib oplax-monoidal coherence
  \(\mathtt{Functor.OplaxMonoidal.comp\_}\delta\), once the composite presheaf pullback
  \(\mathtt{pullback}\,\varphi'_{h \circ f}\) is identified with the composite functor
  \(\mathtt{pullback}\,\varphi'_{f} \,\circ\, \mathtt{pullback}\,\varphi'_{h}\) through the
  Mathlib presheaf coherence \(\mathtt{PresheafOfModules.pullbackComp}\). This
  identification rests on the ring-map reconciliation
  \[
    (\mathtt{toRingCatSheafHom}\,(h \circ f)).\mathrm{hom}
      \;=\; \varphi'_{f} \;\mathbin{;}\;
      (\mathtt{Opens.map}\,f.\mathrm{base})^{\mathrm{op}}.\mathtt{whiskerLeft}\,\varphi'_{h},
  \]
  which is \emph{definitional}: although the two sides are presented in functor categories
  that nominally differ by the equality \(\mathtt{Opens.map}\,(h.\mathrm{base} \circ
  f.\mathrm{base}) = \mathtt{Opens.map}\,f.\mathrm{base} \,\circ\, \mathtt{Opens.map}\,
  h.\mathrm{base}\) (\(\mathtt{Opens.map\_comp}\)), this equality and the underlying
  scheme-composition data already hold up to definitional equality at default
  transparency, so the reconciliation is definitional
  (\(\mathtt{toRingCatSheafHom\_comp\_hom\_reconcile}\)). The genuine content of Sq2 is
  therefore not the reconciliation but the monoidality of \(\mathtt{pullbackComp}\) itself,
  isolated as Sq2b below.

  \emph{Sq2b (monoidality of \(\mathtt{pullbackComp}\) --- the genuinely new ingredient).}
  The decomposition in Sq2 reduces the \(\delta\)-core to a single structural fact, absent
  from Mathlib: the connecting isomorphism
  \[
    \mathtt{pullbackComp}\,\varphi'_{f}\,\varphi'_{h} \;:\;
      \mathtt{pullback}\,\varphi'_{f} \,\circ\, \mathtt{pullback}\,\varphi'_{h}
      \;\cong\; \mathtt{pullback}\,
        \bigl(\varphi'_{f} \;\mathbin{;}\;
          (\mathtt{Opens.map}\,f.\mathrm{base})^{\mathrm{op}}.\mathtt{whiskerLeft}\,
          \varphi'_{h}\bigr)
  \]
  is a \emph{monoidal} natural isomorphism between the oplax-monoidal functors
  \(\bigl(\mathtt{pullback}\,\varphi'_{f} \circ \mathtt{pullback}\,\varphi'_{h},\,
  \mathtt{comp\_}\delta\bigr)\) and \(\bigl(\mathtt{pullback}\,(\varphi'_{f}\mathbin{;}
  \cdots),\,\delta\bigr)\). Writing \(\delta\) for the oplax tensorator of the single
  composite pullback and \(\mathtt{comp\_}\delta = (\mathtt{pullback}\,\varphi'_{h})
  .\mathrm{map}\,\delta^{\varphi'_{f}} \mathbin{;} \delta^{\varphi'_{h}}\) for the
  composite tensorator of the two-step pullback, monoidality is the identity
  \[
    \delta_{M,N}
      \;=\; (\mathtt{pullbackComp}\,\varphi'_{f}\,\varphi'_{h}).\mathrm{inv}_{M\otimes N}
      \;\mathbin{;}\; (\mathtt{comp\_}\delta)_{M,N} \;\mathbin{;}\;
      \bigl((\mathtt{pullbackComp}\,\varphi'_{f}\,\varphi'_{h}).\mathrm{hom}_{M}
        \otimes (\mathtt{pullbackComp}\,\varphi'_{f}\,\varphi'_{h}).\mathrm{hom}_{N}\bigr).
  \]
  This is precisely the \(\delta\)-twin of the unit coherence
  \cref{lem:pullbackObjUnitToUnit_comp}, and it is proved by the \emph{same} mate-calculus
  argument with the unit comparison \(\eta\) replaced by the tensorator \(\delta\). The
  pivot is that, for the canonical oplax structure on a left adjoint
  (\(\mathtt{leftAdjointOplaxMonoidal}\), here
  \cref{lem:presheaf_pullback_oplaxmonoidal}), \(\delta\) is \emph{itself} an adjunction
  transpose: \(\delta_{M,N} = \mathtt{homEquiv}^{-1}\bigl((\eta_{M}\otimes\eta_{N})
  \mathbin{;} \mu_{M,N}\bigr)\), where \(\mu\) is the lax tensorator of the
  \(\mathtt{pushforward}\). Consequently the transpose-and-inject opening that closes the
  unit coherence ports verbatim: applying \(\mathtt{homEquiv}.\mathrm{injective}\) for the
  pullback--pushforward adjunction and transporting \(\mathtt{pullbackComp}\) to
  \(\mathtt{pushforwardComp}\) across the adjunctions via
  \(\mathtt{conjugateEquiv\_pullbackComp\_inv}\) (with \(\mathtt{unit\_conjugateEquiv}\)
  and \(\mathtt{Adjunction.comp\_unit\_app}\)) reduces the goal to the lax-\(\mu\)
  composition coherence of \(\mathtt{PresheafOfModules.pushforward}\) across
  \(\mathtt{pushforwardComp}\). This last residual --- isolated as the lemma
  \(\mathtt{pushforwardComp\_lax\_}\mu\) --- is the \emph{right-adjoint twin} of the
  \(\delta\)-coherence: it is precisely the assertion that the pseudofunctoriality
  isomorphism \(\mathtt{pushforwardComp}\) is \emph{monoidal} for the lax tensorators
  \(\mu\), i.e.\ that comparing the lax-\(\mu\) structure of the single composite
  \(\mathtt{pushforward}\) with the composite of the two individual lax-\(\mu\) structures
  agrees across \(\mathtt{pushforwardComp}\). Two points must be stated precisely.

  First, the mate-calculus reduction proceeds as follows: the transpose-and-inject
  opening, the \(\delta = \mathtt{homEquiv}^{-1}\bigl((\eta\otimes\eta)\mathbin{;}\mu\bigr)\)
  pivot, the \(\mathtt{conjugateEquiv\_pullbackComp\_inv}\) transport across the adjunctions
  (with \(\mathtt{unit\_conjugateEquiv}\) and \(\mathtt{Adjunction.comp\_unit\_app}\)), and
  the two-argument \(\mathtt{tensorHom}\)/\(\delta_{\mathrm{natural}}\) bookkeeping --- which
  follows the template of Mathlib's
  \(\mathtt{CategoryTheory.Adjunction.isMonoidal\_comp}\) (the canonical ``a composite of
  monoidal adjunctions is monoidal'') --- together reduce Sq2b to the single residual
  \(\mathtt{pushforwardComp\_lax\_}\mu\). This is the content of
  \cref{lem:pullbackObjUnitToUnit_comp}'s \(\delta\)-mirror.

  Second, this residual \(\mathtt{pushforwardComp\_lax\_}\mu\) is discharged by a short
  \emph{sectionwise pure-tensor collapse}, not by a change-of-rings calculation. On the
  strongly-monoidal pushforward objects the pushforward tensorator \(\mu\) reduces
  \emph{definitionally} to the lighter \(\mathtt{restrictScalars}\) tensorator \(\mu\) (the
  identity \(\mathtt{pushforward\_}\mu\mathtt{\_eq}\), which holds by reflexivity); after
  this reduction, evaluated on a pure tensor \(m \otimes n\) every
  \(\mathtt{restrictScalars}\) tensorator is the identity (Mathlib's
  \(\mathtt{ModuleCat.restrictScalars\_}\mu\mathtt{\_tmul}\)), so both sides of the
  interchange collapse to the same pure tensor \(m \otimes n\); no change-of-rings
  construction is required. This closes the Sq2b obligation.

  Crucially, this sub-step is stated at the \(\mathtt{PresheafOfModules}\) level --- about
  \(\mathtt{pullback}\,\varphi'\) with \(\varphi'\) the ring map
  \((\mathcal{S}_{0}\mathbin{;}\mathtt{forget}_{2}) \to f^{\mathrm{op}}\mathbin{;}
  (\mathcal{R}_{0}\mathbin{;}\mathtt{forget}_{2})\) carried by \(\mathtt{pullbackTensorMap}\)
  --- rather than at the scheme level, where the composite factorisation is immediate from
  \(\mathtt{pullbackComp}\) and the ring-map reconciliation is definitional.

  The observation above --- that \(\mathtt{pullbackTensorMap}\) is not an adjunction
  transpose --- concerns the \emph{full} four-fold composite and does not constrain
  Sq2b: the inner tensorator \(\delta\) genuinely \emph{is} a transpose
  (\(\mathtt{homEquiv}^{-1}\) of \((\eta\otimes\eta)\mathbin{;}\mu\)), so for Sq2b
  specifically the \(\eta\to\delta\) mirror of \cref{lem:pullbackObjUnitToUnit_comp} is the
  correct route.

  \emph{Sq1 (sheafification \(\leftrightarrow\) pullback) --- the sole open ingredient.} The
  composition coherence of \(\mathtt{SheafOfModules.sheafificationCompPullback}\) across
  \(h \circ f\) --- the analog of \(\mathtt{pullbackComp}\) for the
  ``sheafify-then-pullback'' natural isomorphism --- is \emph{absent from Mathlib} and is
  supplied as a named, private project sub-lemma
  \(\mathtt{sheafificationCompPullback\_comp}\). For a presheaf \(P\) over \(X\), writing
  \(a_X, a_Z\) for the sheafifications and \(\mathtt{sheafCompPb} =
  \mathtt{SheafOfModules.sheafificationCompPullback}\), it asserts
  \[
    ((\mathtt{sheafCompPb}\,(h \circ f)).\mathrm{app}\,P).\mathrm{hom}
      \;=\;
      (\mathtt{pullbackComp}\,h\,f).\mathrm{inv}_{a_X P} \;\mathbin{;}\;
      h^{*}\bigl(((\mathtt{sheafCompPb}\,f).\mathrm{app}\,P).\mathrm{hom}\bigr) \;\mathbin{;}\;
      ((\mathtt{sheafCompPb}\,h).\mathrm{app}\,
        (\mathtt{pullback}\,\varphi'_{f}\,P)).\mathrm{hom} \;\mathbin{;}\;
      a_Z.\mathrm{map}\bigl(
        (\mathtt{PresheafOfModules.pullbackComp}\,\varphi'_{f}\,\varphi'_{h}).\mathrm{hom}_P
      \bigr).
  \]
  It is proved by the \(\mathtt{leftAdjointUniq}\) mate calculus: both \(\mathtt{sheafCompPb}\)
  isomorphisms are \(\mathtt{leftAdjointUniq}\) of composite adjunctions, so transposing the
  identity under the composite adjunction for \(h \circ f\) and evaluating the left-hand side
  by the mate identity \(\mathtt{homEquiv\_leftAdjointUniq\_hom\_app}\) reduces it to a
  unit-level identity; the two \(\mathtt{pullbackComp}\) factors are then transported across
  the adjunction units --- the \(\delta\)-free twin of
  \cref{lem:pullbackObjUnitToUnit_comp}.

  \emph{The reduced tail goal.} The leading pseudofunctoriality factor
  \((\mathtt{pullbackComp}\,h\,f).\mathrm{inv}\) is killed \emph{inline} --- keeping
  \(\mathtt{pullbackComp}\) opaque via the conjugate transport
  \cref{lem:homEquiv_conjugateEquiv_app_prime} and the conjugate collapse
  \(\mathtt{conjugateEquiv\_whiskerLeft} + \mathtt{conjugateEquiv\_pullbackComp\_inv}\), rather
  than distributing the forgetful functor through it (that route is circular). The remaining
  obligation is a unit identity for the two-layer composite adjunction (presheaf
  pullback--pushforward adjunction composed with the sheafification adjunction), the
  \(\mathtt{mateEquiv\_vcomp}\) cocycle interior held open inside
  \cref{lem:sheafificationcomppullback_comp}. Writing \(B_{\varphi}\) for the
  sheaf-level composite adjunction of \cref{lem:leftadjointuniq_app_unit_eta_general} and
  \(\mathtt{PresheafPullback}_f\) for the presheaf-level pullback along \(f\), the assembly
  proceeds in the following ordered steps, mirroring the already-closed
  \cref{lem:pullbackObjUnitToUnit_comp} (whose proof is the \(\delta\)-free analog of this one,
  exactly one sheafification layer down):
  \begin{enumerate}
    \item[(a)] \emph{Strip the outer identity wrapper.} The tail goal carries an outer
      \(\mathtt{restrictScalars}(\mathbf{1})\) change-of-rings wrapper coming from the
      identity ring map; this restriction-of-scalars-along-identity factor is removed first,
      so that the genuine comparison factors are exposed.
    \item[(b)] \emph{Distribute the right-hand sheaf composite under \(\mathtt{forget}\).} The
      right-hand side is the sheaf-level composite carrying the two sub-comparison factors
      \(R_1\) (the \(f\)-layer) and \(R_5\) (the \(h\)-layer). Pushing the forgetful functor
      \(\mathtt{forget}\) (sheaf \(\to\) presheaf) through the composite separates these two
      factors \emph{without} disturbing the left-hand composite unit \(B_{h\circ f}.\mathrm{unit}\),
      which is left intact for the final reassembly.
    \item[(c)] \emph{Recover the sub-comparison units.} Apply
      \cref{lem:leftadjointuniq_app_unit_eta_general}
      (\(\mathtt{leftAdjointUniqUnitEta\_app}\), the \(P\)-general step-1 brick) to rewrite
      \(R_1\) as the composite-adjunction unit \(B_f.\mathrm{unit}.\mathrm{app}\,P\) and
      \(R_5\) as \(B_h.\mathrm{unit}.\mathrm{app}\,(\mathtt{PresheafPullback}_f\,P)\). This is
      precisely the uniqueness-of-left-adjoints characterisation
      \(\mathtt{homEquiv\_leftAdjointUniq\_hom\_app}\), now used in its \(P\)-general form.
    \item[(d)] \emph{Slide the \(\mathtt{pushforwardComp}\) comparison past the recovered
      units.} The sheaf-level comparison \(\mathtt{pushforwardComp}\,h\,f\) is moved across the
      two recovered units \(B_f.\mathrm{unit}\) and \(B_h.\mathrm{unit}\) by its naturality.
    \item[(e)] \emph{Reassemble the composite unit.} Expanding
      \(B_{h\circ f}.\mathrm{unit}\) by the composite-adjunction unit formula
      \(\mathtt{comp\_unit\_app}\) and applying \(\mathtt{unit\_naturality}\) identifies the
      slid expression of step~(d) with the left-hand composite unit, closing the tail.
  \end{enumerate}

  \emph{The precise binding obligation.} The one place where the assembly is not yet a
  mechanical paste is the bridge between the \emph{presheaf}-level left-hand data (the
  presheaf pushforward \(\mathtt{pushforward}^{\mathrm{pre}}\) and the composite unit
  \(B_{h\circ f}.\mathrm{unit}\)) and the \emph{sheaf}-level right-hand factors \(R_1, R_5\)
  wrapped in \(\mathtt{forget}\). Crossing this boundary requires the compatibility
  \[
    \mathtt{forget} \circ \mathtt{pushforward}^{\mathrm{sheaf}}
      \;=\;
    \mathtt{pushforward}^{\mathrm{pre}} \circ \mathtt{forget},
  \]
  interacting with the recovered \(B_f\)- and \(B_h\)-units of step~(c). This compatibility is
  the natural unit to isolate as its own named sub-lemma \emph{before} the assembly: once it
  is in hand, steps (a)--(e) become a mechanical paste, whereas without it the presheaf- and
  sheaf-level pushforwards cannot be aligned at the recovered units. The remainder of the tail
  is bookkeeping around this single compatibility.

  This is the sheafification-laden analog of \cref{lem:pullbackObjUnitToUnit_comp}; unlike
  that lemma it is not definitional sectionwise, because the adjunctions involved are
  composite rather than single transposes. The Sq1 lemma
  \(\mathtt{sheafificationCompPullback\_comp}\) (\cref{lem:sheafificationcomppullback_comp})
  carries this reduction \emph{inline}: it transposes the equation, kills the leading
  \((\mathtt{pullbackComp}\,h\,f).\mathrm{inv}\) factor opaquely via
  \cref{lem:homEquiv_conjugateEquiv_app_prime} and the conjugate collapse, and holds the
  residual unit identity --- the \(\mathtt{mateEquiv\_vcomp}\) cocycle interior --- as its
  single open obligation, with no separate tail sub-lemma.
  \emph{Warning:} transposing the \emph{whole} tail back under the composite adjunction --- as
  opposed to recovering the two sub-units individually via step~(c) and reassembling --- is
  \emph{circular} (verified): it returns the very identity one set out to prove, so the
  step-(c) recovery of the individual \(R_1, R_5\) units is essential.

  \emph{Sq3 (unit-iso transport).} The monoidal-unit sheafification isomorphism
  \(\mathtt{sheafifyTensorUnitIso}\) is carried through the same \(\mathtt{pullbackComp}\)
  identification used for Sq2.

  \emph{Sq4 (the connecting iso).} The scheme-level composition coherence of
  \(\mathtt{pullbackValIso}\) --- which produces the final
  \(\mathtt{tensorObjIsoOfIso}\bigl((\mathtt{pullbackComp}\,h\,f)_M,
  (\mathtt{pullbackComp}\,h\,f)_N\bigr)\) factor by tensoring the \(M\)- and \(N\)-instances
  --- is the \emph{standalone named obligation}
  \cref{lem:pullbackValIso_comp_coherence} (\(\mathtt{pullbackValIso\_comp}\)), absent from
  Mathlib and built separately as its own brick (its decomposition is
  \cref{lem:pullbackValIso_eq_sheafCompPb} and
  \cref{lem:sheafCompPb_counit_comp_coherence}). It is a short corollary of Sq1: since
  \(\mathtt{pullbackValIso}\) factors through \(\mathtt{sheafCompPb}\) (as
  \(\mathtt{sheafCompPb}^{-1}\) followed by the pullback of a counit), its composition
  coherence reduces to that of Sq1 (\cref{lem:sheafificationcomppullback_comp}) together with
  the pseudofunctoriality of the counit
  (\cref{lem:pullbackComp_hom_naturality}). It is consumed inside
  \(\mathtt{pullbackTensorMap\_restrict}\) as the connecting-object boundary where the abstract
  pullback's \(\mathtt{.val}\) presentation meets the helper-\(\mathtt{.obj}\) presentation,
  \emph{not} as a step inlined in the proof of Sq1.

  The proof assembles four squares. Sq1
  (\(\mathtt{sheafificationCompPullback\_comp}\)) is the primary reduction
  (\cref{lem:sheafificationcomppullback_comp}); the monoidality Sq2b
  (\(\mathtt{pullbackComp\_}\delta\) together with \(\mathtt{pushforwardComp\_lax\_}\mu\))
  follows from the mate-calculus reduction above, with the Sq2 ring-map reconciliation
  definitional; and Sq4 (\(\mathtt{pullbackValIso\_comp}\)) reduces to Sq1 via the
  \(\mathtt{pullbackValIso}\) factorisation (\cref{lem:pullbackValIso_comp_coherence}).
  The Sq3 unit-iso transport --- carrying the monoidal-unit
  sheafification isomorphism \(\mathtt{sheafifyTensorUnitIso}\) through \(\mathtt{pullbackComp}\)
  --- follows from the same \(\mathtt{pullbackComp}\) identification.
  The four squares are assembled via the \emph{interleaved paste} that
  fuses the composite-instance left-hand side with the \(f\)-instance/\(h\)-instance right-hand
  side, carried out inline in \(\mathtt{pullbackTensorMap\_restrict}\).
  The adjunction-mate bridge \(\mathtt{conjugateEquiv\_pullbackComp\_inv}\) (relating
  \(\mathtt{pullbackComp}\) for the left adjoints to \(\mathtt{pushforwardComp}\) for the
  right adjoints) is the device used in Sq2b and Sq4 to transport the pseudofunctoriality
  coherence across the adjunctions.

  \emph{The interleaved assembly (steps (1)--(7)).} The four squares do \emph{not} paste
  row-by-row. Because \(S_1^{h}\) acts on
  \(\bigl((\mathtt{pullback}\,f).\mathrm{obj}\,M\bigr) \otimes
  \bigl((\mathtt{pullback}\,f).\mathrm{obj}\,N\bigr)\), \emph{not} on
  \(\mathtt{pullback}\,\varphi'_{f}\,(M \otimes N)\), each square's output must first be slid past
  the next, by the oplax naturality \(\delta_{\mathrm{natural}}\) and the naturality of the
  intervening comparisons, before the four coherences line up --- exactly as in the D1$'$
  naturality paste of \cref{lem:pullback_tensor_map_natural}. Once \(\delta = \mathtt{comp\_}\delta\)
  has been decomposed (Sq2) and \(S_3 = \mathtt{sheafifyTensorUnitIso}\) re-expressed as
  \(a\bigl(\eta \otimes \eta\bigr)\) (\cref{def:sheafify_tensor_unit_iso}), every right-hand factor
  preceding the \(h\)-comparison \(S_1^{h}\) has the shape ``\(h^{*}\) of a sheafified presheaf map'',
  so the comparison \(S_1^{h} = (\mathtt{sheafCompPb}\,h)\) can be slid through it by its
  naturality. The assembly is then the following ordered chase. Write \(a = a_Z\) for the
  target sheafification, \(\mathtt{pb}^{\mathrm{pre}}h\) for the presheaf pullback along \(h\),
  and \(\eta\) for the sheafification-adjunction unit.
  \begin{enumerate}
    \item[(1)] \emph{Slide \(S_1^{h}\) past the \(\delta_f\)-leg.} By the naturality of
      \(\mathtt{sheafCompPb}\,h\) applied to \(\delta_f\), the leading
      \((\mathtt{sheafCompPb}\,h)\mathbin{;} a(\mathtt{pb}^{\mathrm{pre}}h.\mathrm{map}\,\delta_f)\)
      becomes \(h^{*}\!\bigl(a(\delta_f)\bigr) \mathbin{;} (\mathtt{sheafCompPb}\,h)\). The common
      leading factor \(h^{*}\!\bigl(a(\delta_f)\bigr)\) (which also opens the right-hand side as the
      \(h^{*}\)-image of the \(f\)-instance's \(\delta\)-core) is peeled off both sides.
    \item[(2)] \emph{Slide \(S_1^{h}\) past the \(S_3^{f}\)-leg (unit iso).} With \(S_3^{f}\) written
      as \(a(\eta \otimes \eta)\) over the \(f\)-pullbacks, the same naturality of
      \(\mathtt{sheafCompPb}\,h\) moves the right-hand \(S_1^{h}\) leftward past the
      \(h^{*}\)-image of \(S_3^{f}\).
    \item[(3)] \emph{Slide \(S_1^{h}\) past the \(S_4^{f}\)-leg (\(\mathtt{pullbackValIso}\)), then
      pass to a presheaf identity.} A third application of the naturality of
      \(\mathtt{sheafCompPb}\,h\), now on the \(f\)-instance's connecting factor
      \(\mathtt{forget}(\mathtt{pullbackValIso}\,f) \otimes \mathtt{forget}(\mathtt{pullbackValIso}\,f)\),
      lands \(S_1^{h}\) on \(\mathtt{pb}^{\mathrm{pre}}f\,M.\mathrm{val} \otimes
      \mathtt{pb}^{\mathrm{pre}}f\,N.\mathrm{val}\), matching the left-hand \(S_1^{h}\); peeling this
      common comparison leaves a goal lying \emph{entirely} under \(a\), which the functoriality of
      \(a\) collapses to a single presheaf-level identity \(a(\mathrm{BIG}_L) = a(\mathrm{BIG}_R)\),
      reduced to \(\mathrm{BIG}_L = \mathrm{BIG}_R\).
    \item[(4)] \emph{Push \(\delta_h\) across the first \(\mathtt{pb}^{\mathrm{pre}}h\)-leg.} On the
      presheaf identity, the \(h\)-instance tensorator \(\delta_h\) is moved left across the
      \(\mathtt{pb}^{\mathrm{pre}}h\)-image of the \(f\)-leg by the oplax naturality
      \(\mathtt{Functor.OplaxMonoidal.}\delta_{\mathrm{natural}}\).
    \item[(5)] \emph{Push \(\delta_h\) across the second \(\mathtt{pb}^{\mathrm{pre}}h\)-leg.} A second
      application of \(\delta_{\mathrm{natural}}\) brings \(\delta_h\) to a position common to both
      sides; the common \(\delta_h^{\,\mathtt{pb}^{\mathrm{pre}}f M.\mathrm{val},\,
      \mathtt{pb}^{\mathrm{pre}}f N.\mathrm{val}}\) is peeled.
    \item[(6)] \emph{Split into the \(M\)- and \(N\)-legs.} The bifunctoriality of \(\otimes\)
      (\(\mathtt{tensorHom\_comp\_tensorHom}\)) factors both remaining composites as a tensor of an
      \(M\)-component with an \(N\)-component, reducing the identity to the conjunction of the two
      per-leg identities.
    \item[(7)] \emph{Discharge each leg by \cref{lem:cmp_leg}.} The \(M\)-leg is exactly
      \(\mathtt{cmp\_leg}\,h\,f\,M\) and the \(N\)-leg is \(\mathtt{cmp\_leg}\,h\,f\,N\); applying the
      per-leg coherence to each closes the presheaf identity, hence the whole paste.
  \end{enumerate}
  With Sq1, Sq2b, Sq4 and the per-leg core \cref{lem:cmp_leg} in hand, this interleaved
  assembly --- together with the Sq3 unit-iso transport --- fuses the
  four squares into the displayed identity. The base-change-square consequence (the Remark) is then the
  specialisation \(h := j'\) obtained by equating the two factorisations \(j' \mathbin{;} f =
  g \mathbin{;} j\); the iff-clause is immediate, an isomorphism being preserved and reflected
  by the pseudofunctoriality isomorphism it is conjugated by.
\end{proof}

% ---------------------------------------------------------------------------
% The per-leg core of the step-(7) assembly: cmp_leg
%   (presheaf transpose of Sq1; reduces to the leftAdjointUniq core coherence C)
% ---------------------------------------------------------------------------

\begin{lemma}
[Per-leg composition coherence of the sheafification-unit / \(\mathtt{pullbackValIso}\)
 comparison (\(\mathtt{cmp\_leg}\))]
  \label{lem:cmp_leg}
  \lean{AlgebraicGeometry.Scheme.Modules.cmp_leg}
  \uses{lem:pullbackValIso_comp_coherence, lem:sheafificationcomppullback_comp,
        lem:pullbackComp_hom_naturality, def:pullback_val_iso}
  Let \(h : Z \to Y\), \(f : Y \to X\) be composable morphisms of schemes and let
  \(M \in \Scheme.\mathtt{Modules}\,X\). Write \(a_W\) for the sheafification over a scheme \(W\),
  \(\eta\) for the sheafification-adjunction unit (\(\eta_P : P \to a_W P\), at the scheme implied
  by the altitude), \(\mathtt{forget}\) for the forgetful functor
  \(\Scheme.\mathtt{Modules}\,W \to \mathtt{PresheafOfModules}\), \(\mathtt{pb}^{\mathrm{pre}}g\) for
  the presheaf pullback along \(g\), and, for an altitude \(g\) ending at the scheme carrying
  \(a_W\), the per-altitude comparison
  \[
    \mathrm{cmp}_{g}(M)
      \;:=\;
    \eta_{\mathtt{pb}^{\mathrm{pre}}g\,M.\mathrm{val}}
      \;\mathbin{;}\;
    \mathtt{forget}\bigl((\mathtt{pullbackValIso}\,g\,M).\mathrm{hom}\bigr).
  \]
  Then the composite-pullback comparison factors through the \(f\)- and \(h\)-instances,
  conjugated by the presheaf pseudofunctoriality iso
  \(\mathtt{pullbackComp}\,\varphi'_{f}\,\varphi'_{h}\) and the scheme pseudofunctoriality iso
  \(\mathtt{pullbackComp}\,h\,f\): the identity
  \[
    (\mathtt{pullbackComp}\,\varphi'_{f}\,\varphi'_{h}).\mathrm{hom}_{M.\mathrm{val}}
      \;\mathbin{;}\;
    \mathrm{cmp}_{h\circ f}(M)
    \;=\;
    \mathtt{pb}^{\mathrm{pre}}h.\mathrm{map}\bigl(\mathrm{cmp}_{f}(M)\bigr)
      \;\mathbin{;}\;
    \mathrm{cmp}_{h}(f^{*}M)
      \;\mathbin{;}\;
    \mathtt{forget}\bigl((\mathtt{pullbackComp}\,h\,f).\mathrm{hom}.\mathrm{app}\,M\bigr)
  \]
  holds, an equality of presheaf maps
  \(\mathtt{pb}^{\mathrm{pre}}h\bigl(\mathtt{pb}^{\mathrm{pre}}f\,M.\mathrm{val}\bigr)
    \to \mathtt{forget}\bigl((h\circ f)^{*}M\bigr)\). Equivalently, spelt out,
  \[
    (\mathtt{pullbackComp}\,\varphi'_{f}\,\varphi'_{h}).\mathrm{hom}_{M.\mathrm{val}}
      \mathbin{;}
    \eta_{\mathtt{pb}^{\mathrm{pre}}(h\circ f)\,M.\mathrm{val}}
      \mathbin{;}
    \mathtt{forget}\bigl((\mathtt{pullbackValIso}\,(h\circ f)\,M).\mathrm{hom}\bigr)
  \]
  equals
  \[
    \mathtt{pb}^{\mathrm{pre}}h.\mathrm{map}\Bigl(
      \eta_{\mathtt{pb}^{\mathrm{pre}}f\,M.\mathrm{val}} \mathbin{;}
      \mathtt{forget}\bigl((\mathtt{pullbackValIso}\,f\,M).\mathrm{hom}\bigr)\Bigr)
      \mathbin{;}
    \eta_{\mathtt{pb}^{\mathrm{pre}}h((f^{*}M).\mathrm{val})}
      \mathbin{;}
    \mathtt{forget}\bigl((\mathtt{pullbackValIso}\,h\,(f^{*}M)).\mathrm{hom}\bigr)
      \mathbin{;}
    \mathtt{forget}\bigl((\mathtt{pullbackComp}\,h\,f).\mathrm{hom}.\mathrm{app}\,M\bigr).
  \]
  This is precisely the \emph{presheaf transpose} of the sheaf-level Sq4 coherence
  \cref{lem:pullbackValIso_comp_coherence} (equivalently, of Sq1,
  \cref{lem:sheafificationcomppullback_comp}); the \(N\)-leg is the same identity with \(N\) in
  place of \(M\). It is the per-leg core consumed at step~(7) of the assembly of
  \cref{lem:pullback_tensor_map_basechange}.
\end{lemma}

\begin{proof}
  \uses{lem:pullbackValIso_comp_coherence, lem:sheafificationcomppullback_comp,
        lem:pullbackComp_hom_naturality, def:pullback_val_iso}
  \textit{Source: internal categorical construction; no external reference.}

  \emph{Verified prefix (reduction to the core coherence C).} The trailing factor
  \(\mathtt{forget}\bigl((\mathtt{pullbackComp}\,h\,f).\mathrm{hom}.\mathrm{app}\,M\bigr)\) is an
  isomorphism, hence a monomorphism; cancelling it by post-composing both sides with
  its inverse \(\mathtt{forget}\bigl((\mathtt{pullbackComp}\,h\,f).\mathrm{inv}.\mathrm{app}\,M\bigr)\). On the right the adjacent
  \(\mathtt{forget}(\mathtt{pullbackComp}\,h\,f).\mathrm{hom}\mathbin{;}
  \mathtt{forget}(\mathtt{pullbackComp}\,h\,f).\mathrm{inv}\) collapses to the identity by
  functoriality of \(\mathtt{forget}\) and \(\mathrm{hom}\mathbin{;}\mathrm{inv} = \mathbf 1\). On
  the left the grouped factor
  \(\mathtt{forget}\bigl((\mathtt{pullbackValIso}\,(h\circ f)\,M).\mathrm{hom}\bigr) \mathbin{;}
  \mathtt{forget}\bigl((\mathtt{pullbackComp}\,h\,f).\mathrm{inv}.\mathrm{app}\,M\bigr)\) is rewritten,
  again under \(\mathtt{forget}\), by the sheaf-level Sq4 composition coherence
  \cref{lem:pullbackValIso_comp_coherence} (\(\mathtt{pullbackValIso\_comp}\)):
  \[
    (\mathtt{pullbackValIso}\,(h\circ f)\,M).\mathrm{hom}
      \mathbin{;} (\mathtt{pullbackComp}\,h\,f).\mathrm{inv}_M
    \;=\;
    \Psi^{M}_{h,f} \mathbin{;} h^{*}\!\bigl((\mathtt{pullbackValIso}\,f\,M).\mathrm{hom}\bigr),
  \]
  whose left vertical \(\Psi^{M}_{h,f} = a_Z\bigl((\mathtt{pullbackComp}\,\varphi'_{f}
  \,\varphi'_{h}).\mathrm{inv}_{M.\mathrm{val}}\bigr) \mathbin{;}
  ((\mathtt{sheafCompPb}\,h).\mathrm{app}\,(\mathtt{pb}^{\mathrm{pre}}f\,M.\mathrm{val})).\mathrm{inv}\)
  exposes a leading \(\Psi\) whose first factor is the \(a_Z\)-image of the presheaf
  \((\mathtt{pullbackComp}\,\varphi'_{f}\,\varphi'_{h}).\mathrm{inv}\). The opening
  \((\mathtt{pullbackComp}\,\varphi'_{f}\,\varphi'_{h}).\mathrm{hom}\) on the left of the goal is then
  slid through the sheafification-unit naturality of \(\eta\)
  (\(\eta_{P} \mathbin{;} a_Z.\mathrm{map}\,k = k \mathbin{;} \eta_{P'}\) for a presheaf map
  \(k : P \to P'\)), bringing the presheaf \(\mathrm{hom}\) adjacent to the \(a_Z\)-image
  \(\mathrm{inv}\); their composite collapses to the identity (\(\mathrm{hom}\mathbin{;}\mathrm{inv}
  = \mathbf 1\)). This verified prefix reduces the per-leg identity to the single
  \textbf{core coherence C}:
  \[
    \eta_{\mathtt{pb}^{\mathrm{pre}}h(\mathtt{pb}^{\mathrm{pre}}f\,M.\mathrm{val})}
      \;\mathbin{;}\;
    \mathtt{forget}\bigl(((\mathtt{sheafCompPb}\,h).\mathrm{app}\,
      (\mathtt{pb}^{\mathrm{pre}}f\,M.\mathrm{val})).\mathrm{inv}\bigr)
      \;\mathbin{;}\;
    \mathtt{forget}\bigl(h^{*}\!\bigl((\mathtt{pullbackValIso}\,f\,M).\mathrm{hom}\bigr)\bigr)
  \]
  \[
    \;=\;
    \mathtt{pb}^{\mathrm{pre}}h.\mathrm{map}\bigl(\mathrm{cmp}_{f}(M)\bigr)
      \;\mathbin{;}\;
    \eta_{\mathtt{pb}^{\mathrm{pre}}h((f^{*}M).\mathrm{val})}
      \;\mathbin{;}\;
    \mathtt{forget}\bigl((\mathtt{pullbackValIso}\,h\,(f^{*}M)).\mathrm{hom}\bigr),
  \]
  where \(\mathtt{sheafCompPb}\,h = A_h.\mathtt{leftAdjointUniq}\,B_h\) is the
  uniqueness-of-left-adjoints comparison for the two composite adjunctions
  \(A_h = (\text{sheafificationAdj}_Z).\mathrm{comp}\,(\text{pullback--pushforwardAdj}\,\varphi_h)\)
  and \(B_h = (\text{pullback--pushforwardAdj}\,\varphi_h.\mathrm{hom}).\mathrm{comp}\,
  (\text{sheafificationAdj}_Y)\), both presenting the same right adjoint.

  \emph{Core coherence C.} This residual is a unit-level identity for the composite adjunctions,
  the \(\mathtt{leftAdjointUniq}\) mate calculus already used for Sq1
  (\cref{lem:sheafificationcomppullback_comp}). It is discharged by the two Mathlib mate identities
  \(\mathtt{Adjunction.unit\_leftAdjointUniq\_hom\_app}\) (which expresses the unit composed with the
  \(\mathtt{leftAdjointUniq}\) comparison as the unit of the second adjunction) and
  \(\mathtt{Adjunction.leftAdjointUniq\_hom\_app\_counit}\) (which evaluates the comparison against the
  counit). \emph{Hazard.} These two identities do \emph{not} apply directly to the goal as
  presented: the comparison \(\mathtt{sheafCompPb}\,h\) is built from the \emph{plain}
  sheafification unit \(\eta\), whereas the mate identities are stated against \(B_h\)'s
  composite-adjunction unit. One must \emph{first} re-express the plain \(\eta\) as the
  composite-adjunction unit of \(B_h\) via \(\mathtt{Adjunction.comp\_unit\_app}\) --- the
  syntactic-alignment step that the mate identities alone do not perform --- after which
  they collapse both sides to the same composite-adjunction unit, closing C. This reuses verbatim the Sq1 \(\mathtt{leftAdjointUniq}\) template of
  \cref{lem:sheafificationcomppullback_comp}, one sheafification layer down.
\end{proof}

% ---------------------------------------------------------------------------
% Sq4 sub-lemmas: composition coherence of pullbackValIso (the Sq4 chain)
%   - lem:pullbackComp_hom_naturality   (Mathlib: pullback pseudofunctoriality naturality)
%   - lem:pullbackValIso_eq_sheafCompPb (sheafCompPb/counit factorisation of pullbackValIso)
%   - lem:sheafCompPb_counit_comp_coherence (Sq1 corollary for the composite)
%   - lem:pullbackValIso_comp_coherence (Sq4 composition coherence of pullbackValIso)
% ---------------------------------------------------------------------------

\begin{lemma}
[Naturality of the pullback pseudofunctoriality comparison]
  \label{lem:pullbackComp_hom_naturality}
  \lean{CategoryTheory.NatTrans.naturality}
  \mathlibok
  \textit{Provided by Mathlib.}
  For composable \(h : Z \to Y\), \(f : Y \to X\), the pullback pseudofunctoriality
  isomorphism \(\mathtt{pullbackComp}\,h\,f : h^{*} \circ f^{*} \cong (h \circ f)^{*}\) is a
  \emph{natural} isomorphism of functors
  \(\Scheme.\mathtt{Modules}\,X \to \Scheme.\mathtt{Modules}\,Z\). Hence for every morphism
  \(g : A \to B\) in \(\Scheme.\mathtt{Modules}\,X\) its component square commutes. In the
  \(\mathrm{inv}\) orientation consumed by the Sq4 chain (with
  \((\mathtt{pullbackComp}\,h\,f).\mathrm{inv} : (h\circ f)^{*} \Rightarrow h^{*}\circ f^{*}\)),
  \[
    (\mathtt{pullback}\,(h \circ f)).\mathrm{map}\,g \;\mathbin{;}\;
      (\mathtt{pullbackComp}\,h\,f).\mathrm{inv}_B
      \;=\;
    (\mathtt{pullbackComp}\,h\,f).\mathrm{inv}_A \;\mathbin{;}\;
      (\mathtt{pullback}\,h).\mathrm{map}\bigl((\mathtt{pullback}\,f).\mathrm{map}\,g\bigr).
  \]
  This is the naturality square of the natural isomorphism
  \((\mathtt{pullbackComp}\,h\,f).\mathrm{inv}\) (the \(\mathrm{hom}\) orientation is the same
  square with the two horizontal arrows reversed).
\end{lemma}

\begin{lemma}
[\(\mathtt{pullbackValIso}\) as the \(\mathtt{sheafCompPb}\)/counit composite]
  \leanok
  \label{lem:pullbackValIso_eq_sheafCompPb}
  \lean{AlgebraicGeometry.Scheme.Modules.pullbackValIso_eq_sheafCompPb}
  \uses{def:pullback_val_iso, lem:sheafification_comp_pullback_eq_leftadjointuniq,
        def:sheafify_unit_iso}
  For \(f : Y \to X\) and \(M \in \Scheme.\mathtt{Modules}\,X\), the connecting isomorphism
  factors as the inverse sheafification--pullback comparison at the underlying presheaf,
  followed by the pullback of the sheafification counit at the already-sheaf \(M\):
  \[
    \mathtt{pullbackValIso}\,f\,M
      \;=\;
    (\mathtt{sheafificationCompPullback}\,f).\mathrm{symm}.\mathrm{app}\,(M.\mathrm{val})
      \;\mathbin{;}\;
    (\mathtt{pullback}\,f).\mathrm{mapIso}\,(\varepsilon^{X}_{M}),
  \]
  where \(\varepsilon^{X}_{M} : a_X(M.\mathrm{val}) \cong M\) is the sheafification-adjunction
  counit at \(M\) (an isomorphism since \(M\) is a sheaf). We abbreviate
  \(\mathtt{sheafCompPb}\,f := \mathtt{sheafificationCompPullback}\,f\).
\end{lemma}
\begin{proof}
  \leanok
  \uses{def:pullback_val_iso, lem:sheafification_comp_pullback_eq_leftadjointuniq}
  Immediate from the definition of \(\mathtt{pullbackValIso}\) (\cref{def:pullback_val_iso}):
  it \emph{is}, definitionally, this composite. The presheaf
  comparison \(\mathtt{sheafificationCompPullback}\) is the left-adjoint-uniqueness
  isomorphism of \cref{lem:sheafification_comp_pullback_eq_leftadjointuniq}.
\end{proof}

\begin{lemma}
[Composition coherence of the \(\mathtt{sheafCompPb}\)/counit composite]
  \leanok
  \label{lem:sheafCompPb_counit_comp_coherence}
  \lean{AlgebraicGeometry.Scheme.Modules.sheafCompPb_counit_comp_coherence}
  \uses{lem:sheafificationcomppullback_comp, lem:pullbackComp_hom_naturality,
        lem:sheafification_comp_pullback_eq_leftadjointuniq, def:sheafify_unit_iso}
  For composable \(h : Z \to Y\), \(f : Y \to X\) and \(M \in \Scheme.\mathtt{Modules}\,X\),
  the altitude-\((h \circ f)\) composite
  \((\mathtt{sheafCompPb}\,(h\circ f)).\mathrm{symm}.\mathrm{app}\,(M.\mathrm{val}) \mathbin{;}
  (\mathtt{pullback}\,(h\circ f)).\mathrm{map}\,\varepsilon^{X}_{M}\),
  post-composed with \((\mathtt{pullbackComp}\,h\,f).\mathrm{inv}_M\), factors through the
  sheafified presheaf-pullback pseudofunctoriality iso and the \(h^{*}\)-image of the
  altitude-\(f\) composite. The identity is
  \[
    (\mathtt{sheafCompPb}\,(h\circ f)).\mathrm{symm}.\mathrm{app}\,(M.\mathrm{val})
      \;\mathbin{;}\;
    (\mathtt{pullback}\,(h\circ f)).\mathrm{map}\,\varepsilon^{X}_{M}
      \;\mathbin{;}\;
    (\mathtt{pullbackComp}\,h\,f).\mathrm{inv}_M
    \;=\;
    \Psi^{M}_{h,f}
      \;\mathbin{;}\;
    (\mathtt{pullback}\,h).\mathrm{map}\Bigl(
      (\mathtt{sheafCompPb}\,f).\mathrm{symm}.\mathrm{app}\,(M.\mathrm{val})
        \;\mathbin{;}\;
      (\mathtt{pullback}\,f).\mathrm{map}\,\varepsilon^{X}_{M}\Bigr),
  \]
  where the left vertical
  \[
    \Psi^{M}_{h,f}
      \;=\;
    a_Z\!\bigl((\mathtt{pullbackComp}\,\varphi'_{f}\,\varphi'_{h}).\mathrm{inv}_{M.\mathrm{val}}\bigr)
      \;\mathbin{;}\;
    \bigl((\mathtt{sheafCompPb}\,h).\mathrm{app}\,(\mathtt{pb}^{\mathrm{pre}}f\,M.\mathrm{val})\bigr)
      .\mathrm{inv}
  \]
  is the sheafified presheaf-pullback pseudofunctoriality iso (the \(a_Z\)-sheafification of the
  presheaf comparison \(\mathtt{pullbackComp}\,\varphi'_{f}\,\varphi'_{h}\), followed by the
  inverse of \(\mathtt{sheafCompPb}\,h\) at the \emph{presheaf} object
  \(\mathtt{pb}^{\mathrm{pre}}f\,M.\mathrm{val}\)).

  \emph{Note on the factorisation.} The bottom edge is the \(h^{*}\)-image of the altitude-\(f\)
  composite \((\mathtt{sheafCompPb}\,f).\mathrm{symm}.\mathrm{app}(M.\mathrm{val}) \mathbin{;}
  (\mathtt{pullback}\,f).\mathrm{map}\,\varepsilon^{X}_{M} = \mathtt{pullbackValIso}\,f\,M\)
  --- the \(f\)-instance pushed forward by \(h^{*}\), over the presheaf object
  \(\mathtt{pb}^{\mathrm{pre}}f\,M.\mathrm{val}\), with the \(X\)-counit \(\varepsilon^{X}_{M}\).
  It is provably equal to (but \emph{not} the formal form of) the ``geometric'' altitude-\(h\)
  composite \((\mathtt{sheafCompPb}\,h).\mathrm{symm}.\mathrm{app}((f^{*}M).\mathrm{val})
  \mathbin{;} (\mathtt{pullback}\,h).\mathrm{map}\,\varepsilon^{Y}_{f^{*}M}\) over the
  \emph{sheaf} object \((f^{*}M).\mathrm{val}\) with the \(Y\)-counit; the equality is the
  sheafification--counit naturality of \cref{def:pullback_val_iso}.
\end{lemma}
\begin{proof}
  \leanok
  \uses{lem:sheafificationcomppullback_comp, lem:pullbackComp_hom_naturality,
        lem:sheafification_comp_pullback_eq_leftadjointuniq, def:sheafify_unit_iso}
  \textit{Source: internal categorical construction; no external reference.}

  Two ingredients combine, one per ``half'' of the composite.

  \emph{Sheafified-presheaf half (consume Sq1).} The closed Sq1 coherence
  \cref{lem:sheafificationcomppullback_comp} rewrites the altitude-\((h\circ f)\) comparison
  \(\mathtt{sheafCompPb}\,(h\circ f)\) as the composite of the altitude-\(f\) and altitude-\(h\)
  comparisons, conjugated by \(\mathtt{pullbackComp}\,h\,f\) and the strict presheaf-pullback
  pseudofunctor isomorphism \(\mathtt{pb}^{\mathrm{pre}}(h\circ f) \cong
  \mathtt{pb}^{\mathrm{pre}}h \circ \mathtt{pb}^{\mathrm{pre}}f\). This disposes of the
  sheafified-presheaf side and produces the left vertical \(\Psi^{M}_{h,f}\): the
  \(a_Z\)-sheafification of the presheaf comparison
  \((\mathtt{pullbackComp}\,\varphi'_{f}\,\varphi'_{h}).\mathrm{inv}_{M.\mathrm{val}}\),
  followed by the inverse of \(\mathtt{sheafCompPb}\,h\) at the presheaf object
  \(\mathtt{pb}^{\mathrm{pre}}f\,M.\mathrm{val}\). The remaining \(h^{*}\)-image factor on the
  right is then the \(h^{*}\)-pullback of the altitude-\(f\) composite
  \(\mathtt{pullbackValIso}\,f\,M\).

  \emph{Counit half (slide \(\varepsilon\) across \(\mathtt{pullbackComp}\)).} The
  \((h\circ f)\)-counit factor \(\varepsilon^{X}_{M}\) is the \emph{same} isomorphism
  \(a_X(M.\mathrm{val}) \cong M\) over \(X\) (independent of \(f,h\)). Its pullback
  \((\mathtt{pullback}\,(h\circ f)).\mathrm{map}\,\varepsilon^{X}_{M}\) is slid past
  \((\mathtt{pullbackComp}\,h\,f)_M\) by the naturality of the pseudofunctoriality iso
  (\cref{lem:pullbackComp_hom_naturality}, with \(g := \varepsilon^{X}_{M}\)) into
  \((\mathtt{pullback}\,h).\mathrm{map}\bigl((\mathtt{pullback}\,f).\mathrm{map}\,
  \varepsilon^{X}_{M}\bigr)\). The reconciliation of this slid \(X\)-counit with the
  altitude-\(h\) \emph{\(Y\)}-counit \(\varepsilon^{Y}_{f^{*}M}\) is exactly the
  sheafification-counit naturality already packaged in
  \cref{lem:sheafification_comp_pullback_eq_leftadjointuniq} /
  \cref{def:sheafify_unit_iso}: \(\mathtt{pullbackValIso}\,f\,M\) is built precisely to make
  \((\mathtt{pullback}\,f).\mathrm{map}\,\varepsilon^{X}_{M}\) agree with \(\varepsilon^{Y}_{f^{*}M}\)
  through \(\mathtt{sheafCompPb}\,f\).

  Pasting the two reconciled halves gives the claimed commuting square. No transposition is
  needed beyond the Sq1 cocycle of \cref{lem:sheafificationcomppullback_comp}.
\end{proof}

\begin{lemma}
[Sq4 --- composition coherence of \(\mathtt{pullbackValIso}\)]
  \leanok
  \label{lem:pullbackValIso_comp_coherence}
  \lean{AlgebraicGeometry.Scheme.Modules.pullbackValIso_comp}
  \uses{lem:pullbackValIso_eq_sheafCompPb, lem:sheafCompPb_counit_comp_coherence,
        def:pullback_val_iso, lem:sheafificationcomppullback_comp,
        lem:pullbackComp_hom_naturality}
  For composable \(h : Z \to Y\), \(f : Y \to X\) and \(M \in \Scheme.\mathtt{Modules}\,X\),
  the connecting isomorphism \(\mathtt{pullbackValIso}\) satisfies the composition coherence
  relating its \((h\circ f)\)-instance to the \(f\)-instance (transported by \(h^{*}\)), the
  \(h\)-instance (on the pulled-back module \(f^{*}M\)), and the pullback pseudofunctoriality
  iso \(\mathtt{pullbackComp}\,h\,f\): the square
  \[
    \begin{array}{ccc}
      a_Z\bigl(\mathtt{pb}^{\mathrm{pre}}(h\circ f)\,M.\mathrm{val}\bigr)
        & \xrightarrow{\;\;(\mathtt{pullbackValIso}\,(h\circ f)\,M).\mathrm{hom}\;\;}
        & (\mathtt{pullback}\,(h\circ f)).\mathrm{obj}\,M \\[2pt]
      {\scriptstyle \Psi^{M}_{h,f}}\big\downarrow & & \big\downarrow{\scriptstyle (\mathtt{pullbackComp}\,h\,f).\mathrm{inv}_M} \\[2pt]
      (\mathtt{pullback}\,h).\mathrm{obj}\,\bigl(a_Y(\mathtt{pb}^{\mathrm{pre}}f\,M.\mathrm{val})\bigr)
        & \xrightarrow{\;\;h^{*}\bigl((\mathtt{pullbackValIso}\,f\,M).\mathrm{hom}\bigr)\;\;}
        & (\mathtt{pullback}\,h).\mathrm{obj}\,(f^{*}M)
    \end{array}
  \]
  commutes, where the left vertical
  \[
    \Psi^{M}_{h,f}
      \;=\;
    a_Z\!\bigl((\mathtt{pullbackComp}\,\varphi'_{f}\,\varphi'_{h}).\mathrm{inv}_{M.\mathrm{val}}\bigr)
      \;\mathbin{;}\;
    \bigl((\mathtt{sheafCompPb}\,h).\mathrm{app}\,(\mathtt{pb}^{\mathrm{pre}}f\,M.\mathrm{val})\bigr)
      .\mathrm{inv}
  \]
  is the sheafified presheaf-pullback pseudofunctoriality iso (the \(a_Z\)-sheafification of the
  presheaf comparison \(\mathtt{pullbackComp}\,\varphi'_{f}\,\varphi'_{h}\) followed by the inverse
  of \(\mathtt{sheafCompPb}\,h\) at the presheaf object \(\mathtt{pb}^{\mathrm{pre}}f\,M.\mathrm{val}\)),
  and \((\mathtt{pullback}\,h).\mathrm{obj}\,(f^{*}M) =
  (\mathtt{pullback}\,h).\mathrm{obj}\,((\mathtt{pullback}\,f).\mathrm{obj}\,M)\). The bottom edge
  is the \(h^{*}\)-image of the \emph{altitude-\(f\)} comparison \(\mathtt{pullbackValIso}\,f\,M\)
  over the presheaf object \(\mathtt{pb}^{\mathrm{pre}}f\,M.\mathrm{val}\); this equals
  (though differs in form from) the ``geometric'' altitude-\(h\)
  bottom edge \((\mathtt{pullbackValIso}\,h\,(f^{*}M)).\mathrm{hom}\) over the sheaf object
  \((f^{*}M).\mathrm{val}\). As an equation,
  \[
    (\mathtt{pullbackValIso}\,(h\circ f)\,M).\mathrm{hom}
      \;\mathbin{;}\;
    (\mathtt{pullbackComp}\,h\,f).\mathrm{inv}_M
      \;=\;
    \Psi^{M}_{h,f} \;\mathbin{;}\; h^{*}\!\bigl((\mathtt{pullbackValIso}\,f\,M).\mathrm{hom}\bigr).
  \]
  This is the connecting-object boundary where the abstract pullback's \(.\mathrm{val}\)
  presentation meets the helper-\(.\mathrm{obj}\) presentation; tensoring the \(M\)- and
  \(N\)-instances yields the trailing
  \(\mathtt{tensorObjIsoOfIso}\bigl((\mathtt{pullbackComp}\,h\,f)_M,
  (\mathtt{pullbackComp}\,h\,f)_N\bigr).\mathrm{hom}\) factor of the D3$'$ master identity.
\end{lemma}
\begin{proof}
  \leanok
  \uses{lem:pullbackValIso_eq_sheafCompPb, lem:sheafCompPb_counit_comp_coherence,
        def:pullback_val_iso, lem:sheafificationcomppullback_comp,
        lem:pullbackComp_hom_naturality}
  \textit{Source: internal categorical construction; no external reference.}

  A three-step corollary.
  \begin{enumerate}
    \item[(1)] \emph{Unfold.} By \cref{lem:pullbackValIso_eq_sheafCompPb}, rewrite the
      altitude-\((h\circ f)\) occurrence of \(\mathtt{pullbackValIso}\) as its
      \(\mathtt{sheafCompPb}\)/counit composite
      \((\mathtt{sheafCompPb}\,\varphi).\mathrm{symm}.\mathrm{app}(-) \mathbin{;}
       (\mathtt{pullback}\,\varphi).\mathrm{map}\,\varepsilon\). The square to prove becomes a
      statement purely about the \(\mathtt{sheafCompPb}\)/counit composites.
    \item[(2)] \emph{Apply the coherence.} \cref{lem:sheafCompPb_counit_comp_coherence}
      discharges precisely this commuting square: its sheafified-presheaf half is the closed Sq1
      coherence \cref{lem:sheafificationcomppullback_comp}, and its counit half slides
      \(\varepsilon^{X}_{M}\) across \((\mathtt{pullbackComp}\,h\,f)_M\) by
      \cref{lem:pullbackComp_hom_naturality}, identifying the left vertical with
      \(\Psi^{M}_{h,f}\).
    \item[(3)] \emph{Refold.} Re-apply \cref{lem:pullbackValIso_eq_sheafCompPb} in reverse to
      repackage the altitude-\(f\) composite back as
      \(h^{*}(\mathtt{pullbackValIso}\,f\,M)\); the altitude-\(h\) comparison
      \(\mathtt{sheafCompPb}\,h\) and the sheafified presheaf pseudofunctoriality iso remain as
      the left vertical \(\Psi^{M}_{h,f}\), closing the square.
    \end{enumerate}
  The only substantive step is the \(\Psi^{M}_{h,f}\) connecting-object bookkeeping of
  step~(2) (the \(.\mathrm{val}\)/\(.\mathrm{obj}\) boundary).
\end{proof}

\begin{lemma}
[Pullback commutes with $\otimes$ on locally-trivial pairs (D4$'$)]
  \label{lem:pullback_tensor_iso_loctriv}
  \lean{AlgebraicGeometry.Scheme.Modules.pullbackTensorIsoOfLocallyTrivial}
  % NOTE: \lean{} pin target absent from Lean source as of iter-271 (verified by grep).
  % This is the LIVE D4' build target (the committed loc-triv route) — a forward reference
  % to a not-yet-formalized declaration, NOT a stale rename. Once the prover lands the decl
  % under this exact name, sync_leanok resolves the pin automatically.
  \uses{def:IsLocallyTrivial, lem:pullback_tensor_map, lem:pullback_tensor_map_natural,
        lem:pullback_tensor_iso_unit, lem:pullback_tensor_map_basechange,
        lem:isiso_of_isiso_restrict, lem:tensorobj_preserves_locally_trivial,
        lem:IsLocallyTrivial_pullback}
  Let \(f : Y \to X\) be a morphism of schemes and let
  \(M, N \in \Scheme.\mathtt{Modules}\,X\) be \emph{locally trivial}
  (\cref{def:IsLocallyTrivial}). Then the sheaf-level comparison map
  \(\delta_{\mathrm{sheaf}} = \mathtt{pullbackTensorMap}\) of
  \cref{lem:pullback_tensor_map} is an isomorphism:
  \[
    f^*(M \otimes_X N)
      \;\xrightarrow{\;\sim\;}\;
    (f^*M) \otimes_Y (f^*N),
  \]
  natural in \(M\) and \(N\). This is the comparison the relative Picard consumer needs;
  it is established by promoting the already-built map \(\delta_{\mathrm{sheaf}}\) to an
  isomorphism through a chart-chase, with no comparison for arbitrary modules and no
  filtered-colimit machinery.
  % SOURCE: [Stacks Project], Modules on Ringed Spaces, Lemma
  % \ref{lemma-tensor-product-pullback} (pullback commutes with the tensor product,
  % functorially) and Lemma \ref{lemma-pullback-locally-free} (the pullback of a locally
  % free module is locally free --- the geometric content licensing the chart-chase)
  % (read from references/stacks-modules.tex, L2392-2400 and L2112-2118).
  % SOURCE QUOTE:
  % "\begin{lemma}
  % \label{lemma-tensor-product-pullback}
  % Let $f : (X, \mathcal{O}_X) \to (Y, \mathcal{O}_Y)$ be
  % a morphism of ringed spaces. Let $\mathcal{F}$, $\mathcal{G}$
  % be $\mathcal{O}_Y$-modules. Then
  % $f^*(\mathcal{F} \otimes_{\mathcal{O}_Y} \mathcal{G})
  % = f^*\mathcal{F} \otimes_{\mathcal{O}_X} f^*\mathcal{G}$
  % functorially in $\mathcal{F}$, $\mathcal{G}$.
  % \end{lemma}"
  % SOURCE QUOTE:
  % "\begin{lemma}
  % \label{lemma-pullback-locally-free}
  % Let $f : (X, \mathcal{O}_X) \to (Y, \mathcal{O}_Y)$
  % be a morphism of ringed spaces. If $\mathcal{G}$ is
  % a locally free $\mathcal{O}_Y$-module, then
  % $f^*\mathcal{G}$ is a locally free $\mathcal{O}_X$-module.
  % \end{lemma}"
  \textit{Source: [Stacks Project], Modules on Ringed Spaces, Lemmas
  \texttt{lemma-tensor-product-pullback} and \texttt{lemma-pullback-locally-free}:
  pullback commutes with \(\otimes\) functorially, and the pullback of a locally free
  (locally trivial) module is again locally free; on locally-trivial pairs the comparison
  is therefore an isomorphism, checkable on a trivialising cover.}
\end{lemma}

\begin{proof}
  \uses{def:IsLocallyTrivial, lem:pullback_tensor_map, lem:pullback_tensor_map_natural,
        lem:pullback_tensor_iso_unit, lem:pullback_tensor_map_basechange,
        lem:isiso_of_isiso_restrict, lem:tensorobj_preserves_locally_trivial,
        lem:IsLocallyTrivial_pullback}
  \textit{Source: [Stacks Project], Lemmas \texttt{lemma-tensor-product-pullback} and
  \texttt{lemma-pullback-locally-free}.}

  By local triviality of \(M\) and \(N\) (\cref{def:IsLocallyTrivial}) choose a common
  affine open cover \(\{U_i\}\) of \(X\) on which both restrict to the trivial bundle,
  \(M|_{U_i} \cong \mathcal{O}_{U_i}\) and \(N|_{U_i} \cong \mathcal{O}_{U_i}\) --- the
  common refinement of the two trivialising covers, exactly as in the proof that the
  tensor product of locally-trivial modules is locally trivial
  (\cref{lem:tensorobj_preserves_locally_trivial}). The preimages
  \(\{f^{-1}U_i\}\) form an open cover of \(Y\). It suffices, by the
  restriction-detects-isomorphism criterion \cref{lem:isiso_of_isiso_restrict}, to show
  that the canonical map \(\delta_{\mathrm{sheaf}}(M,N)\) restricts to an isomorphism on
  each \(f^{-1}U_i\); the canonical global map being shown locally-iso is exactly what
  lets us conclude (a non-canonical patchwork of local isomorphisms would not glue).

  Fix \(i\) and write \(g_i : f^{-1}U_i \to U_i\) for the restriction of \(f\), with open
  immersions \(j_i : U_i \hookrightarrow X\) and \(j_i' : f^{-1}U_i \hookrightarrow Y\),
  so \(f \circ j_i' = j_i \circ g_i\). The base-change coherence
  \cref{lem:pullback_tensor_map_basechange} identifies the restriction
  \((j_i')^*\delta^{f}_{\mathrm{sheaf}}(M,N)\) with the comparison
  \(\delta^{g_i}_{\mathrm{sheaf}}(j_i^*M, j_i^*N)\) for \(g_i\) on the restricted modules,
  and reduces the claim to: \(\delta^{g_i}_{\mathrm{sheaf}}(j_i^*M, j_i^*N)\) is an
  isomorphism. On \(U_i\) the restricted modules are trivial, \(j_i^*M \cong
  \mathcal{O}_{U_i}\) and \(j_i^*N \cong \mathcal{O}_{U_i}\); by naturality of
  \(\delta_{\mathrm{sheaf}}\) in both arguments (\cref{lem:pullback_tensor_map_natural})
  these isomorphisms transport \(\delta^{g_i}_{\mathrm{sheaf}}(j_i^*M, j_i^*N)\) to the
  comparison on the unit pair \(\delta^{g_i}_{\mathrm{sheaf}}(\mathcal{O}_{U_i},
  \mathcal{O}_{U_i})\). That comparison is an isomorphism by
  \cref{lem:pullback_tensor_iso_unit} --- the unit-pair case, where the structure-sheaf
  comparison \(\mathtt{pullbackUnitIso}\) (an isomorphism for all \(g_i\)) and the unitors
  do the work. Hence \(\delta^{g_i}_{\mathrm{sheaf}}(j_i^*M, j_i^*N)\) is an isomorphism,
  so is its base-change identification \((j_i')^*\delta^{f}_{\mathrm{sheaf}}(M,N)\), and
  thus \(\delta^{f}_{\mathrm{sheaf}}(M,N)\) restricts to an isomorphism on each
  \(f^{-1}U_i\). (Both sides are indeed locally trivial on this cover --- \(f^*M, f^*N,
  f^*(M\otimes_X N)\) are locally trivial by \cref{lem:IsLocallyTrivial_pullback} and
  \cref{lem:tensorobj_preserves_locally_trivial} --- which is what makes the chart-local
  comparison a map between copies of the structure sheaf.) By
  \cref{lem:isiso_of_isiso_restrict}, \(\delta_{\mathrm{sheaf}}(M,N)\) is a global
  isomorphism, natural in \(M, N\) by \cref{lem:pullback_tensor_map_natural}.
\end{proof}

\begin{lemma}
[An invertible module is locally trivial (locally free of rank $1$)]
  \label{lem:isinvertible_implies_locallytrivial}
  \lean{AlgebraicGeometry.Scheme.Modules.IsInvertible.isLocallyTrivial}
  % NOTE: \lean{} pin target absent from Lean source as of iter-271 (verified by grep).
  % Forward reference to a not-yet-formalized (and, per the OFF-path NOTE below,
  % off-critical-path) declaration — NOT a stale rename; no existing Lean decl to repin to.
  \uses{def:scheme_modules_isinvertible, def:IsLocallyTrivial, lem:stalk_tensor_commutation}
  % NOTE: OFF the critical path AND unnecessary. The committed route to
  % lem:isinvertible_pullback is the LOCALLY-TRIVIAL comparison lem:pullback_tensor_iso_loctriv;
  % the relative Picard consumer carries IsLocallyTrivial directly (its carrier OnProduct =
  % { L // IsLocallyTrivial L }), so this forward bridge IsInvertible ⇒ IsLocallyTrivial is
  % never needed on-path. Only the EASY reverse direction IsLocallyTrivial ⇒ IsInvertible
  % (lem:tensorobj_inverse_invertible / exists_tensorObj_inverse) is consumed. Retained solely
  % as a future ingredient for the A.2.c Quot-embedding (invertible ⇒ locally free of rank 1).
  % Still Mathlib-scale: (1) stalk-invertibility plumbing turning M⊗N≅𝒪_X into a stalkwise
  % invertible module, (2) finite-presentation spreading-out for SheafOfModules on a scheme
  % ("M f.p. + M_x free ⇒ M≅𝒪 on an affine nbhd"), present in Mathlib only at the
  % CommRing/LocalizedModule level. Do not build it for the pullback corollary.
  Let \(X\) be a scheme and \(M \in \Scheme.\mathtt{Modules}\,X\). If \(M\) is
  \(\otimes\)-invertible (\cref{def:scheme_modules_isinvertible}), then \(M\) is locally
  trivial: every point of \(X\) has an open (affine) neighbourhood \(U\) on which
  \(M|_U \cong \mathcal{O}_U\). In symbols
  \[
    \mathtt{IsInvertible}(M)
      \;\Longrightarrow\;
    \mathtt{LineBundle.IsLocallyTrivial}\,M.
  \]
  This implication is recorded for later use --- the Quot-embedding of A.2.c --- and is
  \emph{not} on the path to \cref{lem:isinvertible_pullback}. That corollary is obtained
  from the locally-trivial comparison \cref{lem:pullback_tensor_iso_loctriv}, whose
  locally-trivial hypotheses are supplied directly by the consumer's carrier (and, for the
  inverse object, by the easy reverse implication \cref{lem:tensorobj_inverse_invertible});
  this forward bridge is never invoked.
  % SOURCE: [Stacks Project], Modules on Ringed Spaces, Lemma
  % \ref{lemma-invertible-is-locally-free-rank-1} (invertible modules and locally free
  % rank-1 modules coincide when all stalks are local rings)
  % (read from references/stacks-modules.tex, L4159-4165).
  % SOURCE QUOTE:
  % "\begin{lemma}
  % \label{lemma-invertible-is-locally-free-rank-1}
  % Let $(X, \mathcal{O}_X)$ be a ringed space. Any locally free
  % $\mathcal{O}_X$-module of rank $1$ is invertible.
  % If all stalks $\mathcal{O}_{X, x}$ are local rings, then
  % the converse holds as well (but in general this is not the case).
  % \end{lemma}"
  % SOURCE: [Stacks Project], Modules on Ringed Spaces, Lemma
  % \ref{lemma-invertible} (invertibility characterised by a tensor inverse; an
  % invertible module is locally a direct summand of a finite free module)
  % (read from references/stacks-modules.tex, L4066-4079).
  % SOURCE QUOTE:
  % "\begin{lemma}
  % \label{lemma-invertible}
  % Let $(X, \mathcal{O}_X)$ be a ringed space. Let $\mathcal{L}$
  % be an $\mathcal{O}_X$-module. Equivalent are
  % \begin{enumerate}
  % \item $\mathcal{L}$ is invertible, and
  % \item there exists an $\mathcal{O}_X$-module $\mathcal{N}$
  % such that
  % $\mathcal{L} \otimes_{\mathcal{O}_X} \mathcal{N} \cong \mathcal{O}_X$.
  % \end{enumerate}
  % In this case $\mathcal{L}$ is locally a direct summand of a finite free
  % $\mathcal{O}_X$-module and the module $\mathcal{N}$ in (2) is isomorphic to
  % $\SheafHom_{\mathcal{O}_X}(\mathcal{L}, \mathcal{O}_X)$.
  % \end{lemma}"
  \textit{Source: [Stacks Project], Modules on Ringed Spaces, Lemmas
  \texttt{lemma-invertible} and \texttt{lemma-invertible-is-locally-free-rank-1}: on a
  ringed space whose stalks are local rings, an invertible module is locally free of
  rank \(1\); and an invertible module is finitely presented (locally a direct summand of
  a finite free module).}
\end{lemma}

% SOURCE QUOTE PROOF: (read from references/stacks-modules.tex, L4186-4198).
% "Assume all stalks $\mathcal{O}_{X, x}$ are local rings and $\mathcal{L}$
% invertible. In the proof of Lemma \ref{lemma-invertible}
% we have seen that $\mathcal{L}_x$ is an invertible
% $\mathcal{O}_{X, x}$-module for all $x \in X$. Since
% $\mathcal{O}_{X, x}$ is local, we see that
% $\mathcal{L}_x \cong \mathcal{O}_{X, x}$
% (More on Algebra, Section \ref{more-algebra-section-picard}).
% Since $\mathcal{L}$ is of finite presentation by
% Lemma \ref{lemma-invertible} we conclude that $\mathcal{L}$
% is locally free of rank $1$ by
% Lemma \ref{lemma-finite-presentation-stalk-free}."
\begin{proof}
  \uses{def:scheme_modules_isinvertible, lem:stalk_tensor_commutation}
  \textit{Source: [Stacks Project], Lemma \texttt{lemma-invertible-is-locally-free-rank-1}
  (converse direction), using \texttt{lemma-invertible}.}

  By \(\mathtt{IsInvertible}(M)\) (\cref{def:scheme_modules_isinvertible}) there is a
  witness \(N\) and an isomorphism \(\psi : M \otimes_X N \xrightarrow{\sim}
  \mathcal{O}_X\); by Stacks \texttt{lemma-invertible} this already forces \(M\) to be
  finitely presented (locally a direct summand of a finite free \(\mathcal{O}_X\)-module).
  Fix \(x \in X\). Taking stalks of \(\psi\) and commuting the stalk past the tensor
  product --- the varying-ring stalk-tensor identification
  \((M \otimes_X N)_x \cong M_x \otimes_{\mathcal{O}_{X,x}} N_x\) (\(\mathtt{stalkTensorIso}\)) --- gives
  \[
    M_x \otimes_{\mathcal{O}_{X,x}} N_x \;\cong\; \mathcal{O}_{X,x},
  \]
  so \(M_x\) is an invertible module over the \emph{local} ring \(\mathcal{O}_{X,x}\)
  (the scheme hypothesis on \(X\) supplies local stalks). An invertible module over a
  local ring is free of rank \(1\), so \(M_x \cong \mathcal{O}_{X,x}\). Since \(M\) is
  finitely presented and its stalk at \(x\) is free, the standard finite-presentation
  spreading-out (Stacks \texttt{lemma-finite-presentation-stalk-free}) extends the free
  stalk to a free neighbourhood: there is an open \(U \ni x\) with \(M|_U \cong
  \mathcal{O}_U\). As \(x\) was arbitrary and any such \(U\) may be shrunk to an affine
  open, \(M\) is locally trivial.
\end{proof}

\begin{lemma}
[Pullback preserves $\otimes$-invertibility for a general scheme morphism]
  \label{lem:isinvertible_pullback}
  \lean{AlgebraicGeometry.Scheme.Modules.IsInvertible.pullback}
  % NOTE: \lean{} pin target absent from Lean source as of iter-271 (verified by grep).
  % Forward reference to a not-yet-formalized declaration (the three-line consequence
  % below, pending its loc-triv input) — NOT a stale rename; no existing Lean decl to repin to.
  \uses{def:scheme_modules_isinvertible, lem:pullback_tensor_iso_loctriv,
        lem:pullback_unit_iso, def:IsLocallyTrivial}
  % NOTE: Three-line Stacks proof (lemma-pullback-invertible). Witness f^*N; the iso is
  % (pullbackTensorIsoOfLocallyTrivial)⁻¹ ≫ (pullback f).mapIso e ≫ pullbackUnitIso. Depends on
  % the LOCALLY-TRIVIAL comparison lem:pullback_tensor_iso_loctriv and lem:pullback_unit_iso
  % (in hand). The loc-triv hypotheses on M and its inverse N are supplied directly: the
  % consumer carries IsLocallyTrivial, and the inverse N may be taken locally trivial (its dual,
  % via the easy reverse exists_tensorObj_inverse). NO forward bridge IsInvertible ⇒ loc-triv,
  % NO general comparison, NO concrete inverse-image build.
  Let \(f : Y \to X\) be a morphism of schemes and let
  \(M \in \Scheme.\mathtt{Modules}\,X\) be \(\otimes\)-invertible
  (\cref{def:scheme_modules_isinvertible}) with a tensor inverse \(N\) and isomorphism
  \(e : M \otimes_X N \cong \mathcal{O}_X\). If both \(M\) and \(N\) are locally trivial
  (\cref{def:IsLocallyTrivial}) --- as holds for every line bundle, and in particular for
  the modules the relative Picard carrier presents --- then the pullback \(f^*M\) is again
  \(\otimes\)-invertible, with inverse \(f^*N\):
  \[
    \mathtt{IsInvertible}(M)
      \;\Longrightarrow\;
    \mathtt{IsInvertible}\bigl((\mathtt{Scheme.Modules.pullback}\,f).\mathrm{obj}\,M\bigr).
  \]
  % SOURCE: [Stacks Project], Modules on Ringed Spaces, Lemma
  % \ref{lemma-pullback-invertible} (the pullback of an invertible module is
  % invertible) (read from references/stacks-modules.tex, L4142-4147).
  % SOURCE QUOTE:
  % "\begin{lemma}
  % \label{lemma-pullback-invertible}
  % Let $f : (X, \mathcal{O}_X) \to (Y, \mathcal{O}_Y)$ be a
  % morphism of ringed spaces. The pullback $f^*\mathcal{L}$ of an
  % invertible $\mathcal{O}_Y$-module is invertible.
  % \end{lemma}"
  \textit{Source: [Stacks Project], Modules on Ringed Spaces, Lemma
  \texttt{lemma-pullback-invertible}: the pullback \(f^*\mathcal{L}\) of an
  invertible \(\mathcal{O}\)-module along a morphism of ringed spaces is
  invertible.}
\end{lemma}

% SOURCE QUOTE PROOF: (read from references/stacks-modules.tex, L4149-4157).
% "\begin{proof}
% By Lemma \ref{lemma-invertible}
% there exists an $\mathcal{O}_Y$-module $\mathcal{N}$ such that
% $\mathcal{L} \otimes_{\mathcal{O}_Y} \mathcal{N} \cong \mathcal{O}_Y$.
% Pulling back we get
% $f^*\mathcal{L} \otimes_{\mathcal{O}_X} f^*\mathcal{N} \cong \mathcal{O}_X$
% by Lemma \ref{lemma-tensor-product-pullback}.
% Thus $f^*\mathcal{L}$ is invertible by Lemma \ref{lemma-invertible}.
% \end{proof}"
\begin{proof}
  \uses{def:scheme_modules_isinvertible, lem:pullback_tensor_iso_loctriv,
        lem:pullback_unit_iso, def:IsLocallyTrivial}
  \textit{Source: [Stacks Project], Lemma \texttt{lemma-pullback-invertible}.}
  This is the Stacks proof in the project's notation, using the locally-trivial
  comparison \cref{lem:pullback_tensor_iso_loctriv}.

  By \(\mathtt{IsInvertible}(M)\) (\cref{def:scheme_modules_isinvertible}) there is a
  witness \(N \in \Scheme.\mathtt{Modules}\,X\) together with an isomorphism
  \(e : M \otimes_X N \xrightarrow{\sim} \mathcal{O}_X\); by hypothesis \(M\) and \(N\)
  are locally trivial (\cref{def:IsLocallyTrivial}). We take the pullback witness
  \((\mathtt{pullback}\,f).\mathrm{obj}\,N = f^*N\) and exhibit \(f^*M\) as
  \(\otimes\)-invertible by the composite
  \[
    (f^*M) \otimes_Y (f^*N)
      \;\xrightarrow[\sim]{\;(\text{\cref{lem:pullback_tensor_iso_loctriv}})^{-1}\;}\;
    f^*(M \otimes_X N)
      \;\xrightarrow[\sim]{\;f^*e\;}\;
    f^*\mathcal{O}_X
      \;\xrightarrow[\sim]{\;\mathtt{pullbackUnitIso}\;}\;
    \mathcal{O}_Y,
  \]
  whose three arrows are, in order: the inverse of the locally-trivial pullback--tensor
  isomorphism \cref{lem:pullback_tensor_iso_loctriv} (applicable because \(M, N\) are
  locally trivial); the image \((\mathtt{pullback}\,f).\mathtt{mapIso}\,e\) of the
  invertibility witness \(e\) under the pullback functor; and the structure-sheaf
  comparison \(\mathtt{pullbackUnitIso}\) identifying \(f^*\mathcal{O}_X\) with
  \(\mathcal{O}_Y\) (\cref{lem:pullback_unit_iso}). Each is an isomorphism --- the first
  by the locally-trivial strong-monoidality of pullback, the second as the image of an
  isomorphism under a functor, the third by \cref{lem:pullback_unit_iso} --- so the
  composite is an isomorphism \(f^*M \otimes_Y f^*N \cong \mathcal{O}_Y\). Hence \(f^*N\)
  is a tensor inverse of \(f^*M\), and \(\mathtt{IsInvertible}(f^*M)\) holds.
\end{proof}

\section{The invertibility-carrier Picard group}
\label{sec:tensorobj_pic_carrier}

The group law of \cref{lem:tensorobj_isoclass_commgroup} is carried on the
\emph{locally-trivial} predicate \(\mathtt{IsLocallyTrivial}\), and on that carrier
the existence of a tensor inverse for each class is an extra obligation
(\cref{lem:tensorobj_inverse_invertible}): one must \emph{manufacture} the dual
object \(L^{-1} = \mathcal{H}om_{\mathcal{O}_X}(L, \mathcal{O}_X)\) and prove the
contraction \(L \otimes_X L^{-1} \cong \mathcal{O}_X\), which is the whole content of
the dual-infrastructure build of \cref{sec:tensorobj_dual_infra}.

This section carries the group law on a \emph{different} predicate for which that
obligation evaporates: the \(\otimes\)-invertibility predicate
\(\mathtt{IsInvertible}\) of \cref{def:scheme_modules_isinvertible},
\[
  \mathtt{IsInvertible}(M)
  \;:\equiv\;
  \exists\, N,\ \mathrm{Nonempty}\bigl(M \otimes_X N \cong \mathcal{O}_X\bigr).
\]
On this carrier the inverse of a class \([M]\) is supplied by the \emph{membership
witness itself} --- the very \(N\) whose existence \(\mathtt{IsInvertible}(M)\)
asserts --- so no inverse object need be constructed and no separate
tensor-inverse lemma is consumed. The cost is shifted to a routine
well-definedness fact (\cref{lem:isinvertible_inverse_welldef}: the witness \(N\) is
determined up to isomorphism), proved from the monoidal coherence isomorphisms
alone. The carrier is the scheme-theoretic analogue of Mathlib's
\(\mathtt{CommRing.Pic}\), whose elements are exactly the iso-classes of invertible
modules (Stacks tag 01CX).

\begin{lemma}

\leanok
[Associator for $\otimes_X$ on invertible objects (corollary of the unconditional
associator)]
  \label{lem:tensorobj_assoc_iso_invertible}
  \lean{AlgebraicGeometry.Scheme.Modules.tensorObj_assoc_iso_invertible}
  \uses{def:scheme_modules_tensorobj, def:scheme_modules_isinvertible,
        lem:tensorobj_assoc_iso}
  Let \(X\) be a scheme and let
  \(M, N, P \in \Scheme.\mathtt{Modules}\,X\) be \(\otimes\)-invertible
  (\cref{def:scheme_modules_isinvertible}). There is an isomorphism
  \[
    (M \otimes_X N) \otimes_X P
      \;\xrightarrow{\ \sim\ }\;
    M \otimes_X (N \otimes_X P)
  \]
  in \(\Scheme.\mathtt{Modules}\,X\). The invertibility-carrier group law
  (\cref{thm:pic_commgroup}) consumes only the \emph{existence} of this
  isomorphism, the proposition \(\mathrm{Nonempty}\bigl((M \otimes_X N) \otimes_X P
  \cong M \otimes_X (N \otimes_X P)\bigr)\); no pentagon, triangle, or naturality
  coherence is ever required.
\end{lemma}

\begin{proof}
  \uses{lem:tensorobj_assoc_iso}
  Immediate specialisation of the \emph{unconditional} associator
  \cref{lem:tensorobj_assoc_iso} to the invertible arguments \(M, N, P\); no
  hypothesis on the arguments is used. (The earlier route argued via ``invertible
  \(\Rightarrow\) locally free \(\Rightarrow\) sectionwise flat'' to feed the
  flat-whiskering lemma \cref{lem:flat_whisker_localizer}; that step is \emph{false}
  --- local freeness does not give global sectionwise flatness \(M(U)\) over
  \(\mathcal{O}_X(U)\) for a non-affine open \(U\) --- and is no longer needed, since
  the general associator is itself unconditional, its single open left-whiskering
  obligation \cref{lem:islocallyinjective_whisker_of_W} being closed via the d.2
  stalk--tensor commutation \cref{lem:islocallyinjective_whiskerleft_via_stalk}.)
\end{proof}

\begin{definition}

\leanok
[The invertibility-carrier Picard group \(\Pic X\)]
  \label{def:pic_carrier}
  \lean{AlgebraicGeometry.Scheme.Modules.PicGroup}
  \uses{def:scheme_modules_tensorobj, def:scheme_modules_isinvertible, def:pic_setoid}
  % SOURCE: [Stacks Project], Tag 01CX (Modules on Ringed Spaces, \S Invertible
  % modules), Definition definition-pic (the Picard group as the abelian group of
  % iso-classes of invertible modules under tensor product)
  % (read from references/stacks-modules.tex, L4350-L4357).
  % SOURCE QUOTE:
  % "\begin{definition}
  % \label{definition-pic}
  % Let $(X, \mathcal{O}_X)$ be a ringed space.
  % The {\it Picard group} $\Pic(X)$ of $X$ is the
  % abelian group whose elements are isomorphism classes of
  % invertible $\mathcal{O}_X$-modules, with addition
  % corresponding to tensor product.
  % \end{definition}"
  \textit{Source: [Stacks Project], Tag 01CX, Definition~\texttt{definition-pic}:
  the Picard group \(\Pic(X)\) is the abelian group of isomorphism classes of
  invertible \(\mathcal{O}_X\)-modules, with addition the tensor product.}
  Let \(X\) be a scheme. Consider the subtype of \(\otimes\)-invertible objects
  \[
    \mathtt{Invertible}\,X
    \;:=\;
    \bigl\{\, M \in \Scheme.\mathtt{Modules}\,X
              \,\bigm|\, \mathtt{IsInvertible}\,M \,\bigr\}
  \]
  of \cref{def:scheme_modules_isinvertible}, and equip it with the setoid whose
  relation is ``there exists an isomorphism'',
  \(M \sim M' :\equiv \mathrm{Nonempty}(M \cong M')\) (a proposition-valued
  equivalence relation: reflexive, symmetric and transitive because isomorphisms
  compose and invert). The \emph{Picard group carrier}
  \[
    \Pic X \;:=\; \mathtt{Invertible}\,X \,/\!\sim
  \]
  is the quotient type of isomorphism classes of \(\otimes\)-invertible
  \(\mathcal{O}_X\)-modules. We write \([M] \in \Pic X\) for the class of an
  invertible \(M\). This is the scheme-theoretic analogue of Mathlib's
  \(\mathtt{CommRing.Pic}\,R\), carried on \(\mathtt{IsInvertible}\) rather than on
  \(\mathtt{IsLocallyTrivial}\); on a scheme the two predicates agree (Stacks tag
  0B8M), but only the invertibility form carries the inverse witness existentially.
\end{definition}

\begin{lemma}

\leanok
[\(\otimes\)-invertibility is closed under tensor product]
  \label{lem:isinvertible_tensor}
  \lean{AlgebraicGeometry.Scheme.Modules.IsInvertible.tensorObj}
  \uses{def:scheme_modules_isinvertible, lem:tensorobj_assoc_iso_invertible,
        lem:tensorobj_comm_iso, lem:tensorobj_unit_iso}
  % SOURCE: [Stacks Project], Tag 01CR (Modules on Ringed Spaces, \S Invertible
  % modules), Lemma lemma-constructions-invertible (item 1: the tensor product of
  % two invertible modules is invertible)
  % (read from references/stacks-modules.tex, L4200-L4213).
  % SOURCE QUOTE:
  % "\begin{lemma}
  % \label{lemma-constructions-invertible}
  % Let $(X, \mathcal{O}_X)$ be a ringed space.
  % \begin{enumerate}
  % \item If $\mathcal{L}$, $\mathcal{N}$ are invertible
  % $\mathcal{O}_X$-modules, then so is
  % $\mathcal{L} \otimes_{\mathcal{O}_X} \mathcal{N}$.
  % \item If $\mathcal{L}$ is an invertible $\mathcal{O}_X$-module, then so is
  % $\SheafHom_{\mathcal{O}_X}(\mathcal{L}, \mathcal{O}_X)$ and the evaluation map
  % $\mathcal{L} \otimes_{\mathcal{O}_X}
  % \SheafHom_{\mathcal{O}_X}(\mathcal{L}, \mathcal{O}_X) \to \mathcal{O}_X$
  % is an isomorphism.
  % \end{enumerate}
  % \end{lemma}"
  \textit{Source: [Stacks Project], Tag 01CR,
  Lemma~\texttt{lemma-constructions-invertible} (item 1): the tensor product of two
  invertible \(\mathcal{O}_X\)-modules is again invertible.}
  Let \(X\) be a scheme and \(M, M' \in \Scheme.\mathtt{Modules}\,X\). If
  \(\mathtt{IsInvertible}\,M\) and \(\mathtt{IsInvertible}\,M'\), then
  \(\mathtt{IsInvertible}(M \otimes_X M')\).
\end{lemma}

\begin{proof}
  \uses{def:scheme_modules_isinvertible, lem:tensorobj_assoc_iso_invertible,
        lem:tensorobj_comm_iso, lem:tensorobj_unit_iso}
  Let \(N\) be a tensor inverse of \(M\) (so \(M \otimes_X N \cong \mathcal{O}_X\))
  and \(N'\) a tensor inverse of \(M'\) (so \(M' \otimes_X N' \cong \mathcal{O}_X\)),
  furnished by the two invertibility witnesses. We claim \(N \otimes_X N'\) is a
  tensor inverse of \(M \otimes_X M'\). Indeed, reassociating and commuting the four
  factors with the associator \cref{lem:tensorobj_assoc_iso_invertible} and the
  braiding \cref{lem:tensorobj_comm_iso},
  \[
    (M \otimes_X M') \otimes_X (N \otimes_X N')
    \;\cong\;
    (M \otimes_X N) \otimes_X (M' \otimes_X N')
    \;\cong\;
    \mathcal{O}_X \otimes_X \mathcal{O}_X
    \;\cong\;
    \mathcal{O}_X,
  \]
  the middle step substituting the two witness isomorphisms and the last being the
  unitor \cref{lem:tensorobj_unit_iso}. Hence \(N \otimes_X N'\) witnesses
  \(\mathtt{IsInvertible}(M \otimes_X M')\). Every step is consumed only as the
  existence of an isomorphism, so no coherence identity is required.
\end{proof}

\begin{lemma}

\leanok
[The structure sheaf is \(\otimes\)-invertible]
  \label{lem:isinvertible_unit}
  \lean{AlgebraicGeometry.Scheme.Modules.isInvertible_unit}
  \uses{def:scheme_modules_isinvertible, lem:tensorobj_unit_iso}
  % SOURCE: [Stacks Project], Tag 01CR (Modules on Ringed Spaces, \S Invertible
  % modules), Lemma lemma-invertible (the "$\exists\,\mathcal{N}$ with
  % $\mathcal{L} \otimes \mathcal{N} \cong \mathcal{O}_X$" characterisation of
  % invertibility)
  % (read from references/stacks-modules.tex, L4066-L4079).
  % SOURCE QUOTE:
  % "\begin{lemma}
  % \label{lemma-invertible}
  % Let $(X, \mathcal{O}_X)$ be a ringed space. Let $\mathcal{L}$
  % be an $\mathcal{O}_X$-module. Equivalent are
  % \begin{enumerate}
  % \item $\mathcal{L}$ is invertible, and
  % \item there exists an $\mathcal{O}_X$-module $\mathcal{N}$
  % such that
  % $\mathcal{L} \otimes_{\mathcal{O}_X} \mathcal{N} \cong \mathcal{O}_X$.
  % \end{enumerate}
  % In this case $\mathcal{L}$ is locally a direct summand of a finite free
  % $\mathcal{O}_X$-module and the module $\mathcal{N}$ in (2) is isomorphic to
  % $\SheafHom_{\mathcal{O}_X}(\mathcal{L}, \mathcal{O}_X)$.
  % \end{lemma}"
  \textit{Source: [Stacks Project], Tag 01CR, Lemma~\texttt{lemma-invertible}: an
  \(\mathcal{O}_X\)-module is invertible iff there exists \(\mathcal{N}\) with
  \(\mathcal{L} \otimes_{\mathcal{O}_X} \mathcal{N} \cong \mathcal{O}_X\).}
  Let \(X\) be a scheme. The structure sheaf
  \(\mathcal{O}_X = \mathtt{SheafOfModules.unit}\,X.\mathtt{ringCatSheaf}\) is
  \(\otimes\)-invertible: \(\mathtt{IsInvertible}\,\mathcal{O}_X\).
\end{lemma}

\begin{proof}
  \uses{def:scheme_modules_isinvertible, lem:tensorobj_unit_iso}
  Take \(N := \mathcal{O}_X\) as the witness. The right unitor
  \cref{lem:tensorobj_unit_iso} provides
  \(\mathcal{O}_X \otimes_X \mathcal{O}_X \cong \mathcal{O}_X\), which is precisely
  the isomorphism required by \(\mathtt{IsInvertible}\,\mathcal{O}_X\).
\end{proof}

\begin{lemma}

\leanok
[The tensor inverse is determined up to isomorphism]
  \label{lem:isinvertible_inverse_welldef}
  \lean{AlgebraicGeometry.Scheme.Modules.IsInvertible.inverse_unique}
  \uses{def:scheme_modules_isinvertible, lem:tensorobj_unit_iso,
        lem:tensorobj_assoc_iso_invertible, lem:tensorobj_comm_iso}
  % SOURCE: [Stacks Project], Tag 01CR (Modules on Ringed Spaces, \S Invertible
  % modules), Lemma lemma-invertible (the inverse $\mathcal{N}$ is unique up to
  % isomorphism: it is isomorphic to the dual)
  % (read from references/stacks-modules.tex, L4066-L4079).
  % SOURCE QUOTE:
  % "\begin{lemma}
  % \label{lemma-invertible}
  % Let $(X, \mathcal{O}_X)$ be a ringed space. Let $\mathcal{L}$
  % be an $\mathcal{O}_X$-module. Equivalent are
  % \begin{enumerate}
  % \item $\mathcal{L}$ is invertible, and
  % \item there exists an $\mathcal{O}_X$-module $\mathcal{N}$
  % such that
  % $\mathcal{L} \otimes_{\mathcal{O}_X} \mathcal{N} \cong \mathcal{O}_X$.
  % \end{enumerate}
  % In this case $\mathcal{L}$ is locally a direct summand of a finite free
  % $\mathcal{O}_X$-module and the module $\mathcal{N}$ in (2) is isomorphic to
  % $\SheafHom_{\mathcal{O}_X}(\mathcal{L}, \mathcal{O}_X)$.
  % \end{lemma}"
  \textit{Source: [Stacks Project], Tag 01CR, Lemma~\texttt{lemma-invertible}: when
  \(\mathcal{L}\) is invertible the module \(\mathcal{N}\) with
  \(\mathcal{L} \otimes_{\mathcal{O}_X} \mathcal{N} \cong \mathcal{O}_X\) is
  isomorphic to the dual \(\SheafHom_{\mathcal{O}_X}(\mathcal{L}, \mathcal{O}_X)\),
  hence unique up to isomorphism.}
  Let \(X\) be a scheme and \(M, N, N' \in \Scheme.\mathtt{Modules}\,X\). If
  \(M \otimes_X N \cong \mathcal{O}_X\) and \(M \otimes_X N' \cong \mathcal{O}_X\),
  then \(N \cong N'\). Consequently the assignment \([M] \mapsto [N]\), sending an
  invertible class to the class of its tensor inverse, is well-defined on
  \(\Pic X\).
\end{lemma}

\begin{proof}
  \uses{def:scheme_modules_isinvertible, lem:tensorobj_unit_iso,
        lem:tensorobj_assoc_iso_invertible, lem:tensorobj_comm_iso}
  This is the standard ``inverse-of-inverse'' argument in a commutative monoid, run
  with isomorphisms in place of equalities. Chain the unitor
  \cref{lem:tensorobj_unit_iso}, the two hypotheses, the associator
  \cref{lem:tensorobj_assoc_iso_invertible} and the braiding
  \cref{lem:tensorobj_comm_iso}:
  \[
    N
    \;\cong\; N \otimes_X \mathcal{O}_X
    \;\cong\; N \otimes_X (M \otimes_X N')
    \;\cong\; (N \otimes_X M) \otimes_X N'
    \;\cong\; \mathcal{O}_X \otimes_X N'
    \;\cong\; N',
  \]
  where the second step substitutes \(\mathcal{O}_X \cong M \otimes_X N'\), the
  third is the associator, the fourth uses the braiding
  \(N \otimes_X M \cong M \otimes_X N \cong \mathcal{O}_X\), and the outer two are
  the unitors. Hence \(N \cong N'\). For well-definedness on classes, note that if
  \(M \cong M'\) then a tensor inverse of \(M\) is also a tensor inverse of \(M'\)
  (tensor with the isomorphism), so the inverse class depends only on \([M]\).
\end{proof}

\begin{theorem}

\leanok
[The invertibility-carrier Picard group is abelian]
  \label{thm:pic_commgroup}
  \lean{AlgebraicGeometry.Scheme.Modules.picCommGroup}
  \uses{def:pic_carrier, def:pic_mul, def:pic_inv, lem:isinvertible_tensor,
        lem:isinvertible_unit, lem:isinvertible_inverse_welldef, lem:tensorobj_unit_iso,
        lem:tensorobj_comm_iso, lem:tensorobj_assoc_iso_invertible}
  % SOURCE: [Stacks Project], Tag 01CX (Modules on Ringed Spaces, \S Invertible
  % modules), Definition definition-pic + Lemma lemma-constructions-invertible
  % (the iso-classes of invertibles form an abelian group under tensor product)
  % (read from references/stacks-modules.tex, L4200-L4213 and L4350-L4357).
  % SOURCE QUOTE:
  % "\begin{lemma}
  % \label{lemma-constructions-invertible}
  % Let $(X, \mathcal{O}_X)$ be a ringed space.
  % \begin{enumerate}
  % \item If $\mathcal{L}$, $\mathcal{N}$ are invertible
  % $\mathcal{O}_X$-modules, then so is
  % $\mathcal{L} \otimes_{\mathcal{O}_X} \mathcal{N}$.
  % \item If $\mathcal{L}$ is an invertible $\mathcal{O}_X$-module, then so is
  % $\SheafHom_{\mathcal{O}_X}(\mathcal{L}, \mathcal{O}_X)$ and the evaluation map
  % $\mathcal{L} \otimes_{\mathcal{O}_X}
  % \SheafHom_{\mathcal{O}_X}(\mathcal{L}, \mathcal{O}_X) \to \mathcal{O}_X$
  % is an isomorphism.
  % \end{enumerate}
  % \end{lemma}
  % [\dots]
  % \begin{definition}
  % \label{definition-pic}
  % Let $(X, \mathcal{O}_X)$ be a ringed space.
  % The {\it Picard group} $\Pic(X)$ of $X$ is the
  % abelian group whose elements are isomorphism classes of
  % invertible $\mathcal{O}_X$-modules, with addition
  % corresponding to tensor product.
  % \end{definition}"
  \textit{Source: [Stacks Project], Tag 01CX, Definition~\texttt{definition-pic}
  and Lemma~\texttt{lemma-constructions-invertible}: the isomorphism classes of
  invertible \(\mathcal{O}_X\)-modules form an abelian group under tensor product,
  with unit \(\mathcal{O}_X\).}
  Let \(X\) be a scheme. The carrier \(\Pic X\) of \cref{def:pic_carrier} is an
  abelian group, with
  \[
    [M] \cdot [M'] := [M \otimes_X M'],
    \qquad
    1 := [\mathcal{O}_X],
    \qquad
    [M]^{-1} := [N] \ \text{for any witness } M \otimes_X N \cong \mathcal{O}_X.
  \]
\end{theorem}

\begin{proof}
  \uses{def:pic_carrier, def:pic_mul, def:pic_inv, lem:isinvertible_tensor,
        lem:isinvertible_unit, lem:isinvertible_inverse_welldef, lem:tensorobj_unit_iso,
        lem:tensorobj_comm_iso, lem:tensorobj_assoc_iso_invertible}
  Each group-operation datum and each axiom is supplied by a single
  existence-of-isomorphism, read as an equality of classes in \(\Pic X\).

  \emph{The operations are well-defined.} The product \([M] \cdot [M'] := [M
  \otimes_X M']\) lands in \(\Pic X\) because \(\otimes_X\) preserves invertibility
  (\cref{lem:isinvertible_tensor}), and it is independent of representatives by the
  bifunctoriality of \(\otimes_X\)
  (\cref{lem:scheme_modules_tensorobj_functoriality}): an isomorphism in either
  argument tensors to an isomorphism, hence an equality of classes. The unit
  \([\mathcal{O}_X]\) lies in \(\Pic X\) by \cref{lem:isinvertible_unit}. The
  inverse \([M] \mapsto [N]\), \(N\) any membership witness of
  \(\mathtt{IsInvertible}\,M\), is well-defined by
  \cref{lem:isinvertible_inverse_welldef} --- the witness is determined up to
  isomorphism, so its class depends only on \([M]\). Crucially, \emph{no inverse
  object is constructed}: the inverse is read off the existential carried by
  \(\mathtt{IsInvertible}\).

  \emph{The axioms.} Each axiom is the image, under passage to iso-classes, of one
  coherence isomorphism of \(\otimes_X\):
  \[
    \mathtt{one\_mul} \leftarrow \langle \lambda \rangle,
    \quad
    \mathtt{mul\_one} \leftarrow \langle \rho \rangle,
    \quad
    \mathtt{mul\_assoc} \leftarrow \langle \alpha \rangle,
    \quad
    \mathtt{mul\_comm} \leftarrow \langle \beta \rangle,
    \quad
    \mathtt{mul\_left\_inv} \leftarrow \langle \varepsilon \rangle,
  \]
  where the unit laws are the unitors \cref{lem:tensorobj_unit_iso}, associativity
  is the invertible-case associator \cref{lem:tensorobj_assoc_iso_invertible},
  commutativity is the braiding \cref{lem:tensorobj_comm_iso}, and the left-inverse
  law \([N] \cdot [M] = [N \otimes_X M] = [\mathcal{O}_X] = 1\) is the witness
  isomorphism \(\varepsilon : N \otimes_X M \cong \mathcal{O}_X\) itself (obtained
  from \(M \otimes_X N \cong \mathcal{O}_X\) by the braiding). Each axiom asserts
  only \(\mathrm{Nonempty}(\cdots \cong \cdots)\), so no pentagon, triangle, or
  hexagon coherence is invoked; commutativity makes the left-inverse law suffice
  for a two-sided inverse.

  This is the invertibility-carrier analogue of
  \cref{lem:tensorobj_isoclass_commgroup}, differing in exactly one respect: the
  left-inverse field is discharged by the predicate's own membership witness rather
  than by the manufactured dual of \cref{lem:tensorobj_inverse_invertible}. The
  whole dual-infrastructure obligation of \cref{sec:tensorobj_dual_infra} is
  thereby off this group's critical path.
\end{proof}

\section{Consumer: the \(\AddCommGroup\) instance for the relative Picard sheaf}
\label{sec:tensorobj_consumer_acg}

The substrate of \cref{sec:tensorobj_substrate,sec:tensorobj_onproduct_lift}
supplies all the data missing from \cref{lem:rel_pic_sharp_groupoid}'s Lean
encoding. The consumer instance below packages the abelian-group law
\(\bigl[\,\mathcal{L}\,\bigr] + \bigl[\,\mathcal{L}'\,\bigr]
 := \bigl[\,\mathcal{L} \otimes_{C \times_S T} \mathcal{L}'\,\bigr]\) on the
relative Picard quotient via the substrate.

\begin{theorem}


[Abelian-group instance on the relative Picard quotient via \(\Scheme.\mathtt{Modules}.\mathtt{tensorObj}\)]
  \label{thm:rel_pic_addcommgroup_via_tensorobj}
  \lean{AlgebraicGeometry.Scheme.PicSharp.addCommGroup_via_tensorObj}
  % NOTE: \lean{} pin target absent from Lean source as of iter-271 (verified by grep).
  % Forward reference to a not-yet-formalized declaration (the real instance currently in
  % Lean is `PicSharp.addCommGroup`; the `_via_tensorObj` upgrade is not yet built) — NOT a
  % stale rename; no existing Lean decl to repin to under this name.
  % NOTE (carrier pivot, deferred): the \uses below still points at the old
  % locally-trivial group law lem:tensorobj_isoclass_commgroup; this is DEFERRED and
  % will be repointed to thm:pic_commgroup (the invertibility-carrier group, inverse
  % free) once that carrier group is formalized. Left unchanged this iter.
  \uses{lem:tensorobj_lift_onproduct,
        lem:pullback_compatible_with_tensorobj,
        def:scheme_modules_isinvertible,
        lem:tensorobj_isoclass_commgroup,
        lem:tensorobj_assoc_iso,
        lem:tensorobj_unit_iso,
        lem:tensorobj_comm_iso,
        thm:relative_pic_quotient_well_defined,
        lem:rel_pic_sharp_groupoid}

  % SOURCE: [Kleiman], "The Picard scheme", \S 2 (FGA Explained Ch.~9 \S 9.2),
  % Definitions df:aPf + df:Pfs (the absolute and relative Picard functors are
  % abelian-group-valued; the relative one is defined as a quotient of abelian
  % groups)
  % (read from references/kleiman-picard-src/kleiman-picard.tex, L1274-L1318).
  % SOURCE QUOTE:
  % "\begin{dfn}\label{df:aPf}
  %  The {\it absolute Picard functor} $\Pic_X$ is the functor from  the
  %  category of (locally Noetherian) $S$-schemes $T$ to the category of
  % abelian groups defined by the formula
  % 	$$\Pic_X(T):=\Pic(X_T).$$
  % \end{dfn}
  %  [\dots]
  % \begin{dfn}\label{df:Pfs}
  %   The {\it relative Picard functor} $\Pic_{X/S}$ is defined by
  % 	$$\Pic_{X/S}(T):=\Pic(X_T)\big/\Pic(T).$$
  %  Denote its associated sheaves in the Zariski, \'etale, and fppf
  % topologies  by
  % 	$$\Pic_{(X/S)\zar},\ \Pic_{(X/S)\et},\
  % 	\Pic_{(X/S)\fppf}.$$
  %  \end{dfn}"
  \textit{Source: [Kleiman], ``The Picard scheme'', \S 2,
  Defs.~df:aPf + df:Pfs (the target category is the category of abelian
  groups; the relative Picard functor is a quotient of abelian groups).}
  Let \(C/k\) be a smooth proper geometrically integral curve over a
  field \(k\), let \(\pi_C : C \to \Spec k\) be the structure morphism,
  and let \(T\) be a \(k\)-scheme with structure morphism
  \(\pi_T : T \to \Spec k\). The set
  \[
    \Pic^{\sharp}_{C/k}(T) \;\;:=\;\; \Pic(C \times_k T) \,/\, \pi_T^*\Pic(T)
  \]
  of \cref{thm:relative_pic_quotient_well_defined} carries a canonical
  abelian-group structure with addition
  \(\bigl[\,\mathcal{L}\,\bigr] + \bigl[\,\mathcal{L}'\,\bigr]
   := \bigl[\,\mathcal{L} \otimes_{C \times_k T} \mathcal{L}'\,\bigr]\),
  neutral element \(\bigl[\,\mathcal{O}_{C \times_k T}\,\bigr]\), and
  inverse
  \(-\bigl[\,\mathcal{L}\,\bigr] := \bigl[\,\mathcal{L}^{-1}\,\bigr]\),
  where \(\otimes_{C \times_k T}\) and \((-)^{-1}\) are the operations
  supplied by the substrate of \cref{def:scheme_modules_tensorobj}
  restricted to the \(\texttt{LineBundle.OnProduct}\,\pi_C\,\pi_T\) carrier
  via \cref{lem:tensorobj_lift_onproduct}. With this structure the quotient
  map
  \(\Pic(C \times_k T) \twoheadrightarrow \Pic^{\sharp}_{C/k}(T)\) is a
  surjective group homomorphism with kernel exactly \(\pi_T^*\Pic(T)\),
  in agreement with \cref{lem:rel_pic_sharp_groupoid}.
\end{theorem}

\begin{proof}
  \uses{lem:tensorobj_lift_onproduct,
        lem:pullback_compatible_with_tensorobj,
        def:scheme_modules_isinvertible,
        lem:tensorobj_isoclass_commgroup,
        thm:relative_pic_quotient_well_defined,
        lem:rel_pic_sharp_groupoid}
  Write \(\Pic(C \times_k T)\) for the set of isomorphism classes of objects
  of \(\texttt{LineBundle.OnProduct}\,\pi_C\,\pi_T\) (line bundles on
  \(C \times_k T\)), with the operation
  \([L] + [L'] := [L \otimes_{C \times_k T} L']\). This operation is
  well-defined on isomorphism classes: an isomorphism \(L \cong L_0\) tensors
  to an isomorphism \(L \otimes L' \cong L_0 \otimes L'\) by functoriality of
  \(\otimes\) (\cref{lem:scheme_modules_tensorobj_functoriality}), and
  symmetrically in the second argument. The abelian-group structure on
  \(\Pic(C \times_k T)\) is then exactly the \emph{units group of the
  commutative monoid of \(\otimes\)-iso-classes}
  (\cref{lem:tensorobj_isoclass_commgroup}): multiplication
  \([L]\cdot[L'] = [L \otimes_{C\times_k T} L']\) is well-defined and
  commutative-monoidal by the coherence isomorphisms
  \cref{lem:tensorobj_assoc_iso,lem:tensorobj_unit_iso,lem:tensorobj_comm_iso},
  and an iso-class is a unit precisely when its carrier is
  \(\otimes\)-invertible (\cref{def:scheme_modules_isinvertible}), the inverse
  being carried by the predicate. Each axiom is a \emph{proposition} --- the
  existence of an isomorphism --- so no monoidal-coherence datum (pentagon,
  triangle, hexagon) is invoked; the \(\AddCommGroup\) on
  \(\Pic(C \times_k T)\) is the additive form of this units group. This mirrors
  Mathlib's own Picard group: \(\mathtt{CommRing.Pic}\,R\) carries its
  \(\mathtt{CommGroup}\) (\(\mathtt{instCommGroupPic}\),
  \texttt{Mathlib.RingTheory.PicardGroup}) transported from
  \(\mathrm{Units}\bigl(\mathrm{Skeleton}(\mathtt{ModuleCat}\,R)\bigr)\); here
  the same \(\mathrm{Units}(\cdots)\) shape is transported to the scheme-level
  \(\otimes_X\). (The scheme-level Picard group is genuinely absent from
  Mathlib --- \texttt{Mathlib.RingTheory.PicardGroup} records ``connect to
  invertible sheaves'' as an open task --- so the project supplies it lightly.)
  The invertibility-predicate carrier is the project's \(\mathtt{IsInvertible}\)
  (\cref{def:scheme_modules_isinvertible}), the scheme-level analogue of
  Mathlib's \(\mathtt{Module.Invertible}\); the geometric predicate
  \(\mathtt{IsLocallyTrivial}\) is retained separately and is connected to
  \(\mathtt{IsInvertible}\) only by an off-critical-path equivalence (only the
  geometric \(\mathtt{IsInvertible} \Leftarrow \mathtt{IsLocallyTrivial}\)-style
  direction is nontrivial, and it is not needed here).

  By \cref{lem:pullback_compatible_with_tensorobj}, the pullback functor
  \(\pi_T^{*} : \texttt{LineBundle}\,T \to
   \texttt{LineBundle.OnProduct}\,\pi_C\,\pi_T\) carries tensor products to
  tensor products and the unit to the unit, so the induced map on isomorphism
  classes \(\pi_T^{*} : \Pic(T) \to \Pic(C \times_k T)\) is a group
  homomorphism, and its image is a subgroup.

  By \cref{thm:relative_pic_quotient_well_defined}, the underlying set of
  \(\Pic^{\sharp}_{C/k}(T)\) is the set of cosets of this subgroup. The
  quotient of an abelian group by a subgroup is itself an abelian group, via
  the standard \(\mathtt{QuotientAddGroup}\) construction, and the quotient
  map is a surjective group homomorphism whose kernel is exactly the subgroup
  (Stacks tag 01CR for the Picard-group special case; the general
  quotient-of-abelian-groups statement is in any algebra textbook). Putting
  the pieces together gives the \(\AddCommGroup\) instance on
  \(\Pic^{\sharp}_{C/k}(T)\) of the statement, and identifies it with the
  structure of \cref{lem:rel_pic_sharp_groupoid}.
\end{proof}

\section{Out of scope}
\label{sec:tensorobj_out_of_scope}

The following are downstream of this chapter and live in dedicated
sibling chapters; this chapter does \emph{not} touch them:
\begin{itemize}
  \item \textbf{The sheafified relative Picard functor and its abelian-group target}
    (\cref{def:rel_pic_etale_sheafification,
              thm:rel_pic_etale_sheaf_group_structure}) --- once
    \cref{thm:rel_pic_addcommgroup_via_tensorobj} is established, the
    \(\AddCommGroup\)-valued presheaf, its \'etale sheafification, and
    the sheafification preserving the abelian-group target close
    against it (chapter \cref{chap:Picard_RelPicFunctor}); none of those
    declarations is the responsibility of this chapter.
  \item \textbf{Representability of \(\Pic^\sharp_{(C/k)\et}\) by a scheme}
    (A.2.c, \cref{chap:Picard_FGAPicRepresentability}) --- consumes the
    full \(\AddCommGroup\)-valued sheafified presheaf;
    representability is downstream of \cref{thm:rel_pic_etale_sheaf_group_structure}.
  \item \textbf{Identity component \(\Pic^0_{C/k}\)} (A.3,
    \cref{chap:Picard_IdentityComponent}) --- once representability is
    in hand, the identity component is cut out by Hilbert polynomials;
    that is independent of the tensor-product substrate of this chapter.
  \item \textbf{Albanese universal property} (A.4,
    \cref{chap:Albanese_AlbaneseUP}) --- consumes the full Route~A stack
    downstream of A.3.
  \item \textbf{Mathlib upstream PR of \(\Scheme.\mathtt{Modules}.\mathtt{tensorObj}\)} ---
    the substrate of this chapter is project-side, intentionally; a
    parallel Mathlib upstream PR is welcome but not on the critical
    path. Once the substrate ships in Mathlib, the
    \texttt{TensorObjSubstrate.lean} file can be slimmed to a single
    import + namespace re-export, but no downstream Lean file changes.
\end{itemize}

\section{Internal-consistency check}
\label{sec:tensorobj_consistency_check}

The Lean-pinned declaration blocks of this chapter form a connected
\texttt{\textbackslash uses\{\dots\}} chain whose leaves all resolve
either inside this chapter, inside chapter
\cref{chap:Picard_LineBundlePullback} (A.1.b on disk), or inside
chapter \cref{chap:Picard_RelPicFunctor} (A.1.c on disk):
\begin{itemize}
  \item \cref{def:scheme_modules_tensorobj} introduces the substrate
    operation and depends on no prior block.
  \item \cref{lem:scheme_modules_tensorobj_functoriality} uses
    \cref{def:scheme_modules_tensorobj}.
  \item \cref{rem:scheme_modules_monoidal_off_path} uses
    \cref{def:scheme_modules_tensorobj,
          lem:scheme_modules_tensorobj_functoriality, lem:whisker_of_W}; it
    records that, under route~(e), the full monoidal category is \emph{cheap}
    (the output of \(\mathtt{LocalizedMonoidal}\)) and that the group law still
    consumes only existence-of-isomorphism propositions.
  \item \cref{lem:restrictscalars_laxmonoidal} uses
    \cref{def:scheme_modules_tensorobj}; it is a self-contained
    \(\mathtt{CommRingCat}\)-level supplement, off the critical path and
    \emph{not} an ingredient of \cref{lem:tensorobj_restrict_iso} --- nothing
    downstream depends on it.
  \item \cref{def:scheme_modules_isinvertible} (the \(\otimes\)-invertibility
    predicate, carrying the inverse) uses \cref{def:scheme_modules_tensorobj}.
  \item The route~(e) whisker lemmas
    \cref{lem:islocallyinjective_whisker_of_W} (local injectivity of
    \(F \triangleleft g\) for \(g \in J.W\): the PRIMARY route assumes \(F\) locally
    trivial and proves it sheaf-level via \cref{lem:tensorobj_restrict_iso} and the
    unitor, with no stalks; the FALLBACK route, for arbitrary \(F\), is the
    stalkwise argument with residual ingredients d.1/d.2 --- the same (d.2)
    stalk--tensor commutation \cref{lem:stalk_tensor_commutation} that the \emph{live}
    associator obligation \cref{lem:islocallyinjective_whiskerleft_via_stalk} consumes,
    hence on the critical path) and
    \cref{lem:whisker_of_W} (its closed left/right packaging) supply the two
    fields of the instance \cref{lem:jw_ismonoidal}. The latter applies Mathlib's
    \(\mathtt{LocalizedMonoidal}\) to the already-monoidal
    \(\mathtt{PresheafOfModules}\,\mathcal{O}_X\) and the sheafification localizer
    to give the full \(\MonoidalCategory\) structure on
    \(\Scheme.\mathtt{Modules}\,X\); it is stated-but-unformalized, pending the
    sole open proof obligation.
  \item The coherence isomorphisms \cref{lem:tensorobj_unit_iso}
    (left/right unitors) and \cref{lem:tensorobj_comm_iso} (braiding) are the cheap
    \(\mathtt{sheafification.mapIso}\) of the presheaf-level coherence data. The
    associator \cref{lem:tensorobj_assoc_iso} is the sheafification transport of the
    presheaf-level associator, whose single open left-whiskering obligation
    \cref{lem:islocallyinjective_whiskerleft_via_stalk} is closed
    \emph{unconditionally} --- for an arbitrary whiskering object, with no flatness and
    no local triviality --- by the d.2 stalk--tensor commutation
    \cref{lem:stalk_tensor_commutation} (\cref{sec:tensorobj_stalk_tensor}).
    None uses \(\mathtt{MonoidalClosed}\) or sectionwise flatness; the associator's
    pinned \(\mathtt{IsLocallyTrivial}\) hypotheses are vestigial (the construction is
    unconditional).
  \item \cref{lem:isiso_sheafification_map_of_W} (sheafification sends a
    morphism whose underlying map is in \(J.W\) to an isomorphism) depends on no
    prior block; it is a closed go/no-go bridge, retained as a supplement (the
    route~(e) associator is the API-derived one, so this bridge is no longer the
    associator's load-bearing step).
  \item \cref{lem:flat_whisker_localizer} (left/right whiskering of a
    sheafification-localizer morphism by a flat object stays in the localizer
    class) uses \cref{def:scheme_modules_tensorobj}; it is a valid closed
    standalone result, superseded on the critical path by the d.2 route of
    \cref{sec:tensorobj_stalk_tensor}
    (\cref{lem:islocallyinjective_whiskerleft_via_stalk}).
  \item \cref{lem:tensorobj_isoclass_commgroup} (the commutative group on
    locally-trivial \(\otimes\)-iso-classes) uses
    \cref{def:scheme_modules_tensorobj,
          def:scheme_modules_isinvertible,
          lem:scheme_modules_tensorobj_functoriality,
          lem:tensorobj_lift_onproduct,
          lem:tensorobj_inverse_invertible,
          lem:tensorobj_assoc_iso,
          lem:tensorobj_unit_iso,
          lem:tensorobj_comm_iso}. This is the group-law engine: the PRIMARY
    construction is the commutative group on locally-trivial iso-classes à la
    \(\mathtt{Module.Invertible}\) / \(\mathtt{CommRing.Pic}\), needing only
    \cref{lem:tensorobj_restrict_iso} and \cref{lem:tensorobj_inverse_invertible}
    on top of the assembled locally-trivial associator/unitors/braiding; the
    \(\mathrm{Units}(\mathrm{Skeleton}(\Scheme.\mathtt{Modules}\,X))\) presentation
    via \cref{lem:jw_ismonoidal} is retained as the route~(e) fallback.
  \item Under the PRIMARY locally-trivial group law, \cref{lem:tensorobj_restrict_iso}
    (consumed by the PRIMARY proof of \cref{lem:islocallyinjective_whisker_of_W}) and
    \cref{lem:tensorobj_inverse_invertible} (in the \(\uses\) of
    \cref{lem:tensorobj_isoclass_commgroup}, supplying inverses) are ON the critical
    path. The remaining supplements
    \cref{lem:restrictscalars_laxmonoidal,lem:tensorobj_preserves_locally_trivial}
    use \cref{def:scheme_modules_tensorobj} and are auxiliary. The group law rests on
    \cref{lem:tensorobj_isoclass_commgroup} and \cref{def:scheme_modules_isinvertible}.
  \item \cref{lem:tensorobj_lift_onproduct} restricts \(\otimes\) to the
    locally-trivial carrier \(\texttt{LineBundle.OnProduct}\) on \(C \times_S T\),
    using \cref{def:scheme_modules_tensorobj,lem:tensorobj_preserves_locally_trivial}
    (the Lean \(\mathtt{tensorObjOnProduct}\) applies
    \(\mathtt{tensorObj\_isLocallyTrivial}\) directly).
  \item \cref{lem:pullback_compatible_with_tensorobj} uses
    \cref{lem:tensorobj_lift_onproduct} and implicitly
    \cref{def:pullback_along_projection} from
    \cref{chap:Picard_LineBundlePullback}; it is \emph{not} blocked on
    \cref{lem:tensorobj_restrict_iso}.
  \item \cref{thm:rel_pic_addcommgroup_via_tensorobj} uses
    \cref{lem:tensorobj_lift_onproduct,
          lem:pullback_compatible_with_tensorobj,
          def:scheme_modules_isinvertible,
          lem:tensorobj_isoclass_commgroup,
          thm:relative_pic_quotient_well_defined,
          lem:rel_pic_sharp_groupoid}; the last two are in
    \cref{chap:Picard_LineBundlePullback} (A.1.b) and
    \cref{chap:Picard_RelPicFunctor} (A.1.c) respectively.
\end{itemize}
All \(\uses\) targets resolve either in this chapter, in
\cref{chap:Picard_LineBundlePullback}, or in
\cref{chap:Picard_RelPicFunctor}. No declaration cross-references a
chapter that does not yet exist on disk.

\section{Sheaf internal-hom of \(\mathcal{O}_X\)-modules (the \(\otimes\)-inverse
  infrastructure)}
\label{sec:tensorobj_dual_infra}

The PRIMARY group law (\cref{lem:tensorobj_isoclass_commgroup}) needs a
\(\otimes\)-inverse for every line bundle, supplied by
\cref{lem:tensorobj_inverse_invertible} as the dual
\(L^{-1} = \mathcal{H}om_{\mathcal{O}_X}(L, \mathcal{O}_X)\). The standard route
constructs that dual as a special case of the \emph{sheaf internal hom} of
\(\mathcal{O}_X\)-modules; over a fixed ring the internal hom of \(R\)-modules is the
linear-coyoneda module \(\mathtt{ihom}\,M\,N = (M \to N)\), with dual
\(\mathtt{ihom}\,M\,\mathbf{1} = M^\vee\), but for the \emph{varying} structure sheaf
no such
internal-hom / monoidal-closed structure exists at the
\(\mathtt{PresheafOfModules}\), \(\mathtt{SheafOfModules}\), or general-categorical
level. This section decomposes the construction of the missing object into
individually-formalizable sub-steps, following the slice presentation of the internal
hom; each block names the intended Lean declaration and its dependencies. The
specialisation that closes \cref{lem:tensorobj_inverse_invertible} is recorded in
\cref{rem:dual_discharges_inverse}; an alternative descent-theoretic route is noted in
\cref{rem:dual_via_stack}.

\emph{Status: deferred, off the group's critical path.} The blocks of this section
--- the dual object \(\mathtt{dual}\) and its restriction compatibility
\cref{lem:dual_restrict_iso}, the local-triviality of the dual
\cref{lem:dual_isLocallyTrivial}, the gluing engine
\cref{lem:sheafofmodules_hom_of_local_compat}, and the discharge remark
\cref{rem:dual_discharges_inverse} --- together constitute the \emph{deferred}
\(\mathtt{IsInvertible} \Leftrightarrow \mathtt{IsLocallyTrivial}\) bridge: on a
scheme an invertible \(\mathcal{O}_X\)-module is locally free of rank one, and
conversely (Stacks tag 0B8M), so the two carriers agree, and the locally-trivial
carrier's missing inverse (\cref{lem:tensorobj_inverse_invertible}) is recovered
precisely from the dual built here.
% SOURCE: [Stacks Project], Tag 0B8M (Modules on Ringed Spaces, \S Invertible
% modules), Lemma lemma-invertible-is-locally-free-rank-1
% (read from references/stacks-modules.tex, L4159-L4165).
% SOURCE QUOTE:
% "\begin{lemma}
% \label{lemma-invertible-is-locally-free-rank-1}
% Let $(X, \mathcal{O}_X)$ be a ringed space. Any locally free
% $\mathcal{O}_X$-module of rank $1$ is invertible.
% If all stalks $\mathcal{O}_{X, x}$ are local rings, then
% the converse holds as well (but in general this is not the case).
% \end{lemma}"
However, the relative Picard \emph{group law} is now carried on the
\(\otimes\)-invertibility predicate \(\mathtt{IsInvertible}\)
(\cref{sec:tensorobj_pic_carrier}, \cref{thm:pic_commgroup}), where the inverse is
the membership witness and \emph{no} dual object is constructed. This section is
therefore needed only by a downstream consumer that specifically requires the
locally-trivial form of an inverse --- not by the group law. The blocks are retained
intact for that future use.

A structural point governs the whole construction. The naive rule
\(U \mapsto \Hom_{\mathcal{O}_X(U)}(M(U), N(U))\) is \emph{contravariant} in the
restriction maps of \(M\) (a restriction \(M(U) \to M(V)\) for \(V \subseteq U\)
induces precomposition \(\Hom(M(V), N(V)) \to \Hom(M(U), \cdot)\) in the wrong
direction), so it is \emph{not} a presheaf in the covariant sense and cannot be
sheafified like the tensor product \cref{def:scheme_modules_tensorobj}. The remedy is
the \emph{slice} formula \(U \mapsto (M|_U \to N|_U)\) --- the module of morphisms of
\emph{restricted} modules over the restricted ring --- whose presheaf functoriality
(restrict a global-over-\(U\) morphism to a global-over-\(V\) morphism, for
\(V \subseteq U\)) \emph{is} covariant. Constructing that slice presheaf, and
proving it satisfies the sheaf condition, is the substance of the build.

\begin{definition}

\leanok
[The global-ring module on morphisms of presheaves of modules (value module)]
  \label{def:presheaf_internal_hom_value}
  \lean{PresheafOfModules.InternalHom.homModule}
  \uses{def:scheme_modules_tensorobj}
  % SOURCE: [Stacks Project], "Modules on Ringed Spaces", \S Internal Hom
  % (\texttt{section-internal-hom}, tag area 01CM)
  % (read from references/stacks-modules.tex, L3508--L3514).
  % SOURCE QUOTE:
  % "abelian groups. In addition, given an element
  % $\varphi \in \Hom_{\mathcal{O}_X|_U}(\mathcal{F}|_U, \mathcal{G}|_U)$
  % and a section $f \in \mathcal{O}_X(U)$ then we can define
  % $f\varphi \in \Hom_{\mathcal{O}_X|_U}(\mathcal{F}|_U, \mathcal{G}|_U)$
  % by either precomposing with multiplication by $f$ on $\mathcal{F}|_U$
  % or postcomposing with multiplication by $f$ on $\mathcal{G}|_U$ (it gives
  % the same result). Hence we in fact get a sheaf of $\mathcal{O}_X$-modules."
  \textit{Source: [Stacks Project], ``Modules on Ringed Spaces'', \S Internal Hom
  (tag area 01CM): the section ring acts on a morphism of restricted modules by
  post-composition with multiplication by the section.}
  Let \(R\) be a presheaf of commutative rings on a base category \(\mathcal{C}\)
  that has a terminal object \(T\), and let \(M, N \in \mathtt{PresheafOfModules}\,R\).
  The abelian group of morphisms \(M \longrightarrow N\) of presheaves of \(R\)-modules
  carries a canonical \(R(T)\)-module structure --- the \emph{morphism module}
  \(\mathtt{homModule}\,M\,N\). A global scalar \(f \in R(T)\) acts on a morphism
  \(\varphi : M \to N\) by post-composition with multiplication by \(f\),
  \[
    f \cdot \varphi \;:=\; m_f^{N} \circ \varphi,
  \]
  where \(m_f^{N} : N \to N\) is the \emph{multiplication-by-\(f\)} endomorphism of
  \(N\): the endomorphism that, on each object \(U \in \mathcal{C}\), is multiplication
  by the restriction \(f|_U \in R(U)\) of \(f\) along the unique morphism \(U \to T\)
  to the terminal object. The module axioms hold because the assignment
  \(f \mapsto m_f^{N}\) is a ring homomorphism \(R(T) \to \mathrm{End}(N)\) into the
  endomorphism ring of \(N\) --- a product \(f \cdot g\) goes to the \emph{composite}
  \(m_f^{N} \circ m_g^{N}\) of the two multiplication endomorphisms, the order in which
  the factors compose being reversed by the fact that the scalar acts by
  \emph{post}-composition --- and composition of morphisms is biadditive in the
  preadditive category \(\mathtt{PresheafOfModules}\,R\). Distributivity and the
  associativity of the action then follow from bilinearity of composition together with
  the homomorphism property of \(f \mapsto m_f^{N}\). This is the genuinely new core of
  the construction: Mathlib supplies the \emph{fixed}-ring internal hom
  \(\mathtt{ihom}\,M\,N = (M \to N)\) over a single base ring, but nothing equipping the
  morphism group with the module structure for the \emph{varying} structure presheaf at
  the \(\mathtt{PresheafOfModules}\) level.
\end{definition}

\begin{definition}

\leanok
[Slice specialisation of the morphism module]
  \label{def:presheaf_internal_hom_slice_value}
  \lean{PresheafOfModules.InternalHom.internalHomObjModule}
  \uses{def:presheaf_internal_hom_value}
  Specialise \cref{def:presheaf_internal_hom_value} to a slice category. For an object
  (open) \(U\) of \(\mathcal{C} = \Opens X\), restricting along the forgetful functor
  \(\mathtt{Over.forget}\,U : \mathcal{C}/U \to \mathcal{C}\) --- the open-immersion
  restriction \(M \mapsto M|_U\) --- carries \(M, N\) to presheaves of modules
  \(M|_U, N|_U\) over the restricted structure ring on the slice \(\mathcal{C}/U\). The
  slice \(\mathcal{C}/U\) has the terminal object \(\mathtt{Over.mk}\,(\mathrm{id}_U)\),
  at which the restricted ring takes the value \(R(U)\); hence
  \cref{def:presheaf_internal_hom_value} makes the morphism group
  \(M|_U \longrightarrow N|_U\) an \(R(U)\)-module, the
  \(\mathtt{internalHomObjModule}\). This \(R(U)\)-module \((M|_U \to N|_U)\) is exactly
  the value \(\mathcal{H}om(M, N)(U)\) of \cref{def:presheaf_internal_hom}.
\end{definition}

\begin{definition}

\leanok
[Presheaf internal hom of \(\mathcal{O}_X\)-modules]
  \label{def:presheaf_internal_hom}
  \lean{PresheafOfModules.InternalHom.internalHom}
  \uses{def:scheme_modules_tensorobj, def:presheaf_internal_hom_value,
        def:presheaf_internal_hom_slice_value, lem:presheaf_internal_hom_restriction}
  % SOURCE: [Stacks Project], "Modules on Ringed Spaces", \S Internal Hom
  % (\texttt{section-internal-hom}, tag area 01CM)
  % (read from references/stacks-modules.tex, L3500--L3524).
  % SOURCE QUOTE:
  % "Let $(X, \mathcal{O}_X)$ be a ringed space.
  % Let $\mathcal{F}$, $\mathcal{G}$ be $\mathcal{O}_X$-modules.
  % Consider the rule
  % $$
  % U \longmapsto \Hom_{\mathcal{O}_X|_U}(\mathcal{F}|_U, \mathcal{G}|_U).
  % $$
  % It follows from the discussion in Sheaves, Section
  % \ref{sheaves-section-glueing-sheaves} that this is a sheaf of
  % abelian groups. In addition, given an element
  % $\varphi \in \Hom_{\mathcal{O}_X|_U}(\mathcal{F}|_U, \mathcal{G}|_U)$
  % and a section $f \in \mathcal{O}_X(U)$ then we can define
  % $f\varphi \in \Hom_{\mathcal{O}_X|_U}(\mathcal{F}|_U, \mathcal{G}|_U)$
  % by either precomposing with multiplication by $f$ on $\mathcal{F}|_U$
  % or postcomposing with multiplication by $f$ on $\mathcal{G}|_U$ (it gives
  % the same result). Hence we in fact get a sheaf of $\mathcal{O}_X$-modules.
  % We will denote this sheaf
  % $\SheafHom_{\mathcal{O}_X}(\mathcal{F}, \mathcal{G})$.
  % There is a canonical ``evaluation'' morphism
  % $$
  % \mathcal{F}
  % \otimes_{\mathcal{O}_X}
  % \SheafHom_{\mathcal{O}_X}(\mathcal{F}, \mathcal{G})
  % \longrightarrow
  % \mathcal{G}.
  % $$"
  \textit{Source: [Stacks Project], ``Modules on Ringed Spaces'', \S Internal Hom
  (tag area 01CM).}
  Let \(R\) be a presheaf of commutative rings on \(\Opens X\) and let
  \(M, N \in \mathtt{PresheafOfModules}\,R\). The \emph{presheaf internal hom}
  \(\mathcal{H}om(M, N)\) is the presheaf of \(R\)-modules whose value on an open
  \(U\) is the module of morphisms of restricted modules
  \[
    \mathcal{H}om(M, N)(U)
      \;:=\;
    \mathtt{ModuleCat.of}\;\bigl(R(U)\bigr)\;\bigl(M|_U \longrightarrow N|_U\bigr),
  \]
  the \(R(U)\)-module of \(R|_U\)-linear morphisms \(M|_U \to N|_U\) of presheaves of
  modules on \(U\), formed via the open-immersion restriction \(U \hookrightarrow X\).
  This value module is the slice specialisation
  \cref{def:presheaf_internal_hom_slice_value} of the morphism module
  \cref{def:presheaf_internal_hom_value}.

  Explicitly, \(\mathcal{H}om(M, N)\) is the presheaf of \(R\)-modules assembled from
  three pieces. \emph{(a) Value modules.} On each open \(U\) the value is the
  \(R(U)\)-module \((M|_U \to N|_U)\) of \cref{def:presheaf_internal_hom_value} /
  \cref{def:presheaf_internal_hom_slice_value}. \emph{(b) Restriction maps.} For an
  inclusion \(V \subseteq U\) the restriction map
  \(\mathcal{H}om(M, N)(U) \to \mathcal{H}om(M, N)(V)\) is the further-restriction map
  \cref{lem:presheaf_internal_hom_restriction}, \(\varphi \mapsto \varphi|_V\), which
  is additive and semilinear over the ring restriction \(R(g) : R(U) \to R(V)\)
  induced by \(g : V \to U\) --- the compatibility datum that the presheaf-of-modules
  assembly (\(\mathtt{PresheafOfModules.mk}\), via \(\mathtt{ofPresheaf}\)) consumes.
  \emph{(c) Functoriality.} Restriction along the identity \(U \to U\) is the identity
  map, and restriction along a composite \(V \xrightarrow{g} U \xrightarrow{h} W\) of
  inclusions \(V \subseteq U \subseteq W\) (the morphism \(g : V \to U\) followed by
  \(h : U \to W\) in \(\mathcal{C} = \Opens X\)) is the composite of the two
  restriction maps \(\mathcal{H}om(M,N)(W) \to \mathcal{H}om(M,N)(U) \to
  \mathcal{H}om(M,N)(V)\); both identity and composition are inherited from the
  functoriality of further restriction, since \(\mathtt{Over.map}\) is a functor. These three pieces
  make \(\mathcal{H}om(M, N)\) a genuine, covariant presheaf of \(R\)-modules; this
  assembled presheaf is the construction target \(\mathtt{PresheafOfModules.internalHom}\).

  The slice formula is forced by contravariance: the pointwise rule
  \(U \mapsto \Hom_{R(U)}(M(U), N(U))\) varies \emph{against} the restriction maps of
  \(M\) and so is not a covariant presheaf, whereas the module of morphisms of
  \emph{restricted} objects \(M|_U \to N|_U\) restricts covariantly. Thus, unlike the
  tensor product \cref{def:scheme_modules_tensorobj} --- which is covariant in the
  restriction maps and sheafifies directly --- the internal hom must be built through
  the slice presentation, and there is no presheaf-level dual to sheafify naively.
\end{definition}

\begin{lemma}

\leanok
[Restriction map of the presheaf internal hom]
  \label{lem:presheaf_internal_hom_restriction}
  \lean{PresheafOfModules.InternalHom.restrictionMap}
  \uses{def:presheaf_internal_hom_value, def:presheaf_internal_hom_slice_value}
  % SOURCE: [Stacks Project], "Modules on Ringed Spaces", \S Internal Hom
  % (\texttt{section-internal-hom}, tag area 01CM)
  % (read from references/stacks-modules.tex, L3502--L3508).
  % SOURCE QUOTE:
  % "Consider the rule
  % $$
  % U \longmapsto \Hom_{\mathcal{O}_X|_U}(\mathcal{F}|_U, \mathcal{G}|_U).
  % $$
  % It follows from the discussion in Sheaves, Section
  % \ref{sheaves-section-glueing-sheaves} that this is a sheaf of
  % abelian groups."
  \textit{Source: [Stacks Project], ``Modules on Ringed Spaces'', \S Internal Hom
  (tag area 01CM): the rule \(U \mapsto \Hom(\mathcal{F}|_U, \mathcal{G}|_U)\), with its
  further-restriction maps, is a (pre)sheaf of abelian groups.}
  Let \(g : V \to U\) be a morphism of \(\mathcal{C} = \Opens X\) (an inclusion
  \(V \subseteq U\) of opens), and let \(M, N \in \mathtt{PresheafOfModules}\,R\). There
  is a restriction map
  \[
    \mathcal{H}om(M, N)(U) \longrightarrow \mathcal{H}om(M, N)(V),
    \qquad \varphi \longmapsto \varphi|_V,
  \]
  the \(\mathtt{restrictionMap}\), sending a morphism of restricted modules
  \(M|_U \to N|_U\) to its further restriction \(M|_V \to N|_V\). Concretely, the
  functor \(\mathtt{Over.map}\,g : \mathcal{C}/V \to \mathcal{C}/U\) realises
  ``restrict from \(U\) to \(V\)'', and \(\varphi|_V\) is \(\varphi\) pulled back along
  it --- the further restriction of the same morphism of presheaves of modules. This
  map is additive, and it is \emph{semilinear} over the ring restriction
  \(R(g) : R(U) \to R(V)\): for \(f \in R(U)\),
  \[
    (f \cdot \varphi)\big|_V \;=\; R(g)(f) \cdot \bigl(\varphi|_V\bigr).
  \]
  This semilinearity is the compatibility datum that the presheaf assembly of
  \cref{def:presheaf_internal_hom} (through \(\mathtt{PresheafOfModules.ofPresheaf}\) /
  \(\mathtt{mk}\)) requires.
\end{lemma}

\begin{proof}
  \uses{def:presheaf_internal_hom_value}
  \leanok
  Further restriction along \(\mathtt{Over.map}\,g\) is functorial, so it carries a
  morphism \(M|_U \to N|_U\) to a morphism \(M|_V \to N|_V\) and respects identities and
  composites; in particular the induced map on morphism groups is additive. For the
  semilinearity, recall that the global-scalar action of
  \cref{def:presheaf_internal_hom_value} acts on \(\varphi\) by post-composition with
  the multiplication-by-\(f\) endomorphism \(m_f^{N}\); restricting that endomorphism
  from \(U\) to \(V\) is multiplication by \(f|_V = R(g)(f)\) on \(N|_V\). Since further
  restriction commutes with composition of morphisms, restricting \(f \cdot \varphi =
  m_f^{N} \circ \varphi\) yields \(m_{R(g)(f)}^{N|_V} \circ (\varphi|_V) = R(g)(f) \cdot
  (\varphi|_V)\), which is the asserted compatibility.
\end{proof}

\begin{definition}

\leanok
[Presheaf dual against the structure presheaf]
  \label{def:presheaf_dual}
  \lean{PresheafOfModules.dual}
  \uses{def:presheaf_internal_hom}
  With \(R\) and \(M\) as in \cref{def:presheaf_internal_hom}, the
  \emph{presheaf dual} of \(M\) is the internal hom into the unit (structure)
  presheaf,
  \[
    M^\vee \;:=\; \mathcal{H}om(M, R),
  \]
  i.e.\ \(M^\vee(U) = \mathtt{ModuleCat.of}\,(R(U))\,(M|_U \to R|_U)\), the
  \(R(U)\)-module of \(R|_U\)-linear functionals on \(M|_U\). This is the
  presheaf-of-modules analogue of the linear dual
  \(\mathtt{Module.Dual}\,R\,M = (M \to R)\) that serves as the \(\otimes\)-inverse of
  an invertible module over a fixed ring; here the fixed-ring dual is taken
  open-by-open and assembled by the slice presentation of
  \cref{def:presheaf_internal_hom}.
\end{definition}

\begin{lemma}

\leanok
[Evaluation/contraction of the internal hom]
  \label{lem:internal_hom_eval}
  \lean{PresheafOfModules.internalHomEval}
  % NOTE (iter-224): \texttt{PresheafOfModules.internalHomEval} is CLOSED axiom-clean
  % (\(\{\texttt{propext}, \texttt{Classical.choice}, \texttt{Quot.sound}\}\)). The full
  % morphism assembles the per-open contraction \texttt{internalHomEvalApp}
  % (the \(R(U)\)-bilinear \((M|_U) \otimes_{R(U)} (M|_U \to R|_U) \to R(U)\),
  % \(s \otimes \varphi \mapsto \varphi(s)\)) sectionwise; naturality follows from the
  % evaluate/restrict-commutation argument in the proof body. The iter-222/223 ``whnf
  % heartbeat-bomb'' obstacle was a stale Mathlib-version artifact and is gone.
  \uses{def:presheaf_dual, def:scheme_modules_tensorobj}
  % SOURCE: [Stacks Project], "Modules on Ringed Spaces", \S Internal Hom
  % (\texttt{section-internal-hom}, tag area 01CM)
  % (read from references/stacks-modules.tex, L3517--L3524).
  % SOURCE QUOTE:
  % "There is a canonical ``evaluation'' morphism
  % $$
  % \mathcal{F}
  % \otimes_{\mathcal{O}_X}
  % \SheafHom_{\mathcal{O}_X}(\mathcal{F}, \mathcal{G})
  % \longrightarrow
  % \mathcal{G}.
  % $$"
  \textit{Source: [Stacks Project], ``Modules on Ringed Spaces'', \S Internal Hom
  (tag area 01CM).}
  With \(R\) and \(M\) as above, there is a canonical \emph{evaluation} (contraction)
  morphism of presheaves of \(R\)-modules
  \[
    \mathrm{ev}_M \;\colon\; M \otimes_R M^\vee \;\longrightarrow\; R,
    \qquad s \otimes \phi \longmapsto \phi(s),
  \]
  the counit of the tensor--hom adjunction
  \(\mathcal{H}om(M \otimes_R -, R) \cong \mathcal{H}om(-, M^\vee)\) specialised at
  \(R\). On each open \(U\) it is the open-by-open contraction
  \((M|_U) \otimes_{R|_U} (M|_U \to R|_U) \to R|_U\), \(s \otimes \phi \mapsto
  \phi(s)\), and these are compatible with restriction, so they assemble to a morphism
  of presheaves of modules. The tensor product is the presheaf tensor underlying
  \cref{def:scheme_modules_tensorobj}.
\end{lemma}

\begin{proof}
  \uses{def:presheaf_dual, def:scheme_modules_tensorobj, lem:presheaf_internal_hom_restriction}
  \leanok
  On a fixed open \(U\), the value \(M^\vee(U) = (M|_U \to R|_U)\) is the module of
  \(R|_U\)-linear functionals on \(M|_U\), and the evaluation
  \(s \otimes \phi \mapsto \phi(s)\) is the standard contraction
  \((M|_U) \otimes_{R(U)} (M|_U \to R|_U) \to R(U)\) bilinear over \(R(U)\), hence a
  well-defined \(R(U)\)-linear map out of the tensor product. For \(V \subseteq U\)
  the square relating the \(U\)- and \(V\)-contractions commutes because evaluating
  then restricting equals restricting both arguments then evaluating
  (\(\phi(s)|_V = (\phi|_V)(s|_V)\)); hence the open-by-open contractions are natural
  in the restriction maps and define a morphism of presheaves of \(R\)-modules.
  Concretely, the morphism is defined by specifying, on each open \(U\), the per-open
  contraction \((M|_U) \otimes_{R(U)} (M|_U \to R|_U) \to R(U)\), \(s \otimes \phi \mapsto \phi(s)\).
  The single remaining datum is the naturality square, which reduces to the
  evaluate/restrict-commutation identity \(\phi(s)|_V = (\phi|_V)(s|_V)\) above: the
  restriction \(\phi|_V\) of a functional inside the dual \(M^\vee\) is its restriction
  to the open \(V\), and restriction is functorial, so evaluating the restricted
  functional on the restricted section agrees with restricting the value of the
  evaluation. This is the same restriction-functoriality argument already used to make
  the internal-hom restriction maps functorial in
  \cref{lem:presheaf_internal_hom_restriction}. This is the counit of the tensor--hom
  adjunction of \cref{def:presheaf_internal_hom} specialised to the unit presheaf.
\end{proof}

\begin{lemma}


\leanok
[The internal hom is a sheaf; the sheaf-level dual]
  \label{lem:internal_hom_isSheaf}
  \lean{AlgebraicGeometry.Scheme.Modules.dual}
  \uses{def:presheaf_dual, lem:internal_hom_eval}
  % SOURCE: [Stacks Project], "Modules on Ringed Spaces", \S Internal Hom
  % (\texttt{section-internal-hom}, tag area 01CM)
  % (read from references/stacks-modules.tex, L3502--L3514).
  % SOURCE QUOTE:
  % "Consider the rule
  % $$
  % U \longmapsto \Hom_{\mathcal{O}_X|_U}(\mathcal{F}|_U, \mathcal{G}|_U).
  % $$
  % It follows from the discussion in Sheaves, Section
  % \ref{sheaves-section-glueing-sheaves} that this is a sheaf of
  % abelian groups. ... Hence we in fact get a sheaf of $\mathcal{O}_X$-modules."
  \textit{Source: [Stacks Project], ``Modules on Ringed Spaces'', \S Internal Hom
  (tag area 01CM).}
  Let \(X\) be a scheme, \(R = X.\mathtt{ringCatSheaf}\), and let
  \(M, N \in \Scheme.\mathtt{Modules}\,X\) with \(N\) a sheaf of modules. Then the
  presheaf internal hom \(\mathcal{H}om(M, N)\) of \cref{def:presheaf_internal_hom}
  satisfies the sheaf condition: for an open cover \(\{U_i\}\) of \(U\), a family of
  morphisms \(M|_{U_i} \to N|_{U_i}\) agreeing on the overlaps
  \(U_i \cap U_j\) glues to a unique morphism \(M|_U \to N|_U\), because \(N\) is a
  sheaf and morphisms of sheaves of modules are determined and gluable section-wise.
  Hence \(\mathcal{H}om(M, N)\) descends to an object of \(\Scheme.\mathtt{Modules}\,X\).
  Specialising to \(N = \mathcal{O}_X = \mathtt{SheafOfModules.unit}\,R\) gives the
  \emph{sheaf-level dual}
  \[
    \mathtt{dual}\,M \;:=\; \mathcal{H}om_{\mathcal{O}_X}(M, \mathcal{O}_X)
    \;\in\; \Scheme.\mathtt{Modules}\,X,
  \]
  and the evaluation \cref{lem:internal_hom_eval} descends to a morphism of sheaves of
  modules \(M \otimes_X \mathtt{dual}\,M \to \mathcal{O}_X\).
\end{lemma}

\begin{proof}
  \uses{def:presheaf_dual, lem:internal_hom_eval}
  Fix an open \(U\) and an open cover \(\{U_i\}\) of \(U\). A morphism
  \(M|_U \to N|_U\) is a section of \(\mathcal{H}om(M, N)\) over \(U\); restricting it
  to the \(U_i\) gives a compatible family. Conversely, given morphisms
  \(\varphi_i : M|_{U_i} \to N|_{U_i}\) with \(\varphi_i|_{U_i \cap U_j} =
  \varphi_j|_{U_i \cap U_j}\), section-wise gluing in the sheaf \(N\) produces, for
  each open \(W \subseteq U\) and section \(s \in M(W)\), a well-defined section of
  \(N(W)\) by gluing the \(\varphi_i(s|_{W \cap U_i})\) (these agree on overlaps by
  hypothesis, and \(N\) is a sheaf so the glued section is unique). The resulting
  assignment is \(R(W)\)-linear and natural in \(W\), hence a morphism
  \(M|_U \to N|_U\) restricting to the \(\varphi_i\); uniqueness is the separatedness
  of \(N\). This is exactly the sheaf condition for \(\mathcal{H}om(M, N)\), so it is a
  sheaf of \(\mathcal{O}_X\)-modules. Taking \(N = \mathcal{O}_X\) yields
  \(\mathtt{dual}\,M\), and the evaluation morphism of \cref{lem:internal_hom_eval},
  being a morphism of presheaves into the sheaf \(\mathcal{O}_X\), is a morphism of
  the corresponding sheaves of modules.
\end{proof}

\begin{lemma}


\leanok
[Base change along a ring isomorphism commutes with the linear dual (H2$'$)]
  \label{lem:restrictscalars_ringiso_dualequiv}
  \lean{restrictScalarsRingIsoDualEquiv}
  \uses{lem:restrictscalars_ringiso_tensorequiv}
  Let \(e : R \simeq S\) be an isomorphism of commutative rings, and let \(M\) be an
  \(S\)-module. Restriction of scalars along \(e\) commutes with the formation of the
  linear dual: the canonical map
  \[
    \mathtt{restrictScalars}_e\bigl(M \to_S S\bigr)
      \;\xrightarrow{\ \sim\ }\;
    \bigl(\mathtt{restrictScalars}_e M \to_R R\bigr),
    \qquad \varphi \;\longmapsto\; e^{-1} \circ \varphi,
  \]
  is an \(R\)-linear equivalence, with inverse \(\psi \mapsto e \circ \psi\). On both
  sides the \(R\)-module structure is the \(S\)-module structure of \(M\) pulled back
  along \(e\), so that \(r \cdot m = e(r) \cdot m\): the left-hand side is the
  \(S\)-linear dual \((M \to_S S)\) regarded as an \(R\)-module by restriction, and the
  right-hand side is the \(R\)-linear dual of the restricted module
  \(\mathtt{restrictScalars}_e M\) into the restricted ground ring. This is the dual
  analogue of \cref{lem:restrictscalars_ringiso_tensorequiv}: it is to the internal hom
  / dual exactly what that lemma is to the tensor product. It is the \emph{H2$'$
  ingredient} of the restriction-compatibility isomorphism of the C-bridge
  \cref{lem:dual_isLocallyTrivial} --- the ring-iso/dual commutation that, composed with
  the open-immersion pushforward comparison H1, identifies the dual of a restriction
  with the restriction of the dual.
\end{lemma}

\begin{proof}

\leanok
  The two assignments \(\varphi \mapsto e^{-1} \circ \varphi\) and
  \(\psi \mapsto e \circ \psi\) are mutually inverse, since \(e \circ e^{-1} =
  \mathrm{id}_S\) and \(e^{-1} \circ e = \mathrm{id}_R\) as ring homomorphisms, so the
  two composites are the identity on functionals. Each assignment is \(R\)-linear:
  additivity is immediate, and for a scalar \(r \in R\) the defining relation
  \(r \cdot m = e(r) \cdot m\) of the pulled-back module structure makes each linearity
  check a single rewrite --- \(e^{-1}\) (resp.\ \(e\)) intertwines the \(R\)-action on
  the source with the \(R\)-action on the target precisely because \(e\) and \(e^{-1}\)
  are inverse ring homomorphisms identifying \(R\) with \(S\). The forward map is
  therefore an \(R\)-linear bijection, i.e.\ an \(R\)-linear equivalence; packaging the
  two one-sided maps with the inverse identities gives
  \(\mathtt{restrictScalarsRingIsoDualEquiv}\).
\end{proof}

\begin{definition}

\leanok
[Precomposition equivalence on dual sections]
  \label{def:presheaf_dual_precomp_equiv}
  \lean{PresheafOfModules.dualPrecompEquiv}
  \uses{def:presheaf_dual}
  Let \(R\) be a presheaf of commutative rings and let \(e : M \xrightarrow{\sim} M'\)
  be an isomorphism of presheaves of \(R\)-modules. On a fixed open \(U\), precomposition
  with \(e\) gives an \(R(U)\)-linear equivalence between the dual sections
  \[
    \mathtt{dualPrecompEquiv}\,e\,U \;\colon\;
      \bigl(M'|_U \to R|_U\bigr) \;\xrightarrow{\ \sim\ }\; \bigl(M|_U \to R|_U\bigr),
      \qquad \varphi \;\longmapsto\; \varphi \circ e|_U,
  \]
  with inverse \(\psi \mapsto \psi \circ (e^{-1})|_U\). This is the section-level core of
  the contravariant functoriality of the presheaf dual \(M \mapsto M^\vee\)
  (\cref{def:presheaf_dual}) in isomorphisms: precomposing a functional on \(M'\) with
  the comparison isomorphism produces a functional on \(M\), \(R(U)\)-linearly and
  invertibly.
\end{definition}

\begin{definition}

\leanok
[The presheaf dual carries an iso \(M \cong M'\) to an iso \(M'^\vee \cong M^\vee\)]
  \label{def:presheaf_dual_iso_of_iso}
  \lean{PresheafOfModules.dualIsoOfIso}
  \uses{def:presheaf_dual, def:presheaf_dual_precomp_equiv}
  An isomorphism \(e : M \xrightarrow{\sim} M'\) of presheaves of \(R\)-modules induces a
  canonical isomorphism of presheaf duals
  \[
    \mathtt{dualIsoOfIso}\,e \;\colon\; M'^\vee \;\xrightarrow{\ \sim\ }\; M^\vee
  \]
  --- note the reversal of direction, the dual being contravariant --- assembled
  open-by-open from the precomposition equivalences
  \(\mathtt{dualPrecompEquiv}\,e\,U\) of \cref{def:presheaf_dual_precomp_equiv}.
  Naturality in \(U\), i.e.\ compatibility of these sectionwise equivalences with
  restriction, is the routine check that precomposition by \(e\) commutes with
  restricting functionals.
\end{definition}

\begin{definition}


\leanok
[The sheaf-level dual carries an iso \(M \cong M'\) to an iso \(\mathtt{dual}\,M' \cong \mathtt{dual}\,M\)]
  \label{def:scheme_modules_dual_iso_of_iso}
  \lean{AlgebraicGeometry.Scheme.Modules.dualIsoOfIso}
  \uses{lem:internal_hom_isSheaf, def:presheaf_dual_iso_of_iso}
  Let \(X\) be a scheme and \(e : M \xrightarrow{\sim} M'\) an isomorphism in
  \(\Scheme.\mathtt{Modules}\,X\). Then there is a canonical isomorphism of sheaf-level
  duals
  \[
    \mathtt{dualIsoOfIso}\,e \;\colon\; \mathtt{dual}\,M' \;\xrightarrow{\ \sim\ }\;
      \mathtt{dual}\,M,
  \]
  obtained by sheafifying the presheaf isomorphism \(\mathtt{dualIsoOfIso}\,e\) of
  \cref{def:presheaf_dual_iso_of_iso} through the sheaf-level dual of
  \cref{lem:internal_hom_isSheaf}. It is the dual (contravariant) analogue of the
  functoriality of the tensor product in isomorphisms: where an isomorphism of a tensor
  factor induces an isomorphism of tensor products, an isomorphism \(M \cong M'\)
  induces an isomorphism of duals in the reverse direction.
\end{definition}

\begin{lemma}

\leanok
[Open-immersion slice-site sheaf equivalence (the shared slice bridge)]
  \label{lem:open_immersion_slice_sheaf_equiv}
  \lean{AlgebraicGeometry.Scheme.Modules.overSliceSheafEquiv}
  % SOURCE: Mathlib code-location provenance (this is NOT an external literature
  % citation, so it carries no verbatim mathematical SOURCE QUOTE). This lemma
  % completes the documented TODO at Mathlib/Topology/Sheaves/Over.lean, lines
  % 19--22, reproduced here for orientation:
  %   "show that both functors of the equivalence overEquivalence U are continuous
  %    and induce an equivalence between Sheaf ((Opens.grothendieckTopology X).over U) A
  %    and Sheaf (Opens.grothendieckTopology U) A for any category A."
  Let \(X\) be a scheme and \(U \subseteq X\) an open subset, with associated open
  immersion \(\iota : U \hookrightarrow X\). The reindexing equivalence
  \(\mathtt{overEquivalence}\,U : \mathtt{Over}\,U \simeq \mathtt{Opens}\,\widetilde{U}\)
  --- between the slice of \(\mathtt{Opens}\,X\) over \(U\) and the opens of the subspace
  \(\widetilde{U}\) --- upgrades to an \emph{equivalence of sheaf categories}
  \[
    \mathtt{Sheaf}\bigl((\mathtt{Opens.grothendieckTopology}\,X).\mathtt{over}\,U\bigr)\,A
      \;\;\simeq\;\;
    \mathtt{Sheaf}\bigl(\mathtt{Opens.grothendieckTopology}\,U\bigr)\,A
  \]
  for any target category \(A\), compatible with the module-pullback comparison
  \(\mathtt{restrictFunctorIsoPullback}\). This completes the Mathlib TODO recorded at
  \texttt{Topology/Sheaves/Over.lean} (lines 19--22). It is the \emph{shared root}
  underlying both the restriction-compatibility isomorphism for the dual
  (\cref{lem:dual_isLocallyTrivial}) and the morphism-descent / gluing engine
  (\cref{lem:sheafofmodules_hom_of_local_compat}): both reduce to identifying a
  slice-indexed family of sections over \(\mathtt{Over}\,U\) with an
  open-immersion-indexed family over \(\mathtt{Opens}\,\widetilde{U}\), and that
  identification is exactly this equivalence.
\end{lemma}

\begin{proof}

\leanok
  The functor \(\mathtt{overEquivalence}\,U\) realises \(\mathtt{Over}\,U\) as a
  \emph{dense subsite} of \(\mathtt{Opens}\,\widetilde{U}\): every open of the subspace
  \(\widetilde{U}\) is covered by down-sets \(\{V : V \le U\}\), and these down-sets are
  precisely the objects of the slice \(\mathtt{Over}\,U\). One therefore applies the
  dense-subsite transfer \(\mathtt{Functor.IsDenseSubsite.sheafEquiv}\): a continuous,
  cocontinuous, dense inclusion of sites induces an equivalence of the associated sheaf
  categories. The continuity and cover-lifting hypotheses are supplied by the slice-site
  instances of \(\mathtt{Sites/Over}\) (the Grothendieck topology
  \(J.\mathtt{over}\,U\) and its comparison with the topology pulled back along
  \(\iota\)).

  The over-category coherence that is laborious for a general site --- the pseudofunctor
  identities \(\mathtt{Over.mapId}\), \(\mathtt{Over.mapComp}\) and the associativity
  \(\mathtt{overMapPullback\_assoc}\) --- is here \emph{automatic}: because
  \(\mathtt{Opens}\,X\) is a thin poset, every hom-set is a subsingleton, so each such
  square commutes by \(\mathtt{Subsingleton.elim}\) (and its heterogeneous form
  \(\mathtt{Subsingleton.helim}\) where source and target objects differ only up to a
  propositional equality of opens). The sole sectionwise content is the down-set
  identification
  \[
    \iota_{!}\bigl(\iota^{-1}(V)\bigr) \;=\; V
    \qquad (V \le U),
  \]
  i.e.\ the image under the open immersion of the preimage of an open contained in \(U\)
  is that open again. Transporting sections along this equality and invoking the
  dense-subsite equivalence yields the stated sheaf equivalence, compatible with
  \(\mathtt{restrictFunctorIsoPullback}\) by the same thinness argument.
\end{proof}

\begin{lemma}


\leanok
[The per-section dual transport \(\mathtt{sliceDualTransport}\)]
  \label{lem:slice_dual_transport}
  \lean{AlgebraicGeometry.Scheme.Modules.sliceDualTransport}
  \uses{lem:internal_hom_isSheaf, lem:restrictscalars_ringiso_dualequiv,
        def:dual_unit_ring_swap, lem:isiso_eps_restrictscalars_appiso,
        def:dual_unit_iso_gen, lem:image_preimage_of_le, def:dual_unit_ring_swap_hom,
        lem:isiso_eps_restrictscalars_appiso_hom, lem:slice_dual_transport_inv,
        def:dual_unit_ring_swap_inv, lem:dual_unit_ring_swap_comp_inv,
        lem:dual_unit_ring_swap_inv_comp}
  Let \(f : Y \hookrightarrow X\) be an open immersion of schemes and let
  \(M \in \Scheme.\mathtt{Modules}\,X\). Write \(\alpha\) for the open-immersion
  structure-presheaf natural transformation with components
  \(\alpha_U = (f.\mathtt{appIso}\,U).\mathrm{inv}\), a \emph{\(\mathtt{CommRingCat}\)}-level
  ring isomorphism, and \(\beta = \mathtt{whiskerRight}\,\alpha\,
  (\mathtt{forget}_2\,\mathtt{CommRingCat}\,\mathtt{RingCat})\) for its
  \(\mathtt{RingCat}\)-level shadow, with component
  \(\beta_V \colon \mathcal{O}_Y(V) \xrightarrow{\sim} \mathcal{O}_X(fV)\)
  (\(fV := f.\mathtt{opensFunctor}(V)\)). Then for every open \(V \subseteq Y\) there is an
  \(\mathcal{O}_Y(V)\)-linear isomorphism
  \[
    \bigl((\mathtt{pushforward}\,\beta).\mathtt{obj}\,(\mathtt{dual}\,M.\mathtt{val})\bigr)(V)
      \;\xrightarrow{\ \sim\ }\;
    \bigl(\mathtt{dual}\,((\mathtt{pushforward}\,\beta).\mathtt{obj}\,M.\mathtt{val})\bigr)(V),
  \]
  realised as a single \(\mathtt{LinearEquiv.toModuleIso}\) that packages \emph{both} leg (A)
  --- the slice-Hom base-change reindexing across \(f.\mathtt{opensFunctor}\) --- and leg (B)
  --- the unit codomain ring-iso swap. The naturality of the section family in \(W\) has
  \emph{two} distinct parts: the underlying base-morphism uniqueness in
  \((\mathtt{Over}\,V.\mathrm{unop})^{\mathrm{op}}\) is \(\mathtt{Subsingleton.elim}\) (a thin
  poset has at most one inclusion between objects), but the accompanying equation of
  \(\mathcal{O}_Y(V)\)-module maps is a genuine, separate obligation: it is the
  \(\varepsilon\)-naturality of \(\mathtt{restrictScalars}\) along the structure-ring iso
  \(\beta_W\), which \(\mathtt{Subsingleton.elim}\) does \emph{not} discharge.
\end{lemma}

\begin{proof}
  \uses{lem:internal_hom_isSheaf, lem:restrictscalars_ringiso_dualequiv,
        def:dual_unit_ring_swap, lem:isiso_eps_restrictscalars_appiso,
        def:dual_unit_iso_gen, lem:image_preimage_of_le, def:dual_unit_ring_swap_hom,
        lem:isiso_eps_restrictscalars_appiso_hom, lem:slice_dual_transport_inv,
        def:dual_unit_ring_swap_inv, lem:dual_unit_ring_swap_comp_inv,
        lem:dual_unit_ring_swap_inv_comp, def:restrict_scalars_lax_epsilon}
        \leanok
  The isomorphism is produced from an \(\mathcal{O}_Y(V)\)-linear equivalence between the
  left-hand module --- the \(\mathcal{O}_Y(V)\)-module inherited from
  \(\mathtt{pushforward}\,\beta\) --- and the right-hand \(\mathtt{internalHomObjModule}\),
  assembled via \(\mathtt{LinearEquiv.toModuleIso}\), with both module structures
  identified explicitly.

  \emph{Leg (A): the slice-Hom reindexing, built categorically.} A dual section
  \(\varphi\) of the left-hand value is a morphism \(\mathtt{restr}_{fV}\,M.\mathtt{val} \to
  \mathtt{restr}_{fV}\,\mathbf{1}_X\) over \(\mathtt{Over}\,(fV)\). Its image is the dual
  section over \(\mathtt{Over}\,V\) whose component at \(W \le V\) (with image
  \(fW := f.\mathtt{opensFunctor}(W)\)) is the categorical image
  \[
    (\mathtt{restrictScalars}\,\beta_W).\mathrm{map}\,
      \bigl(\varphi_{(\mathtt{Over.mk}\,(f.\mathtt{opensFunctor}(W \le V)))}\bigr),
  \]
  i.e.\ the component of \(\varphi\) at the \(f\)-image of \(W\), reindexed across
  \(f.\mathtt{opensFunctor}\) by restriction of scalars along \(\beta_W\), and then
  post-composed with the leg-(B) codomain swap below.

  \emph{Leg (B): the unit codomain ring-iso swap, via the lax-monoidal unit \(\varepsilon\).}
  The codomain of the reindexed component is the restriction of scalars of the unit section,
  which must be carried to the unit section over \(Y\). This is exactly the inverse of the
  lax-monoidal unit comparison of the restriction-of-scalars functor,
  \[
    \mathtt{codomainMap} \;:=\;
      \mathtt{inv}\bigl(\varepsilon(\mathtt{restrictScalars}\,g)\bigr),
      \qquad g := (f.\mathtt{appIso}\,W).\mathrm{inv}.\mathrm{hom},
  \]
  (\cref{def:dual_unit_ring_swap})
  where \(g\) is taken at the \emph{\(\mathtt{CommRingCat}\)} level. The unit \(\varepsilon\)
  is invertible because \(g\) is a ring isomorphism
  (\cref{lem:isiso_eps_restrictscalars_appiso}),
  so \(\varepsilon(\mathtt{restrictScalars}\,g)\) is an isomorphism. The presheaf-level
  identification underlying the unit handling is the dual-of-unit isomorphism for the unit
  section (\cref{def:dual_unit_iso_gen}).

  \emph{Inverse.} The inverse reindexing depends on a set-theoretic down-set fact: because
  \(f\) is an open immersion, \(fV \subseteq \mathrm{range}\,f\), so every open
  \(W'' \le fV\) already lies in the image of \(f.\mathtt{opensFunctor}\), namely
  \(W'' = f.\mathtt{opensFunctor}(f^{-1}W'')\). This is the open-image down-set coverage of an
  \(\mathtt{IsOpenImmersion}\): an open contained in the (open) image of the immersion is the
  functor-image of its own preimage. The project records it as the down-set identity
  \(\iota_!(\iota^{-1}V) = V\) for \(V \le W\) (the helper
  \(\mathtt{image\_preimage\_of\_le}\) (\cref{lem:image_preimage_of_le})),
  and it is exactly the lemma the inverse reindexing relies on to define the components on the
  whole down-set of \(fV\).

  With this in hand the inverse \(\mathtt{PresheafOfModules.Hom}\) mirrors the forward
  component formula, with two changes: it uses \((f.\mathtt{appIso}\,W'').\mathrm{hom}\) in
  place of its inverse, and its \(\varepsilon\)-codomain swap is
  \(\mathtt{dualUnitRingSwapHom} = \mathtt{inv}\bigl(\varepsilon(\mathtt{restrictScalars}\,
  (f.\mathtt{appIso}\,W'').\mathrm{hom})\bigr)\) (\cref{def:dual_unit_ring_swap_hom})
  --- that is, the \emph{inverse} of \(\varepsilon\) in the \(\mathrm{hom}\)-direction (where
  \(\varepsilon\) is invertible by
  \cref{lem:isiso_eps_restrictscalars_appiso_hom}), \emph{not}
  \(\varepsilon\) of the \(\mathrm{inv}\)-direction \(\beta\). (The reverse reindex turns an
  \(\mathcal{O}_X(fW'')\)-section map into an \(\mathcal{O}_Y(W'')\)-map, so it must restrict
  along the \(\mathrm{hom}\)-direction ring iso, whose codomain swap is therefore \(\mathtt{inv}\,
  \varepsilon\) of \emph{that} functor.) The whole family is reindexed by the inverse of the
  down-set bijection \(W'' \leftrightarrow f^{-1}W''\) above. The round-trip identities \(\mathtt{left\_inv}\)
  and \(\mathtt{right\_inv}\) then close component-wise by \(\mathtt{Iso.inv\_hom\_id}\) and
  \(\mathtt{Iso.hom\_inv\_id}\) of \(f.\mathtt{appIso}\) (the forward/inverse \(\beta\)-swaps
  cancel) together with the cancellation of the down-set bijection
  \(\iota_!(\iota^{-1}V) = V\) against its inverse.

  For additivity and \(\mathcal{O}_Y(V)\)-linearity, one must be careful about \emph{which}
  module structure carries the left-hand value. The additive and scalar structure on a dual
  section of the left-hand value
  \(\bigl((\mathtt{pushforward}\,\beta).\mathtt{obj}\,(\mathtt{dual}\,M.\mathtt{val})\bigr)(V)\)
  is the \(\mathcal{O}_Y(V)\)-module structure on the internal-hom object
  (\cref{def:presheaf_internal_hom_slice_value}) --- the pointwise module structure in which the
  sum and scalar multiple of dual sections are formed at the terminal object \(\mathrm{id}_{fV}\)
  of the slice over \(fV\). This agrees with, but is \emph{a priori} a distinct presentation of,
  the additive and scalar structure of the underlying presheaf morphism (the pointwise add and
  scalar action on morphisms of presheaves of modules). The verification of additivity and
  \(\mathcal{O}_Y(V)\)-linearity of the transport therefore proceeds in three distinct steps, the
  first of which is a required first move rather than an afterthought:
  \begin{enumerate}
    \item[(i)] \emph{Identify the two module presentations.} The internal-hom module structure's
      addition and scalar action are identified with the underlying pointwise add and scalar action
      on presheaf morphisms. This is a definitional identification of the two
      \(\mathcal{O}_Y(V)\)-module presentations of the same group of dual sections: the
      internal-hom operations, formed at the terminal object of the slice, are precisely the
      pointwise morphism-level operations of the presheaf morphism they package. Until this
      identification is made, the additivity and scalar-compatibility obligations of the transport
      cannot be read off the change-of-rings compatibility, because their sums and scalar multiples
      are taken in the \emph{internal-hom} structure, not the morphism structure.
    \item[(ii)] \emph{Match the X-side and Y-side scalars by the \(\beta\)-naturality ring
      identity.} The scalar acting on the X-side (the pushforward-restricted scalar \(s\) by
      which \(\mathtt{restrictScalars}\,\beta_W\) multiplies the reindexed component) and the
      scalar acting on the Y-side (the \(\mathcal{O}_Y(V)\)-scalar \(c\) of the internal-hom
      action) are \emph{a priori} elements of different rings and cannot be compared directly.
      The bridge is the ring identity
      \[
        s \;=\; (\beta.\mathtt{app}\,W').\mathrm{hom}\, c,
      \]
      obtained from \(\mathtt{InternalHom.termRingMap\_naturality}\) --- which records how the
      internal-hom term-ring map transports the scalar action across the slice --- together with
      the naturality of \(\beta\) on the thin poset \(\mathrm{Opens}\,Y\). Without this identity
      the two sides of the scalar-compatibility equation are not even of matching type, so it
      must be supplied before the change-of-rings compatibility of step (iii) can be invoked.
      The restricted scalar action is made explicit by
      \(\mathtt{ModuleCat.restrictScalars.smul\_def'}\), which exhibits the restricted smul
      as the original smul along \(\beta_W\).
    \item[(iii)] \emph{Apply the change-of-rings compatibility on the morphism level.} Once the
      operations are expressed pointwise on the underlying morphisms and the scalars matched by
      (ii), additivity and
      \(\mathcal{O}_Y(V)\)-linearity follow from the post-composition compatibility of the
      change-of-rings reindexing: the functor \(\mathtt{restrictScalars}\,\beta_W\) preserves sums
      and scalar multiples of morphisms, and the structure scalar is intertwined by the ring
      isomorphism \(\beta_W\). This last is the presheaf-level shadow of
      \cref{lem:restrictscalars_ringiso_dualequiv}.
  \end{enumerate}

  \emph{Naturality of the section family in \(W\): the two-leg horizontal paste.} The
  naturality square --- that the two routes around any \(W' \le W \le V\) square in
  \((\mathtt{Over}\,V.\mathrm{unop})^{\mathrm{op}}\) agree --- is the \emph{horizontal paste} of
  two independent legs, and it is a mistake to discharge the whole square by
  \(\mathtt{Subsingleton.elim}\) alone.

  \emph{leg-A --- the slice-Hom base-change reindex.} The first leg is the naturality of the
  reindexed dual section \(\varphi\) itself, transported across the \(\mathtt{Over}\)-slice
  morphism induced by \(f.\mathtt{opensFunctor}\) (the \(\mathtt{homLocalSection}\)-style reindex
  of leg (A)). Concretely it is \(\varphi.\mathtt{naturality}\) pushed along the slice map, with
  the underlying \emph{base} morphisms --- the indexing inclusions of
  \((\mathtt{Over}\,(fV).\mathrm{unop})^{\mathrm{op}}\) --- pinned by
  \(\mathtt{Subsingleton.elim}\), since \(\mathrm{Opens}\,X\) is a thin poset and there is at
  most one inclusion between any two objects. This leg carries \emph{no} \(\varepsilon\)
  content: it is pure base-change reindexing of the source section, and \(\mathtt{Subsingleton.
  elim}\) settles exactly its share --- the uniqueness of the base morphisms --- and nothing
  more.

  \emph{leg-B --- the \(\varepsilon\)-square.} The second leg is the naturality identity of the
  lax-monoidal unit \(\varepsilon\) of \(\mathtt{restrictScalars}\), packaged as the natural
  transformation \(\mathtt{restrictScalarsLax}\varepsilon\)
  (\cref{def:restrict_scalars_lax_epsilon}), whose \(\mathtt{NatTrans}\) naturality field is
  \emph{exactly} the square asserting that \(\varepsilon\) commutes with the restriction maps of
  \emph{both} the source and the target presheaf along the structure-ring iso \(\beta_W\). This
  is the genuine, separate module-map obligation that \(\mathtt{Subsingleton.elim}\) does
  \emph{not} settle: once leg-A's base morphisms are pinned, the residual is an equation of
  \(\mathcal{O}_Y(V)\)-\emph{module} maps, and that equation is precisely the
  \(\varepsilon\)-naturality square post-composed with the definition of the codomain swap
  \(\mathtt{dualUnitRingSwap} = \mathtt{inv}\,\varepsilon(\mathtt{restrictScalars}\,\beta_W)\)
  (\cref{def:dual_unit_ring_swap}).
  % NOTE: the Lean helper that delivers leg-B is
  % `PresheafOfModules.restrictScalarsLaxε`.

  \emph{Assembly discipline --- what makes the paste close.} The two legs compose to the
  required square only if leg-B's equation is \emph{derived from the \(\varepsilon\)-naturality
  identity itself}, never re-asserted by an independent hand-written expression. That is: the
  monoidal-unit subterm \(\mathbf 1\) and the restriction-of-scalars functor subterm appearing
  in leg-B must be \emph{inherited from} the components of
  \(\mathtt{restrictScalarsLax}\varepsilon\) --- read off its naturality field, whose components
  are the lax units
  \(\mathtt{Functor.LaxMonoidal.}\varepsilon(\mathtt{ModuleCat.restrictScalars}\,
  (\alpha.\mathrm{app}\,X).\mathrm{hom})\) --- and never re-introduced by spelling out a fresh
  \(\varepsilon(\mathtt{restrictScalars}\,\dots)\) by hand. When leg-B is sourced this way the
  entire \(\varepsilon\)/unit/restriction content of the two routes matches on the nose, and the
  \emph{only} residual discrepancy is a reconciliation of two \emph{definitionally equal but
  syntactically distinct spellings}:
  \begin{itemize}
    \item the presheaf-level versus sheaf-level restriction-of-scalars functor --- two
      presentations of the same change-of-rings functor along \(\beta_W\); and
    \item two elaborations of the monoidal unit \(\mathbf 1\) (\(\mathtt{tensorUnit}\)) --- as it
      elaborates in the dual-section codomain versus as it elaborates inside the
      \(\varepsilon\)-component.
  \end{itemize}
  Each of these is a \emph{defeq leaf}: the discrepancy is discharged by congruence on that
  leaf, NOT by supplying any missing natural transformation or monoidal structure. We emphasise
  that this is \emph{bounded definitional bookkeeping} --- a finite reconciliation of two
  spellings of the same term --- and not a structural gap: every piece of \(\varepsilon\), unit,
  and restriction-of-scalars data is already present in
  \(\mathtt{restrictScalarsLax}\varepsilon\), so nothing remains to be \emph{constructed}, only
  to be \emph{recognised as equal}.

  \emph{The round-trips \(\mathtt{left\_inv}\) and \(\mathtt{right\_inv}\).} The two round-trips
  of the equivalence \(\mathtt{sliceDualTransport}\) reduce, by \(\mathtt{hom\_ext}\) working
  section-by-section over the slice, to per-component telescopes that \emph{bypass the
  \(\varepsilon\)-square entirely}: a round-trip composes a forward component with its reverse,
  and the leg-(B) swaps cancel against one another \emph{before} any \(\varepsilon\)-commutation
  with restriction maps is ever invoked. Each telescope then collapses by the
  \(\varepsilon\)-cancellation identities already proved --- the swap/reverse cancellations
  \(\mathtt{dualUnitRingSwap} \circ \mathtt{dualUnitRingSwapInv} = \mathrm{id}\) and
  \(\mathtt{dualUnitRingSwapInv} \circ \mathtt{dualUnitRingSwap} = \mathrm{id}\)
  (\cref{lem:dual_unit_ring_swap_comp_inv}, \cref{lem:dual_unit_ring_swap_inv_comp}) ---
  together with the \(\mathtt{appIso}\) hom/inv cancellation
  (\(\mathtt{Iso.hom\_inv\_id}\) and \(\mathtt{Iso.inv\_hom\_id}\) of \(f.\mathtt{appIso}\), by
  which the forward \(\beta_W^{-1}\) and reverse \(\beta_W\) ring-iso swaps annihilate) and the
  \(\mathtt{eqToHom}\)/\(\mathtt{restrictScalarsId'App}\) identity collapse that absorbs the
  \(\mathtt{eqToHom}\) transports across the cross-fiber down-set relabel
  \(W'' \leftrightarrow f^{-1}W''\). No new infrastructure is required: this is the same
  \(\varepsilon\)-cancellation and identity-collapse used to verify the \(\mathtt{app}\) field,
  applied once more under the section-by-section \(\mathtt{hom\_ext}\).
\end{proof}

\subsection{Dual-inverse lane: the slice-transport helper declarations}
\label{subsec:dual_inverse_helpers}

These are the individual ingredients out of which \cref{lem:slice_dual_transport} and its
inverse leg are assembled: the leg-(B) unit-ring swaps, their invertibility and cancellation
identities, and the presheaf-level dual-of-unit identification. Each is a self-contained
internal construction.

\begin{lemma}\leanok
[Leg-(B) unit \(\varepsilon\) is invertible, \(\mathrm{inv}\)-direction]
  \label{lem:isiso_eps_restrictscalars_appiso}
  \lean{AlgebraicGeometry.Scheme.Modules.isIso_ε_restrictScalars_appIso}
  \uses{def:restrict_scalars_lax_epsilon}
  Let \(f : Y \hookrightarrow X\) be an open immersion and \(W' \subseteq Y\) an open. The
  lax-monoidal unit \(\varepsilon\) of the change-of-rings functor
  \(\mathtt{restrictScalars}\,\bigl((f.\mathtt{appIso}\,W').\mathrm{inv}.\mathrm{hom}\bigr)\),
  taken along the \(\mathtt{CommRingCat}\)-level structure-ring isomorphism, is an
  isomorphism.
\end{lemma}
\begin{proof}
  \uses{def:restrict_scalars_lax_epsilon}
  \leanok
  Proved directly in Lean --- the underlying ring map
  \((f.\mathtt{appIso}\,W').\mathrm{inv}.\mathrm{hom}\) is an isomorphism, hence bijective, so
  its restriction-of-scalars unit \(\varepsilon\) (\cref{def:restrict_scalars_lax_epsilon}) is
  invertible by the bijectivity criterion for \(\varepsilon\).
\end{proof}

\begin{lemma}\leanok
[Leg-(B) unit \(\varepsilon\) is invertible, \(\mathrm{hom}\)-direction]
  \label{lem:isiso_eps_restrictscalars_appiso_hom}
  \lean{AlgebraicGeometry.Scheme.Modules.isIso_ε_restrictScalars_appIso_hom}
  \uses{def:restrict_scalars_lax_epsilon}
  In the same setting, the lax-monoidal unit \(\varepsilon\) of
  \(\mathtt{restrictScalars}\,\bigl((f.\mathtt{appIso}\,W').\mathrm{hom}.\mathrm{hom}\bigr)\)
  --- the \(\mathrm{hom}\)-direction of the structure-ring iso --- is likewise an
  isomorphism. This is the variant powering the reverse (\(\mathtt{invFun}\)) leg, which
  reindexes a \(\mathcal{O}_Y\)-section map back to a \(\mathcal{O}_X\)-map.
\end{lemma}
\begin{proof}
  \uses{def:restrict_scalars_lax_epsilon}
  \leanok
  Proved directly in Lean --- identical to \cref{lem:isiso_eps_restrictscalars_appiso} with
  the \(\mathrm{hom}\)-direction (also bijective) ring map.
\end{proof}

\begin{lemma}\leanok
[Element action of \(\mathrm{inv}\,\varepsilon\) for restriction of scalars along a bijection]
  \label{lem:eps_inv_apply}
  \lean{AlgebraicGeometry.Scheme.Modules.εInv_apply}
  \uses{def:restrict_scalars_lax_epsilon, lem:isiso_eps_restrictscalars_appiso}
  Let \(g : R \to S\) be a bijective \(\mathtt{CommRingCat}\)-homomorphism, so that the
  lax-monoidal unit \(\varepsilon\) of \(\mathtt{restrictScalars}\,g\) is an isomorphism
  (\cref{lem:isiso_eps_restrictscalars_appiso}). Then on elements the inverse unit acts as the
  inverse ring bijection: for \(s \in S\),
  \[
    \bigl(\mathrm{inv}\,(\varepsilon(\mathtt{restrictScalars}\,g))\bigr)(s)
      \;=\;
    (\mathtt{RingEquiv.ofBijective}\,g)^{-1}(s).
  \]
  This is the reusable element-level ingredient for every \(\varepsilon\)-swap cancellation in the
  dual-inverse lane (the round-trips and the naturality \(\varepsilon\)-squares), where the swaps
  \(\mathtt{dualUnitRingSwap}\) etc.\ are all such \(\mathrm{inv}\,\varepsilon\) maps.
\end{lemma}
\begin{proof}
  \uses{def:restrict_scalars_lax_epsilon}
  \leanok
  Proved directly in Lean. From \(\varepsilon \mathbin{;} \mathrm{inv}\,\varepsilon = \mathbf{1}\)
  applied at \(r = (\mathtt{RingEquiv.ofBijective}\,g)^{-1}(s)\), the underlying-element identity
  \(\mathtt{ModuleCat.restrictScalars\_\eta}\) (the underlying map of \(\varepsilon\) is \(g\)) gives
  \((\mathrm{inv}\,\varepsilon)\bigl(g((\mathtt{RingEquiv.ofBijective}\,g)^{-1}(s))\bigr)
  = (\mathtt{RingEquiv.ofBijective}\,g)^{-1}(s)\); rewriting \(g \circ g^{-1} = \mathrm{id}\) by
  \(\mathtt{apply\_symm\_apply}\) yields the claim.
\end{proof}

\begin{definition}\leanok
[Leg-(B) codomain unit ring-iso swap, \(\mathrm{inv}\)-direction]
  \label{def:dual_unit_ring_swap}
  \lean{AlgebraicGeometry.Scheme.Modules.dualUnitRingSwap}
  \uses{lem:isiso_eps_restrictscalars_appiso}
  The codomain unit ring-iso swap
  \[
    \mathtt{dualUnitRingSwap}\,f\,W'
      \;:=\; \mathtt{inv}\bigl(\varepsilon(\mathtt{restrictScalars}\,
        (f.\mathtt{appIso}\,W').\mathrm{inv}.\mathrm{hom})\bigr)
    \;\colon\;
    (\mathtt{restrictScalars}\,(f.\mathtt{appIso}\,W')^{-1})(\mathbf 1_X(fW'))
      \;\longrightarrow\; \mathbf 1_Y(W'),
  \]
  the inverse of the (invertible, \cref{lem:isiso_eps_restrictscalars_appiso}) lax-monoidal
  unit \(\varepsilon\), carrying the restriction-of-scalars of the \(X\)-unit section to the
  \(Y\)-unit section. It is the leg-(B) ingredient of \cref{lem:slice_dual_transport}.
\end{definition}
\begin{proof}
  \uses{lem:isiso_eps_restrictscalars_appiso}
  \leanok
  Proved directly in Lean --- the categorical inverse of an isomorphism.
\end{proof}

\begin{definition}\leanok
[Leg-(B) unit ring-iso swap, reverse (\(\varepsilon\) itself)]
  \label{def:dual_unit_ring_swap_inv}
  \lean{AlgebraicGeometry.Scheme.Modules.dualUnitRingSwapInv}
  \uses{lem:isiso_eps_restrictscalars_appiso}
  The reverse map
  \(\mathtt{dualUnitRingSwapInv}\,f\,W' := \varepsilon(\mathtt{restrictScalars}\,
  (f.\mathtt{appIso}\,W').\mathrm{inv}.\mathrm{hom})\) --- the lax-monoidal unit
  \(\varepsilon\) \emph{itself} (not its inverse) --- from \(\mathbf 1_Y(W')\) back to the
  restriction-of-scalars of the \(X\)-unit section. It is the two-sided inverse of
  \cref{def:dual_unit_ring_swap}.
\end{definition}
\begin{proof}
  \uses{lem:isiso_eps_restrictscalars_appiso}
  \leanok
  Proved directly in Lean --- the unit \(\varepsilon\), invertible by
  \cref{lem:isiso_eps_restrictscalars_appiso}.
\end{proof}

\begin{definition}\leanok
[Reverse leg codomain unit ring-iso swap, \(\mathrm{hom}\)-direction]
  \label{def:dual_unit_ring_swap_hom}
  \lean{AlgebraicGeometry.Scheme.Modules.dualUnitRingSwapHom}
  \uses{lem:isiso_eps_restrictscalars_appiso_hom}
  The \(\mathrm{hom}\)-direction codomain swap
  \[
    \mathtt{dualUnitRingSwapHom}\,f\,W'
      \;:=\; \mathtt{inv}\bigl(\varepsilon(\mathtt{restrictScalars}\,
        (f.\mathtt{appIso}\,W').\mathrm{hom}.\mathrm{hom})\bigr)
    \;\colon\;
    (\mathtt{restrictScalars}\,(f.\mathtt{appIso}\,W').\mathrm{hom})(\mathbf 1_Y(W'))
      \;\longrightarrow\; \mathbf 1_X(fW'),
  \]
  the inverse of the (invertible, \cref{lem:isiso_eps_restrictscalars_appiso_hom}) unit
  \(\varepsilon\) in the \(\mathrm{hom}\)-direction. It is the leg-(B) ingredient of the
  reverse transport \cref{lem:slice_dual_transport_inv}.
\end{definition}
\begin{proof}
  \uses{lem:isiso_eps_restrictscalars_appiso_hom}
  \leanok
  Proved directly in Lean --- the categorical inverse of an isomorphism.
\end{proof}

\begin{lemma}


\leanok
[Presheaf-level relabel unit \(\varepsilon\) is invertible]
  \label{lem:isiso_eps_restrictscalars_presheafmap}
  \lean{AlgebraicGeometry.Scheme.Modules.isIso_ε_restrictScalars_presheafMap}
  \uses{def:restrict_scalars_lax_epsilon}
  \textit{Source: internal categorical construction; no external reference.}
  Let \(f : Y \hookrightarrow X\) be an open immersion. The lax-monoidal unit
  \(\varepsilon\) of \(\mathtt{restrictScalars}\) along the presheaf-level structure-ring map
  induced by the propositional down-set identity
  \(f.\mathtt{opensFunctor}(f^{-1}(W')) = W'\) is an isomorphism. This is the
  presheaf-map analogue of \cref{lem:isiso_eps_restrictscalars_appiso}, supplying the
  unit-codomain transport of the reverse leg
  \cref{lem:slice_dual_transport_inv}.
\end{lemma}
\begin{proof}
  \uses{def:restrict_scalars_lax_epsilon}
  \leanok
  Proved directly in Lean --- the underlying presheaf structure-ring map is an isomorphism,
  hence bijective, so its restriction-of-scalars unit \(\varepsilon\)
  (\cref{def:restrict_scalars_lax_epsilon}) is invertible by the bijectivity criterion.
\end{proof}

\begin{definition}


\leanok
[Cross-fiber unit-codomain relabel swap]
  \label{def:unit_relabel_swap}
  \lean{AlgebraicGeometry.Scheme.Modules.unitRelabelSwap}
  \uses{lem:isiso_eps_restrictscalars_presheafmap}
  \textit{Source: internal categorical construction; no external reference.}
  The unit-codomain transport
  \(\mathtt{unitRelabelSwap} := \mathtt{inv}\,\varepsilon\) realizing the cross-fiber relabel
  of the reverse component \(\mathtt{sliceDualTransportInv}.\mathtt{app}\): it carries the
  restriction-of-scalars of the unit section across the propositional down-set identity
  \(f.\mathtt{opensFunctor}(f^{-1}(W')) = W'\), as the inverse of the (invertible,
  \cref{lem:isiso_eps_restrictscalars_presheafmap}) lax-monoidal unit \(\varepsilon\). It is
  the presheaf-map mirror of \cref{def:dual_unit_ring_swap_hom}.
\end{definition}
\begin{proof}
  \uses{lem:isiso_eps_restrictscalars_presheafmap}
  \leanok
  Proved directly in Lean --- the categorical inverse of an isomorphism.
\end{proof}

\begin{lemma}\leanok
[Round-trip of the swap pair, \(\mathrm{inv}\circ\mathrm{swap}\)]
  \label{lem:dual_unit_ring_swap_comp_inv}
  \lean{AlgebraicGeometry.Scheme.Modules.dualUnitRingSwap_comp_dualUnitRingSwapInv}
  \uses{def:dual_unit_ring_swap, def:dual_unit_ring_swap_inv}
  \(\mathtt{dualUnitRingSwap}\,f\,W' \circ \mathtt{dualUnitRingSwapInv}\,f\,W' =
  \mathrm{id}\): the swap followed by its reverse is the identity.
\end{lemma}
\begin{proof}
  \uses{def:dual_unit_ring_swap, def:dual_unit_ring_swap_inv}
  \leanok
  Proved directly in Lean --- the \(\mathtt{Iso.inv\_hom\_id}\) cancellation of
  \(\varepsilon\) with its inverse.
\end{proof}

\begin{lemma}\leanok
[Round-trip of the swap pair, \(\mathrm{swap}\circ\mathrm{inv}\)]
  \label{lem:dual_unit_ring_swap_inv_comp}
  \lean{AlgebraicGeometry.Scheme.Modules.dualUnitRingSwapInv_comp_dualUnitRingSwap}
  \uses{def:dual_unit_ring_swap, def:dual_unit_ring_swap_inv}
  \(\mathtt{dualUnitRingSwapInv}\,f\,W' \circ \mathtt{dualUnitRingSwap}\,f\,W' =
  \mathrm{id}\): the reverse followed by the swap is the identity.
\end{lemma}
\begin{proof}
  \uses{def:dual_unit_ring_swap, def:dual_unit_ring_swap_inv}
  \leanok
  Proved directly in Lean --- the \(\mathtt{Iso.hom\_inv\_id}\) cancellation of
  \(\varepsilon\) with its inverse.
\end{proof}

\begin{definition}\leanok
[Section equivalence for the dual of the unit]
  \label{def:unit_dual_section_equiv}
  \lean{PresheafOfModules.unitDualSectionEquiv}
  \uses{def:eval_lin, def:global_smul, def:presheaf_internal_hom_slice_value}
  For a presheaf of modules over a base ring presheaf \(R_0\) and an object \(X\), the
  \(R_0(X)\)-linear equivalence
  \[
    \bigl(\mathtt{restr}_X\,\mathbf 1 \to \mathtt{restr}_X\,\mathbf 1\bigr)
      \;\simeq_{R_0(X)}\; R_0(X)
  \]
  identifying endomorphisms of the (restricted) monoidal unit with the ground ring, via
  evaluation at the global section \(1\); the inverse is multiplication by a global scalar
  (\cref{def:global_smul}).
\end{definition}
\begin{proof}
  \uses{def:eval_lin, def:global_smul, def:presheaf_internal_hom_slice_value}
  \leanok
  Proved directly in Lean. The forward map is \(\mathtt{evalLin}\) at \(1\)
  (\cref{def:eval_lin}); the inverse is the global-scalar action \(\mathtt{globalSMul}\)
  (\cref{def:global_smul}). The substantive field is \(\mathtt{left\_inv}\): every
  endomorphism of the unit is multiplication by its value at \(1\), which follows from the
  \(\varphi\)-naturality of the section toward the terminal object of the slice. The module
  structure on the hom-space is the slice internal-hom structure
  (\cref{def:presheaf_internal_hom_slice_value}).
\end{proof}

\begin{definition}\leanok
[Presheaf dual of the unit is the unit (general base)]
  \label{def:dual_unit_iso_gen}
  \lean{PresheafOfModules.dualUnitIsoGen}
  \uses{def:unit_dual_section_equiv, def:presheaf_dual, lem:internal_hom_eval}
  The presheaf-level isomorphism
  \(\mathtt{dual}\,\mathbf 1 \cong \mathbf 1\) of presheaves of modules over a general
  single-universe base ring presheaf \(R_0\), assembled sectionwise from
  \cref{def:unit_dual_section_equiv} with the evaluation-at-\(1\) naturality.
\end{definition}
\begin{proof}
  \uses{def:unit_dual_section_equiv, def:presheaf_dual, lem:internal_hom_eval}
  \leanok
  Proved directly in Lean. The sectionwise components are the section equivalences
  \cref{def:unit_dual_section_equiv} packaged by \(\mathtt{isoMk}\); the naturality square is
  the evaluation-at-\(1\) naturality of the internal-hom evaluation
  (\cref{lem:internal_hom_eval}) specialised at \(M = \mathbf 1\).
\end{proof}

\begin{definition}\leanok
[Scheme-level presheaf dual of the unit is the unit]
  \label{def:presheaf_dual_unit_iso}
  \lean{AlgebraicGeometry.Scheme.Modules.presheafDualUnitIso}
  \uses{def:dual_unit_iso_gen}
  The scheme-level instance \(\mathtt{dual}\,\mathbf 1 \cong \mathbf 1\) of the
  presheaf-of-modules dual-of-unit isomorphism, obtained by specialising
  \cref{def:dual_unit_iso_gen} at \(R_0 = Y.\mathtt{presheaf}\). It is the presheaf core of
  the third leg \(\mathtt{dual\_unit\_iso}\) (\cref{lem:dual_unit_iso}).
\end{definition}
\begin{proof}
  \uses{def:dual_unit_iso_gen}
  \leanok
  Proved directly in Lean --- a direct instance of \cref{def:dual_unit_iso_gen}.
\end{proof}

\begin{lemma}\leanok
[Reverse slice dual transport (the \(\mathtt{invFun}\) of \cref{lem:slice_dual_transport})]
  \label{lem:slice_dual_transport_inv}
  \lean{AlgebraicGeometry.Scheme.Modules.sliceDualTransportInv}
  \uses{def:dual_unit_ring_swap_hom, lem:image_preimage_of_le,
        lem:isiso_eps_restrictscalars_appiso_hom, lem:restrictscalars_ringiso_dualequiv,
        def:presheaf_dual_unit_iso, def:unit_relabel_swap}
  \textit{Source: internal categorical construction; no external reference.}
  In the setting of \cref{lem:slice_dual_transport}, with \(f : Y \hookrightarrow X\) an open
  immersion, \(M \in \Scheme.\mathtt{Modules}\,X\), \(V \subseteq Y\) an open and \(\beta\)
  the open-immersion structure-ring morphism, \emph{and assuming the \(\beta\)-compatibility
  hypothesis \(\mathtt{h\beta} : \forall P,\ (\beta.\mathtt{app}\,P).\mathrm{hom} \circ
  (f.\mathtt{appIso}\,P).\mathrm{hom}.\mathrm{hom} = \mathrm{id}\)} (the ring identity that
  collapses the double restriction-of-scalars in the reverse component; it holds for the
  open-immersion \(\beta\) of \cref{lem:slice_dual_transport} but is false for an arbitrary
  \(\beta\), so it is carried explicitly as a hypothesis), there is a reverse transport: a
  dual section
  \(\psi : \mathtt{restr}_V\,((\mathtt{pushforward}\,\beta)\,M.\mathtt{val}) \to
  \mathtt{restr}_V\,\mathbf 1_Y\) over \(\mathtt{Over}\,V\) is sent to an \(X\)-slice dual
  section \(\mathtt{restr}_{fV}\,M.\mathtt{val} \to \mathtt{restr}_{fV}\,\mathbf 1_X\) over
  \(\mathtt{Over}\,(fV)\), where \(fV = f.\mathtt{opensFunctor}(V)\). This is the inverse leg
  of the forward equivalence \cref{lem:slice_dual_transport}.
\end{lemma}
\begin{proof}
  \uses{def:dual_unit_ring_swap_hom, lem:image_preimage_of_le,
        lem:isiso_eps_restrictscalars_appiso_hom, lem:restrictscalars_ringiso_dualequiv,
        def:presheaf_dual_unit_iso, def:unit_relabel_swap}
        \leanok
  We build the underlying \(\mathtt{PresheafOfModules.Hom}\) component-by-component over the
  \(X\)-slice \(\mathtt{Over}\,(fV)\), mirroring the forward component formula of
  \cref{lem:slice_dual_transport} with two changes. The component is a four-leg
  cross-fiber composite
  \(M.\mathtt{val}.\mathrm{map}(\mathtt{eqToHom}) \circ
  (\mathtt{restrictScalars}\,\rho).\mathrm{map}(\mathtt{collapse}\circ\mathtt{core}) \circ
  \mathtt{unitRelabelSwap}\): the outer \(\mathtt{restrictScalars}\,\rho\) conjugation carries
  the semilinear relabel across the ring fibers (a single \(\circ\)-chain cannot express it),
  \(\mathtt{collapse}\) reduces the double restriction-of-scalars to the identity via the
  \(\mathtt{h\beta}\) ring identity, and \(\mathtt{unitRelabelSwap}\)
  (\cref{def:unit_relabel_swap}) is the unit-codomain transport.

  \emph{The component at \(W''\).} Fix \(W'' \le fV\) and set \(W' := W''.\mathrm{left}\) and
  \(P := f^{-1}(W')\) for its preimage. Because \(f\) is an open immersion and
  \(fV \subseteq \mathrm{range}\,f\), the open \(W' \le fV\) lies in the image of
  \(f.\mathtt{opensFunctor}\), so the down-set identity
  \(f.\mathtt{opensFunctor}(P) = W'\) holds --- but only \emph{propositionally}
  (\cref{lem:image_preimage_of_le}); this is the single set-theoretic fact the inverse
  reindexing depends on. The component of the reverse section at \(W''\) is the \(X\)-slice
  mirror of the forward component:
  \[
    M.\mathtt{val}(\mathtt{eqToHom})
      \;\circ\;
    (\mathtt{restrictScalars}\,(f.\mathtt{appIso}\,P).\mathrm{hom}).\mathrm{map}
      \bigl(\psi_{(\mathtt{Over.mk}\,(P \le V))}\bigr)
      \;\circ\;
    \mathtt{dualUnitRingSwapHom}\,f\,P,
  \]
  conjugated on source and target by the \(\mathtt{eqToHom}\) transports arising from the
  propositional equality \(f.\mathtt{opensFunctor}(P) = W'\) (the mirror of
  \(\mathtt{homLocalSection}\)). Compared with the forward formula, this map uses
  \((f.\mathtt{appIso}\,P).\mathrm{hom}\) in place of its inverse --- because the reverse
  reindex turns a \(\mathcal{O}_Y(P)\)-section map into a \(\mathcal{O}_X(fP)\)-map --- and
  its codomain swap is \(\mathtt{dualUnitRingSwapHom}\) (\cref{def:dual_unit_ring_swap_hom}),
  the inverse of \(\varepsilon\) in the \(\mathrm{hom}\)-direction
  (\cref{lem:isiso_eps_restrictscalars_appiso_hom}).

  \emph{Additivity and \(\mathcal{O}_Y(V)\)-linearity} follow exactly as in the forward leg
  (steps (i)--(iii) of \cref{lem:slice_dual_transport}): identify the two module
  presentations, match the X- and Y-side scalars by the \(\beta\)-naturality ring identity,
  and apply the change-of-rings compatibility, the module-level shadow of
  \cref{lem:restrictscalars_ringiso_dualequiv}.

  \emph{Naturality} of the component family across morphisms of \((\mathtt{Over}\,(fV))^{\mathrm{op}}\)
  reduces, after the thin-poset \(\mathtt{Subsingleton.elim}\) on the indexing morphisms, to
  the \(\varepsilon\)-naturality of \(\mathtt{restrictScalars}\) along the structure-ring iso,
  exactly as in the forward leg. The unit handling is the presheaf dual-of-unit
  identification \cref{def:presheaf_dual_unit_iso}. The whole family is reindexed by the
  inverse of the down-set bijection \(W'' \leftrightarrow f^{-1}(W'')\) supplied by
  \cref{lem:image_preimage_of_le}.
\end{proof}

\begin{lemma}




\leanok
[Restriction commutes with the dual along an open immersion (the C-bridge)]
  \label{lem:dual_restrict_iso}
  \lean{AlgebraicGeometry.Scheme.Modules.dual_restrict_iso}
  \uses{lem:internal_hom_isSheaf, lem:restrictscalars_ringiso_dualequiv,
        lem:slice_dual_transport}
  % SOURCE: [Stacks Project], "Modules on Ringed Spaces", \S Internal Hom,
  % Lemma~\texttt{lemma-pullback-internal-hom} (tag area 01CM)
  % (read from references/stacks-modules.tex, L3710--L3721).
  % SOURCE QUOTE:
  % "Let $f : (X, \mathcal{O}_X) \to (Y, \mathcal{O}_Y)$ be a morphism
  % of ringed spaces. Let $\mathcal{F}$, $\mathcal{G}$ be $\mathcal{O}_Y$-modules.
  % If $\mathcal{F}$ is finitely presented and $f$ is flat,
  % then the canonical map
  % $$
  % f^*\SheafHom_{\mathcal{O}_Y}(\mathcal{F}, \mathcal{G})
  % \longrightarrow
  % \SheafHom_{\mathcal{O}_X}(f^*\mathcal{F}, f^*\mathcal{G})
  % $$
  % is an isomorphism."
  \textit{Source: [Stacks Project], ``Modules on Ringed Spaces'', \S Internal Hom,
  Lemma~\texttt{lemma-pullback-internal-hom} (tag area 01CM).}
  Let \(f : Y \hookrightarrow X\) be an open immersion of schemes, with image the open
  \(U \subseteq X\), and let \(M \in \Scheme.\mathtt{Modules}\,X\). Then there is a
  canonical isomorphism of \(\mathcal{O}_Y\)-modules
  \[
    (\mathtt{dual}\,M)\big|_f \;\xrightarrow{\ \sim\ }\; \mathtt{dual}\bigl(M|_f\bigr),
  \]
  natural in \(M\), between the restriction of the sheaf-level dual
  (\cref{lem:internal_hom_isSheaf}) and the dual of the restriction. (Only opens
  \(V \subseteq U\) intervene downstream, where the comparison is needed.)
\end{lemma}

\begin{proof}
  \uses{lem:internal_hom_isSheaf, lem:restrictscalars_ringiso_dualequiv,
        lem:slice_dual_transport}
        \leanok
  This is the classical internal-hom base-change isomorphism
  \(f^*\mathcal{H}\!om_{\mathcal{O}_X}(M, \mathcal{O}_X) \cong
  \mathcal{H}\!om_{\mathcal{O}_Y}(f^*M, f^*\mathcal{O}_X)\) of
  Lemma~\texttt{lemma-pullback-internal-hom}, specialised to the structure sheaf
  \(\mathcal{G} = \mathcal{O}_X\). For an open immersion the morphism \(f\) is flat and
  restriction \((-)|_f\) is the exact pullback \(f^*\), so the finite-presentation and
  flatness hypotheses of the general ringed-space statement are automatic; this is the
  \emph{favorable} (open-immersion) case of the base-change comparison, in which the map
  is an isomorphism for \emph{every} \(M\), not only the finitely-presented ones.

  We emphasise at the outset that this is a genuine \emph{natural isomorphism}, not a
  definitional equality, and --- crucially --- that the dual is \emph{not sectionwise}.
  Unlike the tensor product, whose value is computed fibrewise
  (\((M \otimes N)(V) = M(V) \otimes N(V)\)), the dual evaluates to a Hom over an entire
  down-set: for an open \(V \subseteq U\),
  \[
    (\mathtt{dual}\,M)(V)
      \;=\; \mathcal{H}om\bigl(M|_V,\ \mathcal{O}|_V\bigr)\big|_{\mathtt{Over}\,V},
  \]
  a module of morphisms of presheaves over the whole slice \(\mathtt{Over}\,V \subseteq
  \mathrm{Opens}\,Y\), \emph{not} a single fibre \(M(V) \to \mathcal{O}(V)\). Consequently
  the assertion that restriction along the open immersion \(f\) commutes with the dual does
  \emph{not} reduce to a single ground-ring reconciliation. The proof proceeds in two
  stages: a preliminary adjunction-uniqueness rewrite (stage~H1) that re-presents the
  pulled-back dual in pushforward form, followed by the resolution of the resulting
  \emph{residual} into \emph{two} legs, carried out objectwise in the open \(V\) (with
  \(V \subseteq U\)) and assembled over the slice with poset-thin naturality.

  \emph{Stage H1 --- adjunction-uniqueness rewrite.} The pulled-back dual is initially
  presented in the \(\mathtt{pullback}\,\varphi\)-form inherited from the abstract
  definition of the restriction functor, and no sectionwise comparison can begin against
  that opaque form. We first rewrite it into a \(\mathtt{pushforward}\,\beta\)-form: the
  restriction-along-an-open-immersion pullback and the pushforward along the structure-ring
  map \(\beta\) are both left adjoints to one and the same functor, so by uniqueness of left
  adjoints --- the adjunction \(\mathtt{pushforwardPushforwardAdj}\) composed with
  \(\mathtt{leftAdjointUniq}\) --- they are canonically isomorphic, and this isomorphism
  transports the pulled-back dual into pushforward form. After H1 the residual obligation is
  \emph{exactly}
  \[
    (\mathtt{pushforward}\,\beta).\mathtt{obj}\,(\mathtt{dual}\,M.\mathtt{val})
      \;\xrightarrow{\ \sim\ }\;
    \mathtt{dual}\bigl((\mathtt{pushforward}\,\beta).\mathtt{obj}\,M.\mathtt{val}\bigr),
  \]
  i.e.\ ``\(\mathtt{pushforward}\,\beta\) commutes with the dual''. It is this residual ---
  \emph{not} the original Step-4 isomorphism --- that legs (A) and (B) below resolve; legs
  (A) and (B) are the content remaining \emph{after} the H1 rewrite, not a standalone
  two-leg split of the whole step.

  \emph{Leg (A) --- slice-site Hom base-change (Beck--Chevalley).} The domain presheaf must
  first be transported across the open immersion. Because \(f\) is an open immersion, its
  direct-image functor on open sets is fully faithful with image
  the down-set \(\{W : W \le U\}\); this carries the restricted domain
  \(\mathtt{restr}_V(f_*M)\) over \(\mathtt{Over}\,V \subseteq \mathrm{Opens}\,Y\) to the
  restricted domain \(\mathtt{restr}_{fV}\,M\) over \(\mathtt{Over}\,(fV) \subseteq
  \mathrm{Opens}\,X\), and hence identifies the two Hom-modules computed over these matched
  slices. This is a genuine slice-site Hom base-change --- a Beck--Chevalley square on the
  internal hom across the fully-faithful open-immersion functor, restricted to the down-set
  of \(fV\) --- \emph{not} a fibrewise identity.

  \emph{Leg (B) --- ground-ring reconciliation.} The codomain of the internal hom is in
  both cases the structure sheaf. Once leg (A) has matched the domains, the unit codomain
  is reconciled along the open-immersion ring isomorphism
  \(\beta_V = (\mathtt{toRingCatSheafHom}\,f).\mathrm{hom}_V \colon
  \mathcal{O}_X(fV) \xrightarrow{\sim} \mathcal{O}_Y(V)\): restriction of scalars along the
  ring \emph{isomorphism} \(\beta_V\) is an equivalence, and its ModuleCat-level
  shadow on duals is exactly \cref{lem:restrictscalars_ringiso_dualequiv}, carrying the
  \(\mathcal{O}_X(fV)\)-linear dual onto the \(\mathcal{O}_Y(V)\)-linear dual,
  \(R(V)\)-linearly and invertibly. This is the dual analogue, at the module level, of the
  ring-isomorphism tensor equivalence (\cref{lem:restrictscalars_ringiso_strongmonoidal}).

  \emph{The combined leg-(A)\(\circ\)(B) atom \(\mathtt{sliceDualTransport}\).} Both legs are
  realised together by a single \(\mathcal{O}_Y(V)\)-linear isomorphism
  \(\mathtt{sliceDualTransport}\,f\,M\,V\) (\cref{lem:slice_dual_transport}), built
  \emph{by hand} --- it is \emph{not} obtained by consuming a sheaf-level slice equivalence.
  This is essential: the equivalence of \emph{sheaves of modules} over an open immersion
  compares the restriction functors and the units, but carries \emph{no} internal-hom or
  dual content; the assertion realized by leg (A) --- that the dual commutes with slice
  reindexing along the open-immersion direct-image functor \(f_{!}\) on opens
  (\(f.\mathtt{opensFunctor}\)) --- cannot be supplied by such an object without a
  monoidal-closed structure on sheaves of modules, which this development deliberately does
  \emph{not} carry (cf.\ \cref{rem:scheme_modules_monoidal_off_path}). The atom must
  therefore be constructed directly.

  Concretely (\cref{lem:slice_dual_transport}), the forward component of
  \(\mathtt{sliceDualTransport}\,f\,M\,V\) at \(W \le V\) is the categorical image
  \((\mathtt{restrictScalars}\,\beta_W).\mathrm{map}\,(\varphi_{\bullet})\) of the dual-section
  component of \(\varphi\) at the \(f\)-image of \(W\), reindexed across
  \(f.\mathtt{opensFunctor}\) by restriction of scalars along \(\beta_W\), post-composed with
  the leg-(B) codomain swap \(\mathtt{inv}(\varepsilon(\mathtt{restrictScalars}\,\beta_W))\).
  Leg (A) is built categorically via \((\mathtt{restrictScalars}\,\beta_W).\mathrm{map}\):
  a genuine change-of-rings functor acts here, so the reindexing must be applied as a
  functor map rather than pointwise. Leg (B)'s unit comparison \(\varepsilon\) is invertible
  because \(\beta_W\) is a ring isomorphism, as detailed in \cref{lem:slice_dual_transport}.
  Because \(\mathrm{Opens}\,Y\) is a thin poset, the
  naturality squares commute by \(\mathtt{Subsingleton.elim}\). The forward and inverse maps
  identify the slice-Hom value over \(\mathtt{Over}\,(fV)\) with the slice-Hom value over
  \(\mathtt{Over}\,V\), \(\mathcal{O}_Y(V)\)-linearly.

  The Step-4 residual is then \(\mathtt{PresheafOfModules.isoMk}\) applied \emph{directly} to
  the family \(V \mapsto \mathtt{sliceDualTransport}\,f\,M\,V\) --- which already packages
  \emph{both} legs (A) and (B) --- with the outer naturality square thin-poset-trivial. No
  separate leg-(B) step is interposed in \(\mathtt{isoMk}\): the per-section atom of
  \cref{lem:slice_dual_transport} carries the full leg-(A)\(\circ\)(B) content. The outer
  \(\mathtt{isoMk}\) naturality square follows from the per-section family naturality
  \(\mathtt{sliceDualTransport.naturality}\) of \cref{lem:slice_dual_transport}: the
  \(\varepsilon\)-naturality is established first at the level of individual sections, and the
  outer presheaf-morphism naturality of \(\mathtt{isoMk}\) then follows.

  The residual of stage~H1 (``\(\mathtt{pushforward}\,\beta\) commutes with the dual'') is
  the sectionwise composite of leg (A) followed by leg (B), packaged over the slice with
  poset-thin naturality; the full Step-4 isomorphism is then the H1 isomorphism followed by
  this composite. The slice-reindexing coherence is \emph{light}, because \(\mathrm{Opens}\)
  is a thin poset and the relevant squares commute by \(\mathtt{Subsingleton.elim}\),
  whereas H1 and legs (A) and (B) carry the substantive content.

  It is worth recording \emph{why} the strong-monoidal restriction
  (\cref{lem:presheaf_pushforward_adj_substrate}) has no direct dual analogue here. For
  the tensor product the transport \emph{collapsed to leg (B) alone}: because \(\otimes\)
  is sectionwise, pushing a tensor forward reduced, section by section, to a single
  ground-ring tensor equivalence, with no domain slice left to re-index. The dual's
  non-sectionwise nature removes that collapse and forces the additional leg (A): the
  domain \(\mathtt{restr}_V(f_*M)\) genuinely lives over a varying slice and must be
  transported before the ground ring can be reconciled.

  We caution that the naive ``fixed-value-category slice equivalence'' of
  \cref{lem:open_immersion_slice_sheaf_equiv} is \emph{not} applicable to leg (A). That instrument
  operates in a fixed value category over a fixed base ring, whereas the residual of
  leg (A) lives at the \emph{presheaf} level, with a per-section ring \(\mathcal{O}_Y(V)\)
  that varies with \(V\) and a finer slicing than the sheaf-level equivalence resolves;
  leg (A) must therefore be built directly at the presheaf-of-modules / varying-ring level.

  Composing legs (A) and (B) gives, for each \(V \subseteq U\), an
  \(\mathcal{O}_Y(V)\)-linear isomorphism between the two values; assembling these into a
  presheaf-of-modules natural isomorphism and sheafifying through the sheaf-level dual of
  \cref{lem:internal_hom_isSheaf} resolves the H1 residual. Precomposing with the stage-H1
  rewrite then yields the stated isomorphism
  \((\mathtt{dual}\,M)|_f \cong \mathtt{dual}(M|_f)\).

  \emph{The inverse-uniqueness shortcut is not viable.} One might hope to derive the
  isomorphism from the already-established \cref{lem:tensorobj_restrict_iso} by uniqueness of
  monoidal inverses --- the dual being the \(\otimes\)-inverse of an invertible module and
  restriction commuting with \(\otimes\) --- sidestepping leg (A). The abstract statements of
  this idiom (\(\mathtt{hasRightDualOfEquivalence}\), \(\mathtt{rightDualIso}\)) require,
  however, four structures the present setting lacks: a registered \(\mathtt{MonoidalCategory}\)
  instance on \(\mathtt{PresheafOfModules}\) (the monoidal structure here is carried entirely
  by hand), a \emph{strong} (not merely lax/oplax) monoidal restriction functor, restriction
  being an \emph{equivalence} (false for a non-surjective open immersion, which is a
  localization), and an exact pairing \(\mathtt{ExactPairing}\,M\,(\mathtt{dual}\,M)\) (which
  exists only for dualizable \(M\), whereas the lemma is stated for general \(M\)). Even
  restricted to the invertible consumer this would require building a coevaluation and
  verifying the zig-zag identities --- strictly \emph{more} infrastructure than leg (A) --- so
  the construction commits to legs (A) and (B) above.
\end{proof}

\begin{lemma}\leanok
[Dual of the structure sheaf]
  \label{lem:dual_unit_iso}
  \lean{AlgebraicGeometry.Scheme.Modules.dual_unit_iso}
  \uses{lem:internal_hom_eval, lem:internal_hom_isSheaf, lem:tensorobj_unit_iso}
  Let \(Y\) be a scheme. The sheaf-level dual of the structure sheaf is canonically the
  structure sheaf itself,
  \[
    \mathtt{dual}\,\mathcal{O}_Y \;\xrightarrow{\ \sim\ }\; \mathcal{O}_Y,
  \]
  obtained by sheafifying the presheaf-level ``evaluation at \(1\)'' isomorphism
  \(\mathcal{H}om(\mathbf{1}, \mathbf{1}) \cong \mathbf{1}\): every endomorphism of the
  unit presheaf is multiplication by its value at the global section \(1\), so evaluation
  at \(1\) identifies the internal hom of the unit into itself with the unit.
\end{lemma}

\begin{proof}
  \uses{lem:internal_hom_eval, lem:internal_hom_isSheaf, def:presheaf_dual_unit_iso}
  \leanok
  At the presheaf level the dual of the unit \(\mathbf{1} = \mathcal{O}_Y\) is the internal
  hom \(\mathcal{H}om(\mathbf{1}, \mathbf{1})\), whose value on an open \(U\) is the module
  of \(\mathcal{O}(U)\)-linear endomorphisms of \(\mathbf{1}|_U\). Each such endomorphism is
  determined by its value at the global section \(1\) and equals multiplication by that
  value; hence \emph{evaluation at \(1\)} is itself the sectionwise isomorphism
  \(\mathcal{H}om(\mathbf{1}, \mathbf{1}) \cong \mathbf{1}\) of presheaves of modules, with
  inverse the scalar-action map \(r \mapsto (\,\cdot\,r)\). (No left unitor is needed: the
  identification is the direct evaluation-at-\(1\) / global-scalar-multiplication
  isomorphism, not the composite through \(\mathbf{1} \otimes \mathbf{1} \cong \mathbf{1}\).)
  Sheafifying through the sheaf-level dual of \cref{lem:internal_hom_isSheaf} yields
  \(\mathtt{dual}\,\mathcal{O}_Y \cong \mathcal{O}_Y\). The scheme-level alias of this
  presheaf evaluation-at-\(1\) isomorphism is
  \(\mathtt{presheafDualUnitIso}\) (\cref{def:presheaf_dual_unit_iso}).
\end{proof}

\begin{lemma}




\leanok
[The dual of a line bundle is a line bundle]
  \label{lem:dual_isLocallyTrivial}
  \lean{AlgebraicGeometry.Scheme.Modules.dual_isLocallyTrivial}
  \uses{lem:internal_hom_isSheaf, lem:dual_restrict_iso, lem:dual_unit_iso,
        def:scheme_modules_dual_iso_of_iso, lem:restrictscalars_ringiso_dualequiv}
  % SOURCE: [Stacks Project], Tag 01CR (Modules on Ringed Spaces, \S Invertible
  % modules), Lemma~lemma-constructions-invertible (item~2)
  % (read from references/stacks-modules.tex, L4200--L4213).
  % SOURCE QUOTE:
  % "\item If $\mathcal{L}$ is an invertible $\mathcal{O}_X$-module, then so is
  % $\SheafHom_{\mathcal{O}_X}(\mathcal{L}, \mathcal{O}_X)$ and the evaluation map
  % $\mathcal{L} \otimes_{\mathcal{O}_X}
  % \SheafHom_{\mathcal{O}_X}(\mathcal{L}, \mathcal{O}_X) \to \mathcal{O}_X$
  % is an isomorphism."
  \textit{Source: [Stacks Project], Tag~01CR,
  Lemma~\texttt{lemma-constructions-invertible} (item~2).}
  Let \(X\) be a scheme and \(L \in \Scheme.\mathtt{Modules}\,X\) with
  \(\mathtt{LineBundle.IsLocallyTrivial}\,L\). Then the sheaf-level dual
  \(\mathtt{dual}\,L\) of \cref{lem:internal_hom_isSheaf} is again locally trivial of
  rank one, i.e.\ \(\mathtt{LineBundle.IsLocallyTrivial}\,(\mathtt{dual}\,L)\).
\end{lemma}

\begin{proof}
  \uses{lem:internal_hom_isSheaf, lem:dual_restrict_iso, lem:dual_unit_iso,
        def:scheme_modules_dual_iso_of_iso, lem:restrictscalars_ringiso_dualequiv}
        \leanok
  Local triviality of \(L\) furnishes, around each point, an affine open \(U\) with a
  trivialisation \(L|_U \cong \mathcal{O}_U\); writing \(f : U \hookrightarrow X\) for the
  associated open immersion, it suffices to exhibit an isomorphism
  \((\mathtt{dual}\,L)|_U \cong \mathcal{O}_U\). This is the following three-step chain:
  \[
    (\mathtt{dual}\,L)\big|_f
      \;\xrightarrow{\ \sim\ }\; \mathtt{dual}\bigl(L|_f\bigr)
      \;\xrightarrow{\ \sim\ }\; \mathtt{dual}\,\mathcal{O}_U
      \;\xrightarrow{\ \sim\ }\; \mathcal{O}_U.
  \]
  \begin{enumerate}
    \item The first isomorphism is the restriction-compatibility isomorphism
      \cref{lem:dual_restrict_iso}: restriction along the open immersion \(f\) commutes
      with the formation of the dual. This is the C-bridge proper; its substantive content
      is the sectionwise ring-iso transport of the module structure
      (\cref{lem:restrictscalars_ringiso_dualequiv}), as detailed there.
    \item The second isomorphism is the dual applied to the trivialisation
      \(L|_U \cong \mathcal{O}_U\), via the contravariant functoriality of the sheaf-level
      dual in isomorphisms (\(\mathtt{dualIsoOfIso}\),
      \cref{def:scheme_modules_dual_iso_of_iso}): an isomorphism \(L|_U \cong \mathcal{O}_U\)
      induces \(\mathtt{dual}(L|_U) \cong \mathtt{dual}\,\mathcal{O}_U\).
    \item The third isomorphism trivialises the dual of the unit,
      \(\mathtt{dual}\,\mathcal{O}_U \cong \mathcal{O}_U\) (\cref{lem:dual_unit_iso}):
      the dual of the structure sheaf is the structure sheaf itself, the comparison being
      evaluation at \(1 \in \mathcal{O}_U\), under which
      \(\mathcal{H}om_{\mathcal{O}_U}(\mathcal{O}_U, \mathcal{O}_U) \cong \mathcal{O}_U\).
  \end{enumerate}
  Composing the three exhibits \((\mathtt{dual}\,L)|_U \cong \mathcal{O}_U\). Hence
  \(\mathtt{dual}\,L\) is locally trivial of rank one, as required by item~2 of Stacks
  Lemma~\texttt{lemma-constructions-invertible} (tag 01CR).
\end{proof}

\begin{lemma}


\leanok
[Local isomorphisms of \(\mathcal{O}_X\)-modules glue to a global isomorphism]
  \label{lem:isiso_of_isiso_restrict}
  \lean{AlgebraicGeometry.Scheme.Modules.isIso_of_isIso_restrict}
  Let \(X\) be a scheme and \(\varphi \colon M \to N\) a morphism in
  \(\Scheme.\mathtt{Modules}\,X = \mathtt{SheafOfModules}\,X.\mathtt{ringCatSheaf}\).
  Suppose that for every point \(x \in X\) there is an open neighbourhood \(U_x\) of
  \(x\) such that the restriction of \(\varphi\) to \(U_x\) --- the image
  \((\mathtt{restrictFunctor}\,(U_x).\mathtt{i}).\mathtt{map}\,\varphi\) under the
  restriction functor along the open immersion \((U_x).\mathtt{i}\colon U_x \to X\) ---
  is an isomorphism in \(\Scheme.\mathtt{Modules}\,U_x\). Then \(\varphi\) is an
  isomorphism.
\end{lemma}

\begin{proof}



  Being an isomorphism for a morphism of sheaves of \(\mathcal{O}_X\)-modules can be
  checked on the underlying morphism of sheaves of abelian groups, and there it can be
  checked stalkwise. Fix a point \(x \in X\). Choose a preimage \(x' \in U_x\) of \(x\)
  under the open immersion \((U_x).\mathtt{i}\). By hypothesis the restriction
  \((\mathtt{restrictFunctor}\,(U_x).\mathtt{i}).\mathtt{map}\,\varphi\) is an
  isomorphism, hence so is its image under any functor --- in particular under the
  stalk-at-\(x'\) functor on the underlying ab-sheaves. The natural isomorphism
  ``restriction commutes with stalks'' (\(\mathtt{restrictStalkNatIso}\)) identifies
  this with the stalk of (the underlying ab-sheaf morphism of) \(\varphi\) at the image
  point \(x\); transporting along it shows that stalk is an isomorphism. As this holds
  at every \(x\), and a morphism of sheaves of abelian groups all of whose stalks are
  isomorphisms is itself an isomorphism (stalkwise iso-detection,
  \(\mathtt{isIso\_of\_stalkFunctor\_map\_iso}\) on \(\mathtt{TopCat.Sheaf}\)), the
  underlying ab-sheaf morphism of \(\varphi\) is an isomorphism. Finally the forgetful
  functor \(\Scheme.\mathtt{Modules.toPresheaf}\) reflects isomorphisms, so \(\varphi\)
  itself is an isomorphism.
\end{proof}

\begin{definition}



\leanok
[Scheme-level linearity wrapper \(\mathtt{homMk}\) (A-bridge step (ii))]
  \label{def:scheme_modules_homMk}
  \lean{AlgebraicGeometry.Scheme.Modules.homMk}
  Let \(X\) be a scheme and \(M, N \in \Scheme.\mathtt{Modules}\,X\). Given a morphism
  \(g \colon M.\mathtt{val}.\mathtt{presheaf} \to N.\mathtt{val}.\mathtt{presheaf}\) of
  the underlying presheaves of abelian groups, together with a hypothesis \(hg\) that
  \(g\) is \(\mathcal{O}_X\)-linear sectionwise, \(\mathtt{homMk}\,g\,hg\) is the
  morphism \(M \to N\) in \(\Scheme.\mathtt{Modules}\,X\) that this data determines. It
  is a thin wrapper that packages \(\mathtt{PresheafOfModules.homMk}\) at the scheme
  level: the sectionwise-linearity datum promotes the ab-group morphism to a genuine
  morphism of sheaves of \(\mathcal{O}_X\)-modules. Its \(\mathtt{@[simp]}\) companion
  \(\mathtt{toPresheaf\_map\_homMk}\) records that forgetting back recovers \(g\), i.e.\
  \((\mathtt{toPresheaf}\,X).\mathtt{map}\,(\mathtt{homMk}\,g\,hg) = g\). This is step
  (ii) of the A-bridge \cref{lem:sheafofmodules_hom_of_local_compat}: once the
  underlying ab-sheaf morphism has been glued, \(\mathtt{homMk}\) re-attaches the
  \(\mathcal{O}_X\)-linearity to produce the global morphism of modules.
\end{definition}

\begin{lemma}



\leanok
  [The load-bearing local section \(\mathtt{localSection}\) of the hom-sheaf]
  \label{lem:scheme_modules_hom_local_section}
  \lean{AlgebraicGeometry.Scheme.Modules.homLocalSection}
  \uses{lem:open_immersion_slice_sheaf_equiv, lem:image_preimage_of_le}
  Let \(X\) be a scheme, \(M, N \in \Scheme.\mathtt{Modules}\,X\), \(\{U_i\}\) an open
  cover, and \(f_i \colon M|_{U_i} \to N|_{U_i}\) a local morphism. Write
  \(\mathcal{H}\) for the sheaf of abelian-group morphisms from the underlying presheaf of
  \(M\) to that of \(N\) (i.e.\ \(\mathtt{presheafHom}\) of the two underlying presheaves
  of abelian groups). Then \(f_i\) determines a local section
  \(\mathtt{localSection}\,i \in \mathcal{H}(U_i)\): its
  component at a pair \((V, h : V \to U_i)\) is the given component of \(f_i\), conjugated
  by the canonical \(\mathtt{eqToHom}\) identifications arising from the open-set equality
  \(\iota_i(\iota_i^{-1}(V)) = V\) valid for \(V \subseteq U_i\). The non-trivial content is
  the \emph{naturality condition}: these \(\mathtt{eqToHom}\)-conjugated components commute
  across the morphisms of the slice category over \(U_i\) --- precisely the
  section-direction slice of \cref{lem:open_immersion_slice_sheaf_equiv}, with coherence
  automatic on the thin poset of opens of \(X\) by \(\mathtt{Subsingleton.elim}\). This
  is the standalone sub-lemma the gluing of \cref{lem:sheafofmodules_hom_of_local_compat}
  is built upon.
\end{lemma}

\begin{proof}
  \uses{lem:open_immersion_slice_sheaf_equiv, lem:image_preimage_of_le}
  \leanok
  Fix \(i\). The section \(\mathtt{localSection}\,i \in \mathcal{H}(U_i)\) is a morphism of
  the restrictions of the underlying abelian-group presheaves of \(M\) and \(N\), indexed
  over the slice category \(\mathtt{Over}\,U_i\). We define its component at a pair
  \((V, h : V \to U_i)\) by conjugating the given component of \(f_i\): writing
  \(V' = \iota_i^{-1}(V)\) for the preimage of \(V\) under the open immersion
  \(\iota_i : U_i \to X\), take the section map \((f_i).\mathtt{app}(V') \colon M(V') \to N(V')\)
  of the underlying ab-presheaf morphism of \(f_i\), and conjugate it on the source by the
  restriction \(M(\mathtt{eqToHom})\) and on the target by \(N(\mathtt{eqToHom})\) along the
  down-set identity \(\iota_i\bigl(\iota_i^{-1}(V)\bigr) = V\) (valid for \(V \subseteq U_i\),
  \cref{lem:image_preimage_of_le}).
  The result is a morphism \(M(V) \to N(V)\).

  The substantive content is the \emph{naturality} of this assignment across a morphism
  \(\varphi\) of the slice \(\mathtt{Over}\,U_i\), say from a pair \((B)\) to a pair \((A)\)
  with underlying inclusion \(B \subseteq A\). Unwinding the two composite routes, one must
  show --- separately on the source presheaf \(M\) and the target presheaf \(N\) --- that the
  \(\varphi\)-restriction followed by the \(\mathtt{eqToHom}\) at \(B\) equals the
  \(\mathtt{eqToHom}\) at \(A\) followed by the \(\kappa\)-restriction, where
  \(\kappa \colon \iota_i^{-1}(B) \to \iota_i^{-1}(A)\) is the induced inclusion of
  preimages. Each of these two equations is an equality of parallel morphisms between the
  same pair of opens in the thin poset \((\mathtt{Opens}\,X)^{\mathrm{op}}\), so it holds by
  \(\mathtt{Subsingleton.elim}\); combined with the naturality of the underlying ab-presheaf
  morphism of \(f_i\) across \(\kappa\), the two routes agree. This is precisely the
  section-direction slice of the comparison packaged in
  \cref{lem:open_immersion_slice_sheaf_equiv}.
\end{proof}

\begin{lemma}\leanok
[Down-set image identity \(\iota_!(\iota^{-1}V) = V\)]
  \label{lem:image_preimage_of_le}
  \lean{AlgebraicGeometry.Scheme.Modules.image_preimage_of_le}
  For opens \(V \le W\) of a scheme \(X\), the image under the open immersion \(W.\iota\) of
  the preimage of \(V\) is \(V\) again: \(W.\iota\,{}_!(W.\iota^{-1}(V)) = V\). This is the
  equality powering the \(\mathtt{eqToHom}\)-conjugations of \(\mathtt{homLocalSection}\)
  (\cref{lem:scheme_modules_hom_local_section}) and of the reverse slice transport
  (\cref{lem:slice_dual_transport_inv}).
\end{lemma}
\begin{proof}

\leanok
  Proved directly in Lean --- the image of the preimage is the intersection with the
  (open) range of the immersion, which equals \(V\) since \(V \le W\) lies in that range.
\end{proof}

\begin{definition}\leanok
[Top section to global hom]
  \label{def:top_section_to_hom}
  \lean{AlgebraicGeometry.Scheme.Modules.topSectionToHom}
  Let \(F, G\) be presheaves of abelian groups on a topological space \(X\). A section of
  \(\mathtt{presheafHom}\,F\,G\) over the terminal open \(\top\) determines a global
  morphism \(F \to G\): since \(\top\) is terminal in \(\mathtt{Opens}\,X\), the value at
  \(\top\) restricts to a full compatible family of sections, which
  \(\mathtt{presheafHomSectionsEquiv}\) identifies with a morphism \(F \to G\). This is
  sub-step (b) of \cref{lem:sheafofmodules_hom_of_local_compat}.
\end{definition}
\begin{proof}

\leanok
  Proved directly in Lean --- restrict the top section along \(W \le \top\) and pass through
  the sections equivalence \(\mathtt{presheafHomSectionsEquiv}\).
\end{proof}

\begin{lemma}\leanok
[Sectionwise value of \(\mathtt{topSectionToHom}\)]
  \label{lem:top_section_to_hom_app}
  \lean{AlgebraicGeometry.Scheme.Modules.topSectionToHom_app}
  \uses{def:top_section_to_hom}
  At an open \(W\), the morphism \(\mathtt{topSectionToHom}\,s\)
  (\cref{def:top_section_to_hom}) evaluates to the \(\mathtt{Over.mk}\,(W \le \top)\)-component
  of the top section \(s\).
\end{lemma}
\begin{proof}
  \uses{def:top_section_to_hom}
  \leanok
  Proved directly in Lean --- unfolding \(\mathtt{topSectionToHom}\) and the
  \(\mathtt{presheafHom}\) restriction at the slice object \(\mathtt{Over.mk}\,(W \le \top)\).
\end{proof}

\begin{lemma}




\leanok
[Compatible local morphisms of \(\mathcal{O}_X\)-modules glue to a global morphism]
  \label{lem:sheafofmodules_hom_of_local_compat}
  \lean{AlgebraicGeometry.Scheme.Modules.homOfLocalCompat}
  \uses{def:scheme_modules_homMk, lem:open_immersion_slice_sheaf_equiv,
        lem:scheme_modules_hom_local_section, def:top_section_to_hom,
        lem:top_section_to_hom_app, lem:image_preimage_of_le}
  Let \(X\) be a scheme, \(M, N \in \Scheme.\mathtt{Modules}\,X\), and \(\{U_i\}\) an
  open cover of \(X\). Suppose given, for each \(i\), a morphism
  \(f_i \colon M|_{U_i} \to N|_{U_i}\) in \(\Scheme.\mathtt{Modules}\,U_i\) such that
  for all \(i, j\) the restrictions of \(f_i\) and \(f_j\) to \(U_i \cap U_j\) agree.
  Then there is a unique global morphism \(f \colon M \to N\) in
  \(\Scheme.\mathtt{Modules}\,X\) whose restriction to each \(U_i\) is \(f_i\).
\end{lemma}

\begin{proof}
  \uses{def:scheme_modules_homMk, lem:open_immersion_slice_sheaf_equiv,
        lem:scheme_modules_hom_local_section, def:top_section_to_hom,
        lem:top_section_to_hom_app, lem:image_preimage_of_le}
        \leanok
  A morphism of sheaves of \(\mathcal{O}_X\)-modules is, concretely, a morphism of the
  underlying sheaves of abelian groups that is in addition \(\mathcal{O}_X\)-linear
  sectionwise. We build \(f\) in two steps.

  \emph{(i) Glue the underlying ab-sheaf morphism.} Forget \(M\) and \(N\) to their
  underlying sheaves of abelian groups. The morphisms-between-them presheaf
  \(\mathcal{H}\) (\(\mathtt{presheafHom}\) of the underlying presheaves of abelian groups
  of \(M\) and \(N\)), whose value at an open \(W\) is the abelian
  group of morphisms \(M|_W \to N|_W\) of presheaves of abelian groups, is itself a
  \emph{sheaf of types} --- this is the ``hom into a sheaf is a sheaf'' principle
  \(\mathtt{Presheaf.IsSheaf.hom}\), which consumes the sheaf condition of \(N\).
  Regarding \(\mathcal{H}\) as a \(\mathtt{TopCat.Sheaf}\), its universal gluing
  property \(\mathtt{existsUnique\_gluing}\) amalgamates any compatible family of local
  sections of \(\mathcal{H}\) into a unique global section. This route is preferred over
  the raw sieve-level amalgamation \(\mathtt{presheafHom\_isSheafFor}\), which carries
  strictly more bookkeeping.

  The load-bearing datum is, for each \(i\), the \emph{local section}
  \(\mathtt{localSection}\,i \in \mathcal{H}(U_i)\)
  of \cref{lem:scheme_modules_hom_local_section},
  manufactured from \(f_i\). Its component at a pair \((V, h : V \to U_i)\) is the given
  component of \(f_i\), conjugated by the canonical \(\mathtt{eqToHom}\) identifications
  that arise from the open-set equality \(\iota_i\bigl(\iota_i^{-1}(V)\bigr) = V\) valid
  for \(V \subseteq U_i\) (the image under the open immersion \(\iota_i : U_i \to X\) of
  the preimage of \(V\) is \(V\) again). The genuine coherence risk of the whole
  construction is concentrated in the \emph{naturality condition} of
  \(\mathtt{localSection}\,i\): the requirement that these
  \(\mathtt{eqToHom}\)-conjugated components of \(f_i\) commute across the morphisms of
  the slice category \(\mathtt{Over}\,U_i\). This is the one sub-piece to be
  built and verified first, as a standalone lemma.

  This \(\mathtt{localSection}\) naturality is precisely the \emph{section-direction
  slice} of the comparison packaged in \cref{lem:open_immersion_slice_sheaf_equiv}:
  converting the local datum \(f_i\), which lives over the open-immersion pullback
  \(\mathtt{Opens}\,\widetilde{U_i}\), into a section of \(\mathcal{H}\) indexed over the
  slice \(\mathtt{Over}\,U_i\) is the one-sided realisation of the open-immersion
  \(\leftrightarrow\) slice transport. As there, the \(\mathtt{eqToHom}\)-conjugation is
  controlled by the down-set identity \(\iota_i(\iota_i^{-1}(V)) = V\) for
  \(V \subseteq U_i\), and the coherence is automatic on the thin poset
  \(\mathtt{Opens}\,X\) by \(\mathtt{Subsingleton.elim}\). This A-engine uses the
  sheaf-site equivalence \cref{lem:open_immersion_slice_sheaf_equiv}; the C-bridge
  \cref{lem:dual_isLocallyTrivial}, by contrast, uses a per-open slice comparison at
  the presheaf level over the \emph{varying} section ring \(\mathcal{O}_X(V)\), and so
  does \emph{not} reduce to the same sheaf-site root --- the two are independent and
  must be built separately.

  Of the remaining four sub-steps, three are mechanical and one --- the cocycle bridge
  (a) --- carries the genuine content. We spell it out.

  \emph{(a) The cocycle/\(\mathtt{IsCompatible}\) condition on the family
  \((\mathtt{localSection}\,i)_i\).} The condition \(\mathtt{IsCompatible}\) asks that for
  every pair of indices \(i, j\) the two restrictions of \(\mathtt{localSection}\,i\) and
  \(\mathtt{localSection}\,j\) to the overlap coincide as sections of the hom-sheaf
  \(\mathcal{H}\) over \(U_i \sqcap U_j\), i.e.\ an honest equality in the single abelian
  group \(\mathcal{H}(U_i \sqcap U_j)\). We must therefore identify the overlap-agreement
  hypothesis on the family \((f_i)_i\) that furnishes this. The correct hypothesis --- and
  the only one a caller can actually produce --- is the \emph{sectionwise} one:

  \begin{quote}
    for all indices \(i, j\) and every open \(V\) with \(V \subseteq U_i\) and
    \(V \subseteq U_j\), the two section maps \((f_i).\mathtt{app}(V)\) and
    \((f_j).\mathtt{app}(V)\) agree as morphisms in the single, fixed abelian-group hom-type
    \(M(V) \to N(V)\), the value of the hom-presheaf \(\mathcal{H} = \mathcal{H}om(M, N)\) at
    \(V\).
  \end{quote}

  This is exactly the data a caller has at hand: two local morphisms that literally agree on
  every open contained in both \(U_i\) and \(U_j\).

  It is worth contrasting this with a naive alternative that does \emph{not} work. One might
  try to compare the two double-restrictions \((f_i)|_{U_i \cap U_j}\) and
  \((f_j)|_{U_i \cap U_j}\) directly, as morphisms reached by the two slice-restriction
  routes. But these two morphisms do not live in a common type at all: each is the image of
  \(f\) under the pullback-of-modules functor along a \emph{different} slice-restriction
  map, and the resulting source and target objects are sheafifications of pullback presheaves
  along different morphisms. Such objects are only \emph{isomorphic}, through a nontrivial
  comparison isomorphism, and are \emph{not} propositionally equal. A heterogeneous
  comparison of the two double-restrictions therefore cannot be reduced to an honest
  equation --- there is no object equality to transport along, so no
  \(\mathtt{Subsingleton.elim}\) of the thin poset and no heterogeneous-to-homogeneous
  elimination applies --- and, crucially, no caller can produce such a datum in the first
  place. The sectionwise hypothesis sidesteps this entirely by comparing only the section
  maps, which live in the fixed group \(M(V) \to N(V)\).

  With the sectionwise hypothesis in hand the passage to \(\mathtt{IsCompatible}\) is
  \emph{direct}. Restrict \(\mathtt{localSection}\,i\) to the overlap \(U_i \sqcap U_j\). By
  \cref{lem:scheme_modules_hom_local_section}, the component of \(\mathtt{localSection}\,i\)
  at an open \(V \subseteq U_i\) is precisely \((f_i).\mathtt{app}(V')\), where
  \(V' = \iota_i^{-1}(V)\), conjugated by the canonical \(\mathtt{eqToHom}\) identifications
  arising from the down-set identity \(\iota_i\bigl(\iota_i^{-1}(V)\bigr) = V\); restricting
  to the overlap simply evaluates this component at the opens \(V \subseteq U_i \sqcap U_j\),
  and likewise for \(j\). Thus both restrictions, evaluated at such a \(V\), are conjugates
  of \((f_i).\mathtt{app}(V')\) and \((f_j).\mathtt{app}(V')\) by \(\mathtt{eqToHom}\)
  morphisms, all lying in the fixed group \(M(V) \to N(V)\). The sectionwise hypothesis
  equates the two middle terms \((f_i).\mathtt{app}\) and \((f_j).\mathtt{app}\) directly.
  The only remaining identification is that the two \(\mathtt{eqToHom}\)-conjugation routes
  --- built from the two down-set identities, one for \(U_i\) and one for \(U_j\) --- agree;
  since both are parallel morphisms in the thin poset \((\mathtt{Opens}\,X)^{\mathrm{op}}\),
  they coincide by \(\mathtt{Subsingleton.elim}\). Hence the two restricted local sections
  agree in \(\mathcal{H}(U_i \sqcap U_j)\), which is exactly \(\mathtt{IsCompatible}\).

  Note the division of labour: the substantive content is the sectionwise agreement of the
  \(f_i\), an equation in fixed abelian-group hom-types that a caller can supply; the
  \(\mathtt{eqToHom}\)-conjugation and the \(\mathtt{Subsingleton.elim}\) collapse of the
  thin-poset routes are the only mechanical identifications, and they are legitimate
  precisely because they act on the thin poset of opens of \(X\), never on the genuinely
  distinct \(\mathtt{Ab}\)-hom-set section maps.

  Two of the remaining sub-steps are mechanical: (b) gluing via
  \(\mathtt{existsUnique\_gluing}\) and converting the amalgamated section over
  \(\mathrm{op}\,\top\) into a genuine morphism \(g\) of the underlying presheaves of
  \(M\) and \(N\) --- the
  \(\top\)-section-to-morphism passage being the global-sections equivalence
  \(\mathtt{presheafHomSectionsEquiv}\) / \(\mathtt{sheafHomSectionsEquiv}\) of the
  hom-sheaf \(\mathcal{H}\) (which, since \(\top\) is terminal in \(\mathtt{Opens}\,X\),
  lands the glued section as an honest morphism of presheaves), \emph{not} a hand-unfold
  of the \(\top\)-section; and (d) the recovery that the restriction of \(g\) to
  each \(U_i\) is the ab-sheaf morphism underlying \(f_i\).

  Sub-step (c) --- the sectionwise \(\mathcal{O}_X\)-linearity check that promotes the
  glued morphism \(g\) of underlying ab-sheaves to a morphism of \(\mathcal{O}_X\)-modules
  --- is, by contrast, \emph{not} mechanical. It is the point at which the several module
  structures in play (the native section-ring action of \(\mathcal{O}_X\), and the
  restriction-of-scalars action induced by the open immersions) must be reconciled. The
  required identity \(g(r \cdot x) = r \cdot g(x)\), checked sectionwise over an open
  \(P\), splits into three legs.
  \begin{itemize}
    \item \emph{(c.i) The \(M\)-leg (source semilinearity).} Restriction of a section
      against a structure-ring scalar is discharged by the native, \(\Gamma\)-module-level
      functoriality identity \(\mathtt{Scheme.Modules.map\_smul}\,M\): for a restriction
      map \(i\) one has
      \(M(i)(r \cdot x) = X(i)(r) \cdot M(i)(x)\) --- an equation phrased directly over the
      section ring \(\mathcal{O}_X\), with \emph{no} restriction-of-scalars artifact. This
      leg closes cleanly.
    \item \emph{(c.ii) The \(f\)-leg (the genuine obstacle).} The section component
      \((f_i).\mathtt{app}(P)\) is, by construction, linear over the
      \emph{restriction-of-scalars} action carried by \((M|_{\iota_i})(P)\) --- restriction
      of scalars along the structure-ring map of the open immersion \(\iota_i\) --- whereas
      the \(f\)-leg of the goal feeds in the \emph{native}
      \(\mathcal{O}_X(\mathrm{image})\)-module action. These two actions are
      propositionally equal but \emph{not} definitionally equal. The reconciliation is the
      genuine content of (c): for an open immersion the structure-ring comparison
      isomorphism is the identity --- \(\mathtt{appIso}(\iota_i) = \mathtt{Iso.refl}\)
      --- so the ring map along which scalars are restricted is the identity \(\mathbf{1}\);
      and restriction of scalars along the \emph{identity} ring map returns the original
      action,
      \(c \cdot_{\,\mathrm{restr}\,\mathbf{1}} x = \mathbf{1}(c) \cdot x = c \cdot x\),
      identified by \(\mathtt{ModuleCat.restrictScalars.smul\_def}\) together with
      \(\mathtt{restrictScalarsId'App}\). This identity-ring-map smul bridge --- the
      localization observation that restriction of scalars along an isomorphism that
      happens to be the identity reproduces the native module action --- is the substantive
      mathematics of sub-step (c).
    \item \emph{(c.iii) The \(N\)-leg (target semilinearity).} Once the bridge of (c.ii) is
      in place, the target leg is discharged by the native identity
      \(\mathtt{Scheme.Modules.map\_smul}\,N\), exactly as in (c.i); the scalars transported
      across the two legs then reconcile because the two open-set \(\mathtt{eqToHom}\) maps
      arising from the down-set identity \(\iota_i\bigl(\iota_i^{-1}(W)\bigr) = W\) compose
      to the identity on the overlap.
  \end{itemize}

  \emph{(ii) Promote to \(\mathcal{O}_X\)-linear.} The linearity equation
  \(g(r \cdot m) = r \cdot g(m)\) is a sectionwise identity between two sections of a
  separated presheaf. It holds on each \(U_i\), because there \(g\) restricts to the
  module morphism \(f_i\); since the two sides agree on the cover \(\{U_i\}\) and the
  ambient presheaf is separated, they agree globally. Packaging \(g\) together with this
  global linearity through the scheme-level wrapper \cref{def:scheme_modules_homMk}
  (\(\mathtt{homMk}\)) yields a morphism \(f \colon M \to N\) in
  \(\Scheme.\mathtt{Modules}\,X\); its restriction to each \(U_i\) is \(f_i\) and its
  uniqueness are both inherited from the ab-sheaf amalgamation of step (i).

  This gluing engine is substantially larger than a routine glue: it is dominated by the
  \(\mathtt{localSection}\) naturality condition and the convert-and-recover bookkeeping, and
  the naturality is the section-direction slice of
  \cref{lem:open_immersion_slice_sheaf_equiv} --- so its true cost is bounded below by
  building that sheaf-site equivalence, not by the surrounding assembly. It is therefore to be
  scoped as a multi-piece build --- with \(\mathtt{localSection}\), together with its
  naturality, built as a standalone verified lemma first --- rather than a one-shot.
  The step computes no tensor stalk --- it is purely a statement about gluing morphisms
  of sheaves --- and so it does not re-enter the stalk-tensor commutation.
\end{proof}

\begin{remark}
[How the dual infrastructure discharges the \(\otimes\)-inverse]
  \label{rem:dual_discharges_inverse}
  \uses{lem:tensorobj_inverse_invertible, lem:dual_isLocallyTrivial,
        lem:sheafofmodules_hom_of_local_compat, lem:isiso_of_isiso_restrict,
        lem:tensorobj_restrict_iso, lem:tensorobj_unit_iso}
  Once \cref{def:presheaf_internal_hom,def:presheaf_dual,%
  lem:internal_hom_isSheaf,lem:dual_isLocallyTrivial} are available,
  \cref{lem:tensorobj_inverse_invertible} is discharged with no further missing object:
  take \(L^{-1} := \mathtt{dual}\,L\) (\cref{lem:internal_hom_isSheaf}), which is a line
  bundle by \cref{lem:dual_isLocallyTrivial}. The global contraction \(\varepsilon_L
  \colon L \otimes_X L^{-1} \to \mathcal{O}_X\) is \emph{not} obtained by sheafifying
  the presheaf-level evaluation; it is instead built by \emph{gluing the canonical local
  contraction morphisms}. On each trivialising open \(U\), the
  restriction-compatibility iso \cref{lem:tensorobj_restrict_iso} carries
  \((L \otimes_X L^{-1})|_U\) to \(\mathcal{O}_U \otimes_U \mathcal{O}_U\), and the
  canonical local contraction \(\mathcal{O}_U \otimes_U \mathcal{O}_U \to \mathcal{O}_U\)
  of the free rank-one module with its dual is the left unitor
  \cref{lem:tensorobj_unit_iso}, an isomorphism. These local contractions are canonical,
  hence agree on overlaps, and so glue --- by the A-bridge
  \cref{lem:sheafofmodules_hom_of_local_compat} --- to a single global morphism
  \(\varepsilon_L \colon L \otimes_X L^{-1} \to \mathcal{O}_X\). Since ``being an
  isomorphism'' is local on \(X\) and \(\varepsilon_L\) restricts to an isomorphism on
  each \(U\), it is a global isomorphism \(L \otimes_X L^{-1} \xrightarrow{\sim}
  \mathcal{O}_X\) by the B-connector \cref{lem:isiso_of_isiso_restrict}.

  This descent construction deliberately avoids the forbidden
  ``sheafify-the-evaluation'' route --- which would re-enter the abandoned
  tensor-stalk / sheafification-unit whiskering commutation --- and computes no tensor
  stalk anywhere. It is the realization the ``infrastructure-blocked'' note in
  \cref{lem:tensorobj_inverse_invertible} points to.
\end{remark}

\begin{remark}
[Alternative route via object-level descent (fallback)]
  \label{rem:dual_via_stack}
  Should the slice-and-sheafify route of
  \cref{def:presheaf_internal_hom,lem:internal_hom_isSheaf} bottom out, the dual can
  instead be assembled by \emph{object-level descent}: glue the local trivial duals
  \(\mathcal{O}_{U_i}\) along the inverse-transpose transition isomorphisms
  \(g_{ij}^{-\mathsf{T}}\) of \(L\) into a global \(L^{-1}\) on \(X\). Abstractly this
  is effective descent for sheaves of modules on the small Zariski site --- the
  essential-surjectivity of the descent-data functor of an \(\mathtt{IsStack}\)
  pseudofunctor \(U \mapsto (\mathcal{O}_U\text{-modules})\). This is the more
  principled and reusable construction but the larger build (it requires constructing
  the restriction pseudofunctor and proving the stack/effective-descent property),
  and is recorded only as a fallback; the slice internal-hom route above is the
  intended one. It is the same morphism-/object-level descent family on which the
  associator re-route of \cref{lem:tensorobj_assoc_iso} also waits.
\end{remark}

\subsection{Internal-hom and restrict-scalars helper declarations}
\label{subsec:tensorobj_internalhom_helpers}

The following are the internal Lean helper declarations assembling the morphism
module \cref{def:presheaf_internal_hom_value}, the slice internal hom
\cref{def:presheaf_internal_hom}, its restriction maps
\cref{lem:presheaf_internal_hom_restriction}, the presheaf dual
\cref{def:presheaf_dual}, the evaluation morphism \cref{lem:internal_hom_eval}, the
lax-monoidal restrict-scalars structure \cref{lem:restrictscalars_laxmonoidal}, the
ring-iso base-change equivalence \cref{lem:restrictscalars_ringiso_tensorequiv} and
the pushforward adjunction \cref{lem:presheaf_pushforward_adj_substrate}. Each is
proved directly in Lean and carries the dependency edges of the construction.

\begin{definition}\leanok
[Canonical ring map from a terminal object]
  \label{def:term_ring_map}
  \lean{PresheafOfModules.InternalHom.termRingMap}
  \uses{def:presheaf_internal_hom_value}
  For a presheaf of commutative rings \(R\) on a base category with terminal object
  \(T\) and an object \(Y\), the unique morphism \(Y \to T\) induces the ring map
  \(R(T) \to R(Y)\) along which a global scalar acts on the value at \(Y\).
\end{definition}
\begin{proof}
  \uses{def:presheaf_internal_hom_value}
  \leanok
  Proved directly in Lean, as \(R\) applied to the opposite of the terminal morphism.
\end{proof}

\begin{lemma}\leanok
[Naturality of the terminal ring map]
  \label{lem:term_ring_map_naturality}
  \lean{PresheafOfModules.InternalHom.termRingMap_naturality}
  \uses{def:term_ring_map}
  The restriction map \(R(g) : R(X) \to R(Y)\) along \(g : X \to Y\) carries
  \(\mathtt{termRingMap}_X\,f\) to \(\mathtt{termRingMap}_Y\,f\), so the images of a
  global scalar \(f \in R(T)\) are compatible with restriction.
\end{lemma}
\begin{proof}
  \uses{def:term_ring_map}
  \leanok
  Proved directly in Lean, from uniqueness of morphisms into the terminal object and
  functoriality of \(R\).
\end{proof}

\begin{lemma}\leanok
[The terminal ring map at the terminal object is the identity]
  \label{lem:term_ring_map_terminal}
  \lean{PresheafOfModules.InternalHom.termRingMap_terminal}
  \uses{def:term_ring_map}
  At the terminal object itself the induced ring map \(R(T) \to R(T)\) is the
  identity, since the unique \(T \to T\) is \(\mathrm{id}_T\).
\end{lemma}
\begin{proof}
  \uses{def:term_ring_map}
  \leanok
  Proved directly in Lean, using \(\mathtt{IsTerminal.hom\_ext}\) and
  \(R.\mathtt{map\_id}\).
\end{proof}

\begin{definition}\leanok
[Multiplication by a global scalar as an endomorphism]
  \label{def:global_smul}
  \lean{PresheafOfModules.InternalHom.globalSMul}
  \uses{def:term_ring_map, lem:term_ring_map_naturality}
  For \(f \in R(T)\), the endomorphism \(m_f^{N} : N \to N\) of a presheaf of modules
  \(N\) that at each object \(Y\) is multiplication by \(\mathtt{termRingMap}_Y\,f \in
  R(Y)\); its naturality is the heart of the global-scalar action on morphism modules.
\end{definition}
\begin{proof}
  \uses{def:term_ring_map, lem:term_ring_map_naturality}
  \leanok
  Proved directly in Lean: the naturality square commutes by
  \(\mathtt{termRingMap\_naturality}\) and the semilinearity of the restriction maps of
  \(N\).
\end{proof}

\begin{lemma}\leanok
[Sectionwise action of the global-scalar endomorphism]
  \label{lem:global_smul_hom_apply}
  \lean{PresheafOfModules.InternalHom.globalSMul_hom_apply}
  \uses{def:global_smul}
  At an object \(Y\), \(m_f^{N}\) acts on a section as multiplication by
  \(\mathtt{termRingMap}_Y\,f\).
\end{lemma}
\begin{proof}
  \uses{def:global_smul}
  \leanok
  Proved directly in Lean, by definitional unfolding.
\end{proof}

\begin{lemma}\leanok
[The global-scalar endomorphism at one is the identity]
  \label{lem:global_smul_one}
  \lean{PresheafOfModules.InternalHom.globalSMul_one}
  \uses{def:global_smul}
  \(m_1^{N} = \mathrm{id}_N\).
\end{lemma}
\begin{proof}
  \uses{def:global_smul}
  \leanok
  Proved directly in Lean, from \(\mathtt{map\_one}\) and \(1 \cdot s = s\).
\end{proof}

\begin{lemma}\leanok
[The global-scalar endomorphism at zero is zero]
  \label{lem:global_smul_zero}
  \lean{PresheafOfModules.InternalHom.globalSMul_zero}
  \uses{def:global_smul}
  \(m_0^{N} = 0\).
\end{lemma}
\begin{proof}
  \uses{def:global_smul}
  \leanok
  Proved directly in Lean, from \(\mathtt{map\_zero}\) and \(0 \cdot s = 0\).
\end{proof}

\begin{lemma}\leanok
[Additivity of the global-scalar endomorphism]
  \label{lem:global_smul_add}
  \lean{PresheafOfModules.InternalHom.globalSMul_add}
  \uses{def:global_smul}
  \(m_{f+g}^{N} = m_f^{N} + m_g^{N}\).
\end{lemma}
\begin{proof}
  \uses{def:global_smul}
  \leanok
  Proved directly in Lean, from \(\mathtt{map\_add}\) and \((f+g) \cdot s = f \cdot s +
  g \cdot s\).
\end{proof}

\begin{lemma}\leanok
[Multiplicativity of the global-scalar endomorphism]
  \label{lem:global_smul_mul}
  \lean{PresheafOfModules.InternalHom.globalSMul_mul}
  \uses{def:global_smul}
  \(m_{fg}^{N} = m_g^{N} \circ m_f^{N}\): ring multiplication becomes composition,
  the order reversed because the action is by post-composition.
\end{lemma}
\begin{proof}
  \uses{def:global_smul}
  \leanok
  Proved directly in Lean, from \(\mathtt{map\_mul}\) and iterated scalar
  multiplication.
\end{proof}

\begin{definition}\leanok
[Restriction of a presheaf of modules to the over-category]
  \label{def:internal_hom_restr}
  \lean{PresheafOfModules.InternalHom.restr}
  \uses{def:presheaf_internal_hom_value}
  The restriction \(M|_U\) of the slice internal-hom formula: the pushforward of \(M\)
  along the over-category projection \(\mathtt{Over.forget}\,U\), whose base ring at
  the terminal \(\mathtt{Over.mk}\,(\mathrm{id}_U)\) is \(R(U)\).
\end{definition}
\begin{proof}
  \uses{def:presheaf_internal_hom_value}
  \leanok
  Proved directly in Lean, as \(\mathtt{PresheafOfModules.pushforward}_0\) along
  \(\mathtt{Over.forget}\,U\).
\end{proof}

\begin{lemma}\leanok
[Component agreement up to heterogeneous equality]
  \label{lem:hom_app_heq}
  \lean{PresheafOfModules.InternalHom.hom_app_heq}
  \uses{lem:presheaf_internal_hom_restriction}
  The component of a morphism of presheaves of modules at two equal objects agrees up
  to heterogeneous equality.
\end{lemma}
\begin{proof}
  \uses{lem:presheaf_internal_hom_restriction}
  \leanok
  Proved directly in Lean by substituting the object equality.
\end{proof}

\begin{lemma}\leanok
[Identity functoriality of the restriction map]
  \label{lem:restriction_map_id}
  \lean{PresheafOfModules.InternalHom.restrictionMap_id}
  \uses{lem:presheaf_internal_hom_restriction}
  Restricting along the identity \(U \to U\) is the identity on morphism modules.
\end{lemma}
\begin{proof}
  \uses{lem:presheaf_internal_hom_restriction}
  \leanok
  Proved directly in Lean, using \(\mathtt{Over.mapId}\) and the heterogeneous
  component agreement \cref{lem:hom_app_heq}.
\end{proof}

\begin{lemma}\leanok
[Composition functoriality of the restriction map]
  \label{lem:restriction_map_comp}
  \lean{PresheafOfModules.InternalHom.restrictionMap_comp}
  \uses{lem:presheaf_internal_hom_restriction}
  Restricting along \(h \circ g\) is restricting along \(g\) then along \(h\).
\end{lemma}
\begin{proof}
  \uses{lem:presheaf_internal_hom_restriction}
  \leanok
  Proved directly in Lean, using \(\mathtt{Over.mapComp}\) and
  \cref{lem:hom_app_heq}.
\end{proof}

\begin{lemma}\leanok
[The restriction map respects composition of morphisms]
  \label{lem:restriction_map_comp_hom}
  \lean{PresheafOfModules.InternalHom.restrictionMap_comp_hom}
  \uses{lem:presheaf_internal_hom_restriction}
  Restriction carries a composite \(\varphi \circ \psi\) of morphisms of restricted
  modules to the composite of the restrictions.
\end{lemma}
\begin{proof}
  \uses{lem:presheaf_internal_hom_restriction}
  \leanok
  Proved directly in Lean, by functoriality of the pushforward.
\end{proof}

\begin{lemma}\leanok
[Additivity of the restriction map]
  \label{lem:restriction_map_add}
  \lean{PresheafOfModules.InternalHom.restrictionMap_add}
  \uses{lem:presheaf_internal_hom_restriction}
  The restriction map is additive in the morphism it restricts.
\end{lemma}
\begin{proof}
  \uses{lem:presheaf_internal_hom_restriction}
  \leanok
  Proved directly in Lean, sectionwise.
\end{proof}

\begin{lemma}\leanok
[The restriction map preserves zero]
  \label{lem:restriction_map_zero}
  \lean{PresheafOfModules.InternalHom.restrictionMap_zero}
  \uses{lem:presheaf_internal_hom_restriction}
  The restriction map sends the zero morphism to the zero morphism.
\end{lemma}
\begin{proof}
  \uses{lem:presheaf_internal_hom_restriction}
  \leanok
  Proved directly in Lean, sectionwise.
\end{proof}

\begin{lemma}\leanok
[Restriction commutes with the global-scalar endomorphism]
  \label{lem:restriction_map_global_smul}
  \lean{PresheafOfModules.InternalHom.restrictionMap_globalSMul}
  \uses{lem:presheaf_internal_hom_restriction, def:global_smul, lem:term_ring_map_naturality}
  Restricting the multiplication-by-\(r\) endomorphism of \(N|_U\) along \(g : V \to U\)
  yields the multiplication-by-\(R(g)(r)\) endomorphism of \(N|_V\).
\end{lemma}
\begin{proof}
  \uses{lem:presheaf_internal_hom_restriction, def:global_smul, lem:term_ring_map_naturality}
  \leanok
  Proved directly in Lean: both sides act sectionwise by scalar multiplication and the
  scalars agree by \(\mathtt{termRingMap\_naturality}\).
\end{proof}

\begin{lemma}\leanok
[Semilinearity of the restriction map]
  \label{lem:restriction_map_smul}
  \lean{PresheafOfModules.InternalHom.restrictionMap_smul}
  \uses{lem:presheaf_internal_hom_restriction, lem:restriction_map_comp_hom, lem:restriction_map_global_smul}
  For \(f : X \to Y\), a global scalar \(r \in R(X)\) and \(m\),
  \(\mathtt{restrictionMap}\,(r \cdot m) = R(f)(r) \cdot \mathtt{restrictionMap}\,m\) ---
  the compatibility datum consumed by the presheaf assembly.
\end{lemma}
\begin{proof}
  \uses{lem:presheaf_internal_hom_restriction, lem:restriction_map_comp_hom, lem:restriction_map_global_smul}
  \leanok
  Proved directly in Lean, from \(\mathtt{restrictionMap\_comp\_hom}\) and
  \(\mathtt{restrictionMap\_globalSMul}\).
\end{proof}

\begin{definition}\leanok
[Restriction map as an additive group homomorphism]
  \label{def:restriction_map_add_hom}
  \lean{PresheafOfModules.InternalHom.restrictionMapAddHom}
  \uses{lem:presheaf_internal_hom_restriction, lem:restriction_map_add, lem:restriction_map_zero}
  The restriction map packaged as an additive group homomorphism on the morphism
  groups, using its additivity and zero-preservation.
\end{definition}
\begin{proof}
  \uses{lem:presheaf_internal_hom_restriction, lem:restriction_map_add, lem:restriction_map_zero}
  \leanok
  Proved directly in Lean, bundling \(\mathtt{restrictionMap\_zero}\) and
  \(\mathtt{restrictionMap\_add}\).
\end{proof}

\begin{definition}\leanok
[Underlying abelian-group presheaf of the internal hom]
  \label{def:internal_hom_presheaf}
  \lean{PresheafOfModules.InternalHom.internalHomPresheaf}
  \uses{lem:presheaf_internal_hom_restriction, def:restriction_map_add_hom, lem:restriction_map_id, lem:restriction_map_comp}
  The \(\mathbf{Ab}\)-valued presheaf \(U \mapsto (M|_U \to N|_U)\) with the
  further-restriction maps, prior to the \(R(U)\)-module refinement.
\end{definition}
\begin{proof}
  \uses{lem:presheaf_internal_hom_restriction, def:restriction_map_add_hom, lem:restriction_map_id, lem:restriction_map_comp}
  \leanok
  Proved directly in Lean; functoriality is \(\mathtt{restrictionMap\_id}\) /
  \(\mathtt{restrictionMap\_comp}\).
\end{proof}

\begin{definition}\leanok
[Evaluation functional of a dual section]
  \label{def:eval_lin}
  \lean{PresheafOfModules.evalLin}
  \uses{def:presheaf_dual}
  At an object \(X\), the \(R(X)\)-linear map \(M(X) \to R(X)\) obtained by evaluating
  a dual section \(\varphi\) at its terminal component.
\end{definition}
\begin{proof}
  \uses{def:presheaf_dual}
  \leanok
  Proved directly in Lean, as the underlying linear map of the terminal component of
  \(\varphi\).
\end{proof}

\begin{lemma}\leanok
[Additivity of the evaluation functional]
  \label{lem:eval_lin_add}
  \lean{PresheafOfModules.evalLin_add}
  \uses{def:eval_lin}
  \(\mathtt{evalLin}\) is additive in the dual section \(\varphi\).
\end{lemma}
\begin{proof}
  \uses{def:eval_lin}
  \leanok
  Proved directly in Lean, the categorical \(\mathrm{Hom}\)-addition being applied
  sectionwise.
\end{proof}

\begin{lemma}\leanok
[Linearity of the evaluation functional]
  \label{lem:eval_lin_smul}
  \lean{PresheafOfModules.evalLin_smul}
  \uses{def:eval_lin, lem:term_ring_map_terminal}
  \(\mathtt{evalLin}\) is \(R(X)\)-linear in \(\varphi\), using the morphism-module
  action and the terminal ring map.
\end{lemma}
\begin{proof}
  \uses{def:eval_lin, lem:term_ring_map_terminal}
  \leanok
  Proved directly in Lean, via \(c \cdot \varphi = \varphi \circ m_c\) and
  \(\mathtt{termRingMap\_terminal}\).
\end{proof}

\begin{definition}\leanok
[Open-by-open evaluation/contraction map]
  \label{def:internal_hom_eval_app}
  \lean{PresheafOfModules.internalHomEvalApp}
  \uses{def:eval_lin, lem:eval_lin_add, lem:eval_lin_smul}
  At an object \(X\), the \(R(X)\)-bilinear contraction \((M(X)) \otimes_{R(X)}
  (M|_X \to R|_X) \to R(X)\), \(s \otimes \varphi \mapsto \varphi(s)\).
\end{definition}
\begin{proof}
  \uses{def:eval_lin, lem:eval_lin_add, lem:eval_lin_smul}
  \leanok
  Proved directly in Lean, lifting the bilinear contraction through the tensor
  product; bilinearity is \(\mathtt{evalLin\_add}\) / \(\mathtt{evalLin\_smul}\).
\end{proof}

\begin{lemma}\leanok
[Value of the contraction on a simple tensor]
  \label{lem:internal_hom_eval_app_tmul}
  \lean{PresheafOfModules.internalHomEvalApp_tmul}
  \uses{def:internal_hom_eval_app}
  The contraction sends \(s \otimes \varphi\) to \(\mathtt{evalLin}\,\varphi\,s\).
\end{lemma}
\begin{proof}
  \uses{def:internal_hom_eval_app}
  \leanok
  Proved directly in Lean, by definitional unfolding.
\end{proof}

\begin{lemma}\leanok
[Further restriction along the canonical over-morphism]
  \label{lem:restr_map_hom_mk}
  \lean{PresheafOfModules.restr_map_homMk}
  \uses{def:internal_hom_restr}
  For \(f : X \to Y\), further-restricting \(M|_{X}\) along the canonical over-category
  morphism is definitionally \(M.\mathtt{map}\,f\).
\end{lemma}
\begin{proof}
  \uses{def:internal_hom_restr}
  \leanok
  Proved directly in Lean, by definitional unfolding (extracted to keep its
  \(\mathtt{whnf}\) cost local).
\end{proof}

\begin{lemma}\leanok
[Sectionwise applied form of the pushforward natural transformation]
  \label{lem:pushforward_nat_trans_app_app_apply}
  \lean{PresheafOfModules.pushforwardNatTrans_app_app_apply}
  \uses{lem:presheaf_pushforward_adj_substrate}
  The component of \(\mathtt{pushforwardNatTrans}\,\varphi\,\alpha\) at \(M\), \(U\)
  acts on a section as \(M.\mathtt{map}\,(\alpha_U)^{\mathrm{op}}\).
\end{lemma}
\begin{proof}
  \uses{lem:presheaf_pushforward_adj_substrate}
  \leanok
  Proved directly in Lean, by definitional unfolding.
\end{proof}

\begin{lemma}\leanok
[Sectionwise form of the pushforward congruence (forward)]
  \label{lem:pushforward_congr_hom_app_app}
  \lean{PresheafOfModules.pushforwardCongr_hom_app_app}
  \uses{lem:presheaf_pushforward_adj_substrate}
  The forward component of \(\mathtt{pushforwardCongr}\,e\) acts as the identity on
  sections.
\end{lemma}
\begin{proof}
  \uses{lem:presheaf_pushforward_adj_substrate}
  \leanok
  Proved directly in Lean, by substituting the equality \(e\).
\end{proof}

\begin{lemma}\leanok
[Sectionwise form of the pushforward congruence (inverse)]
  \label{lem:pushforward_congr_inv_app_app}
  \lean{PresheafOfModules.pushforwardCongr_inv_app_app}
  \uses{lem:presheaf_pushforward_adj_substrate}
  The inverse component of \(\mathtt{pushforwardCongr}\,e\) acts as the identity on
  sections.
\end{lemma}
\begin{proof}
  \uses{lem:presheaf_pushforward_adj_substrate}
  \leanok
  Proved directly in Lean, by substituting the equality \(e\).
\end{proof}

\begin{definition}\leanok
[Lax-monoidal unit of restrict-scalars]
  \label{def:restrict_scalars_lax_epsilon}
  \lean{PresheafOfModules.restrictScalarsLaxε}
  \uses{lem:restrictscalars_laxmonoidal}
  % NOTE (iter-306, review): this declaration EXISTS and is axiom-clean —
  % `PresheafOfModules.restrictScalarsLaxε` at PresheafInternalHom.lean:290, with a proved
  % `NatTrans.naturality` field — and is DIRECTLY IMPORTED by DualInverse.lean (line 7), hence in
  % scope at the DUAL naturality sorries. The iter-306 DUAL prover wrongly declared an
  % "architectural wall" claiming this NatTrans is absent (it ran `lean_local_search` -> empty and
  % missed the imported decl). The 3 DUAL naturality sorries should USE this, not rebuild it.
  The lax-monoidal unit \(\varepsilon\) of \(\mathtt{restrictScalars}\,\alpha\),
  assembled sectionwise from the \(\mathbf{ModuleCat}\)-level lax units.
\end{definition}
\begin{proof}
  \uses{lem:restrictscalars_laxmonoidal}
  \leanok
  Proved directly in Lean, sectionwise from
  \(\mathtt{ModuleCat.restrictScalars}\)'s lax unit.
\end{proof}

\begin{definition}\leanok
[Lax-monoidal tensorator of restrict-scalars]
  \label{def:restrict_scalars_lax_mu}
  \lean{PresheafOfModules.restrictScalarsLaxμ}
  \uses{lem:restrictscalars_laxmonoidal}
  The lax-monoidal tensorator \(\mu\) of \(\mathtt{restrictScalars}\,\alpha\),
  assembled sectionwise from the \(\mathbf{ModuleCat}\)-level lax tensorators.
\end{definition}
\begin{proof}
  \uses{lem:restrictscalars_laxmonoidal}
  \leanok
  Proved directly in Lean, sectionwise from
  \(\mathtt{ModuleCat.restrictScalars}\)'s lax tensorator.
\end{proof}

\begin{lemma}\leanok
[Ring-iso base-change equivalence on a simple tensor]
  \label{lem:restrictscalars_ringiso_tensorequiv_apply_tmul}
  \lean{restrictScalarsRingIsoTensorEquiv_apply_tmul}
  \uses{lem:restrictscalars_ringiso_tensorequiv}
  The forward map of \(\mathtt{restrictScalarsRingIsoTensorEquiv}\,e\) sends
  \(a \otimes_R b\) to \(a \otimes_S b\).
\end{lemma}
\begin{proof}
  \uses{lem:restrictscalars_ringiso_tensorequiv}
  \leanok
  Proved directly in Lean, unfolding the linear equivalence on a simple tensor.
\end{proof}

\subsection{Vestigial / off-path helper declarations}
\label{subsec:tensorobj_vestigial_helpers}

The following are the off-critical-path Lean helper declarations supporting the
flat-whiskering localizer stability \cref{lem:flat_whisker_localizer}, the
flatness-free whiskering \cref{lem:whisker_of_W}, the \(R_x\)-linear stalk map
\cref{lem:stalk_linear_map}, the stalkwise localizer characterisation
\cref{lem:W_implies_stalkwise_iso} and the shared slice-site equivalence
\cref{lem:open_immersion_slice_sheaf_equiv}. Each is proved directly in Lean.

\begin{lemma}\leanok
[Sectionwise left-whiskering on a simple tensor]
  \label{lem:to_presheaf_whisker_left_app_tmul}
  \lean{PresheafOfModules.toPresheaf_whiskerLeft_app_tmul}
  \uses{lem:flat_whisker_localizer, def:scheme_modules_tensorobj}
  The underlying additive map of \((F \triangleleft g)_X\) sends \(a \otimes b\) to
  \(a \otimes g_X(b)\).
\end{lemma}
\begin{proof}
  \uses{lem:flat_whisker_localizer, def:scheme_modules_tensorobj}
  \leanok
  Proved directly in Lean, as the action of \(\mathtt{LinearMap.lTensor}\) on a simple
  tensor.
\end{proof}

\begin{lemma}\leanok
[Sectionwise left-whiskering is a tensor map]
  \label{lem:to_presheaf_whisker_left_app_apply}
  \lean{PresheafOfModules.toPresheaf_whiskerLeft_app_apply}
  \uses{lem:to_presheaf_whisker_left_app_tmul}
  The underlying additive map of \((F \triangleleft g)_X\) is
  \(\mathtt{LinearMap.lTensor}\,(F(X))\,g_X\).
\end{lemma}
\begin{proof}
  \uses{lem:to_presheaf_whisker_left_app_tmul}
  \leanok
  Proved directly in Lean, identifying the whiskered map with the left-tensor map.
\end{proof}

\begin{lemma}\leanok
[Local surjectivity is preserved by left-whiskering]
  \label{lem:is_locally_surjective_whisker_left}
  \lean{PresheafOfModules.isLocallySurjective_whiskerLeft}
  \uses{lem:to_presheaf_whisker_left_app_tmul, lem:flat_whisker_localizer}
  If \(g\) is locally surjective then so is \(F \triangleleft g\), with no flatness
  hypothesis (pure right-exactness of \(\otimes\)).
\end{lemma}
\begin{proof}
  \uses{lem:to_presheaf_whisker_left_app_tmul, lem:flat_whisker_localizer}
  \leanok
  Proved directly in Lean, lifting a local preimage of \(b\) to \(a \otimes b\).
\end{proof}

\begin{lemma}\leanok
[Local injectivity is preserved by flat left-whiskering]
  \label{lem:is_locally_injective_whisker_left_of_flat}
  \lean{PresheafOfModules.isLocallyInjective_whiskerLeft_of_flat}
  \uses{lem:to_presheaf_whisker_left_app_apply, lem:flat_whisker_localizer}
  For sectionwise-flat \(F\), if \(g\) is locally injective then so is
  \(F \triangleleft g\).
\end{lemma}
\begin{proof}
  \uses{lem:to_presheaf_whisker_left_app_apply, lem:flat_whisker_localizer}
  \leanok
  Proved directly in Lean, via \(\mathtt{Module.Flat.lTensor\_exact}\) on the kernel
  exact sequence and local injectivity of \(g\) on each simple tensor.
\end{proof}

\begin{lemma}\leanok
[Flat right-whiskering preserves the localizer]
  \label{lem:w_whisker_right_of_flat}
  \lean{PresheafOfModules.W_whiskerRight_of_flat}
  \uses{lem:flat_whisker_localizer}
  For sectionwise-flat \(F\) and \(g \in J.W\), the right-whiskered \(g \triangleright F\)
  again lies in \(J.W\).
\end{lemma}
\begin{proof}
  \uses{lem:flat_whisker_localizer}
  \leanok
  Proved directly in Lean, conjugating the left-whiskered statement by the braiding.
\end{proof}

\begin{lemma}\leanok
[Flatness-free right-whiskering preserves the localizer]
  \label{lem:w_whisker_right_of_w}
  \lean{PresheafOfModules.W_whiskerRight_of_W}
  \uses{lem:whisker_of_W}
  For arbitrary \(F\) and a locally bijective \(g \in J.W\), the right-whiskered
  \(g \triangleright F\) again lies in \(J.W\).
\end{lemma}
\begin{proof}
  \uses{lem:whisker_of_W}
  \leanok
  Proved directly in Lean, conjugating the flatness-free left-whiskered statement by
  the braiding.
\end{proof}

\begin{lemma}\leanok
[Germ characterisation of the linear stalk map]
  \label{lem:stalk_linear_map_germ}
  \lean{PresheafOfModules.stalkLinearMap_germ}
  \uses{lem:stalk_linear_map}
  On the germ of a section \(s\) over \(U\), \(\mathtt{stalkLinearMap}\,g\,x\) is the
  germ of \(g_U(s)\).
\end{lemma}
\begin{proof}
  \uses{lem:stalk_linear_map}
  \leanok
  Proved directly in Lean, the defining naturality of the stalk map.
\end{proof}

\begin{lemma}\leanok
[A stalkwise-iso morphism is bijective on stalks]
  \label{lem:stalk_linear_map_bijective_of_isiso}
  \lean{PresheafOfModules.stalkLinearMap_bijective_of_isIso}
  \uses{lem:stalk_linear_map}
  If the underlying \(\mathbf{Ab}\)-stalk map of \(g\) at \(x\) is an isomorphism then
  \(\mathtt{stalkLinearMap}\,g\,x\) is bijective.
\end{lemma}
\begin{proof}
  \uses{lem:stalk_linear_map}
  \leanok
  Proved directly in Lean, from \(\mathtt{ConcreteCategory.bijective\_of\_isIso}\).
\end{proof}

\begin{definition}\leanok
[Linear stalk equivalence of a stalkwise-iso morphism]
  \label{def:stalk_linear_equiv_of_isiso}
  \lean{PresheafOfModules.stalkLinearEquivOfIsIso}
  \uses{lem:stalk_linear_map, lem:stalk_linear_map_bijective_of_isiso}
  When the \(\mathbf{Ab}\)-stalk map of \(g\) at \(x\) is an isomorphism,
  \(\mathtt{stalkLinearMap}\,g\,x\) bundles to an \(R_x\)-linear equivalence
  \(M_x \xrightarrow{\sim} N_x\).
\end{definition}
\begin{proof}
  \uses{lem:stalk_linear_map, lem:stalk_linear_map_bijective_of_isiso}
  \leanok
  Proved directly in Lean, via \(\mathtt{LinearEquiv.ofBijective}\).
\end{proof}

\begin{lemma}\leanok
[Injective on stalks implies locally injective]
  \label{lem:is_locally_injective_of_injective_stalk}
  \lean{PresheafOfModules.isLocallyInjective_of_injective_stalk}
  \uses{lem:W_implies_stalkwise_iso}
  For a morphism of presheaves valued in a concrete category over a topological space,
  stalkwise injectivity implies local injectivity for the topological Grothendieck
  topology.
\end{lemma}
\begin{proof}
  \uses{lem:W_implies_stalkwise_iso}
  \leanok
  Proved directly in Lean, through \(\mathtt{germ\_eq}\): agreeing germs agree on a
  neighbourhood, which covers each point.
\end{proof}

\begin{lemma}\leanok
[Locally injective implies injective on stalks]
  \label{lem:injective_stalk_of_is_locally_injective}
  \lean{PresheafOfModules.injective_stalk_of_isLocallyInjective}
  \uses{lem:W_implies_stalkwise_iso}
  The converse: a locally injective morphism of presheaves over a topological space
  induces injective stalk maps.
\end{lemma}
\begin{proof}
  \uses{lem:W_implies_stalkwise_iso}
  \leanok
  Proved directly in Lean, via \(\mathtt{germ\_eq}\) and the covering of the equalizer
  sieve.
\end{proof}

\begin{lemma}\leanok
[Pointwise cover-correspondence across the open embedding]
  \label{lem:over_equiv_image_cover_iff}
  \lean{AlgebraicGeometry.Scheme.Modules.overEquiv_image_cover_iff}
  \uses{lem:open_immersion_slice_sheaf_equiv}
  A sieve \(S\) on \(W : \mathtt{Opens}\,\widetilde{U}\) is a covering sieve in the
  subspace topology iff its pushforward along the open-embedding image functor
  \(\widetilde{U} \hookrightarrow X\) is a covering sieve in \(X\).
\end{lemma}
\begin{proof}
  \uses{lem:open_immersion_slice_sheaf_equiv}
  \leanok
  Proved directly in Lean: both sides are the pointwise neighbourhood-cover condition,
  matched across the injection \(\mathtt{Subtype.val}\).
\end{proof}

\begin{definition}\leanok
[Dense-subsite datum for the slice-site equivalence]
  \label{def:over_equiv_inverse_is_dense_subsite}
  \lean{AlgebraicGeometry.Scheme.Modules.overEquivInverseIsDenseSubsite}
  \uses{lem:open_immersion_slice_sheaf_equiv, lem:over_equiv_image_cover_iff}
  The inverse of \(\mathtt{overEquivalence}\,U\) exhibits the opens of the subspace
  \(\widetilde{U}\) as a dense subsite of the slice site over \(U\); the only
  non-formal field is the pointwise cover-correspondence.
\end{definition}
\begin{proof}
  \uses{lem:open_immersion_slice_sheaf_equiv, lem:over_equiv_image_cover_iff}
  \leanok
  Proved directly in Lean, discharging \(\mathtt{functorPushforward\_mem\_iff}\) by
  \(\mathtt{overEquiv\_image\_cover\_iff}\).
\end{proof}

\section{Substrate helper declarations (wired lean-aux blocks)}
\label{sec:tensorobj_substrate_helpers}

This section gives a blueprint entry to each of the project-local helper declarations
of \texttt{TensorObjSubstrate.lean} that the substrate construction and the
pullback--tensor monoidality machinery invoke internally. Each had been written only
as an inline mid-proof \(\mathtt{\backslash lean}\) reference of an existing block,
which leaves it unregistered as a graph node; here each receives its own \(\mathtt{\backslash label}\)
and its \(\mathtt{\backslash uses}\) edges are read off its Lean body, and the consumers
above are wired to depend on them. All but one are sorry-free; the single residual is the
\(\mathtt{mateEquiv\_vcomp}\) cocycle interior held open inside the proof of
\cref{lem:sheafificationcomppullback_comp}, whose surrounding reduction (transpose, inline
R0-kill, conjugate collapse) is given in full.

\begin{definition}\leanok
[Substrate \(\otimes\) respects a pair of isomorphisms]
  \label{def:tensorobjisoofiso}
  \lean{AlgebraicGeometry.Scheme.Modules.tensorObjIsoOfIso}
  \uses{def:scheme_modules_tensorobj}
  For \(M \cong M'\) and \(N \cong N'\) in \(\Scheme.\mathtt{Modules}\,X\), sheafifying the
  presheaf-of-modules tensor of the underlying presheaf isomorphisms produces an
  isomorphism \(M \otimes_X N \cong M' \otimes_X N'\). This is the functor-of-isomorphisms
  form of the substrate tensor, used to build the middle-four interchange and the
  well-definedness of the multiplication on iso-classes.
\end{definition}
\begin{proof}
  \uses{def:scheme_modules_tensorobj}
  Proved directly in Lean: it is \(\mathtt{sheafification}.\mathrm{mapIso}\) applied to the
  monoidal \(\mathtt{tensorIso}\) of the two forgetful images of the given isomorphisms.
\end{proof}

\begin{definition}\leanok
[\(\mathcal{O}_X \otimes_X \mathcal{O}_X \cong \mathcal{O}_X\)]
  \label{def:tensorobj_unit_self_iso}
  \lean{AlgebraicGeometry.Scheme.Modules.tensorObj_unit_iso}
  \uses{def:scheme_modules_tensorobj}
  The substrate tensor of the structure sheaf with itself is the structure sheaf:
  \(\mathtt{tensorObj}\,\mathcal{O}_X\,\mathcal{O}_X \cong \mathcal{O}_X\), where
  \(\mathcal{O}_X = \mathtt{SheafOfModules.unit}\,X.\mathtt{ringCatSheaf}\). It is the unit
  isomorphism consumed by the invertibility constructions.
\end{definition}
\begin{proof}
  \uses{def:scheme_modules_tensorobj}
  Proved directly in Lean: the sheafification of the presheaf left unitor
  \(\lambda_{\mathbf{1}}\), composed with the sheafification-adjunction counit at the
  already-sheaf unit.
\end{proof}

\begin{definition}\leanok
[Middle-four interchange for \(\otimes_X\)]
  \label{def:tensorobj_middlefour}
  \lean{AlgebraicGeometry.Scheme.Modules.tensorObj_middleFour}
  \uses{lem:tensorobj_assoc_iso, lem:tensorobj_comm_iso, def:tensorobjisoofiso}
  For arbitrary \(A,B,C,D \in \Scheme.\mathtt{Modules}\,X\) there is an isomorphism
  \((A \otimes B) \otimes (C \otimes D) \cong (A \otimes C) \otimes (B \otimes D)\),
  reassociating the four factors. It is the bookkeeping isomorphism used to show
  \(\otimes\)-invertibility is closed under tensor product.
\end{definition}
\begin{proof}
  \uses{lem:tensorobj_assoc_iso, lem:tensorobj_comm_iso, def:tensorobjisoofiso}
  Proved directly in Lean: an assembly of the unconditional associator
  \cref{lem:tensorobj_assoc_iso} and the braiding \cref{lem:tensorobj_comm_iso} through
  \(\mathtt{tensorObjIsoOfIso}\); no coherence pentagon is consumed.
\end{proof}

\begin{definition}\leanok
[Refining a trivialisation to a smaller open]
  \label{def:restrict_iso_unit_of_le}
  \lean{AlgebraicGeometry.Scheme.Modules.restrictIsoUnitOfLE}
  \uses{lem:pullback_unit_iso}
  If \(M\) is trivialised on an open \(U\) (i.e. \(M|_U \cong \mathcal{O}_U\)) then it is
  trivialised on every open \(W \le U\): \(M|_W \cong \mathcal{O}_W\).
\end{definition}
\begin{proof}
  \uses{lem:pullback_unit_iso}
  Proved directly in Lean: factor \(W \hookrightarrow X\) through \(W \hookrightarrow U\),
  transport the given trivialisation along the open immersion via the restriction/pullback
  comparison functors, and identify the restricted unit with the unit through the
  pullback unit comparison \cref{lem:pullback_unit_iso} (an isomorphism because the
  open immersion's \(\mathtt{Opens.map}\) is \(\mathtt{Final}\)).
\end{proof}

\begin{lemma}\leanok
[Underlying presheaf map of a promoted module morphism]
  \label{lem:topresheaf_map_hommk}
  \lean{AlgebraicGeometry.Scheme.Modules.toPresheaf_map_homMk}
  \uses{def:scheme_modules_tensorobj}
  The underlying \(\mathtt{Ab}\)-presheaf morphism of a module morphism promoted from a
  sectionwise-\(\mathcal{O}_X\)-linear presheaf map \(g\) is \(g\) itself. This is the
  round-trip identity used by the gluing descent of
  \cref{lem:tensorobj_inverse_invertible}, where local contraction isomorphisms are glued
  to a global module morphism.
\end{lemma}
\begin{proof}
  \uses{def:scheme_modules_tensorobj}
  Proved directly in Lean by \(\mathtt{rfl}\): promotion wraps \(\mathtt{PresheafOfModules.homMk}\,g\)
  and the underlying-presheaf projection unwraps it.
\end{proof}

\begin{lemma}\leanok
[Unit comparison is an iso from a finality hypothesis]
  \label{lem:isiso_pbu_of_final}
  \lean{AlgebraicGeometry.Scheme.Modules.isIso_pbu_of_final}
  For \(g : Y \to X\) with \((\mathtt{Opens.map}\,g.\mathtt{base}).\mathtt{Final}\), the
  Mathlib unit comparison \(\mathtt{pullbackObjUnitToUnit}\,g\) is an isomorphism.
\end{lemma}
\begin{proof}
  Proved directly in Lean by \(\mathtt{inferInstance}\): the lone finality hypothesis is
  isolated at a clean site so the Mathlib \(\mathtt{OfFinal}\) instance synthesises
  without colliding with other in-scope \(\mathtt{Final}\)/\(\mathtt{IsIso}\) hypotheses.
\end{proof}

\begin{definition}\leanok
[Bundled iso form of the unit comparison]
  \label{def:pullback_obj_unit_to_unit_iso}
  \lean{AlgebraicGeometry.Scheme.Modules.pullbackObjUnitToUnitIso}
  \uses{lem:isiso_pbu_of_final}
  For \(g : Y \to X\) with \((\mathtt{Opens.map}\,g.\mathtt{base}).\mathtt{Final}\), the
  isomorphism \(g^*\mathcal{O}_X \cong \mathcal{O}_Y\) bundling the unit comparison
  \(\mathtt{pullbackObjUnitToUnit}\,g\) of \cref{lem:isiso_pbu_of_final}. Handing out the
  bundled \(\mathtt{Iso}\) keeps downstream reasoning at the \(\mathtt{Iso}\) level.
\end{definition}
\begin{proof}
  \uses{lem:isiso_pbu_of_final}
  Proved directly in Lean: \(\mathtt{asIso}\) of the comparison with the explicit witness
  \cref{lem:isiso_pbu_of_final}.
\end{proof}

\begin{lemma}\leanok
[\(\mathtt{hom}\) of the bundled unit comparison]
  \label{lem:pullback_obj_unit_to_unit_iso_hom}
  \lean{AlgebraicGeometry.Scheme.Modules.pullbackObjUnitToUnitIso_hom}
  \uses{def:pullback_obj_unit_to_unit_iso}
  The \(\mathtt{hom}\) component of \(\mathtt{pullbackObjUnitToUnitIso}\,g\) is the bare
  unit comparison \(\mathtt{pullbackObjUnitToUnit}\,g\).
\end{lemma}
\begin{proof}
  \uses{def:pullback_obj_unit_to_unit_iso}
  Proved directly in Lean by \(\mathtt{rfl}\) (\(\mathtt{asIso}\) exposes its underlying map).
\end{proof}

\begin{definition}\leanok
[Sheafification--monoidality reconciliation brick]
  \label{def:sheafify_tensor_unit_iso}
  \lean{AlgebraicGeometry.Scheme.Modules.sheafifyTensorUnitIso}
  \uses{lem:isiso_sheafification_map_of_W, lem:whisker_of_W}
  For presheaves of modules \(P,Q\) on \(X\), sheafifying the presheaf tensor \(P \otimes Q\)
  agrees with sheafifying the tensor of the underlying presheaves of their sheafifications.
  This is the ``sheafification is monoidal'' reconciliation bridging a presheaf-level
  tensorator with the substrate \(\otimes_X\).
\end{definition}
\begin{proof}
  \uses{lem:isiso_sheafification_map_of_W, lem:whisker_of_W}
  Proved directly in Lean, exactly as in \cref{lem:tensorobj_assoc_iso}: whisker the
  sheafification unit \(\eta\) (a \(J.W\)-morphism) on each side and invert under
  sheafification via \cref{lem:isiso_sheafification_map_of_W} together with the
  flatness-free whiskering stability \cref{lem:whisker_of_W}.
\end{proof}

\begin{lemma}\leanok
[\(\mathtt{hom}\) of the reconciliation brick as a whisker composite]
  \label{lem:sheafify_tensor_unit_iso_hom_eq}
  \lean{AlgebraicGeometry.Scheme.Modules.sheafifyTensorUnitIso_hom_eq}
  \uses{def:sheafify_tensor_unit_iso}
  The \(\mathtt{hom}\) of \cref{def:sheafify_tensor_unit_iso} equals the sheafification of
  \(\eta_P \rhd Q\) followed by the sheafification of \((aP).\mathrm{val} \lhd \eta_Q\),
  stated on the common \(\circ\,\mathtt{forget}_2\) carrier so the two whisker factors can
  be merged.
\end{lemma}
\begin{proof}
  \uses{def:sheafify_tensor_unit_iso}
  Proved directly in Lean by \(\mathtt{rfl}\), stripping the carrier cast.
\end{proof}

\begin{lemma}\leanok
[\(\mathtt{hom}\) of the reconciliation brick as one \(\mathtt{tensorHom}\)]
  \label{lem:sheafify_tensor_unit_iso_hom_eq_prime}
  \lean{AlgebraicGeometry.Scheme.Modules.sheafifyTensorUnitIso_hom_eq'}
  \uses{lem:sheafify_tensor_unit_iso_hom_eq, def:sheafify_tensor_unit_iso}
  The \(\mathtt{hom}\) of \cref{def:sheafify_tensor_unit_iso} equals the sheafification of
  the single \(\mathtt{tensorHom}\) \(\eta_P \otimes \eta_Q\) of the two sheafification-unit
  components --- keeping every term in one monoidal instance so naturality reduces to plain
  bifunctoriality.
\end{lemma}
\begin{proof}
  \uses{lem:sheafify_tensor_unit_iso_hom_eq}
  Proved directly in Lean: merge the two whisker factors of
  \cref{lem:sheafify_tensor_unit_iso_hom_eq} by \(\mathtt{\leftarrow Functor.map\_comp}\)
  and identify with \(\mathtt{tensorHom\_def}\).
\end{proof}

\begin{definition}\leanok
[Codomain unit reconciliation \(a_Y\,\mathbf{1} \cong \mathcal{O}_Y\)]
  \label{def:sheafify_unit_iso}
  \lean{AlgebraicGeometry.Scheme.Modules.sheafifyUnitIso}
  The sheafification counit identifies the sheafified presheaf monoidal unit
  \(a_Y.\mathrm{obj}\,\mathbf{1}\) with the sheaf-level structure module
  \(\mathcal{O}_Y\). It is the right-hand vertical of the \(\eta\)-bridge square of
  \cref{lem:eta_bridge_unit_square}.
\end{definition}
\begin{proof}
  Proved directly in Lean: the unit \(\mathbf{1} = (\mathtt{unit}\,Y).\mathrm{val}\) is
  already a sheaf, so the sheafification-adjunction counit at it is an isomorphism, taken
  via \(\mathtt{asIso}\).
\end{proof}

\begin{definition}\leanok
[Sheafified presheaf-pullback \(\cong\) abstract pullback on a value]
  \label{def:pullback_val_iso}
  \lean{AlgebraicGeometry.Scheme.Modules.pullbackValIso}
  \uses{lem:sheafification_comp_pullback_eq_leftadjointuniq}
  For \(f : Y \to X\) and \(M \in \Scheme.\mathtt{Modules}\,X\), sheafifying the
  presheaf-level pullback of the underlying presheaf \(M.\mathrm{val}\) recovers the
  abstract \(\mathcal{O}\)-module pullback \((\mathtt{pullback}\,f).\mathrm{obj}\,M\). It
  reconciles the right-hand side of \(\mathtt{pullbackTensorMap}\) with the substrate
  tensor of the pullbacks.
\end{definition}
\begin{proof}
  \uses{lem:sheafification_comp_pullback_eq_leftadjointuniq}
  Proved directly in Lean: the inverse of
  \(\mathtt{sheafificationCompPullback}\) (\cref{lem:sheafification_comp_pullback_eq_leftadjointuniq})
  at \(M.\mathrm{val}\), composed with the pullback of the sheafification counit at the
  already-sheaf \(M\).
\end{proof}

\begin{definition}\leanok
[Topological inverse image \(\mathtt{pullback0}\)]
  \label{def:pullback0}
  \lean{AlgebraicGeometry.Scheme.Modules.pullback0}
  \uses{lem:pushforward0_is_right_adjoint}
  The left adjoint of the fixed-ring pushforward \(\mathtt{pushforward}_0\,F\,R\); on
  underlying presheaves it is the left Kan extension along \(F^{\mathrm{op}}\). A
  project-local carrier of the pullback Lan decomposition.
\end{definition}
\begin{proof}
  \uses{lem:pushforward0_is_right_adjoint}
  Proved directly in Lean: \(\mathtt{pushforward}_0\,F\,R\) is definitionally a pushforward
  hence a right adjoint (\cref{lem:pushforward0_is_right_adjoint}), and \(\mathtt{pullback0}\)
  is its \(\mathtt{leftAdjoint}\).
\end{proof}

\begin{definition}\leanok
[Adjunction \(\mathtt{pullback0} \dashv \mathtt{pushforward}_0\)]
  \label{def:pullback0_adjunction}
  \lean{AlgebraicGeometry.Scheme.Modules.pullback0Adjunction}
  \uses{def:pullback0, lem:pushforward0_is_right_adjoint}
  The adjunction exhibiting \cref{def:pullback0} as the left adjoint of the fixed-ring
  pushforward.
\end{definition}
\begin{proof}
  \uses{def:pullback0, lem:pushforward0_is_right_adjoint}
  Proved directly in Lean by \(\mathtt{Adjunction.ofIsRightAdjoint}\) of the pushforward
  (\cref{lem:pushforward0_is_right_adjoint}).
\end{proof}

\begin{lemma}\leanok
[The sheafified \(\otimes\) of the two \(\mathtt{pullbackValIso}\) is an iso]
  \label{lem:isiso_sheafify_tensorhom_pullbackvaliso}
  \lean{AlgebraicGeometry.Scheme.Modules.isIso_sheafify_tensorHom_pullbackValIso}
  \uses{def:pullback_val_iso, def:scheme_modules_tensorobj}
  The fourth factor of \(\mathtt{pullbackTensorMap}\) --- the sheafification of the
  \(\mathtt{tensorHom}\) of the two \(\mathtt{pullbackValIso}\) comparisons (one for \(M\),
  one for \(N\)) --- is an isomorphism.
\end{lemma}
\begin{proof}
  \uses{def:pullback_val_iso}
  Proved directly in Lean: it is the \(\mathtt{hom}\) of the sheafification \(\mathtt{mapIso}\)
  of the monoidal \(\mathtt{tensorIso}\) of the two \(\mathtt{pullbackValIso}\) isomorphisms.
\end{proof}

\begin{lemma}\leanok
[Converse of \(\mathtt{isIso\_sheafification\_map\_of\_W}\)]
  \label{lem:w_of_isiso_sheafification}
  \lean{AlgebraicGeometry.Scheme.Modules.W_of_isIso_sheafification}
  \uses{lem:isiso_sheafification_map_of_W}
  A morphism of presheaves of modules whose image under the associated-sheaf functor is an
  isomorphism lies, on underlying additive presheaves, in the sheafification localizer
  \(J.W\).
\end{lemma}
\begin{proof}
  \uses{lem:isiso_sheafification_map_of_W}
  Proved directly in Lean: it is the morphism-property identity
  ``the sheafification functor is the localization at \(J.W.\mathtt{inverseImage}\)''
  (\cref{lem:isiso_sheafification_map_of_W}) read backwards.
\end{proof}

\begin{lemma}\leanok
[\(\delta\)-wrapping on the unit pair from the \(\eta\)-bridge]
  \label{lem:isiso_sheafifydelta_unitpair_of_isiso_sheafifyeta}
  \lean{AlgebraicGeometry.Scheme.Modules.isIso_sheafifyDelta_unitPair_of_isIso_sheafifyEta}
  \uses{lem:w_of_isiso_sheafification, lem:whisker_of_W, lem:isiso_sheafification_map_of_W}
  If the sheafified presheaf unit comparison \(a_Y.\mathrm{map}\,(\eta\,(\mathtt{pullback}\,\varphi'))\)
  is an isomorphism, then so is the sheafified oplax cotensorator \(a_Y.\mathrm{map}\,\delta\)
  on the unit pair \((\mathcal{O},\mathcal{O})\).
\end{lemma}
\begin{proof}
  \uses{lem:w_of_isiso_sheafification, lem:whisker_of_W, lem:isiso_sheafification_map_of_W}
  Proved directly in Lean: the oplax left-unitality
  \(\mathtt{left\_unitality\_hom}\) factors \(\delta\,\mathbf{1}\,\mathbf{1}\) through the
  whiskered \(\eta\) and the (iso) unitor; sheafifying, the whiskered \(\eta\)-factor is an
  iso because a sheafified iso lies in \(J.W\) (\cref{lem:w_of_isiso_sheafification}), \(J.W\)
  is stable under right-whiskering (\cref{lem:whisker_of_W}), and \(J.W\)-morphisms sheafify
  to isos (\cref{lem:isiso_sheafification_map_of_W}).
\end{proof}

\begin{lemma}\leanok
[\(\mathtt{pullbackTensorMap}\) on the unit pair from the \(\eta\)-bridge]
  \label{lem:isiso_pullbacktensormap_unitpair_of_isiso_sheafifyeta}
  \lean{AlgebraicGeometry.Scheme.Modules.isIso_pullbackTensorMap_unitPair_of_isIso_sheafifyEta}
  \uses{lem:isiso_sheafifydelta_unitpair_of_isiso_sheafifyeta, lem:isiso_pullbacktensormap_of_sheafifydelta}
  If the sheafified presheaf unit comparison \(a_Y.\mathrm{map}\,(\eta\,(\mathtt{pullback}\,\varphi'))\)
  is an isomorphism, then \(\mathtt{pullbackTensorMap}\,f\,\mathcal{O}_X\,\mathcal{O}_X\) is an
  isomorphism.
\end{lemma}
\begin{proof}
  \uses{lem:isiso_sheafifydelta_unitpair_of_isiso_sheafifyeta, lem:isiso_pullbacktensormap_of_sheafifydelta}
  Proved directly in Lean: chain the \(\delta\)-wrapping
  \cref{lem:isiso_sheafifydelta_unitpair_of_isiso_sheafifyeta} into the reduction brick
  \cref{lem:isiso_pullbacktensormap_of_sheafifydelta} on \(M = N = \mathcal{O}\).
\end{proof}

\begin{lemma}\leanok
[Composite ring-map reconciliation]
  \label{lem:toringcatsheafhom_comp_hom_reconcile}
  \lean{AlgebraicGeometry.Scheme.Modules.toRingCatSheafHom_comp_hom_reconcile}
  For composable \(h : Z \to Y\), \(f : Y \to X\), the structure ring-presheaf map of the
  composite \(h \circ f\) factors as \(\varphi'_f\) followed by the whiskered \(\varphi'_h\).
  This feeds the presheaf \(\mathtt{pullbackComp}\) into the oplax \(\mathtt{comp\_}\delta\)
  decomposition (Sq2 of \cref{lem:pullback_tensor_map_basechange}).
\end{lemma}
\begin{proof}
  Proved directly in Lean by \(\mathtt{rfl}\): the two sides hold up to definitional
  equality at default transparency.
\end{proof}

\begin{lemma}\leanok
[Sq2b --- monoidality of the presheaf \(\mathtt{pullbackComp}\)]
  \label{lem:pullbackcomp_delta}
  \lean{AlgebraicGeometry.Scheme.Modules.pullbackComp_δ}
  \uses{lem:presheaf_pullback_oplaxmonoidal, lem:presheaf_pushforward_laxmonoidal,
        lem:pushforwardcomp_lax_mu}
  The oplax cotensorator \(\delta\) of the pullback for a composite ring map transports,
  through the pullback pseudofunctor coherence \(\mathtt{pullbackComp}\), into the
  \(\mathtt{Functor.OplaxMonoidal.comp}\) cotensorator of the composite pullback. This is the
  \(\eta \to \delta\) analogue, at the \(\mathtt{PresheafOfModules}\) level, of the unit
  coherence \cref{lem:pullbackObjUnitToUnit_comp}.
\end{lemma}
\begin{proof}
  \uses{lem:presheaf_pullback_oplaxmonoidal, lem:presheaf_pushforward_laxmonoidal}
  Proved directly in Lean by the adjunction-mate calculus: transpose under the
  pullback--pushforward adjunction, rewrite the oplax \(\delta\) as the mate of the lax
  \(\mu\) of the pushforward (\cref{lem:presheaf_pushforward_laxmonoidal}), and use the
  conjugate identity for \(\mathtt{pullbackComp}.\mathrm{inv}\) (here \(\mathtt{pushforwardComp}
  = \mathtt{Iso.refl}\)); the sole residual is the lax-\(\mu\) composition coherence
  \(\mathtt{pushforwardComp\_lax\_}\mu\), discharged by a sectionwise pure-tensor collapse.
\end{proof}

\begin{lemma}\leanok
[Sheaf-level conjugate of \(\mathtt{pullbackComp}.\mathrm{inv}\)]
  \label{lem:sheaf_unit_comp_pushforward_pullbackcomp_inv}
  \lean{AlgebraicGeometry.Scheme.Modules.sheaf_unit_comp_pushforward_pullbackComp_inv}
  % NOTE: this packaged R0-peel is no longer the route the Sq1 lemma
  % `lem:sheafificationcomppullback_comp` takes: that lemma now kills R0 INLINE via
  % `lem:homEquiv_conjugateEquiv_app_prime` (keeping `pullbackComp` opaque). This block is
  % retained as the standalone sheaf-level conjugate building block (its `\lean{}` names a live
  % decl) and may serve an alternative reassembly; it is currently not consumed in-chapter.
  For composable \(h : Z \to Y\), \(f : Y \to X\) and \(Q \in \Scheme.\mathtt{Modules}\,X\),
  the unit of the composite-pullback adjunction post-composed with the pushforward of
  \(\mathtt{pullbackComp}.\mathrm{inv}\) equals the unit of the composite of the \(f\)- and
  \(h\)-adjunctions post-composed with \(\mathtt{pushforwardComp}.\mathrm{hom}\). It is the
  cheap, sheafification-free R0-peel piece of the Sq1 mate calculus.
\end{lemma}
\begin{proof}
  Proved directly in Lean: the sheaf-level instance of \(\mathtt{unit\_conjugateEquiv}\)
  combined with \(\mathtt{conjugateEquiv\_pullbackComp\_inv}\) (the mate of
  \(\mathtt{pullbackComp}.\mathrm{inv}\) is \(\mathtt{pushforwardComp}.\mathrm{hom}\)).
\end{proof}

\begin{lemma}\leanok
[Sq1 --- composition coherence of \(\mathtt{sheafificationCompPullback}\)]
  \label{lem:sheafificationcomppullback_comp}
  \lean{AlgebraicGeometry.Scheme.Modules.sheafificationCompPullback_comp}
  \uses{lem:sheafification_comp_pullback_eq_leftadjointuniq,
        lem:homEquiv_conjugateEquiv_app_prime}
  For composable \(h : Z \to Y\), \(f : Y \to X\) and a presheaf \(P\) over \(X\), the
  sheafification--pullback comparison of \(h \circ f\) factors through the comparisons of
  \(f\) and \(h\), conjugated by the sheaf-level pullback pseudofunctor iso
  \(\mathtt{pullbackComp}\,h\,f\) and the sheafified presheaf pullback pseudofunctor iso.
  This is the \(\mathtt{sheafificationCompPullback}\) twin of \cref{lem:pullbackObjUnitToUnit_comp};
  it is the S1-foundational coherence consumed by \cref{lem:pullback_tensor_map_basechange}.
\end{lemma}
\begin{proof}  \uses{lem:sheafification_comp_pullback_eq_leftadjointuniq,
        lem:homEquiv_conjugateEquiv_app_prime}
  \textit{Source: internal categorical construction; no external reference.}

  \emph{Discharge in the conjugate calculus.} The decisive structural fact is that
  \(\mathtt{leftAdjointUniq}\) \emph{is}, definitionally, a conjugate of the identity: for
  adjunctions \(\mathrm{adj}_1, \mathrm{adj}_2\) to a common right adjoint \(R\),
  \((\mathtt{conjugateEquiv}\,\mathrm{adj}_1\,\mathrm{adj}_2)^{-1}(\mathbf{1}_R)
  = (\mathrm{adj}_1.\mathtt{leftAdjointUniq}\,\mathrm{adj}_2).\mathrm{inv}\) holds by \(\mathtt{rfl}\),
  and \(\mathtt{sheafificationCompPullback}\,\varphi
  = A_\varphi.\mathtt{leftAdjointUniq}\,B_\varphi\)
  (\cref{lem:sheafification_comp_pullback_eq_leftadjointuniq}), with
  \(A_\varphi = (\text{sheafification}_X)\mathbin{.\mathrm{comp}}(\text{pullback--pushforward}_\varphi)\)
  and \(B_\varphi\) its transpose-mate composite. Hence each comparison inverse
  \((\mathtt{sheafificationCompPullback}\,\varphi).\mathrm{inv}\) is
  \(\mathtt{conjugateEquiv}\,A_\varphi\,B_\varphi\) applied to \(\mathbf{1}_{R_\varphi}\), and the
  whole coherence is discharged inside Mathlib's conjugate calculus
  (\(\mathtt{Mathlib.CategoryTheory.Adjunction.Mates}\)) with \emph{no} \(\mathrm{homEquiv}\)
  telescope.

  \emph{Why this dissolves the obstruction.} The telescope route transposes the entire identity
  under the \emph{single} composite adjunction \(A_{h \circ f}\), so every comparison factor must be
  expressed through that one unit/counit; the intermediate-scheme sheafification
  \(\mathtt{sheafAdj}_Y\) carried by \(A_h\) is then neither a factor of \(A_{h \circ f}\) nor free in
  the residual, and would have to be slid in by hand. The fusion law
  \[
    \mathtt{conjugateEquiv}\,\mathrm{adj}_1\,\mathrm{adj}_2\,\alpha \mathbin{;}
    \mathtt{conjugateEquiv}\,\mathrm{adj}_2\,\mathrm{adj}_3\,\beta
      \;=\;
    \mathtt{conjugateEquiv}\,\mathrm{adj}_1\,\mathrm{adj}_3\,(\beta \mathbin{;} \alpha)
    \qquad (\mathtt{conjugateEquiv\_comp})
  \]
  instead keeps each comparison a conjugate \emph{with respect to its own adjunction pair}, and its
  \emph{middle} adjunction \(\mathrm{adj}_2\) is free --- it is exactly the intermediate-scheme
  composite \(A_Y / B_Y\) carrying \(\mathtt{sheafAdj}_Y\). So \(\mathtt{sheafAdj}_Y\) becomes a
  \emph{parameter} of \(\mathrm{adj}_2\), absorbed, never needing to be free in a transposed chain.
  This is the precise structural reason the cocycle closes where the telescope stalled.

  \emph{Steps.}
  \begin{enumerate}
    \item \textbf{Pass to inverses at the natural-transformation level.} Equivalently prove the
      identity of the inverses, dropping \(.\mathrm{app}\,P\) (recovered at the end by
      \(\mathtt{NatTrans.congr\_app}\)); rewrite each
      \((\mathtt{sheafificationCompPullback}\,\varphi).\mathrm{inv}\) as
      \((\mathtt{conjugateEquiv}\,A_\varphi\,B_\varphi)^{-1}(\mathbf{1}_{R_\varphi})\) by the two
      \(\mathtt{rfl}\) bridges above.
    \item \textbf{Apply the injective \(\mathtt{conjugateEquiv}\,A_{h \circ f}\,B_{h \circ f}\)} to
      both sides; the left-hand side, being this equivalence applied to its own
      \(\mathtt{symm}\) of \(\mathbf{1}\), collapses to \(\mathbf{1}\) by
      \(\mathtt{Equiv.apply\_symm\_apply}\).
    \item \textbf{Expose bare conjugates.} On the right-hand side the four factors --- the whiskered
      \((\mathtt{pullbackComp}\,h\,f).\mathrm{inv}\), the \((\mathtt{pullback}\,h)\)-pushed
      \(f\)-comparison, the \(h\)-comparison, and the sheafified presheaf coherence
      \(\delta_{\mathrm{pre}}\) --- are each a \emph{whiskered} conjugate; rewrite each to a whisker
      of a bare conjugate by \(\mathtt{conjugateEquiv\_whiskerLeft}\) /
      \(\mathtt{conjugateEquiv\_whiskerRight}\) (already used in-file at the former R0 step).
    \item \textbf{Fuse.} Apply \(\mathtt{conjugateEquiv\_comp}\) (three times) to fuse the four
      conjugates into a single \(\mathtt{conjugateEquiv}\,A_{h \circ f}\,B_{h \circ f}\) of the
      composite of coherences. Identify the pseudofunctor factors by
      \(\mathtt{conjugateEquiv\_pullbackComp\_inv}\) (the mate of
      \((\mathtt{pullbackComp}\,h\,f).\mathrm{inv}\) on the left adjoints is
      \((\mathtt{pushforwardComp}\,h\,f).\mathrm{hom}\) on the right adjoints) and the presheaf
      strictness \(\mathtt{pushforwardComp} = \mathtt{Iso.refl}\) (so the \(\delta_{\mathrm{pre}}\)
      conjugate is the identity).
    \item \textbf{Close on the strict cocycle.} The residual right-adjoint identity is the
      \(\mathtt{pushforwardComp}\) pseudofunctor cocycle, strict on the presheaf side
      (\(\mathtt{rfl}\)-level). This discharges the goal.
  \end{enumerate}

  The over-general \(\mathtt{mateEquiv\_vcomp}\)/\(\mathtt{TwoSquare}\) machinery
  (\(\mathtt{Mates.lean}\)) only \emph{justifies} \(\mathtt{conjugateEquiv\_comp}\) and is not needed
  directly: the comparisons here are identity-vertical mates, so \(\mathtt{conjugateEquiv\_comp}\) is
  the correct altitude (no explicit verticals to construct). The cautionary partial
  \(\mathtt{leftAdjointUniq\_trans}\) assumes a \emph{single} shared right adjoint and must not be
  forced onto the cross-scheme chain (each comparison has its own \(A/B\) pair) --- that common-\(G\)
  shape is precisely the telescope's trap.
\end{proof}

\begin{lemma}\leanok
[Conjugate transport of a left-functor map through \(\mathrm{homEquiv}\)]
  \label{lem:homEquiv_conjugateEquiv_app_prime}
  \lean{AlgebraicGeometry.Scheme.Modules.homEquiv_conjugateEquiv_app'}
  Let \(\mathrm{adj}_1 : L_1 \dashv R_1\) and \(\mathrm{adj}_2 : L_2 \dashv R_2\) be adjunctions
  between categories \(\mathcal{C}, \mathcal{D}\), let \(\alpha : L_2 \to L_1\) be a natural
  transformation of the left adjoints, and let \(f : L_1.\mathrm{obj}\,c \to d\). Then
  \[
    \mathrm{adj}_2.\mathrm{homEquiv}\,c\,d\,(\alpha.\mathrm{app}\,c \mathbin{;} f)
      \;=\;
    \mathrm{adj}_1.\mathrm{homEquiv}\,c\,d\,f \mathbin{;}
      (\mathtt{conjugateEquiv}\,\mathrm{adj}_1\,\mathrm{adj}_2\,\alpha).\mathrm{app}\,d.
  \]
  This is the device that lets the leading \(\mathtt{pullbackComp}\) factor of the transposed
  Sq1 goal (\cref{lem:sheafificationcomppullback_comp}) be peeled while keeping
  \(\mathtt{pullbackComp}\) \emph{opaque} (a term-mode transport, no rewrite on the iso),
  thereby avoiding the circular forget-distribution route. It is a project-local, \texttt{private}
  verbatim copy of the Mathlib-adjacent declaration
  \(\mathtt{AlgebraicGeometry.Scheme.Modules.homEquiv\_conjugateEquiv\_app}\)
  (reproduced because the chapter where the twin lives is not imported by this file);
  sorry-free.
\end{lemma}
\begin{proof}
  Both sides are computed by \(\mathtt{Adjunction.homEquiv\_unit}\): the left side is
  \(\mathrm{adj}_2.\mathrm{unit}.\mathrm{app}\,c \mathbin{;}
  R_2.\mathrm{map}(\alpha.\mathrm{app}\,c \mathbin{;} f)\) and the right side is
  \(\mathrm{adj}_1.\mathrm{unit}.\mathrm{app}\,c \mathbin{;} R_1.\mathrm{map}\,f \mathbin{;}
  (\mathtt{conjugateEquiv}\,\mathrm{adj}_1\,\mathrm{adj}_2\,\alpha).\mathrm{app}\,d\). Expanding
  \(R_2.\mathrm{map}\) over the composite and using the unit characterisation of the conjugate,
  \(\mathrm{adj}_2.\mathrm{unit}.\mathrm{app}\,c \mathbin{;} R_2.\mathrm{map}(\alpha.\mathrm{app}\,c)
  = \mathrm{adj}_1.\mathrm{unit}.\mathrm{app}\,c \mathbin{;}
  (\mathtt{conjugateEquiv}\,\mathrm{adj}_1\,\mathrm{adj}_2\,\alpha).\mathrm{app}(L_1.\mathrm{obj}\,c)\)
  (the \(\mathtt{unit\_conjugateEquiv}\) identity), followed by naturality of
  \(\mathtt{conjugateEquiv}\,\mathrm{adj}_1\,\mathrm{adj}_2\,\alpha\) at \(f\), identifies the two
  sides.
\end{proof}

\begin{lemma}\leanok
[Y-side sheafification right-triangle for the oplax unit comparison]
  \label{lem:pullback_sheafify_unit_eta_triangle}
  \lean{AlgebraicGeometry.Scheme.Modules.pullbackSheafifyUnitEtaTriangle}
  \uses{def:sheafify_unit_iso}
  For \(f : Y \to X\) with \(\varphi' = (\mathtt{toRingCatSheafHom}\,f).\mathrm{hom}\) and
  \(F = \mathtt{pullback}\,\varphi'\), the sheafification unit at \(F.\mathrm{obj}\,\mathbf{1}\),
  post-composed with the underlying presheaf maps of \(a_Y.\mathrm{map}\,(\eta\,F)\) and
  \(\mathtt{sheafifyUnitIso}.\mathrm{hom}\), recovers the oplax unit comparison \(\eta\,F\). It
  is a substep of the unit square \cref{lem:eta_bridge_unit_square}.
\end{lemma}
\begin{proof}
  \uses{def:sheafify_unit_iso}
  Proved directly in Lean: this is \(\mathtt{Equiv.apply\_symm\_apply}\) for the presheaf
  sheafification adjunction \(\mathtt{homEquiv}\) --- the second factor is
  \(\mathtt{homEquiv}^{-1}(\eta\,F)\), its counit leg being the adjunction counit
  \(\mathtt{sheafifyUnitIso}.\mathrm{hom}\) (\cref{def:sheafify_unit_iso}) at the sheaf
  \(\mathcal{O}_Y\) --- discharged by sheafification-unit naturality and the right triangle.
\end{proof}

\subsection{Change-of-rings, lax-\(\mu\) and \(\mathtt{extendScalars}\) helpers}
\label{subsec:tensorobj_changeofrings_helpers}

This subsection blueprints the change-of-base-ring cluster of
\texttt{TensorObjSubstrate.lean}: the extension-of-scalars functor and its adjunction,
the lax-monoidal structure maps \(\mu\) of \(\mathtt{restrictScalars}\) and
\(\mathtt{pushforward}\), and the monoidality of the pushforward composition cell
(\(\mathtt{pushforwardComp}\)) --- the change-of-rings coherence that is the sole residual
of the presheaf-level Sq2b lemma \cref{lem:pullbackcomp_delta}. Each helper is sorry-free;
its \(\mathtt{\backslash uses}\) edges are read off its Lean body.

\begin{lemma}\leanok
[\(\mathtt{restrictScalars}\) is a right adjoint]
  \label{lem:restrictscalars_is_right_adjoint}
  \lean{AlgebraicGeometry.Scheme.Modules.restrictScalarsIsRightAdjoint}
  For a morphism of ring presheaves \(\varphi : S \to F^{\mathrm{op}} \!\cdot\! R\), the
  restriction-of-scalars functor \(\mathtt{restrictScalars}\,\varphi\) on presheaves of
  modules is a right adjoint. This packages the existence of its left adjoint
  \(\mathtt{extendScalars}\,\varphi\).
\end{lemma}
\begin{proof}
  Proved directly in Lean: \(\mathtt{restrictScalars}\,\varphi\) is definitionally the
  fixed-target pushforward \(\mathtt{pushforward}\,(F = \mathbf{1})\,\varphi\), whose
  right-adjointness is the ambient Mathlib instance; conclude by \(\mathtt{inferInstanceAs}\).
\end{proof}

\begin{lemma}\leanok
[\(\mathtt{pushforward}_0\) is a right adjoint]
  \label{lem:pushforward0_is_right_adjoint}
  \lean{AlgebraicGeometry.Scheme.Modules.pushforward₀IsRightAdjoint}
  The fixed-ring presheaf pushforward \(\mathtt{pushforward}_0\,F\,R\) is a right adjoint;
  this carries the existence of the topological inverse image \(\mathtt{pullback}_0\).
\end{lemma}
\begin{proof}
  Proved directly in Lean: \(\mathtt{pushforward}_0\,F\,R\) is definitionally
  \(\mathtt{pushforward}\,(\mathbf{1}\,(F^{\mathrm{op}} \!\cdot\! R))\) (restriction of
  scalars along the identity is the identity functor), whose right-adjointness is the
  ambient Mathlib instance; conclude by \(\mathtt{inferInstanceAs}\).
\end{proof}

\begin{definition}\leanok
[Extension of scalars \(\mathtt{extendScalars}\)]
  \label{def:extend_scalars}
  \lean{AlgebraicGeometry.Scheme.Modules.extendScalars}
  \uses{lem:restrictscalars_is_right_adjoint}
  For \(\varphi : S \to F^{\mathrm{op}} \!\cdot\! R\), the extension-of-scalars functor
  \(\mathtt{extendScalars}\,\varphi := (\mathtt{restrictScalars}\,\varphi)
  .\mathrm{leftAdjoint}\), the left adjoint of restriction of scalars along \(\varphi\). It
  is the carrier of the strong-monoidal change-of-base on presheaves of modules.
\end{definition}
\begin{proof}
  \uses{lem:restrictscalars_is_right_adjoint}
  Proved directly in Lean: the left adjoint exists because
  \(\mathtt{restrictScalars}\,\varphi\) is a right adjoint
  (\cref{lem:restrictscalars_is_right_adjoint}).
\end{proof}

\begin{definition}\leanok
[The adjunction \(\mathtt{extendScalars} \dashv \mathtt{restrictScalars}\)]
  \label{def:extend_scalars_adjunction}
  \lean{AlgebraicGeometry.Scheme.Modules.extendScalarsAdjunction}
  \uses{def:extend_scalars, lem:restrictscalars_is_right_adjoint}
  The adjunction \(\mathtt{extendScalars}\,\varphi \dashv \mathtt{restrictScalars}\,\varphi\)
  witnessing extension of scalars as left adjoint to restriction of scalars.
\end{definition}
\begin{proof}
  \uses{def:extend_scalars, lem:restrictscalars_is_right_adjoint}
  Proved directly in Lean: \(\mathtt{Adjunction.ofIsRightAdjoint}\) applied to
  \(\mathtt{restrictScalars}\,\varphi\) (\cref{lem:restrictscalars_is_right_adjoint}).
\end{proof}

\begin{lemma}\leanok
[\(\mathtt{restrictScalars}(\mathbf{1})\) on morphisms]
  \label{lem:restrictscalarsid_map}
  \lean{AlgebraicGeometry.Scheme.Modules.restrictScalarsId_map}
  For a morphism \(g\) of presheaves of modules over \(R\),
  \((\mathtt{restrictScalars}\,(\mathbf{1}\,R)).\mathrm{map}\,g = g\): restriction of scalars
  along the identity ring map acts as the identity on morphisms.
\end{lemma}
\begin{proof}
  Proved directly in Lean by \(\mathtt{rfl}\): \(\mathtt{restrictScalars}\,(\mathbf{1})\) is
  definitionally the identity functor, stated as a propositional rewrite so the wrapper can
  be stripped syntactically without a whole-term \(\mathtt{whnf}\).
\end{proof}

\begin{lemma}\leanok
[Sectionwise value of the \(\mathtt{restrictScalars}\) tensorator \(\mu\)]
  \label{lem:restrictscalars_mu_app}
  \lean{AlgebraicGeometry.Scheme.Modules.restrictScalars_μ_app}
  \uses{def:restrict_scalars_lax_mu}
  The lax tensorator \(\mu\) of \(\mathtt{restrictScalars}\,\alpha\), evaluated at a section
  \(W\), is the \(\mathtt{ModuleCat}\) lax \(\mu\) of \(\mathtt{restrictScalars}\,(\alpha_W)\)
  on the sections. It exposes the ambient presheaf \(\mu\) in a form on which the
  pure-tensor rule matches syntactically.
\end{lemma}
\begin{proof}
  \uses{def:restrict_scalars_lax_mu}
  Proved directly in Lean by \(\mathtt{rfl}\): the presheaf tensorator
  \(\mathtt{restrictScalarsLax}\mu\) (\cref{def:restrict_scalars_lax_mu}) is defined
  sectionwise as the \(\mathtt{ModuleCat}\) tensorator.
\end{proof}

\begin{lemma}\leanok
[Pure-tensor value of the \(\mathtt{ModuleCat}\) \(\mathtt{restrictScalars}\) tensorator]
  \label{lem:forget2_restrictscalars_mu_hom_tmul}
  \lean{AlgebraicGeometry.Scheme.Modules.forget₂_restrictScalars_μ_hom_tmul}
  On a pure tensor \(m \otimes_t n\), the \(\mathtt{ModuleCat}\) lax tensorator \(\mu\) of
  \(\mathtt{restrictScalars}\,f\) (with \(\mathtt{forget}_2\)-carrier rings coming from
  presheaves of \emph{commutative} rings) is the identity: \(\mu(m \otimes_t n) =
  m \otimes_t n\).
\end{lemma}
\begin{proof}
  Proved directly in Lean: it is \(\mathtt{ModuleCat.restrictScalars\_}\mu\mathtt{\_tmul}\),
  the Mathlib pure-tensor formula for the change-of-rings tensorator, with the
  \(\mathtt{CommRing}\) instances supplied by the \(\mathtt{CommRingCat}\) carriers.
\end{proof}

\begin{lemma}\leanok
[Pure-tensor value of the presheaf \(\mathtt{restrictScalars}\) tensorator]
  \label{lem:restrictscalars_mu_app_tmul}
  \lean{AlgebraicGeometry.Scheme.Modules.restrictScalars_μ_app_tmul}
  \uses{lem:restrictscalars_mu_app}
  On a pure tensor, \((\mu(\mathtt{restrictScalars}\,\alpha)\,M_1\,M_2)_W
  (m \otimes_t n) = m \otimes_t n\): the presheaf \(\mathtt{restrictScalars}\) tensorator is
  the identity on pure tensors.
\end{lemma}
\begin{proof}
  \uses{lem:restrictscalars_mu_app}
  Proved directly in Lean: expose the \(\mathtt{ModuleCat}\) tensorator by
  \cref{lem:restrictscalars_mu_app}, then apply the Mathlib pure-tensor rule
  \(\mathtt{ModuleCat.restrictScalars\_}\mu\mathtt{\_tmul}\).
\end{proof}

\begin{lemma}\leanok
[Pure-tensor value of the \(\mathtt{pushforward}\)-mapped \(\mathtt{restrictScalars}\) tensorator]
  \label{lem:pushforward_map_restrictscalars_mu_app_tmul}
  \lean{AlgebraicGeometry.Scheme.Modules.pushforward_map_restrictScalars_μ_app_tmul}
  \uses{lem:restrictscalars_mu_app_tmul}
  The outer leg of \cref{lem:pushforwardcomp_lax_mu}: applied to a pure tensor,
  \(((\mathtt{pushforward}\,\varphi).\mathrm{map}\,(\mu(\mathtt{restrictScalars}\,\psi)
  \,N_1\,N_2))_W\) is the identity.
\end{lemma}
\begin{proof}
  \uses{lem:restrictscalars_mu_app_tmul}
  Proved directly in Lean: reindex the section map through \(\mathtt{pushforward\_map\_app\_apply}\)
  (\(\mathtt{pushforward}\,\varphi\) is \(\mathtt{pushforward}_0 \cdot \mathtt{restrictScalars}\,\varphi\),
  so the section map at \(W\) is the \(\mu\) at \(F^{\mathrm{op}}.\mathrm{obj}\,W\)), then collapse
  by \cref{lem:restrictscalars_mu_app_tmul}.
\end{proof}

\begin{lemma}\leanok
[Reduction of the \(\mathtt{pushforward}\) tensorator to the \(\mathtt{restrictScalars}\) \(\mu\)]
  \label{lem:pushforward_mu_eq}
  \lean{AlgebraicGeometry.Scheme.Modules.pushforward_μ_eq}
  \uses{lem:presheaf_pushforward_laxmonoidal, def:restrict_scalars_lax_mu}
  The lax tensorator \(\mu\) of a single \(\mathtt{pushforward}\,\varphi\) equals the lax
  tensorator of the change-of-rings \(\mathtt{restrictScalars}\,\varphi'\) on the
  (strongly-monoidal, \(\mu_{\mathrm{Iso}} = \mathbf{1}\)) reindexed objects
  \(\mathtt{pushforward}_0^{\mathrm{Comm}}\,F\,R_0\). This unfolds the opaque pushforward
  lax-monoidal \(\mu\) into the directly computable \(\mathtt{restrictScalars}\) \(\mu\).
\end{lemma}
\begin{proof}
  \uses{lem:presheaf_pushforward_laxmonoidal, def:restrict_scalars_lax_mu}
  Proved directly in Lean by \(\mathtt{rfl}\): the pushforward lax-monoidal structure
  (\cref{lem:presheaf_pushforward_laxmonoidal}) is the \(\mathtt{Functor.LaxMonoidal.comp}\)
  of the (identity) \(\mathtt{pushforward}_0\) tensorator with the
  \(\mathtt{restrictScalars}\) tensorator (\cref{def:restrict_scalars_lax_mu}), which is
  definitionally the right-hand side.
\end{proof}

\begin{lemma}\leanok
[Sq2b residual --- monoidality of the \(\mathtt{pushforward}\) composition cell]
  \label{lem:pushforwardcomp_lax_mu}
  \lean{AlgebraicGeometry.Scheme.Modules.pushforwardComp_lax_μ}
  \uses{lem:pushforward_mu_eq, lem:restrictscalars_mu_app,
        lem:forget2_restrictscalars_mu_hom_tmul,
        lem:pushforward_map_restrictscalars_mu_app_tmul}
  The change-of-rings coherence ``\(\mathtt{pushforwardComp}\) is monoidal'': the lax
  tensorator \(\mu\) of the composite pushforward
  \(\mathtt{pushforward}\,\psi \cdot \mathtt{pushforward}\,\varphi\) (built by
  \(\mathtt{Functor.LaxMonoidal.comp}\)) agrees with the lax tensorator of the single
  pushforward \(\mathtt{pushforward}\,(\varphi \,;\, F^{\mathrm{op}} \!\lhd\! \psi)\). It is
  the sole genuine residual of the presheaf-level Sq2b lemma \cref{lem:pullbackcomp_delta}.
\end{lemma}
\begin{proof}
  \uses{lem:pushforward_mu_eq, lem:restrictscalars_mu_app,
        lem:forget2_restrictscalars_mu_hom_tmul,
        lem:pushforward_map_restrictscalars_mu_app_tmul}
  Proved directly in Lean by a sectionwise pure-tensor collapse: \(\mathtt{hom\_ext}\) then
  \(\mathtt{tensor\_ext}\), unfold the composite \(\mu\) by \(\mathtt{Functor.LaxMonoidal.comp\_}\mu\),
  lighten each \(\mu(\mathtt{pushforward}\,\_)\) to a \(\mathtt{restrictScalars}\) \(\mu\) by
  \cref{lem:pushforward_mu_eq}, expose the \(\mathtt{ModuleCat}\) \(\mu\) by
  \cref{lem:restrictscalars_mu_app}, and collapse the inner leg by
  \cref{lem:forget2_restrictscalars_mu_hom_tmul} and the outer
  \((\mathtt{pushforward}\,\varphi).\mathrm{map}\) leg by
  \cref{lem:pushforward_map_restrictscalars_mu_app_tmul}; the two collapsed pure tensors are
  then definitionally equal. (The equality is genuinely not \(\mathtt{rfl}\) at the presheaf
  level --- the \(\mathtt{restrictScalars}\) \(\mu\) on a pure tensor is real base-change
  content --- so the atomic-object helpers are applied as explicit arguments and matched by
  \(\mathtt{erw}\) to avoid a \(\mathtt{whnf}\)-explosion.)
\end{proof}

\begin{lemma}\leanok
[Forgetful image of a sheaf-pushforward map]
  \label{lem:forget_map_pushforward_map}
  \lean{AlgebraicGeometry.Scheme.Modules.forget_map_pushforward_map}
  The forgetful functor \(\mathtt{SheafOfModules.forget}\) intertwines the sheaf-level
  \(\mathtt{SheafOfModules.pushforward}\,\varphi\) with the presheaf-level
  \(\mathtt{PresheafOfModules.pushforward}\,\varphi.\mathrm{hom}\): for a morphism \(g\) of
  sheaves of modules, \(\mathtt{forget}.\mathrm{map}\,((\mathtt{pushforward}\,\varphi)
  .\mathrm{map}\,g) = (\mathtt{PresheafOfModules.pushforward}\,\varphi.\mathrm{hom})
  .\mathrm{map}\,(\mathtt{forget}.\mathrm{map}\,g)\). This is the STEP-1 binding
  reconciliation of the Sq1 tail.
\end{lemma}
\begin{proof}
  Proved directly in Lean by \(\mathtt{rfl}\): the sheaf pushforward is built sectionwise
  from the presheaf pushforward and \(\mathtt{forget}\) is the underlying-presheaf
  projection, so the two sides are definitionally equal.
\end{proof}

\subsection{The Picard quotient: setoid, multiplication and inverse}
\label{subsec:tensorobj_pic_quotient_helpers}

This subsection blueprints the three quotient-level helpers underlying the
invertibility-carrier Picard group \cref{def:pic_carrier} and its abelian-group law
\cref{thm:pic_commgroup}: the isomorphism setoid, the tensor-induced multiplication, and
the dual-induced inverse on iso-classes.

\begin{definition}\leanok
[The isomorphism setoid of invertible modules]
  \label{def:pic_setoid}
  \lean{AlgebraicGeometry.Scheme.Modules.picSetoid}
  \uses{def:scheme_modules_isinvertible}
  The setoid on \(\otimes\)-invertible \(\mathcal{O}_X\)-modules with \(M \sim M'\) iff there
  is an isomorphism \(M \cong M'\). The quotient of this setoid is the carrier of the Picard
  group \(\Pic X\).
\end{definition}
\begin{proof}
  \uses{def:scheme_modules_isinvertible}
  Proved directly in Lean: reflexivity, symmetry and transitivity of the
  \(\mathtt{Nonempty}\,(M \cong M')\) relation are witnessed by \(\mathtt{Iso.refl}\),
  \(\mathtt{Iso.symm}\) and \(\mathtt{Iso.trans}\).
\end{proof}

\begin{definition}\leanok
[Multiplication on \(\Pic X\) induced by \(\otimes_X\)]
  \label{def:pic_mul}
  \lean{AlgebraicGeometry.Scheme.Modules.picMul}
  \uses{def:pic_setoid, def:pic_carrier, def:tensorobjisoofiso, lem:isinvertible_tensor}
  The multiplication \([M] \cdot [M'] := [M \otimes_X M']\) on \(\Pic X\), well-defined on
  iso-classes.
\end{definition}
\begin{proof}
  \uses{def:pic_setoid, def:pic_carrier, def:tensorobjisoofiso, lem:isinvertible_tensor}
  Proved directly in Lean by \(\mathtt{Quotient.lift}_2\): the value lands in \(\Pic X\)
  because the tensor of invertibles is invertible (\cref{lem:isinvertible_tensor}), and it
  respects the setoid because \(\mathtt{tensorObjIsoOfIso}\) (\cref{def:tensorobjisoofiso})
  turns a pair of isomorphisms into an isomorphism of tensors.
\end{proof}

\begin{definition}\leanok
[Inverse on \(\Pic X\) induced by the tensor-inverse]
  \label{def:pic_inv}
  \lean{AlgebraicGeometry.Scheme.Modules.picInv}
  \uses{def:pic_setoid, def:pic_carrier, lem:isinvertible_inverse_welldef, lem:tensorobj_comm_iso}
  The inverse class \([M] \mapsto [N]\) on \(\Pic X\), where \(N\) is a chosen tensor-inverse
  of \(M\).
\end{definition}
\begin{proof}
  \uses{def:pic_setoid, def:pic_carrier, lem:isinvertible_inverse_welldef, lem:tensorobj_comm_iso}
  Proved directly in Lean by \(\mathtt{Quotient.lift}\): a chosen inverse witness
  (\(\mathtt{Classical.choose}\) of the invertibility) is itself invertible, with the
  braiding (\cref{lem:tensorobj_comm_iso}) supplying the membership isomorphism; the choice
  is independent of the representative because any two tensor-inverses of \(M\) are
  isomorphic (\cref{lem:isinvertible_inverse_welldef}).
\end{proof}
